\documentclass[12pt]{amsart}
\usepackage{amsmath}
\usepackage{amssymb}
\usepackage{xcolor}
\usepackage{mathrsfs}
\usepackage{stmaryrd}
\usepackage{longtable}
\usepackage{multirow}
\usepackage{makecell}
\usepackage{comment}
\usepackage{tablefootnote}
\usepackage{enumerate}
\usepackage{tikz-cd}
\usepackage{adjustbox}
\usepackage[lowtilde]{url}
\usepackage{hyperref}
%\usepackage{showkeys}
%\usepackage{cite}
%\usepackage{pxfonts}
\renewcommand{\arraystretch}{1.5}

\newcommand{\C}{\mathbb{C}}
\newcommand{\A}{\mathcal{A}}
\newcommand{\B}{\mathcal{B}}
\newcommand{\U}{\mathcal{U}}
\newcommand{\I}{\mathcal{I}}
\newcommand{\J}{\mathcal{J}}
\newcommand{\M}{\mathcal{M}}
\newcommand{\g}{\mathfrak{g}}
\newcommand{\Orb}{\mathbb{O}}
\newcommand{\Walg}{\mathcal{W}}
\newcommand{\Z}{\mathbb{Z}}
\newcommand{\gr}{\operatorname{gr}}
\newcommand{\kf}{\mathfrak{k}}
\newcommand{\tf}{\mathfrak{t}}
\newcommand{\hf}{\mathfrak{h}}
\newcommand{\uf}{\mathfrak{u}}
\newcommand{\lf}{\mathfrak{l}}
\newcommand{\gl}{\mathfrak{gl}}
\newcommand{\su}{\mathfrak{su}}

\newcommand{\HC}{\operatorname{HC}}
\newcommand{\q}{\mathfrak{q}}
\newcommand{\Pb}{\mathbb{P}}
\newcommand{\SL}{\operatorname{SL}}
\newcommand{\SO}{\operatorname{SO}}
\newcommand{\Sp}{\operatorname{Sp}}
\newcommand{\SU}{\operatorname{SU}}
\newcommand{\Spin}{\operatorname{Spin}}
\newcommand{\Str}{\mathcal{O}}
\newcommand{\jet}{J^\infty}
\newcommand{\Spec}{\operatorname{Spec}}
\newcommand{\send}{\mathcal{E}nd}
\newcommand{\Id}{\operatorname{Id}}
\newcommand{\pr}{\operatorname{pr}}

% Shilin's commands

\newcommand{\R}{\mathbb{R}}
\newcommand{\ul}{\underline}
\newcommand{\T}{\mathcal{T}}
\newcommand{\F}{\mathcal{F}}
\newcommand{\TT}{\mathscr{T}}
\newcommand{\LL}{\mathcal{L}}
\newcommand{\Ttwo}{\TT_{\frac{1}{2}}}
\newcommand{\tatp}{\TT^+(\Q, Y)}
\newcommand{\tat}{\TT(\Q,Y)}
\newcommand{\pic}{\mathscr{PA}}
\newcommand{\pay}{\pic_{\kf}(Y)}
\newcommand{\quan}{\mathcal{Q}}
\newcommand{\per}{\operatorname{Per}}
\newcommand{\cm}{\C^\times}
\newcommand{\OO}{\mathcal{O}}
\newcommand{\Q}{\OO_\hbar}
\newcommand{\lp}{\mathfrak{p}}
\newcommand{\z}{\mathfrak{z}}
\newcommand{\rf}{\mathfrak{r}}
\newcommand{\GR}{G_\R}
\newcommand{\Gad}{G_{ad}}
\newcommand{\KR}{K_\R}
\newcommand{\Lh}{\mathcal{L}_\hbar}
\newcommand{\wt}{\widetilde}
\newcommand{\slc}{\mathfrak{sl}(2, \C)}
\newcommand{\xreg}{X^{reg}}
\newcommand{\txreg}{\wt{X}^{reg}}
\newcommand{\spec}{\operatorname{Spec}}
\newcommand{\EE}{\mathcal{E}}
\newcommand{\Ehb}{\EE_\hbar}
\newcommand{\pws}{\C \llbracket \hbar \rrbracket}
\newcommand{\wHC}{\operatorname{wHC}}
\newcommand{\Weyl}{\mathbb{A}}
\newcommand{\slf}{\mathfrak{sl}}
\newcommand{\Dcal}{\mathcal{D}}
\newcommand{\Hom}{\operatorname{Hom}}
\newcommand{\param}{\mathfrak{P}}
\newcommand{\Ecal}{\mathcal{E}}
\newcommand{\GL}{\operatorname{GL}}
\newcommand{\loc}{\operatorname{Loc}}
\newcommand{\coh}{\operatorname{Coh}}
\newcommand{\flag}{\mathcal{B}}

\newcommand{\der}{{\operatorname{Der}}}
\newcommand{\aut}{{\operatorname{Aut}}}
\newcommand{\autd}{\aut(\Dcal)}
\newcommand{\derd}{\der(\Dcal)}
\newcommand{\autder}{\langle \autd, \derd  \rangle}
\newcommand{\GHC}{\langle \mathcal{G}, \mathfrak{G} \rangle}




\newtheorem{Thm}{Theorem}[subsection]
%\newtheorem*{MainThm}{Main Theorem}
\newtheorem{Prop}[Thm]{Proposition}
\newtheorem{Cor}[Thm]{Corollary}
\newtheorem{Lem}[Thm]{Lemma}
\theoremstyle{definition}
\newtheorem{Ex}[Thm]{Example}
\newtheorem{defi}[Thm]{Definition}
\newtheorem{Rem}[Thm]{Remark}
\newtheorem{Conj}[Thm]{Conjecture}
\numberwithin{equation}{section}
\title{On Harish-Chandra modules over quantizations of nilpotent orbits}
\author{Ivan Losev and Shilin Yu}
\oddsidemargin=0cm
\evensidemargin=0cm
\textwidth=16cm
\textheight=224mm
\unitlength=1mm
\begin{document}
\begin{abstract}
Let $G$ be a semisimple algebraic group over the complex numbers and $K$ be a connected reductive group mapping to $G$
so that the Lie algebra of $K$ gets identified with a symmetric subalgebra of $\g$. So we can talk about Harish-Chandra 
$(\g,K)$-modules, where $\g$ is the Lie algebra of $G$.
The goal of this paper is to give a geometric classification of irreducible Harish-Chandra modules with full support 
over the filtered quantizations of the algebras of the form $\C[\Orb]$, where $\Orb$ is a nilpotent orbit in $\g$ with 
codimension of the boundary at least $4$. Namely, we embed the set of isomorphism classes of irreducible Harish-Chandra modules into the set of 
isomorphism classes of irreducible $K$-equivariant suitably twisted local systems on $\Orb\cap \mathfrak{k}^\perp$. We show that
under certain conditions, for example when $K\subset G$ or when $\g\cong \mathfrak{so}_n,\mathfrak{sp}_{2n}$, this embedding is 
in fact a bijection. On the other hand, for $\g=\mathfrak{sl}_n$ and $K=\operatorname{Spin}_n$, the embedding is not 
bijective and we give a description of the image. Finally, we perform a partial classification for exceptional Lie algebras. 
\end{abstract}
\maketitle
\tableofcontents
\section{Introduction}
\subsection{Harish-Chandra modules}
Our base field is $\C$. Let $G$ be a semisimple simply connected algebraic group,
$\g$ be its Lie algebra and $\U$
be the universal enveloping algebra of $\g$. Let $K$ be a connected algebraic group equipped with
a homomorphism to $G$ such that the induced homomorphism of Lie algebras is injective and
identifies the Lie algebra $\kf$ with the symmetric subalgebra $\g^{\sigma}$ for some involution
$\sigma$. Note that we do not require that $K$ is a subgroup of $G$. 

By a {\it Harish-Chandra} (shortly, HC) {$(\g,K)$-module} one means a finite length $\U$-module such that the action of $\kf$ integrates to that of $K$. A primary reason to be interested in
HC modules is that they are related to representations of real forms of $G$ and their nonlinear covers.
Namely, when $K\hookrightarrow G$, the relevant real group is the real form of $G$ corresponding
to the involution $\sigma$, while, in general, we need to consider a cover $\tilde{G}_{\R}$ of $G_\mathbb{R}$ with $\ker[\tilde{G}_{\R}\rightarrow G_\R]\cong \ker[K\rightarrow G]$.

A basic question motivating the present work is: given a primitive ideal $\J$ (=the annihilator of
an irreducible module) in $\U$ classify all irreducible HC $(\g,K)$-modules whose annihilator
coincides with $\J$.

In this paper we will only deal with very special (but also very interesting) ideals $\J$.


\subsection{Quantizations of nilpotent orbits}
Let $\Orb$ be a nilpotent orbit in $\g^*$, $\overline{\Orb}$ be its closure and $\partial\Orb$
be the boundary. Until the end of Introduction, we assume that
$$\operatorname{codim}_{\overline{\Orb}}\partial\Orb\geqslant 4.$$
The algebra of regular functions $\C[\Orb]$ is a graded Poisson algebra. In this case
it makes sense to speak about filtered quantizations of $\C[\Orb]$. These are filtered
algebras $\A$ with an isomorphism $\gr\A\xrightarrow{\sim} \C[\Orb]$ of graded Poisson algebras. We have the following properties.
\begin{itemize}
\item The quantizations are parameterized by $H^2(\Orb,\C)$. This is a special case
of \cite[Theorem 3.4]{orbit}.
\item The action of $G$ lifts from $\C[\Orb]$ to $\A$. The lift possesses a unique quantum comoment map, i.e., a $G$-equivariant algebra homomorphism
$\Phi:U(\g)\rightarrow \A$ lifting the classical comoment map $\g\rightarrow \C[\Orb]$.
See \cite[Section 5.3]{orbit}.
\end{itemize}

In this paper we consider the ideals $\J$ that arise as kernels of the quantum comoment maps
$\Phi:U(\g)\rightarrow \A$.
Here is a reason why one cares about such ideals. Consider the {\it canonical} (in the terminology
from \cite{LMM}) quantization $\A_0$, its parameter is $0\in H^2(\Orb,\C)$. The corresponding
kernel $\J$ is a {\it unipotent ideal} in the terminology of \cite{LMM}. All such ideals are maximal,
\cite[Theorem 8.5.1]{LMM} and \cite[Theorem 1.0.2]{MM}. As argued in \cite[Section 6.3.4]{LMM}, all HC $(\g,K)$-modules annihilated by $\J$
should be considered unipotent. Such modules play a special role: one hopes that they are unitary and the general unitary 
modules are obtained from unipotent ones by operations that are relatively easy to understand. 

On the other hand, assume that $H^2(\Orb,\C)=\{0\}$. All orbits $\Orb$ satisfying this condition
(and the codimension $4$ condition above) are birationally rigid (in the sense of
\cite[Definition 1.2]{orbit}).
As a partial converse, all rigid orbits $\Orb$
satisfy $H^2(\Orb,\C)=\{0\}$ and all but six of them (in exceptional types) satisfy
the codimension 4 condition, see \cite[Lemma 2.4]{reg}. Unipotent modules
for rigid orbits play a special role because they cannot be induced.

\subsection{Main results}
Now we state the main results of the paper. While we do get the classification for all quantizations,  we concentrate on the case when
the quantization of $\C[\Orb]$ is canonical, hence $\J$ is the unipotent ideal associated
with $\Orb$. Let $M$ be an irreducible HC $(\g,K)$-module
annihilated by $\J$.
A result of Vogan, \cite[Theorem 1.2]{Vogan}, says that the associated variety of $M$ is irreducible
hence is the closure of a single $K$-orbit in $\Orb\cap \kf^\perp$, where we write $\kf^\perp$
for the annihilator in $\g^*$ of (the image of) $\kf$. So, let $\Orb_K$ denote the open
orbit in the associated variety. We need to classify all irreducible HC modules $M$ annihilated by $\J$ with associated variety
$\overline{\Orb}_K$.

Pick a good filtration on $M$. One can show the restriction of $\gr M$ to $\Orb_K$ is a (strongly) $K$-equivariant twisted local system with half-canonical twist, see \cite[Theorem 8.7]{Vogan} for a related result, we will suitably generalize it in our setting, see Section \ref{subsec:W-algebra_nilpotent}.  We will also see that the twisted local system is irreducible and non-isomorphic irreducible
modules give rise to non-isomorphic twisted local systems.

Here is the first main result of the paper (Theorem \ref{Thm:main1}).

\begin{Thm}\label{Thm:omnibus}
Suppose that one of the following conditions hold:
\begin{enumerate}
\item $K\hookrightarrow G$,
\item $\g$ is of type $B,C$ or $D$,
\item or $\g=\slf_n$  and $\kf\neq \mathfrak{so}_n$.
\end{enumerate}
Then there is a bijection between
\begin{itemize}
\item[(i)]
the irreducible HC $(\g,K)$-modules annihilated by $\J$
with associated variety $\overline{\Orb}_K$,
\item[(ii)] And irreducible equivariant twisted (by half-canonical twist) local systems
on $\Orb_K$.
\end{itemize}
\end{Thm}

We remark that one can describe (ii) in elementary representation theoretic terms.
We will do so in the main body of the paper.

In the case when $\g=\slf_n$ and $K=\operatorname{Spin}_n$ the map from (i) to (ii)
is not surjective (it is known, see, e.g. \cite[Example 12.4]{Vogan}, that in the case when $n=3$ there are four irreducible twisted local
systems and three irreducible HC modules, this case  is of great importance for the paper). 
Our second main result is a description of the image, see Theorem \ref{Thm:spin_HC_classif_general}.

We will also compute some examples in exceptional types when $K\rightarrow G$
is not injective. In these cases, it is also a general phenomenon that some of the irreducible twisted local systems do not correspond to any irreducible HC modules.

\begin{Rem}
An easy case is when $\operatorname{codim}_{\overline{\Orb}_k}\partial\Orb_K\geqslant 3$.
Here we also have a bijection between (i) and (ii) in the theorem. This was proved
in \cite{LY}. In this paper we will explain an alternative approach based on the
previous results of the first named author.
\end{Rem}

\begin{Rem}
We now compare the results of this paper to classification results
from \cite{BMSZ1,BMSZ2,BMSZ3,BMSZ4}. Compared to these papers, our description is more geometric and uniform, we do not require the unipotent representations we classify to be special. Also, in our approach, conceptually
non-linear groups are not more difficult than linear ones, and exceptional groups are not more
difficult than classical ones (with a few exceptions; also technically non-linear/exceptional groups are somewhat
more complicated). On the other hand, the four papers mentioned above avoid the restrictive codimension restriction
that is crucial for us. Also these papers prove the unitarity of all the special unipotent representations of real classical Lie groups.
\end{Rem}

\subsection{Description of approach}\label{SS_intro_approach}
Recall that for the time being we assume that $\J$ is a unipotent ideal, hence maximal, ideal.


First, let us explain a representation theoretic description of the equivariant
half-canonical twisted local systems on $\Orb_K$. Pick a point $\chi\in \Orb_K$.
Let $K_Q$ denote the reductive part of the stabilizer of $\chi$ in $K$. This is a
reductive group, in general, it is disconnected. Let $\omega_{\Orb_K}$ denote the canonical
bundle on $\Orb_K$, it is $K$-equivariant. So $K_Q$ acts on the fiber $(\omega_{\Orb_K})_\chi$
by a character. Let $\rho_{\omega}\in (\kf_Q^*)^{K_Q}$ be one half of this character.
Let $(K_Q,\rho_\omega)\operatorname{-mod}$ denote the full subcategory in the category
of rational representations of $K_Q$ consisting of all modules, where $\kf_Q$ acts via
$\rho_\omega$ (for the action of $\kf_Q$ obtained by differentiating the $K_Q$-action).
This is a semisimple category. It is equivalent to the category of local systems of interest:
an equivalence is given by taking the fiber of a twisted local system at $\chi$. Below in this section
we will identify the category of twisted local systems of interest with $(K_Q,\rho_\omega)\operatorname{-mod}$. Note that a $K_Q$-representation $V$ in $(K_Q,\rho_\omega)\operatorname{-mod}$ is called \emph{admissible} (with respect to $K$) in \cite[Definition 7.13]{Vogan} and the pair $(\chi, V)$ is called a \emph{(nilpotent) admissible $K$-orbit datum}. We also say that the twisted local system corresponding to $V$ is an \emph{admissible $K$-equivariant vector bundle} over $\Orb_K$. 


Set $\A:=\U/\J$, in fact, since $\J$ is maximal, the algebra $\A$ admits a filtration making it a filtered quantization of $\C[\Orb]$. Let $\HC(\A,K)$ denote the category of all HC $(\g,K)$-modules annihilated by $\J$\footnote{In the general case -- when $\A$ is an arbitrary quantization of $\C[\Orb]$ -- we only have an embedding $\U/\J\hookrightarrow \A$; still
there is a natural bijection between irreducible HC modules over $\U$ with annihilator $\J$ and certain irreducible HC modules over $\A$,
see Proposition \ref{Prop:HC_bijection}}. Let $\HC_{\Orb_K}(\A,K)$  denote the full subcategory of $\HC(\A,K)$ consisting of all modules whose associated variety is contained in $\overline{\Orb}_K$.

The first step in our construction is to produce a fully faithful embedding
$\bullet_{\dagger,\chi}: \HC_{\Orb_K}(\A,K)\xrightarrow{\sim} (K_Q,\rho_\omega)\operatorname{-mod}$, see
Proposition \ref{Prop:adj_properties} for a more general statement.
The construction of this functor was sketched in \cite[Section 6.1]{Wdim}
and the claim that it is a full embedding is proved as the similar claim for
Harish-Chandra bimodules, \cite[Theorem 1.3.1]{HC}. On the level of objects, the functor sends a HC module
$M$ to the twisted local system $(\gr M)|_{\Orb_K}$. This implies that if $M$ is irreducible,
then so is the corresponding twisted local system. Moreover, non-isomorphic HC modules
give rise to non-isomorphic twisted local systems. What remains is to describe
the essential image of $\bullet_{\dagger,\chi}$.

Here our inspiration comes from \cite{HC_symp}. This part is technical so we will describe
it here in a somewhat informal fashion. Any equivariant twisted local system
on $\Orb_K$ can be uniquely quantized to a sheaf of modules over the microlocalization of
$\A$ to $\Orb$. An irreducible twisted local system $\mathcal{E}$
lies in the image of $\bullet_{\dagger,\chi}$
if and only if the corresponding quantizations has nonzero global sections. To check whether
the quantization has global sections, we first push it to the locus in $\overline{\Orb}_K$
obtained by taking the union of $\Orb_K$ and all codimension $2$ orbits. The property
necessary for $\mathcal{E}\in \operatorname{Im}(\bullet_{\dagger,\chi})$ is that this
pushforward is still a {\it coherent} module, see \cite[Section 2.5]{HC_symp}. It is also close to being sufficient so for the purposes
of this section we will treat it as necessary and sufficient.

A key observation of \cite{HC_symp} (in the context of HC bimodules over quantizations of
conical symplectic singularities)  is that the condition that the pushforward of the quantization of
$\mathcal{E}$ is coherent can be tested on slices to codimension $2$ orbit. Namely, let $\Orb'_K\subset \overline{\Orb}_K$ be a codimension $2$ orbits. Let $S'$ be a slice in $\overline{\Orb}$ to a point of $\Orb'$ (this definition works as stated when $\overline{\Orb}$
is normal at the points of $\Orb'$, in general, in requires a modification that we will make in the main body of the paper). 
We can consider the restriction $\mathcal{L}|_{S'}$, this is a twisted local system
on $S'\cap \Orb_K$. The claim that the pushforward of the quantization of $\mathcal{E}$ is coherent
is equivalent to the following condition holding for every codimension
$2$ orbit in $\overline{\Orb}_K$:
\begin{itemize}
\item[(*)] The pushforward of the quantization of $\mathcal{E}|_{S'}$
from $S'\cap \Orb_K$ to $S'\cap \overline{\Orb}_K$ is coherent.
\end{itemize}
All slices $S'$ are known, see \cite{KP1,KP2} for the classical groups, and \cite{FJLS1} for
the exceptional groups. In the setting of Theorem  \ref{Thm:omnibus} we check that (*) is always
satisfied (there is one orbit in type $E_8$, where the slice is too complicated to analyze (*), but
we can use other methods -- Atlas -- in that case). For $\g=\mathfrak{sl}_n$ and $K=\operatorname{Spin}_n$
and in the number of cases in exceptional types, (*) reduces (modulo some caveats)
to understanding the case of $\g=\mathfrak{sl}_3, K=\operatorname{SL}_2$ and the minimal orbit $\Orb$.



{\bf Acknowledgements}: We would like to thank Jeff Adams, Jia-Jun Ma, Lucas Mason-Brown for helpful discussions on computations using Atlas. We are also grateful to Baohua Fu for explaining the papers \cite{FJLS1} and \cite{FJLS2}. The second author wants to thank Binyong Sun for his hospitality and inspiring discussions during his visits to Institute for Advanced Study in Mathematics of Zhejiang University. The work of I.L. has been partially supported by the NSF under grant DMS-2001139. The work of S.Y. has been partially supported by China NSFC grants (Grant No. 12001453 and 12131018) and Fundamental Research Funds for the Central Universities (Grant No. 20720200067 and 20720200071).

%\subsection{Content of the paper}

%\subsection{Notation and conventions}

\section{Harish-Chandra modules over filtered algebras}
The goal of this section is to establish a general framework for working with Harish-Chandra modules.
\subsection{HC $(\A,\theta)$-modules}\label{SS_HC_mod}
Our setting is as follows. The base field is $\C$. Let
\begin{itemize}
\item $d$ be a positive integer,
\item $\A=\bigcup_{i\geqslant 0}\A_{\leqslant i}$ be a filtered associative
algebra,
\item $\epsilon$ be a primitive $d$th root of $1$,
\item and $\zeta$ be a filtration preserving automorphism of $\A$ with $\zeta^d=1$
\end{itemize}
satisfying:
\begin{enumerate}
\item $[\A_{\leqslant i},\A_{\leqslant j}]\subset \A_{\leqslant i+j-d}$,
\item $\gr\A$ (a commutative algebra) is  finitely generated.
\item $\zeta$ acts on the graded component $(\gr\A)_i$ by $\epsilon^i$.
\end{enumerate}

Now let $\theta$ be an anti-involution of $\A$ preserving the filtration and commuting with $\zeta$. Following \cite[Section 6.1]{Wdim},
we define a Harish-Chandra $(\A,\theta)$-module as follows.

\begin{defi}\label{defi:HC} By a Harish-Chandra $(\A,\theta)$-module we mean a finitely generated
$\A$-module $M$ with an automorphism $\zeta$ compatible with the eponymous automorphism
of $\A$ that admits a filtration (called {\it good} for $\theta$),
$M=\bigcup_{j\geqslant 0}M_{\leqslant j}$
with the following properties:
\begin{itemize}
\item[(i)] $\gr M$ is finitely generated over $\gr\A$,
\item[(ii)] For all $a\in \A_{\leqslant i}$ and $m\in M_{\leqslant j}$,
we have $(\theta(a)-a)m\subset M_{\leqslant i+j-d}$.
\item[(iii)] $\zeta$ acts on $(\gr M)_i$ by $\epsilon^i$.
\end{itemize}
\end{defi}

HC $(\A,\theta)$-modules form a full subcategory in the category of $(\A,\Z/d\Z)$-modules, it will be denoted by
$\HC(\A,\theta,\zeta)$. 

Here is a classical example.
\begin{Ex}\label{Ex:HC_classical}
Let $\g$ be a semisimple Lie algebra and let
$\sigma$ be an involution of $\g$. We write $\U$ for the universal
enveloping algebra of $\g$ equipped with the PBW filtration (and $d=1$). It comes with
the unique anti-involution $\theta$ extending $-\sigma$.
Let $\A=\U$ and $\kf=\g^\sigma$. Then, as was pointed out in
\cite[Section 6.1]{Wdim}, a Harish-Chandra $(\U,\theta)$-module is the same
things as a Harish-Chandra $(\g,\kf)$-module, a classical object of
study in Lie representation theory.
\end{Ex}

\begin{Rem}\label{Rem:HC_bimodules}
Let $\A,\zeta$ be as above. Form the algebra $\A\otimes \A^{opp}$ and consider its anti-involution
$\theta$ given by $\theta(x\otimes y)=y\otimes x$. The HC $(\A\otimes \A^{opp},\theta)$-modules will
be called HC $\A$-bimodules.
\end{Rem}


\subsection{Equivariance}
Now suppose that $K$ is a (not necessarily connected) reductive algebraic group that acts on
$\A$ rationally by algebra automorphisms and  with quantum comoment map $\Phi: \mathfrak{k}\rightarrow \A_{\leqslant d}$.
Suppose that this Hamiltonian action is compatible with
the anti-involution $\theta$ in the following way:
\begin{equation}\label{eq:action_compatible}
\theta(ka)=k^{-1}\theta(a),\theta \Phi(x)=-\Phi(x),\forall k\in K, a\in \A, x\in \mathfrak{k}.
\end{equation}
Further, suppose that the $K$-action commutes with $\zeta$ and $\zeta.\Phi(\xi)=\Phi(\xi)$
for all $\xi\in \kf$.

Pick $\kappa\in (\kf^*)^K$.
We can speak about $(K,\kappa)$-equivariant $\A$-modules -- those are $\A$-modules $M$
equipped with a rational $K$-action making the structure map $\A\otimes M\rightarrow M$
equivariant and whose differential is $x\mapsto \Phi(x)-\langle\kappa,x\rangle$.


We can also talk about $(K,\kappa)$-equivariant HC $(\A,\theta)$-modules.

\begin{defi}\label{eq:HC_equiv}
By a  $(K,\kappa)$-equivariant (sometimes to be called twisted equivariant) HC $(\A,\theta)$ we mean a $(K,\kappa)$-equivariant $\A$-module that admits a $K$-stable good filtration.
\end{defi}

The corresponding category will be
denoted by $\HC(\A,\theta,\zeta)^{K,\kappa}$ or simply
$\HC(\A,\theta,\zeta)^K$ if $\kappa=0$ (the morphisms in this category are $K$-equivariant $\A$-linear maps; note that the $\A$-linearity implies the $K$-equivariance if $K$ is connected). We note that if $\kappa_0$
is a character of $K$, then there is a natural category
equivalence between $\HC(\A,\theta,\zeta)^{K,\kappa}$
and $\HC(\A,\theta,\zeta)^{K,\kappa+\kappa_0}$ given by tensoring
with the one-dimensional $K$-representation, where $K$
acts via $\kappa_0$.

For example, when $\A=\U$, $\theta$ is as in Example \ref{Ex:HC_classical},
and $K=(G^\sigma)^\circ$, then a $K$-equivariant HC $(\A,\theta)$-module
is the same thing as a classical HC $(\g,K)$-module.

Now we return to the general situation in the beginning of this section.

\begin{Lem}\label{Lem:twisted_equivar}
Suppose that the derived subgroup $(K,K)$ is simply connected. Then
every simple object in $\HC(\A,\theta,\zeta)$ admits a $K$-action making it
an object of $\HC(\A,\theta,\zeta)^{K,\kappa}$ for some $\kappa$.
\end{Lem}
\begin{proof}
Let $M\in \operatorname{Irr}(\HC(\A,\theta,\zeta))$. By the definition of a HC module,
every vector in $M$ lies in a finite dimensional $\A_{\leqslant d}$-stable
subspace.  The Lie algebra $\kf$
acts on $M$ locally finitely because $\Phi(\kf)\subset \A_{\leqslant d}$. We can decompose $M$ into the
sum $\bigoplus M_{\mu}$, where $\mu$ runs over the set of characters of $\kf$
and $M_{\mu}$ stands for the generalized $\mu$-eigenspace for the center $\mathfrak{z}(\kf)$.
Since $M$ is irreducible it is generated by a highest vector for $\kf$ (and an eigenvector for $\zeta$), say $v$. In particular,
the action of $\mathfrak{z}(\kf)$  on $M$ is diagonalizable. We can shift the
action of $\kf$ by a character, say $\kappa$, of $\kf$ and assume that the weight of
$v$ is a weight for $K$. This gives rise to a $K$-action on $M$.
This action turns $M$ into  a twisted equivariant object in
$\HC(\A,\theta,\zeta)$.
\end{proof}

\subsection{Weakly HC modules}\label{SS_weakly_HC}
Let $\Dcal_\hbar$ be an associative $\C[\hbar]$-algebra with unit. We assume that
$\Dcal_\hbar$ is flat over $\C[\hbar]$ and that $\Dcal_0 = \Dcal_\hbar/ \hbar \Dcal_\hbar$ is a finitely
generated commutative $\C$-algebra.
Let $\theta$ be a $\C[\hbar]$-linear anti-involution on $\Dcal_\hbar$.


Suppose $\A,\zeta,\theta$ are as in Section \ref{SS_HC_mod}.
Define the (modified) Rees algebra for $(\A,\zeta)$ as follows.
For $j\in \Z/d\Z$ let $\A_j$ denote the eigenspace for $\zeta$
with eigenvalue $\epsilon^j$. Then set $R_\hbar(\A):=\bigoplus_{i=0}^\infty
(\A_{\leqslant i}\cap \A_{i+d\Z})\hbar^{i/d}$. This is a graded $\C[\hbar]$-algebra
with $\hbar$ in degree $d$. For $\Dcal_\hbar$ we can take $R_\hbar(\A)$ or its
$\hbar$-adic completion. In both cases the anti-involution $\theta$
of $\A$ gives rise to a $\C[\hbar]$-linear anti-involution of $\Dcal_\hbar$
also denoted by $\theta$.

%For example, take $\Dcal=\U$. For  $\theta$ we take the  anti-involution of $\U$ that on $\g$
%is given by $\theta=-\sigma$.

We write $\Dcal_\hbar^{-\theta}$ for the $-1$ eigenspace of $\theta$.
	We have $\Dcal_\hbar=\Dcal_\hbar^{-\theta}\oplus \Dcal_\hbar^\theta$.
	Note that the following inclusions hold:
	\begin{equation}\label{eq:invol_decomp_properties}
	\begin{split}
	&[\Dcal_\hbar^{-\theta},\Dcal_\hbar^{-\theta}]\subset \Dcal_\hbar^{-\theta}, \\
	&a\in \Dcal_\hbar^{-\theta}, b\in \Dcal_\hbar^\theta\Rightarrow ab+ba\in \Dcal_\hbar^{-\theta}.
	\end{split}
	\end{equation}


Set $\J_\hbar:=\hbar \mathcal{D}_\hbar+ \Dcal_\hbar\Dcal_\hbar^{-\theta}$. This is a two-sided ideal of $\Dcal_\hbar$ and is a Lie algebra under taking the bracket defined by
$[a,b]_\hbar:=\frac{1}{\hbar}(ab-ba)$.

\begin{defi}\label{defi:wHC}
	A {\it weakly HC $(\Dcal_\hbar, \theta)$-module}, is a finitely generated
	$\Dcal_\hbar$-module $M_\hbar$ equipped with
	$\C[\hbar]$-bilinear Lie algebra representation
	$\J_\hbar \times M_\hbar\rightarrow M_\hbar, (a,m)\mapsto a.m$, such that
	\begin{enumerate}[(i)]
		\item $am=\hbar(a.m)$, $\forall\, a \in \J_\hbar$.
		\item $(ba).m = b(a.m)$, $\forall\, a \in \J_\hbar, b \in \Dcal_\hbar$.
		\item $a.(bm)=[a,b]_\hbar m+b(a.m), \forall a\in \J_\hbar, b\in \Dcal_\hbar$.
	\end{enumerate}
A weakly HC $(\Dcal_\hbar,\theta)$-module is called HC if it is flat over $\C[\hbar]$.
\end{defi}

\begin{Rem}\label{Rem:flat_wHC}
Note that for a $\C[\hbar]$-flat $\Dcal_\hbar$-module being HC is not a structure but a condition:
$bm\in \hbar M_\hbar$ for all $b\in \Dcal_{\hbar}^{-\theta},m\in M_\hbar$. Then we can define
$b.m$ as $\hbar^{-1}(bm)$.
\end{Rem}

The category of weakly HC $(\Dcal_\hbar, \theta)$-modules will be denoted by $\wHC(\Dcal_\hbar,\theta)$.

Now suppose that we have a $\C^\times$-action on $\Dcal_\hbar$ by $\C$-algebra automorphisms
such that $t.\hbar=t^d\hbar$ for $t \in \C^\times$. Further suppose that either
\begin{itemize}\item the action is rational or
\item $\Dcal_\hbar$
is complete and separated in the $\hbar$-adic topology and the action is pro-rational.
\end{itemize}
The Rees algebra $R_\hbar(\A)$ satisfies the first condition, while its $\hbar$-adic completion
satisfies the second condition.

\begin{defi}\label{defi:wHC_graded}
A weakly HC $(\Dcal_\hbar,\theta)$-module $M_\hbar$ is {\it graded} if it is equipped with
a (rational/pro-rational) $\C^\times$-action (rational, if the action on $\Dcal_\hbar$ is rational
and pro-rational, if the action on $\Dcal_\hbar$ is pro-rational) such that
the multiplication map $\Dcal_\hbar\otimes M_\hbar\rightarrow M_\hbar$ is
$\C^\times$-equivariant and the Lie algebra action $\J_\hbar
\otimes M_\hbar\rightarrow M_\hbar$ is of degree $-d$.
\end{defi}

The abelian category of graded weakly HC modules will be denoted by $\wHC^{gr}(\Dcal_\hbar,\theta)$.

Now suppose $\Dcal_\hbar=R_\hbar(\A)$ for a filtered algebra $\A$ as in the beginning of this
section.
A connection between $\wHC^{gr}(\Dcal_\hbar,\theta)$ and $\HC(\A,\theta,\zeta)$ is as follows.
Note that we have the functor $\wHC^{gr}(\Dcal_\hbar,\theta)\rightarrow \HC(\A,\theta,\zeta)$
given by $M_\hbar\mapsto M_\hbar/(\hbar-1)M_\hbar$. It is straightforward to check that $M_\hbar/(\hbar-1) M_\hbar$
is indeed a HC $(\A,\theta)$-module: a good filtration on $M_\hbar/(\hbar-1) M_\hbar$
comes from the grading on $M_\hbar$ and the automorphism $\zeta$ is induced by $\epsilon\in \C^\times$.

\begin{Lem}\label{Lem:essential_surjectivity}
The functor $\wHC^{gr}(\Dcal_\hbar,\theta)\rightarrow \HC(\A,\theta,\zeta)$, $M_\hbar\mapsto M_\hbar/(\hbar-1)M_\hbar$, is essentially surjective
and it kills precisely the $\hbar$-torsion modules.
\end{Lem}
\begin{proof}
The description of the kernel is clear.
Let $M \in \HC(\A,\theta,\zeta)$. Pick a good filtration on $M$ as in Definition
\ref{defi:HC}. We can form the (modified) Rees module $R_\hbar(M)$
similarly to what was done for $\A$. By Remark \ref{Rem:flat_wHC}, $R_\hbar(M)$ is an object of
$\wHC^{gr}(\Dcal_\hbar,\theta)$. Of course, $R_\hbar(M)/(\hbar-1)\cong M$.
\end{proof}

We also have a $(K,\kappa)$-equivariant version of weakly HC modules. Let $K$ be an algebraic group which acts rationally or pro-rationally on $\Dcal_\hbar$ by $\C[\hbar]$-linear algebra automorphisms, so that any element $\xi\in \kf$ gives rise to a $\C[\hbar]$-linear endomorphism $\xi_{\Dcal_\hbar}$ of $\Dcal_\hbar$.
Let $\Phi: \kf\rightarrow \Dcal_\hbar$ be a quantum comoment map
meaning that this map is $K$-equivariant and $\hbar^{-d}[\Phi(\xi),\cdot]=\xi_{\Dcal_\hbar}$.
Assume additionally that $\theta$ is $K$-equivariant and $\operatorname{Im} \Phi \subset \J\hbar$.

\begin{defi}\label{defi_wHC_equiv}
Let $M_\hbar\in \wHC(\Dcal_\hbar,\theta)$.
Let $K$ act rationally or pro-rationally on $M_\hbar$ by $\C[\hbar]$-linear automorphisms (so that every element
$\xi\in \kf$ gives rise to a $\C[\hbar]$-linear endomorphism $\xi_{M_\hbar}$ of $M_\hbar$).  We say that
$M_\hbar$ is $(K,\kappa)$-equivariant if
\begin{itemize}
  	\item the multiplication map $\Dcal_\hbar\times M_\hbar\rightarrow M_\hbar$
  and the Lie algebra action map $\J\hbar \times M_\hbar\rightarrow M_\hbar$
  are $K$-equivariant,
  \item $\Phi(\xi).m-\langle\kappa,\xi\rangle=\xi_{M_\hbar}m, \forall \xi\in \kf, \, m \in M_\hbar$.
\end{itemize}
\end{defi}

We denote the corresponding category by $\wHC(\Dcal_\hbar,\theta)^{K,\kappa}$.
Similarly, we can talk about $(K,\kappa)$-equivariant objects in $\wHC^{gr}(\Dcal_\hbar,\theta)$.
Here we assume that the $K$-action commutes with the $\C^\times$-action and that
$\operatorname{Im}\Phi$ lies in the subspace of degree $d$ homogeneous elements with respect to the $\C^\times$-action.

\subsection{Weakly HC modules over formal Weyl algebras}
Here we consider an important special case of weakly HC modules.
Let $V$ be a symplectic vector space. Consider the homogenized Weyl algebra
$\Weyl_\hbar(V)=T(V)[\hbar]/(u\otimes v-v\otimes u-\hbar\omega(u,v))$.

Let $\mathfrak{m}\subset \Weyl_\hbar(V)$ denote the preimage of the maximal ideal of $0$ in $S(V)$
under the epimorphism $\Weyl_\hbar(V)\twoheadrightarrow S(V)$. Let $\Dcal_\hbar$ denote the
$\mathfrak{m}$-adic completion of $\Weyl_\hbar(V)$.

Fix an anti-symplectic involution $\theta$ of $V$. It defines the decomposition
$V=L\oplus L'$, where $L,L'$ are Lagrangian. Namely, $L=V^{-\theta}$ and $L'=V^{\theta}$.
We note that $\C[L][\hbar]$ is a HC $(\Weyl_\hbar, \theta)$-module. Similarly,
$\C[[L,\hbar]]$ is a HC $\Dcal_\hbar$-module.

For a weakly HC $(\Dcal_\hbar,\theta)$-module $M$ we write $M^L$
for the annihilator of $L$ in $M$ (for the Lie algebra action).

\begin{Lem}\label{Lem:wHC_formal_Weyl}
The functors $\bullet^L$ and $\C[[L,\hbar]]\otimes_{\C[[\hbar]]}\bullet$
are mutually quasi-inverse equivalences between $\wHC(\Dcal_\hbar,\theta)$
and $\C[[\hbar]]\operatorname{-mod}$.
\end{Lem}
\begin{proof}
It is easy to see that $(\C[[L,\hbar]]\otimes_{\C[[\hbar]]}\bullet)^L$
is the identity endo-functor of $\C[[\hbar]]\operatorname{-mod}$.

To prove that the other composition is isomorphic to the identity one argues as follows.
Note that $M$ is complete in the $\mathfrak{m}$-adic topology.
Let $p_1,\ldots,p_n$ be a basis in $L'$ and $q_1,\ldots,q_n$ be the
dual basis in $L$. For $m\in M$, we have $q_i.(p_jm)=m+p_j(q_i.m)$.
Using this identity, we finish the proof of this lemma similarly to the analogous claim in
the proof of \cite[Proposition 3.3.1]{HC}.
\end{proof}




%Current section 2.1, 2.2, 2.5.
\section{Harish-Chandra modules over quantizations of nilpotent orbits}
\subsection{Quantizations of nilpotent orbits and their HC modules}\label{SS_quant_nilp}
Let $\Orb\subset \g^*$ be a nilpotent orbit. This is a symplectic algebraic variety.
The group $G$ naturally acts on $\Orb$ and the action is Hamiltonian.

The algebra
$\C[\Orb]$ of regular functions on $\Orb$ is a graded Poisson algebra
that comes with a Hamiltonian action of $G$, the corresponding moment
map $\operatorname{Spec}(\C[\Orb])\rightarrow \g^*$ is induced by the inclusion
$\Orb\hookrightarrow \g^*$. The bracket on
$\C[\Orb]$ has degree $-1$.

\begin{defi}\label{defi:quantization}
Let $A$ be a $\Z_{\geqslant 0}$-graded Poisson algebra with bracket of degree $-d$.
By a {\it filtered quantization} of $A$ we mean a pair $(\Dcal,\iota)$ of
\begin{itemize}
\item
a filtered
associative algebra $\Dcal$ such that $\deg [\cdot,\cdot]\leqslant -d$
\item and a graded Poisson algebra isomorphism $\iota:\gr \Dcal\xrightarrow{\sim} A$.
\end{itemize}
\end{defi}

In what follows in the present section we will place an additional restriction on
$\Orb$. Set $\partial \Orb:=\overline{\Orb}\setminus \Orb$.
We will require that
\begin{equation}\label{eq:codim_condition}
\operatorname{codim}_{\overline{\Orb}}\partial\Orb\geqslant 4.
\end{equation}

The following classical lemma describes the second cohomology $H^2(\Orb,\C)$. Let $e\in \Orb$
and $Q$ be the reductive part of the centralizer of $e$ in $G$.

\begin{Lem}\label{Lem:G_orbit_cohomology}
We have an isomorphism $(\q^*)^Q\xrightarrow{\sim} H^2(\Orb,\C)$ that sends a character of $Q$
to the 1st Chern class of the corresponding line bundle on $\Orb$.
\end{Lem}

To fix the normalization, to a character $\chi:Q\rightarrow \C^\times$ (equivalently,
a character of $Z_G(e)$) we assign the $G$-equivariant line bundle $\mathcal{L}$ on $\Orb$
such that $Z_G(e)$ acts on the fiber of $\mathcal{L}$ at $e\in \Orb$ by $\chi$.

The variety $\operatorname{Spec}(\C[\Orb])$ is singular symplectic, \cite{Beauville}.
In what follows we take a singular symplectic variety $X$ with $\operatorname{codim}_X X^{sing}\geqslant 4$. We also assume that $X$ is conical: meaning that $\C[X]$ admits a $\Z_{\geqslant 0}$-algebra grading
with $\C[X]_0=\{0\}$ and the degree of the Poisson bracket is $-d$ for $d\in \Z_{>0}$.

Filtered quantizations of the algebra $\C[X]$ were classified in \cite[Theorem 3.4]{orbit}
in the general setting. Under the assumption $\operatorname{codim}_X X^{sing}\geqslant 4$, they are parameterized by the points of $\param:=H^2(X^{reg},\C)$ via taking the so called {\it period} of the microlocalization of
the quantization to $X^{reg}$, \cite[Corollary 3.2(1)]{orbit}. We will discuss periods in more detail in
Section \ref{subsec:quan}. Moreover, we can take the universal
quantization $\A_{\param}$ of $\C[X]$, \cite[Corollary 3.2(2)]{orbit}.
This is a filtered $\C[\param]$-algebra with
$\deg \param^*=1$ such that for $\lambda\in \param$, the specialization of $\A_{\param}$
to $\lambda$ is the filtered quantization $\A_\lambda$ of $\C[\Orb]$ corresponding to $\lambda$.
Note that $\A_{\param}^{opp}$ is identified with $\A_{\param}$, where the identification is
$-1$ on $\param^*$. This can be deduced from \cite[Proposition 3.2]{BPW}.
The quantization $\A_\lambda$ of $\C[X]$ corresponding to $\lambda\in \param$
is obtained from $\A_{\param}$ by specializing the parameter to $\lambda$.

Fix a Poisson anti-involution $\theta$ of
$\C[X]$ that preserves the grading. Thanks to the universality property
of $\A_{\param}$ (\cite[Proposition 3.5]{orbit}), $\theta$ lifts to an anti-involution
of $\A_{\param}$ also denoted by $\theta$, compare to the beginning of
\cite[Section 3.7]{orbit}. The action of $\theta$ on $\param^*\subset \A_\param$
comes from the action on $\param=H^2(X^{reg},\C)$.

In particular, we can talk about the category $\HC(\A_{\param},\theta,\zeta)$
of HC $(\A_{\param},\theta)$-modules. For $\lambda\in \param$ we write $\HC(\A_\lambda,\theta)$
for the full subcategory of $\HC(\A_{\param},\theta,\zeta)$ consisting of all modules, where
the action of $\A_{\param}$ factors through $\A_\lambda$. If $\lambda$ is fixed by $\theta$,
then $\theta$ induces an anti-involution of $\A_\lambda$ (again denoted by $\theta$), then
$\HC(\A_\lambda,\theta,\zeta)$ is indeed the category of HC $(\A_\lambda,\theta)$-modules.


For any $M\in \HC(\A_\lambda,\theta,\zeta)$ we can consider its associated variety, to be denoted
by $\operatorname{V}(M)$. By definition, it is the support of the associated graded of
$M$ with respect to any good filtration.

Recall that $X$ has finitely many symplectic leaves, \cite{Kaledin_symplectic}.

\begin{Lem}\label{Lem:lag_antiinvolution}
The following are true:
\begin{enumerate}
\item For any symplectic leaf $\mathcal{L}\subset X$, the intersection $X^\theta\cap \mathcal{L}$
is Lagrangian in $\mathcal{L}$. In particular, $X^\theta$ is an isotropic subvariety in the sense of
\cite{B_ineq}.
\item For any $M\in \HC(\A_\lambda,\theta,\zeta)$, we have $\operatorname{V}(M)\subset X^\theta$. In particular,  $M$ is holonomic in the sense of \cite{B_ineq}.
\item For $M\in \HC(\A_\lambda,\theta,\zeta)$, the following conditions are equivalent:
\begin{itemize}
\item[(a)] $\dim \operatorname{V}(M)<\frac{1}{2}\dim X$,
\item[(b)] $\operatorname{V}(M)\cap X^{reg}=\varnothing$,
\item[(c)] The annihilator of $M$ in $\A_\lambda$ is nonzero.
\end{itemize}
\end{enumerate}
\end{Lem}
\begin{proof}
It is a classical and easy fact that the fixed point locus of an anti-symplectic involution of a
smooth symplectic variety is a Lagrangian subvariety. If $X^\theta\cap \mathcal{L}\neq \varnothing$,
then $\theta$ preserves $\mathcal{L}$ and induces an anti-symplectic involution of $\mathcal{L}$.
So $X^\theta\cap \mathcal{L}$ is Lagrangian, in particular, isotropic. So $X^\theta$ is isotropic
in the sense of \cite{B_ineq}. This proves (1).

The claim that $\operatorname{V}(M)\subset X^\theta$ is an easy consequence of the definition of
a good filtration on a module. Together with (1), this proves (2).

Now we prove (3). (b)$\Rightarrow$(a) follows from (1). (c)$\Rightarrow$(b) follows from
the observation that $V(M)$ is contained in $V(\A_\lambda/\operatorname{Ann}(M))$, while
the latter is the union of symplectic leaves. Finally, (a)$\Rightarrow$(c) follows
from \cite[Theorem 1.2]{B_ineq} (the finiteness of algebraic fundamental groups needed in that
theorem is a result of Namikawa \cite{Namikawa})).
\end{proof}

\begin{defi}\label{defi:HC_full_support}
We consider the category of HC $\A_\lambda$-modules
{\it with full support},  $\overline{\HC}(\A_\lambda,\theta,\zeta)$, defined as the quotient of $\HC(\A_\lambda,\theta,\zeta)$ by the full
subcategory of all modules $M$ such that the equivalent conditions of (3) of Lemma
\ref{Lem:lag_antiinvolution} hold.
\end{defi}

We now get back to the situation when $X=\operatorname{Spec}(\C[\Orb])$.
Note that the action of $G$ on $\A_{\param}$ is Hamiltonian, this is similar to Step 1 of the proof
of \cite[Proposition 4.5]{orbit} as well as the discussion in the beginning of
\cite[Section 5.3]{orbit}. Let $\Phi:\g\rightarrow \A_{\param}$
denote the quantum comoment map. We also write $\Phi$ for the induced homomorphism $\U\rightarrow
\A_{\param}$.

Suppose that the Poisson anti-involution $\theta$ of $\C[\Orb]$ is induced by $-\sigma$, where $\sigma$ is an involution of $\g$ and $\kf:=\g^\sigma$. Let $\kf^\perp$
denote the annihilator of $\kf$ in $\g^*$. Then $X^\theta$ is the preimage of $\kf^\perp$ in $X$.
%It is a classical fact that, for any $G$-orbit $\Orb'\subset \g^*$, the intersection $\Orb'\cap %\kf^\perp$
%is a Lagrangian subvariety in $\Orb'$ and is a union of finitely many $K$-orbits.
In particular, thanks to Lemma \ref{Lem:lag_antiinvolution},
for any $K$-orbit $\Orb_K\subset \Orb\cap \kf^\perp$, we have

\begin{equation}\label{eq:codim_condition_K}
\operatorname{codim}_{\overline{\Orb}_K}\partial\Orb_K\geqslant 2,
\end{equation}
where we write $\partial\Orb_K$ for $\overline{\Orb}_K\setminus \Orb_K$.

\subsection{HC $\A_{\param}$- and $\U$-modules}
Let $\g,\kf,K,\sigma,\theta$ be as in Example \ref{Ex:HC_classical}.
Suppose that $\Orb$ satisfies (\ref{eq:codim_condition}) and intersects $\kf^\perp$, the annihilator
of $\kf$ in $\g^*$. It follows that $\Orb$ is stable under $\sigma$ and hence under $\theta$ as well.
So $\theta$ induces an anti-Poisson involution of $\C[\Orb]$ and also a linear transformation of
$\param=H^2(\Orb,\C)$.
Denote the resulting involution of $\A_{\param}$ also by $\theta$. So we can talk about
HC $(\A_{\param},\theta,\zeta)$-modules and, for $\lambda\in \param$, about HC $(\A_\lambda,\theta)$-modules.

For $\lambda\in \param$, let $\J_\lambda$ denote the kernel of $\U\rightarrow \A_\lambda$. The goal of this section is to prove the following
proposition.

\begin{Prop}\label{Prop:HC_bijection}
There is a natural bijection between
\begin{itemize}
\item irreducible $(K,\kappa)$-equivariant HC $(\A_\lambda,\theta)$-modules of GK dimension
$\frac{1}{2}\dim \Orb$,
\item and irreducible $(K,\kappa)$-equivariant $\U/\J_\lambda$-modules
of GK dimension $\frac{1}{2}\dim \Orb$.
\end{itemize}
\end{Prop}

The proposition follows from the next two lemmas.

\begin{Lem}\label{Lem:HC_bijection1}
There is a natural bijection between
\begin{itemize}
\item irreducible  $(K,\kappa)$-equivariant $\A_\lambda$-modules of GK dimension
$\frac{1}{2}\dim \Orb$,
\item and irreducible  $(K,\kappa)$-equivariant $\U/\J_\lambda$-modules
of GK dimension $\frac{1}{2}\dim \Orb$.
\end{itemize}
\end{Lem}

\begin{Lem}\label{Lem:HC_bijection2}
Every irreducible $(K,\kappa)$-equivariant $\A_\lambda$-module of GK dimension $\frac{1}{2}\dim \Orb$
admits a good filtration (with respect to $\theta$).
\end{Lem}

\begin{proof}[Proof of Lemma \ref{Lem:HC_bijection1}]
Take a simple object  $M\in\A_\lambda\operatorname{-mod}^{K,\kappa}$.
It is a finitely generated $(K,\kappa)$-equivariant $\U/\J_\lambda$-module
because $\A_\lambda$ is a finitely generated $\U$-module.
Note that $\J_\lambda$
intersects the center of $\U$ at a maximal ideal, see the discussion in
the beginning of \cite[Section 5.3]{orbit}. Let $\mathcal{I}$ denote
the (two-sided) ideal in $\U$ generated by this intersection. So $\U/\mathcal{I}$
is a filtered quantization of the nilpotent cone  $\mathcal{N}\subset \g^*$.

Every $(K,\kappa)$-equivariant $\U/\mathcal{I}$-module is supported on
$\kf^\perp\cap \mathcal{N}$. The intersection with every orbit is classically
known to be a Lagrangian subvariety, see also (1) of
Lemma \ref{Lem:lag_antiinvolution}. So $\kf^\perp\cap \mathcal{N}$ is an
isotropic subvariety  in $\mathcal{N}$ and $M$ is a holonomic $\U/\mathcal{I}$-module  in the sense of \cite[Section 1.4]{B_ineq}.
In particular, it  has finite length and has Gelfand-Kirillov dimension
not exceeding $\frac{1}{2}\dim \Orb$, by \cite[Theorems 1.2,1.3]{B_ineq}.
Let $M^0$ be an irreducible $\U$-submodule of $M$. Then $\A_\lambda M^0$
is an irreducible $\A_\lambda$-submodule hence $\A_\lambda M^0=M$.
%Recall that $\A_\lambda$
%is finite over $\U/\J_\lambda$. It follows that
%$$\operatorname{GK-dim}M^0=\operatorname{GK-dim}M,$$
%where we write $\operatorname{GK-dim}$ for the Gelfand-Kirillov dimension.
Also note that the multiplicity of $\A_\lambda$ on $\Orb$ is $1$, while the multiplicity of
$\U/\J_\lambda$ is at least $1$.  Hence $\operatorname{GK-dim}(\A_\lambda/ (\U/\J_\lambda))<\dim \Orb$.
In particular, $\A_\lambda/(\U/\J_\lambda)$ is annihilated by an ideal $\J$
that strictly contains $\J_\lambda$. Since $M$ and hence $M/M^0$ are holonomic $\U/\mathcal{I}$-modules, it follows that
$$\operatorname{GK-dim}(M/M^0)<\frac{1}{2}\dim \Orb,$$
where we write $\operatorname{GK-dim}$ for the Gelfand-Kirillov dimension.

So if $M$ has GK dimension $\frac{1}{2}\dim \Orb$, then $M^0$ is the unique irreducible $\U/\J_\lambda$-submodule of $M$. This gives rise to a
map between the sets of irreducibles with GK dimension $\frac{1}{2}\dim \Orb$:
 $$\operatorname{Irr}_{\frac{1}{2}\dim \Orb}(\A_\lambda\operatorname{-mod}^{K,\kappa})\rightarrow \operatorname{Irr}_{\frac{1}{2}\dim \Orb}(\U/\J_\lambda\operatorname{-mod}^{K,\kappa}),
M\mapsto M^0.$$

Let us prove that this map is a bijection. Take $N\in \operatorname{Irr}_{1/2\dim \Orb}(\U/\J_\lambda\operatorname{-mod}^{K,\kappa})$.
Consider the $\A_\lambda$-module $\A_\lambda\otimes_\U N$.
It is $(K,\kappa)$-equivariant because the action of $K$ on $\A_\lambda$ is Hamiltonian.
We have the map
$N\rightarrow \A_\lambda\otimes_\U N$ induced by $\U/\J_\lambda\hookrightarrow \A_\lambda$.
Similarly to the second paragraph in the proof, the cokernel of this map is annihilated by an ideal strictly containing $\J_\lambda$
and is holonomic. It follows that it is nonzero, hence an embedding. As noted above in the proof, this means that the cokernel has GK
dimension less than $\frac{1}{2}\dim \Orb$.
% So if the map is zero,
%we have $\Hom_{\U/\J_\lambda}(N,\A_\lambda\otimes_\U N)=0$. On the other hand, this Hom
%space is identified with \Hom_{\U/\J_\lambda}(N,\A_\lambda\otimes_\U N)


Let $N'$ be the quotient of $\A_\lambda\otimes_\U N$
by the maximal $\A_\lambda$-submodule with GK dimension $<\frac{1}{2}\dim \Orb$.
The embedding $N\hookrightarrow \A_\lambda\otimes_\U N$ gives rise
to a $\U/\J_\lambda$-module embedding $N\rightarrow N'$.
Therefore $N'=\A_\lambda N$. If an $\A_\lambda$-submodule of $N'$ does not intersect $N$ then it has
GK-dimension $<\frac{1}{2}\dim \Orb$. But $N'$ has no such submodules. So
every $\A_\lambda$-submodule of $N'$ must contain $N$. Therefore $N'$ is simple.

It is easy to see that the maps $M\mapsto M^0$ and $N\mapsto N'$ are mutually inverse.
\end{proof}

\begin{proof}[Proof of Lemma \ref{Lem:HC_bijection2}]
Let $M$ be an irreducible $(K,\kappa)$-equivariant $\A_\lambda$-module.
Let $M=\bigcup_{j}M_{\leqslant j}$ be a filtration such that $\gr M$
is a finitely generated $\C[\Orb]$-module. We claim that this filtration
is good (for $\theta$) if and only if
\begin{itemize}
\item[(*)] $a$ acts by zero on $\gr M$ for every $a\in \C[\Orb]^{-\theta}$,
the $-1$ eigenspace for the action of $\theta$ on $\C[\Orb]$.
\end{itemize}
It is clear that every good filtration for $\theta$ has the desired property.
On the other hand, $\param^*\subset \gr\A_{\param}$ acts by
$0$ on $\gr M$ and the image of $(\gr \A_{\param})^{-\theta}$
in $\C[\Orb]=(\gr \A_{\param})/(\param^*)$ coincides with
$\C[\Orb]^{-\theta}$. So (*) implies the filtration is good
for $\theta$ (recall that in this case $d=1$).

Let $\A_{\lambda,\hbar}$, $\U_\hbar$ denote the Rees algebras of $\A_{\lambda}$, $\U$
and let $\A_{\lambda,\hbar}^{loc}$, $\U_\hbar^{loc}$ denote their microlocalizations
to $\g^*$.  Note that the natural homomorphism $\U_\hbar^{loc}\rightarrow \A_{\lambda,\hbar}^{loc}$
is an epimorphism over $\Orb$. Pick a $K$-stable filtration on $M$ such that
$\gr M$ is finitely generated over $\C[\kf^\perp]$. Note that $\gr M$ is set
theoretically supported on  $\kf^\perp\cap \overline{\Orb}$ and the
restriction to $\Orb$ is, in fact, scheme theoretically supported on
the Lagrangian subvariety $\Orb\cap \kf^\perp$.
The ideals $\langle \C[\Orb]^{-\theta} \rangle$ generated by $\C[\Orb]^{-\theta}$ and the ideal $\C[\Orb]\kf$ in $\C[\Orb]$ have the same localizations
to $\Orb$, the localizations to $\Orb$ of the quotients by these ideals are reduced. It follows that $\C[\Orb]^{-\theta}$ acts by
$0$ on $(\gr M)|_{\Orb}$. Since $\operatorname{codim}_{\overline{\Orb}}\partial \Orb\geqslant 4$,
we have
\begin{equation}\label{eq:codim_K}
\operatorname{codim}_{\overline{\Orb}\cap \kf^\perp} (\partial\Orb\cap \kf^\perp)\geqslant 2.
\end{equation}
Note that $(\gr M)|_{\Orb}$ a $K$-equivariant coherent sheaf on $\Orb\cap \kf^{\perp}$ hence a vector bundle. Therefore, thanks to (\ref{eq:codim_K})
$\Gamma((\gr M)|_{\Orb})$ is finitely generated over $\C[\overline{\Orb}\cap \kf^\perp]$.
It follows that $\Gamma((\gr M)|_{\Orb})$ is finitely generated over $\C[\Orb]$.
Let $M_\hbar^{loc}$ denote the microlocalization of $M_\hbar$ to $\g^*$. It makes to sense 
to restrict $M_\hbar^{loc}$ to $\Orb$. 
The global sections $\Gamma(M_\hbar^{loc}|_{\Orb})$
is finitely generated over the $\hbar$-adic completion of $\A_{\lambda,\hbar}$.
Let $\tilde{M}_\hbar:=\Gamma(M_\hbar^{loc}|_{\Orb})^{fin}$ denote the subspace of $\C^\times$-finite
elements and let $\tilde{M}$ be the specialization of this module to $\hbar=1$.
The latter is a filtered module such that $\gr \tilde{M}$ admits a natural embedding into
$\Gamma((\gr M)|_{\Orb})$. Therefore, $\C[\Orb]^{-\theta}$ acts by $0$
on $\gr\tilde{M}$.

We have a natural filtered homomorphism $M\rightarrow \tilde{M}$. The composition of the associated graded of this homomorphism with the embedding
$\gr\tilde{M}\hookrightarrow \Gamma((\gr M)|_{\Orb})$ is the natural
map $\gr M\rightarrow \Gamma((\gr M)|_{\Orb})$. It follows that $M\rightarrow \tilde{M}$
is an inclusion such that $\tilde{M}/M$ is supported on $\partial\Orb$. In particular,
$M$ inherits a good filtration from $\tilde{M}$. We denote the associated graded module of $M$ with respect to this new filtration by $\gr' M$. By the construction,
$\gr' M\hookrightarrow \gr\tilde{M}$. In particular,
$\C[\Orb]^{-\theta}$ acts by $0$ on $\gr' M$. We have constructed the required
filtration on $M$ finishing the proof of the lemma.
\end{proof}




%Generalities on nilpotent orbits as needed (some of them currently are in Section 6), Sections 2.3, 2.4.
\section{Quasi-classical limits and quantizations}\label{S_qclass_quant}
%Sections 4.2, 4.3, 4.4, 4.6 + whatever is left from relating various versions of weakly HC for $\hbar=0$.


\subsection{Picard Lie algebroids and  twisted local systems }

\subsubsection{Picard algebroids and TDOs} \label{subsec:pic_defn}

Let $X$ be a smooth variety/scheme and $\T_X$ be its tangent sheaf. We first review the notion of Picard algebroid, its basic properties and relations with twisted differential operators and twisted $\mathscr{D}$-modules. We will see in Section \ref{subsec:RH} that modules over Picard algebroids are equivalent to twisted local systems. A standard reference is \cite[Section 2]{BB}. For the defintion and basic properties of general Lie algebroids, see \cite[Section 1.2]{BB}.

\begin{defi} \label{defn:pic}
	A {\it Picard Lie algebroid}, or simply a {\it Picard algebroid} on $X$ is a Lie algebroid $\TT$ with the ($\OO_X$-linear) anchor map $\eta: \TT \to\T_X$ and a central section $1_{\TT}$ of $\ker \eta$, such that $\ker \eta = \OO_X \cdot 1_{\TT}$, i.e., the sequence of sheaves of $\OO_X$-modules,
	\[  0 \to \OO_X \xrightarrow{\iota_\TT} \TT \xrightarrow{\eta} \T_X \to 0\]
	is exact.
\end{defi}

A morphism of Picard algebroids is a morphism of Lie algebroids that preserves the $1_{\TT}$'s. It follows that any morphism of Picard algebroids is in fact an isomorphism, i.e., the category of all Picard algebroids over $X$ forms a groupoid. We write $\pic(X)$ for the set of isomorphsim classes of Picard algebroids over $X$. The Baer sum construction defines  a natural structure of complex vector space on $\pic(X)$. The zero vector  is given by the trivial Picard algebroid $\TT_{\operatorname{triv}} = \OO_X \oplus \T_X$.

For any line bundle $L$ over $X$, we write $\TT(L)$ for the associated Picard algebroid, i.e., the sheaf  over $X$ of local algebraic differential operators  with coefficients in $L$ of order less or equal than $1$. Let $\operatorname{Pic}(X)$ be the Picard group of $X$, then $L \mapsto \TT(L)$ defines a homomorphism of $\TT(\cdot): \operatorname{Pic}(X) \to \pic(X)$ of abelian groups which transfrom tensor product to the Baer sum. We abuse the notation and write $L^r := r \TT(L)$ for $r \in \C$, so that when $r$ is an integer, $L^r$ stand for both the $r$-tensor power of $L$ as a line bundle and its associated Picard algebroid.

Now we recall the classification of Picard algebroids. Let $\TT$ be a Picard algebroid over $X$. Let $H^\bullet_F(X)$ denote the hypercohomology of the first piece $F^1\Omega_X^\bullet$ of the de Rham complex $\Omega_X^\bullet$ with respect to the {\it filtration bête}. In other words,
\[ F^1\Omega_X^\bullet := (0 \to \Omega_X^1 \to \Omega_X^2 \to \cdots) \]
is the algebraic de Rham complex of $X$ with $\Omega_X^0 = \OO_X$ replaced by $0$ (in degree $0$). The natural inclusion $F^1\Omega_X^\bullet \hookrightarrow \Omega_X^\bullet$ of complexes induces a natural long exact sequnce of hypercohomology
\begin{equation}\label{exsq:DR}
	\cdots \to H_F^\bullet(X) \to H^\bullet_{DR}(X) \to H^\bullet(X, \OO_X) \to H_F^{\bullet + 1}(X) \to \cdots.
\end{equation}
One can naturally associate a cohomology class $c_1(\TT) \in H^2_F(X)$ to $\TT$.  %\footnote{This complex is used in \cite{BK}. In \cite{BB}, the class $c_1(\TT)$ was originally defined as a class in $H^2(X, \Omega_X^{\geqslant 1})$ where $\Omega_X^{\geqslant 1} : = (0 \to \Omega_X^1 \to \Omega_X^{\operatorname{cl}})$ is the truncated de Rham complex.}.  
The map $c_1(\cdot): \pic(X) \to H^2_F(X)$ is an isomorphism of vector spaces (\cite[Lemma 2.1.6]{BB}). In particular, we have $c_1(\TT_{\operatorname{triv}}) = 0 \in H^2_F(X)$. The composition of $c_1: \pic(X) \xrightarrow{\sim} H^2_F(X)$ with the natural map $H^2_F(X) \to H^2_{DR}(X)=H^2(X,\C)$ is the {\it topological first Chern class} map, denoted by $c^{\operatorname{top}}_1 : \pic(X) \to H^2(X,\C)$. For any algebraic line bundle $L$, the class $c^{\operatorname{top}}_1(\TT(L))$ is just the usual topological first Chern class of $L$.

Now we discuss a connection of Picard algebroids with sheaves of twisted differential operators.
Given a Picard algebroid $\TT$,
consider the universal enveloping sheaf $\mathscr{U}(\TT)$ of algebras.
Let $\mathscr{D}_\TT$ denote the quotient of $\mathscr{U}(\TT)$ modulo the ideal generated by the central section $1 - 1_\TT$, where $1$ the  identity section of $\mathscr{U}(\TT)$ and $1_\TT$ is the central section of $\TT$. Then $\mathscr{D}(\TT)$ is a {\it sheaf  of twisted differential operators (TDO)} in the sense of \cite[Definition 2.1.1]{BB}. In particular, it is equipped with
\begin{itemize}
	\item
	a morphism of algebras $\iota_{\mathscr{D}}: \OO_X \to \mathscr{D}_{\TT}$, and
	\item
	an increasing exhaustive algebra filtration $\mathscr{D}_{\TT, \leqslant i}$, $i \geqslant0$.
\end{itemize}
Both data are inherited from the similar structures on $\mathscr{U}(\TT)$. The morphism $\iota_{\mathscr{D}}$ identifies $\OO_X$ with $\mathscr{D}_{\TT, \leqslant 0}$ and the obvious morphism from the symmetric algebra $S(\mathscr{D}_{\TT, \leqslant 1} / \mathscr{D}_{\TT, \leqslant 0})$ into $\operatorname{gr} \mathscr{D}_{\TT}$ is an isomorphism of  commutative graded $\OO_X$-algebras.
%Alternatively, this filtration can be defined by induction: $\mathscr{D}_{\TT, \leqslant 0} = \OO_X$ and \[\mathscr{D}_{\TT, \leqslant i} (U) = \{ a \in \mathscr{D}_{\TT}(U) \,|\, r a - a r \in \mathscr{D}_{\TT, \leqslant i-1} (U)  \text{ for any }  r \in  \mathscr{D}_{\TT}(U) \}\] for any open subset $U$ of $X$.

In particular, taking the Lie bracket makes $\mathscr{D}_{\TT, \leqslant 1}$ into a Picard algebroid and the composite map $\TT \hookrightarrow \mathscr{U}(\TT) \to \mathscr{D}_{\TT}$ identifies $\TT$ with $\mathscr{D}_{\TT, \leqslant 1}$ as Picard algebroids. It follows that $\TT \mapsto \mathscr{D}_\TT$ gives an equivalence between the category of Picard algebroids and the category of TDOs. Under this equivalence, the trivial Picard algebroid $\TT_{\operatorname{triv}}$ corresponds to the sheaf $\mathscr{D}_X$ of usual (untwisted) differential operators on $X$.

We proceed to modules.
A $\TT$-module structure on a quasi-coherent $\OO_X$-module $\EE$ is equivalent to a (left) $\mathscr{D}_\TT$-module structure on $\EE$, such that the $\OO_X$-action on $\EE$ via $\iota_{\mathscr{D}}: \OO_X \to \mathscr{D}_{\TT}$ agrees with the $\OO_X$-module structure on $\EE$. Let $\coh(\TT)$ denote the abelian category of all $\TT$-modules over $X$ which are $\OO_X$-coherent and let $\coh(\mathscr{D}_\TT)$ denote the abelian category of all $\mathscr{D}_\TT$-modules over $X$ which are $\OO_X$-coherent. Then we have a canonical equivalence of categories $\coh(\TT) \simeq \coh(\mathscr{D}_\TT)$. It is well-known that a $\mathscr{D}_\TT$-module or a $\TT$-module $\EE$ is coherent as an $\OO_X$-module if and only if it is locally free (\cite[Lemma 2.3.1 (i)]{BB}).  

Assume now that $\EE$ is an algebraic vector bundle of rank $rk(\EE)$ over $X$. A $\TT$-module structure on  $\EE$ induces an isomorphism $\TT \xrightarrow{\sim} \frac{1}{rk(\EE)} \TT(\det(\EE))$ of Picard algebroids, where $\det(\EE) = \bigwedge^{\operatorname{top}} \EE$ is the determinant line bundle of $\EE$ (\cite[Lemma 2.3.1 (ii)]{BB}). Moreover, the $\TT$-module structure induces a projective flat connection on $\EE$, i.e, a flat connection on the projectivization $\mathbb{P}(\EE)$ of $\EE$. Conversely, by \cite[Lemma 2.3.2]{BB} and \cite[Section 2]{BC}, an isomorphism $\TT \xrightarrow{\sim} \frac{1}{rk(\EE)} \TT(\det(\EE))$ and a projective connection on $\EE$ determine the $\TT$-module structure on $\EE$.

This discussion proves the following claim.

\begin{Lem}\label{Lem:algebroid_w_coh_module} If $\TT$ admits a nonzero coherent module, then
	$\TT$ is of the form $r \TT(L)$, where $L$ is an algebraic line bundle and $r \in \mathbb{Q}$.
\end{Lem}


To finish this part we discuss pullbacks.
Let $\varphi:Y \to X$ be a morphism of smooth varieties. For a Picard algebroid, or more generally, a Lie algebroid $\TT$ over $X$, we define the pullback Picard algebroid $\varphi^{\circ} \TT$ as in \cite[Section 1.4.6 and 2.2]{BB}. The pullback map $\TT \mapsto \varphi^{\circ}\TT$ gives a linear map $\varphi^{\circ}: \pic(X) \to \pic(Y)$. Under the identifications $\pic(\cdot) \simeq H^2_F(\cdot)$, $\varphi^{\circ}$ is nothing else but the map $\varphi^*: H^2_F(X) \to H^2_F(Y)$ induced by the morphism $\varphi^* : \varphi^{-1} \Omega^{\geqslant 1}_X \to \Omega^{\geqslant 1}_Y$ of complexes (\cite[Section 2.2]{BB}). We also have standard pullback map $\varphi^*: \operatorname{Pic}(X) \to \operatorname{Pic}(Y)$ of line bundles given by $L \mapsto \varphi^*L$. The pullback maps intertwine the natural maps $\operatorname{Pic}(\cdot) \to \pic(\cdot) \simeq H^2_F(\cdot)$ by construction.

For any $\TT$-module $\M$ over $X$, $\varphi^*\M = \OO_Y \otimes_{\varphi^{-1} \OO_X} \varphi^{-1} \M$ admits naturally the structure of a $\varphi^\circ \TT$-module.
When $\M$ is $\OO_X$-coherent, $\varphi^*\M$ is $\OO_Y$-coherent. Therefore we have a pullback functor $\varphi^*: \coh(\TT) \to \coh(\varphi^\circ \TT)$. In a similar manner, we can pull back TDOs and twisted $D$-modules. Namely, for the TDO $\mathscr{D}_\TT$ assoicated to a Picard algebroid $\TT$ over $X$, we define $\varphi^\circ \mathscr{D}_\TT$ as the sheaf of algebras $\operatorname{Diff}_{\varphi^{-1} \mathscr{D}_\TT} (\varphi^*\mathscr{D}_\TT, \varphi^*\mathscr{D}_\TT)$ of all $\OO_Y$-differential operators on $\varphi^*\mathscr{D}_\TT$ that commutes with the right $\varphi^{-1} \mathscr{D}_\TT$-action. Then $\varphi^\circ \mathscr{D}_\TT$ is a TDO over $Y$ and there is a canonical isomorphism $\varphi^\circ \mathscr{D}_\TT \simeq \mathscr{D}_{\varphi^\circ \TT}$ of TDOs.

\subsubsection{Equivariant setting}
Let $K$ be an algebraic group with Lie algebra $\kf$ and $X$ be a smooth variety equppied with a $K$-action $\alpha: K \times X \to X$. Then $\alpha$ induces an infinitesimal $\kf$-action $d \alpha: \kf \to \Gamma(X, \T_X)$, which is a Lie algebra morphism. Recall the following definitions from \cite[Section 1.8]{BB}.

\begin{defi}\label{defn:pic_action}
	Let $(\LL, \eta)$ be a Lie algebroid over $X$. A {\it (weak) $K$-action} or a {\it $K$-equivariant structure} $\alpha_\LL$ on $\LL$ is a $K$-action on $\LL$ by automorphisms of Lie algebroid which lifts $\alpha$ and makes $\LL$ into a $K$-equivariant coherent sheaf.
	
	A {\it strong $K$-action} on $\LL$ is a pair $(\alpha_\LL, \nu_\kf)$, where $\alpha_\LL$ is a weak $K$-action on $\LL$ and $\nu_\kf: \kf \to \LL$ is a morphism of Lie algebras, such that
	\begin{enumerate}[(i)]
		\item
		$\nu_\kf$ is $K$-equivariant, where $\kf$ is equipped with the adjoint $K$-action;
		\item
		the infinitesimal $\kf$-action $d\alpha_\LL: \kf \to \der(\LL)$ on $\LL$ induced by $\alpha_\LL$ coincides with the adjoint $\kf$-action $\operatorname{ad}_{\nu_\kf}$ induced by $\nu_\kf$. Here
		$\der(\LL)$ is the Lie algebra of derivations $\LL\rightarrow \LL$.
	\end{enumerate}
	
	When $\LL$ is a Picard algebroid $(\TT, \eta, 1_\TT)$, we require that $K$ acts on $\TT$ by automorphisms of Picard algebroid, i.e., additionally $k^* 1_\TT = 1_\TT$ for any $k \in K$.
	
	A morphism between Lie algebroids or Picard algebroids with strong $K$-actions are required to be $K$-equivariant and intertwine the $\nu_\kf$'s. As in the non-equivariant case, the set of all isomorphism classes of strongly $K$-equivariant Picard algebroids over $X$ forms a vector space under Baer sum, which is denoted by $\pic_\kf(X)$.
\end{defi}

\begin{Rem}\label{rem:strong_diff}
	A strong $K$-action on $\TT$ is the same as a $K$-action $\alpha_{\mathscr{D}}$ on the corresponding TDO $\mathscr{D}_\TT$ by algebra automorphisms making the embedding $\mathcal{O}_X\hookrightarrow \mathscr{D}_\TT$ equivariant together with a Lie algebra map $\nu_\kf: \kf \to \mathscr{D}_{\TT, \leqslant 1}$ that satisfies conditions similar to (i) and (ii) above. In this case, we also say that $\mathscr{D}_\TT$ carries a strong $K$-action.
\end{Rem}


\begin{defi}\label{defn:pic_module_action}
	Let $\LL$ be a Lie algebroid over $X$ equipped with a weak $K$-action $\alpha_\LL$. A {\it weakly $K$-equivariant $\LL$-module}, or {\it weak $(\LL, K)$}-module, is an $\LL$-module $\M$ equipped with a $K$-equivariant structure (as an $\OO_X$-module) that is compatible with the $K$-action $\alpha_\LL$ on $\LL$. 	
	
	Suppose now $\LL$ is equipped with a strong $K$-action $(\alpha_\LL, \nu_\kf)$. Given $\kappa \in (\kf^*)^{K}$, a {\it $(K, \kappa)$-equivariant $\LL$-module} is a weak $(\LL, K)$-module $\M$ such that the differential of the $K$-action equals $\nu_\kf - \kappa \cdot \Id_\M$ (cf. Definition \ref{eq:HC_equiv}). A {\it (strong) $(\LL, K)$-module} is a $(K,\kappa)$-equivariant $\LL$-module $\M$ with $\kappa=0$.
	%, i.e., $\M$ is a weak $(\LL, K)$-module $\M$ such that the infinitesimal $\kf$-action on $\M$ coming from the $K$-action coincides with the one coming from $\nu_\kf: \kf \to \LL$ and the $\LL$-module structure on $\M$.
\end{defi}

Consider the case when $\LL$ is a Picard algebroid $\TT$. Let $\coh_K^\kappa(\TT)$ denote the abelian category of all $\OO_X$-coherent $(K,\kappa)$-equivariant $\TT$-modules over $X$, with morphisms being $K$-equivariant morphisms of $\TT$-modules.  In a similar manner, we can talk about weak and strong $(\mathscr{D}_\TT, K)$-modules and $(K, \kappa)$-equivariant $\mathscr{D}_\TT$-modules. Let $\coh_K^\kappa(\mathscr{D}_\TT)$ denote the abelian category of all $(K,\kappa)$-equivariant $\mathscr{D}_\TT$-modules that are $\OO_X$-coherent, with morphisms being $K$-equivariant morphisms of $\mathscr{D}_\TT$-modules. We also have an equivalence $\coh_K^\kappa(\TT) \simeq \coh_K^\kappa(\mathscr{D}_\TT)$ of abelian categories. When $\kappa=0$, we simply write $\coh_K^0(\TT)$ as $\coh_K(\TT)$ and write $\coh_K^{0}(\mathscr{D}_\TT)$ as $\coh_K(\mathscr{D}_\TT) $.
% to denote the abelian category of all strong $(\mathscr{D}_\TT,K)$-modules that are $\OO_X$-coherent.

When the variety $X$ carries a $\cm$-action, we also consider the groupoid of {\it graded Picard algebroids}, i.e., (weakly) $\cm$-equivariant Picard algebroids in the sense of Definition \ref{defn:pic_action}, together with $\cm$-equivariant isomorphisms among them. Let $\pic(X)^{gr}$ denote the vector space of equivalence classes of graded Picard algebroids over $X$, the linear structure of which is defined in the same way as the definition of that of $\pic(X)$. Now suppose $X$ is equipped with a $K \times \cm$-action. A \emph{strongly $K$-equivariant graded Picard algebroid} is a weakly $\cm$-equivariant Picard algebroid $\TT$ equipped with a strong $K$-action commuting with the $\cm$-action, such that the image of the associated Lie algebra morphism $\nu_\kf: \kf \to \Gamma(X, \TT)$ lies in the $\cm$-invariant part of $\Gamma(X, \TT)$. Isomorphisms of two such objects are required to be $\cm$-equivariant and compatible with the strong $K$-actions (in the sense of Defintion \ref{defn:pic_action}).  Let $\pic_\kf(X)^{gr}$ to denote the space of all isomorphism classes of strongly $K$-equivariant graded Picard algebroids. Again $\pic_\kf(X)^{gr}$ admits the structure of a vector space. When $\TT$ is a strongly $K$-equivariant graded Picard algebroid, we always assume a $\TT$-module  $\M$ admits a $\cm$-action which is compatible with the $\cm$-action on $\TT$ and the $\TT$-module structure.
We say that $\M$ is a graded $\TT$-module (or a graded twisted local system).

In what follows we will need the following easy technical lemma.

\begin{Lem}\label{Lem:graded_loc_sys}
Let $\M_1,\ldots,\M_k$ be graded $\TT$-modules that are simple as ordinary $\TT$-modules.
Let $\M$ be a graded $\TT$-module that is isomorphic to the direct sum of $\M_i$
as an ordinary module. Then $\M$ is isomorphic to the direct sum of $\M_i$'s with some shifts
as a graded $\TT$-module.
\end{Lem}
\begin{proof}
For two coherent $\TT$-modules $\M,\M'$, the space $\operatorname{Hom}_{\TT}(\M,\M')$ is finite
dimensional: it is the space of flat sections of the finite rank flat vector bundle $\mathcal{H}om_{\mathcal{O}_X}(\M,\M')$. In particular, the Schur lemma holds for
coherent $\TT$-modules. If $\M,\M'$ are graded $\TT$-modules, then $\operatorname{Hom}_{\TT}(\M,\M')$
has a natural grading. Then $\M\cong \bigoplus_{i=1}^k \operatorname{Hom}_{\TT}(\M_i,\M)\otimes_{\C}\M_i$,
as a graded module.
\end{proof}


\subsubsection{Monodromic sheaves}\label{SSS:monodromic}

For our later applications, it suffices to consider Picard algebroids $\TT$ which admit at least one $\OO_X$-coherent $\TT$-module $\EE$. By Lemma \ref{Lem:algebroid_w_coh_module}, $\TT$ must be of the form $L^r = r \TT(L)$, where $L$ is an algebraic line bundle and $r \in \mathbb{Q}$. Let $L^\times$ denote the total space of $L$ with the zero section excluded and let $\pi: L^\times \to X$ denote the projection map.  Let $H = \cm$, then $L^\times$ is the (right) $H$-torsor over $X$ associated to $L$, so that $L \simeq L^\times \times_{H} \C$. We describe the following alternative equivalent way to define $\TT$-modules via $H$-monodromic $\mathscr{D}$-modules over $L^\times$ as follows. This is a special case of \cite[Section 2.5.2]{BB}, where $H$ can be a torus of any dimension.

For a moment, we treat the general case when $H$ is any torus. Let $\hf$ denote the Lie algebra of $H$ and $\exp_H: \hf \to H$ denote the exponential map for $H$. Then $\hf_{\Z} := \ker(\exp)$ is a lattice in $\hf$ such that there are identifications of groups $\hf/\hf_\Z \simeq H$ and $\hf_\Z \simeq \pi_1(H)$ via the exponential map, where $\pi_1(H)$ is the fundamental group of $H$.  Define the dual torus $H^\vee$ to be $\Hom(\pi_1(H), \cm)$. Then the Lie algebra of $H^\vee$ is naturally identified with the linear dual $\hf^*$ of $\hf$ so that the exponential map $\exp_{H^\vee}: \hf^* \to H^\vee$ sends $\mu \in \hf^*$ to the character $e^{2\pi i \mu}: X \mapsto e^{2\pi i {\langle \mu, X \rangle}}$, where $X \in \hf_{\Z} \simeq \pi_1(H)$. Let $\hf^*_\Z := \{ \mu \,|\, \operatorname{Im}\mu \subset \Z \}$, then $\hf^*_\Z$ is a lattice of $\hf^*$. Then $\exp_{H^\vee}$ induces a Lie group isomorphism $\hf^*/\hf^*_\Z \xrightarrow{\sim} H^\vee$, $\bar{\mu} \mapsto e^{2\pi i \mu}$,  where $\bar{\mu}$ stands for the image of $\mu\in \hf^*$.

Back to the case when $H=\cm$, we identify $\hf$ with $\C$ in the way that the exponential map $\exp_H$ is given by $\exp_H(z) = e^{2 \pi i z}$, $\forall \, z \in \hf = \C$, so that $\hf_\Z = \pi_1(H)$ is identified with $\Z \subset \C$. We also identify $\hf^*$ with $\C$ so that $w \in \C$ corresponds to the linear functional $f_w$ of $\hf$ given by $f_w(z) = w z$, $\forall \, z \in \hf = \C$. In this way $\hf^*_\Z$ is identified with $\Z \subset \C = \hf^*$ and $H^\vee= \Hom(\pi_1(H), \cm) = \Hom(\Z, \cm)$ is identified with $\cm$ via $\phi \mapsto \phi(1)$, $\forall \, \phi \in H^\vee$. The exponential map of $H^\vee$ is hence given by $\exp_{H^\vee} (w) = e^{2 \pi i w}$. In what follows, we regard the complex number $r$ that appears in $\TT = r \TT(L)$ as an element $r \in \hf^*$ via the chosen identification $\hf^* = \C$.


Let $\wt{\mathscr{D}} = \wt{\mathscr{D}}_{L^\times}$ denote the sheaf of (usual) differential operators over $L^\times$. Then $\wt{\mathscr{D}}$ carries a natural $H$-action. Set $\mathscr{D}_\hf := (\pi_*\mathscr{D})^H$ and  $\wt{\mathscr{D}}_\hf := \pi^{-1} \mathscr{D}_\hf$. Then $\wt{\mathscr{D}}_\hf$ is the subseaf of $\wt{\mathscr{D}}$ consisting of all $H$-invariant differential operators. The differential of the right $H$-action on $L^\times$ yields an injective morphism $i_\hf: \hf \to \wt{\mathscr{D}}$ of Lie algebras, whose image lies in $\wt{\mathscr{D}}_\hf$ and hence $\hf$ also embeds into $\mathscr{D}_\hf$. Therefore we have an embedding of $\U\hf = S\hf$ into the center of $\mathscr{D}_\hf$. Regard $r$ as an element of $\hf^* \simeq \C$ as above, then it determines a maximal ideal $I_r$ of $\U\hf = S\hf$ generated by $h - r(h)$ for all $h \in \hf$.  Set $\mathscr{D}_{r,X} (= \mathscr{D}_r) := \mathscr{D}_\hf / I_r \mathscr{D}_\hf$ to be the quotient sheaf of $\mathscr{D}_\hf$ by the two-sided ideal $I_r \mathscr{D}_\hf$. We have a natural isomorphism between $\mathscr{D}_r$ and the sheaf of twisted differential operators $\mathscr{D}_{L^r}$ with twist $L^r = r \TT(L)$.

The $H$-action on $\wt{\mathscr{D}}$ and $i_\hf: \hf \to \wt{\mathscr{D}}$ give a strong $H$-action on $\wt{\mathscr{D}}$. Therefore we can talk about $(H, r)$-equivariant $\wt{\mathscr{D}}$-modules over $L^\times$ (which are also called {\it weakly $H$-equivariant $\wt{\mathscr{D}}$-module with monodromy $r$} in \cite[Section 2.5.2]{BB}). We use  $\coh(\wt{\mathscr{D}})^{r} = \coh_H^r(\wt{\mathscr{D}})$ to denote the abelian category of all $(H, r)$-modules that are $\OO_{L^\times}$-coherent. We then have an equivalence of categories $\coh(\wt{\mathscr{D}})^{r} \xrightarrow{\sim} \coh(\mathscr{D}_r)$ given by $\wt{\M} \mapsto (\pi_* \wt{\M})^H$. The quasi-inverse functor is given by $\M \mapsto \pi^* \M = \wt{\mathscr{D}} \otimes_{\wt{\mathscr{D}}_\hf } \pi^{-1}\M$, where the $\mathscr{D}_r$-module $\M$ is regarded as a $\mathscr{D}_\hf$-module via the quotient homomorphism $\mathscr{D}_\hf \to \mathscr{D}_r$ and hence $\pi^{-1} \M$ can be regarded as a module over $\wt{\mathscr{D}}_\hf = \pi^{-1} \mathscr{D}_\hf$.

Similarly we can regard $X$ and $L^\times$ as analytic manifolds. Denote the analytic versions of relevant sheaves by $\OO_X^{an}$, $\OO_{L^\times}^{an}$, $\mathscr{D}^{an}_{r,X}$, $\wt{\mathscr{D}}^{an}$, etc. Define the categories $\coh(\wt{\mathscr{D}}^{an})^{r} \simeq \coh(\mathscr{D}^{an}_r)$ in a similar manner. We have the analytification functors $(\cdot)^{an}:  \coh(\wt{\mathscr{D}})^{r}  \to \coh(\wt{\mathscr{D}}^{an})^{r}$, $\M \mapsto \M^{an} = \OO^{an}_{L^\times} \otimes_{\OO_{L^\times}} \M$, and $(\cdot)^{an}: \coh(\mathscr{D}_r) \to \coh(\mathscr{D}^{an}_r)$, $\M \mapsto \M^{an} = \OO^{an}_{X} \otimes_{\OO_{X}} \M$. These two functors intertwine the equivalences $\coh(\wt{\mathscr{D}})^{r} \simeq \coh(\mathscr{D}_r)$ and $\coh(\wt{\mathscr{D}}^{an})^{r} \simeq \coh(\mathscr{D}^{an}_r)$.

Note that if we have a weakly $H$-equivariant $\wt{\mathscr{D}}$-module $\wt{\M}$ with monodromy $r$ and $\wt{\M}'$ with monodromy $\nu$, then $\wt{\M} \otimes_{\OO_{L^\times}} \wt{\M}'$ is a naturally $\wt{\mathscr{D}}$-module with monodromy $r + \nu$. Let $\OO(1) := \pi^* L$, then it has a natural $H$-equivaraint structure. Each point $\tilde{x}$ in the fiber $L^\times_x $ at any $x \in X$ corresponds to a linear  isomorphism $\C \simeq L_x$. Therefore there is a canonical trivialization $\OO(1) \simeq \OO_{L^\times}$, via which $\OO_{L^\times}$ becomes a $\wt{\mathscr{D}}$-module. However, this trivialization is not $H$-equivariant and we have $\OO(1) \in \coh(\wt{\mathscr{D}})^{1}$. Suppose $\nu \in \hf^*_\Z=\Z$. Set $\OO(\nu):=\OO(1)^{\otimes \nu}$. Then $\OO(\nu) \in \coh(\wt{\mathscr{D}})^{\nu}$. We then have an equivalence of categories $\coh(\wt{\mathscr{D}})^{r}  \xrightarrow{\sim} \coh(\wt{\mathscr{D}})^{r + \nu}$, $\wt{\M} \mapsto \OO(\nu) \otimes_{\OO_{L^\times}} \wt{\M}$, whose inverse is given by tensoring with $\OO(-\nu)$. The image of $\OO(\nu)$ under the equivalence $\coh(\wt{\mathscr{D}})^{\nu} \xrightarrow{\sim} \coh(\mathscr{D}_\nu)$ is $(\pi_*\OO(\nu))^H = L^{\nu} = L^{\otimes \nu}$. We have an isomorphism of algebras $\mathscr{D}_{r+\nu} \simeq  L^{\nu} \otimes_{\OO_X} \mathscr{D}_{r} \otimes_{\OO_X} L^{-\nu}$ for any $r \in \hf^*$ and $\nu \in \hf^*_\Z$, which induces an equivalence of categories $\coh(\mathscr{D}_r) \xrightarrow{\sim} \coh(\mathscr{D}_{r + \nu})$, $\M \mapsto L^{\nu} \otimes_{\OO_X} \M$, whose inverse is given tensoring with $L^{-\nu}$. Of course these two equivalences intertwine the equivalences $\coh(\wt{\mathscr{D}})^{r} \simeq \coh(\mathscr{D}_r)$ and $\coh(\wt{\mathscr{D}})^{r+\nu} \simeq \coh(\mathscr{D}_{r+\nu})$.


\subsubsection{Riemann-Hilbert correpondence} \label{subsec:RH}
We now turn to monodromic local systems over $L^\times$, considered as an analytic manifold $(L^\times)^{an}$.  Given any $r \in \hf^*$, let $\loc(L^\times)^{\bar{r}}$ denote the category of all local systems with monodromy $e^{2 \pi i r}$ over $L^\times$ (as an analytic manifold). Note that the definition only depends on the image $\bar{r}$ of $r$ in $\hf^*/\hf_\Z^*$.

Assume $X$ is a connected smooth variety, so that $X^{an}$ and $(L^\times)^{an}$ are also connected under the analytic topology. Fix a point $x$ in $X$ and choose any point $\tilde{x}$ in the fiber of $L^\times$ at $x$. Since $L^\times$ is an $H$-torsor, all groups $\pi_1(L^\times, \tilde{x})$ for different choices of $\tilde{x}$ in the fiber of $L^\times$ at a given $x$ are canonically isomorphic. Hence we will write $\pi_1(L^\times, x)$  instead of $\pi_1(L^\times, \tilde{x})$. Moreover, the action on $L^\times$ by $H = \cm$ gives rise to a canonical embedding of $\pi_1(H)$ into $\pi_1(L^\times, x)$. We then have a short exact sequence of groups
\begin{equation}\label{eq:pi_exsq}
	1 \to \pi_1(H) \to \pi_1(L^\times, x) \to \pi_1(X, x) \to 1.
\end{equation}
It is well-known that the category of local systems over $(L^\times)^{an}$ is equivalent to the category of finite dimensional representations of the fundamental group $\pi_1(L^\times, \tilde{x})$ of $L^\times$ based at $\tilde{x}$ by taking the monodromy. Given any $r \in \hf^*=\C$, we say that a local system over $(L^\times)^{an}$ \emph{has monodromy $\bar{r}$} if  $\pi_1(H)$ acts on the corresponding monodromy representation by the character $\bar{r} = e^{2\pi i r} \in H^\vee$. Let $\loc(L^\times)^{\bar{r}}$ denote the category of all local systems over $(L^\times)^{an}$ with monodromy $\bar{r}$.  The Riemann-Hilbert correspondence induces an equivalence of abelian categories $( \cdot )^\nabla : \coh(\wt{\mathscr{D}}^{an}_{L^\times})^r \to \loc(L^\times)^{\bar{r}}$, which sends a $\wt{\mathscr{D}}^{an}$-module $\wt{\M}$ to the sheaf $\wt{\M}^\nabla$ of its flat sections. The inverse functor is given by $\wt{S} \mapsto \OO_{L^\times} \otimes_{\C} \wt{S}$, where the $H$-action on $\OO_{L^\times} \otimes_{\C} \wt{S}$ is the diagonal one. Note that objects in $\loc(L^\times)^{\bar{r}}$ are automatically $H$-equivariant. 

The functor $(\cdot)^{an}:  \coh(\wt{\mathscr{D}})^{r}  \to \coh(\wt{\mathscr{D}}^{an})^{r}$ is not an equivalence in general. On the other hand, let $\coh^{reg}(\wt{\mathscr{D}})^{r}$ denote the full subcategory of $\coh(\wt{\mathscr{D}})^{r}$ consisting of those $\wt{\mathscr{D}}$-modules with regular singularities. Then $(\cdot)^{an}$ restricts to an equivalence of categories $\coh^{reg}(\wt{\mathscr{D}})^{r} \xrightarrow{\sim} \coh(\wt{\mathscr{D}}^{an})^{r}$ and therefore we have the equivalence of categories $\coh^{reg}(\wt{\mathscr{D}})^{r} \xrightarrow{\sim}  \loc(L^\times)^{\bar{r}}$, $\wt{\M} \mapsto (\wt{\M}^{an})^{\nabla}$(\cite{Deligne}, \cite[Remark 2.5.5, (iii)]{BB}). Denote by $\coh^{reg}(\mathscr{D}_r)$ the essential image of $\coh^{reg}(\wt{\mathscr{D}})^{r}$ under the equivalence $\coh(\wt{\mathscr{D}})^{r} \xrightarrow{\sim} \coh(\mathscr{D}_r)$. Then we also have an equivalence  $\coh^{reg}(\mathscr{D}_r) \xrightarrow{\sim} \loc(L^\times)^{\bar{r}}$. Note the latter category depends only on the image of $r$ in $\hf^*/\hf^*_\Z$. The equivalences $\coh(\mathscr{D}_r) \xrightarrow{\sim} \coh(\mathscr{D}_{r + \nu})$ for $\nu \in \hf^*_\Z$ preserve the subcategory of objects with regular singularities and intertwine the equivalences with the category of twisted local systems.


We take the pushout of the short exact sequence \eqref{eq:pi_exsq} along the character $\bar{r}: \pi_1(H) \to \cm$, to obtain a central extension of the group $\pi_1(X,x)$ by $\cm$,
\[  \quad 1 \to \cm \to \pi^{\tilde{r}}_1(X, x) \to \pi_1(X, x) \to 1.  \]
This extension determines a Schur multiplier $\tilde{r} \in H^2(\pi_1(X, x), \cm)$ of $\pi_1(X, x)$. Let $\operatorname{Rep}(\pi^{\tilde{r}}_1(X, x))$ denote the abelian category of all finite dimensional representations $V$ of $\pi^{\tilde{r}}_1(X, x)$ such that $\cm \subset \pi^{\tilde{r}}_1(X, x)$ acts on $V$ by the usual scalar multiplication, i.e., $V$ is a projective representation of $\pi_1(X)$ with Schur multiplier $\tilde{r}$. For any $\wt{S} \in \loc(L^\times)^{\bar{r}}$, we take its fiber at any point $\tilde{x} \in L^\times$ above $x$, then its monodromy representation factors through a representation in $\operatorname{Rep}(\pi^{\tilde{r}}_1(X, x))$. This defines an equivalence of abelian categories $\loc(L^\times)^{\bar{r}} \xrightarrow{\sim} \operatorname{Rep}(\pi^{\tilde{r}}_1(X, x))$.

All the constructions and equivalences of categories have natural functoriality with respect to pullback. Namely, let $\varphi: Y \to X$ be a morphism of smooth varieties. Suppose $L$ is an algebraic line bundle over $X$, then $L':=\varphi^* L$ is a line bundle over $Y$ and $L'^\times$ is the pullback $H$-torsor of $L^\times$ via $\varphi$, together with an $H$-equivariant map $\tilde{\varphi}: L'^\times \to L^\times$ lifting $\varphi$. We then have the pullback functor $\tilde{\varphi}^{-1}: \operatorname{Loc}(L^\times)^{\bar{r}} \to \operatorname{Loc}(L'^\times)^{\bar{r}}$, $\wt{S} \mapsto \tilde{\varphi}^{-1} \wt{S}$. Suppose $X$ and $Y$ are connected as complex
analytic manifolds. Fix a point $y \in Y$ and set $x = \varphi(y)$. We then have the Cartesian diagrams of groups
\begin{equation}\label{diag:pi_groups}
	\begin{tikzcd}
		\pi_1(L'^\times, y) \arrow{r} \arrow[d, "\tilde{\varphi}_*"]  & \pi_1(Y,y)  \arrow[d, "\varphi_*"]   \\
		\pi_1(L^\times, x) \arrow{r}  & \pi_1(X, x)
	\end{tikzcd}	
	\quad \text{and} \quad
	\begin{tikzcd}
		\pi^{\tilde{r}'}_1(Y, y) \arrow{r} \arrow[d, "\tilde{\varphi}_*"]  & \pi_1(Y,y)  \arrow[d, "\varphi_*"]   \\
		\pi^{\tilde{r}}_1(X, x) \arrow{r}  & \pi_1(X, x)
	\end{tikzcd}
\end{equation}
where $\tilde{r}' \in H^2(\pi_1(Y,y), \cm)$ is the Schur multiplier for the $H$-torsor $L'^\times$ and $r$, which coincides with the pullback of $\tilde{r} \in H^2(\pi_1(X,x), \cm)$ via the group homomorphism $\varphi_*: \pi_1(Y,y) \to \pi_1(X,x)$. Therefore we have a functor $\tilde{\varphi}^*: \operatorname{Rep}(\pi^{\tilde{r}}_1(X, x)) \to \operatorname{Rep}(\pi^{\tilde{r}'}_1(Y, y))$, $(V, \rho) \mapsto (V, \rho \circ  \tilde{\varphi}_*)$. The discussion above yields the following proposition.

\begin{Prop}\label{Prop:regular_local_systems}
	There are natural equivalence of categories \[ \coh^{reg}(\mathscr{D}_{r,X}) \xrightarrow{\sim} \coh^{reg}(\wt{\mathscr{D}}_{L^\times})^{r}  \xrightarrow{\sim} \loc(L^\times)^{\bar{r}}  \xrightarrow{\sim} \operatorname{Rep}(\pi^{\tilde{r}}_1(X, x)) \]
	which are functorial with respect to pullback functors, i.e., for any morphism $\varphi: Y \to X$ of smooth varieties which are connected as  complex analytic manifolds, we have natural equivalences of categories illustrated by the diagram
	\begin{equation}\label{diag:pi_functors}
		\begin{tikzcd}
			\coh^{reg}(\mathscr{D}_{r,X}) \arrow[r, "\sim"] \arrow[d, "\tilde{\varphi}^*"] & \coh^{reg}(\wt{\mathscr{D}}_{L^\times})^{r} \arrow[r, "\sim"] \arrow[d, "\tilde{\varphi}^*"]   & \loc(L^\times)^{\bar{r}}  \arrow[r, "\sim"]  \arrow[d, "\varphi^{-1}"]  & \operatorname{Rep}(\pi^{\tilde{r}}_1(X, x)) \arrow[d, "\tilde{\varphi}^*"] \\	
			\coh^{reg}(\mathscr{D}_{r,Y}) \arrow[r, "\sim"]   & \coh^{reg}(\wt{\mathscr{D}}_{L'^\times})^{r} \arrow[r, "\sim"]    & \loc(L'^\times)^{\bar{r}}  \arrow[r, "\sim"]   & \operatorname{Rep}(\pi^{\tilde{r}'}_1(Y, y)). 			
		\end{tikzcd}	
	\end{equation}
\end{Prop}

Now suppose $X$ admits an algebraic action by an algebraic group $K$ and $L$ is a $K$-equivariant line bundle over $X$. We can consider the category $\coh_K(\mathscr{D})^r$ of strongly $K$-equivariant objects in $\coh(\mathscr{D})^r$ and the category $\loc_{K}(L^\times)^{\bar{r}}$ of $K$-equivariant local systems in $\loc(L^\times)^{\bar{r}}$, i.e., local systems that carry $K$-actions which lift the $K$-action on $L^\times$
subject to the usual compatibility conditions. The category $\coh_K(\mathscr{D})^r$ is equivalent to the category $\coh_K(\mathscr{D}_r)$ of strongly $K$-equivariant objects in $\coh(\mathscr{D}_r)$.  Assume additionally that the $K$-action on $X$ has only finitely many orbits. Then $L^\times$ has finitely many $K \times H$-orbits. This implies that all objects in $\coh_K(\mathscr{D})^r$ automatically have regular singularities (\cite[Theorem 11.6.1]{Hotta}). Therefore we have the equivalences of categories $\coh_K(\mathscr{D}_r) \simeq \coh_K(\mathscr{D})^r \simeq  \loc_{K}(L^\times)^{\bar{r}}$.


\subsubsection{Homogeneous case}\label{SSS_homog_case}
Let $G$ be a connected reductive algebraic group, and $H$ be its closed subgroup. Pick $\lambda\in
\operatorname{Hom}(H,\C^\times)\otimes_{\Z}\C = (\hf^*)^H$. Similar to Section \ref{SSS:monodromic}, to $\lambda$ we can assign a strongly $G$-equivariant Picard algebroid $\TT_\lambda$ over $G/H$ (see \cite[Section 2.4.2]{Yu} for details), or equivalently, a sheaf of twisted
differential operators $\mathscr{D}^\lambda_{G/H}$ with a strong $G$-action (see Remark \ref{rem:strong_diff}). In fact, this gives an isomoprhism $(\hf^*)^H \simeq \pic_\g(G/H)$ of vector spaces (see \cite[Proposition 2.4]{Yu}). Pick $\kappa\in \g^{*G}$. We want to classify strongly
$(G,\kappa)$-equivariant coherent $\mathscr{D}^\lambda_{G/H}$-modules. Since the action of $G$
on $G/H$ is transitive, each such module $M$ is a $G$-equivariant coherent sheaf on
$G/H$. Set $x:=1H\in G/H$. So, the fiber $M_x$ is a rational representation of $H$.
The following lemma is quite standard and is a variant of \cite[Lemma 2.7]{Yu}.

\begin{Lem}\label{Lem:homog_classif}
The assignment $M\mapsto M_x$ defines a category equivalence between
\begin{itemize}
\item[(a)] The category of coherent strongly $(G,\kappa)$-equivariant $\mathscr{D}^\lambda_{G/H}$-modules,
\item[(b)] and the category of finite dimensional rational $H$-representations, where
$\mathfrak{h}$ acts by $\lambda-\kappa|_{\mathfrak{h}}$.
\end{itemize}
\end{Lem}
\begin{proof}
Let $\pi$ denote the projection $G\rightarrow G/H, g\mapsto gH$. As in Section \ref{SSS:monodromic},
the functor $\pi^*$ gives an equivalence between (a) and
\begin{itemize}
\item[(c)] The category of coherent strongly $(G\times H, \kappa\times (-\lambda))$-equivariant
$\mathscr{D}_G$-modules (the minus sign appears because the action of $H$ on $G$ is by $h.g=gh^{-1}$).
\end{itemize}
For $\xi\in \g$, let $\xi_G$ (resp., $\xi_M$) denote the vector field on $G$ (resp., $\C$-linear endomorphism
of $M$ induced by $\xi$). For $\eta\in \mathfrak{h}$, let $\eta^r_H,\eta^r_M$ have the similar meaning
for the $H$-action on $G$ from the right, like in (c). 
Note that $\eta_{G,1}=-\eta^r_{H,1}$
for $\eta\in \mathfrak{h}$. The stabilizer of $1$ in $G\times H$ is $H$ embedded diagonally. So the action
of $\mathfrak{h}$ on $M_1$ is by $\eta\mapsto \eta_{M,1}+\eta^r_{M,1}$. Since $\xi_G-\langle\kappa,\xi\rangle=\xi_M$
and $\eta_G+\langle\lambda,\eta\rangle=\eta_M$, we see that $\mathfrak{h}$ acts on $M_1$ by
$\eta\mapsto \langle\lambda-\kappa|_{\mathfrak{h}},\eta\rangle$, so $M\mapsto M_{x}=(\pi^*M)_1$
is indeed a functor from (a) to (b).

It remains to show that every object $N$ in (b) arises as $M_1$. We have a unique $(G,\kappa)$-equivariant
$\mathscr{D}_G$-module with fiber $N$ at $1$. It is easy to see that it is also $(H,-\lambda)$-equivariant.
This finishes the proof.
\end{proof}

\begin{Rem}\label{Rem:homog_analyt_classif}
The same argument shows that the same statement as in Lemma \ref{Lem:homog_classif} 
holds if in (a) we replace $\mathscr{D}^\lambda_{G/H}$ with its analytification.
\end{Rem}

Now we consider a more general situation. Let $K$ be a subgroup of $G$ that acts
on $G/H$ with finitely many orbits, say $Kx_1,\ldots, Kx_\ell$ for $x_i:=g_i H\in G/H$.
Note that $H$ acts trivially on $\operatorname{Hom}(H,\C^\times)_{\Z} \otimes \C$. So conjugating with
$g_i$ gives rise to a well-defined element $\lambda_i\in \operatorname{Hom}(G_{x_i},\C^\times)_{\Z} \otimes \C$.
Let $\Upsilon_i$ denote the set of isomorphism classes of irreducible $K_{x_i}$-modules
where $\kf_{x_i}$ acts by $-(\lambda_i|_{\kf_{x_i}}-\kappa|_{\kf_{x_i}})$.

The following result is proved in the same way as \cite[Theorem 11.6.1,(ii)]{Hotta}.

\begin{Lem}\label{Lem:fin_many_orbit_classif}
The simple (automatically regular) $(K,\kappa)$-equivariant $\mathscr{D}^{\lambda}_{G/H}$-modules
are classified by the elements of the set $\bigsqcup_{i=1}^\ell \Upsilon_i$:
to get from an element of $\Upsilon_i$ to a simple object, take the minimal extension of the twisted local
system on $Kx_i$ corresponding to the element of $\Upsilon_i$ via Lemma \ref{Lem:homog_classif}.
\end{Lem}



\subsection{Deformation quantization of symplectic varieties}\label{subsec:quan}
\subsubsection{Definitions}
We review deformation quantization in the algebraic context from \cite{BK}. Let $X$ be a smooth algebraic variety (over $\C$) equipped with an algebraic symplectic form $\Omega \in H^0(X, \Omega_X^2)$. Then $\Omega$ induces a  Poisson bracket $\{ - , -\}$ on the structure sheaf $\OO_X$ of $X$.

\begin{defi}\label{defi:formal_quan_t}
	A \emph{formal deformation quantization over $\C[[\hbar]]$}, or simply {\it quantization} of a smooth symplectic variety $(X,\Omega)$ is the following datum
	\begin{itemize}
		\item
		a sheaf $\OO_\hbar$ of flat associative $\C[[\hbar]]$-algebras on $X$, complete and separated in the $\hbar$-adic topology, such that the commutator $[a, b] = a \cdot b - b \cdot a \in \hbar \OO_\hbar$ for any $a, b \in \OO_\hbar$, and
		\item
		an isomorphism $\phi: \OO_\hbar / \hbar \OO_\hbar \xrightarrow{\sim} \OO_X$ of sheaves of commutative Poisson algebras, where the Poisson bracket on $ \OO_\hbar / \hbar \OO_\hbar$ is induced by $\hbar^{-1}[a,b]$, $a, b \in \OO_\hbar$.
	\end{itemize}
\end{defi}

There is an obvious definition of equivalent quantizations. Denote by $\quan(X,\Omega)$ or $\quan(X)$ the set of equivalence classes of quantizations of $X$. There is also a natural notion of quantization for a scheme $X$ which is smooth over a scheme $S$ of finite type and is equipped with a relative sympletic form $\Omega \in H^0(X, \Omega^2_{X/S})$, where $\Omega^\bullet_{X/S}$ is the relative de Rham complex. All the constructions and results in this section also work in the general relative setting. See \cite{BK} for details. We will suppress $S$ for most of the time in the rest of the paper.

We recall the following main result from \cite[Theorem 1.8]{BK} which classifies all quantizations under an additional assumption: $X$ is said to be {\it admissible} if the natural homomorphism $H^i_{DR}(X) \to H^i(X, \OO_X)$ is surjective for $i = 1, 2$. This implies that the natural map $H^2_F(X) \to H^2_{DR}(X)$ is injective and its image coincides with the kernel of the map $H^2_{DR}(X) \twoheadrightarrow H^2(X, \OO_X)$.

\begin{Thm}\label{thm:quan_BK}
	Let $X$ be an admissible smooth variety equipped with a symplectic form $\Omega$. Then there is a natural injective map
	\[  \per: \quan(X, \Omega) \hookrightarrow \hbar^{-1} H^2_{DR}(X) [[\hbar]],   \]
	called the {non-commutative period map}. Futhermore, for any $\OO_\hbar \in \quan(X,\Omega)$, we have $\per(\OO_\hbar) \in \hbar^{-1} [\Omega] + H^2_{DR}[[\hbar]]$, where $[\Omega] \in H^2_{DR}(X)$ denotes the cohomology class of the symplectic form $\Omega$. Finally, any choice of linear splitting $P: H^2_{DR}(X) \twoheadrightarrow H^2_F(X)$ of the canonical embedding $H^2_F(X) \to H^2_{DR}(X)$ induces a bijection $\quan(X, \Omega) \xrightarrow{\sim} H^2_F(X)[[\hbar]]$ given by $\OO_\hbar \mapsto P( \per(\OO_\hbar) - \hbar^{-1} [\Omega])$.
\end{Thm}

For our applications, we need to consider \emph{graded} quantizations(see \cite[Section, 2.2]{quant_iso} \cite[Section 3.1]{orbit}  and \cite[Section 3.2]{BPW}). Fix a positive integer $d$.  Suppose that the symplectic variety $X$ admits a rational $\cm$-action, with respect to which the symplectic form $\Omega$ has weight $d$, i.e., the push-forward $z.\Omega = z_* \Omega$ of the form $\Omega$ by the automorphism of $X$ induced by $z \in \cm$ equals the rescaled form $z^d\Omega$. In this case we say that $(X,\Omega)$ is \emph{graded} (of weight $d$).

\begin{defi}\label{defn:graded_quan}
	Suppose $(X, \Omega)$ is a graded symplectic variety of weight $d$. A quantization $\OO_\hbar$ of $(X,\Omega)$  is said to be \emph{graded} if it is equipped with a pro-rational  $\cm$-action by algebra automorphisms such that the isomorphism $\OO_\hbar / \hbar \OO_\hbar \simeq \OO_X$ is $\cm$-equivariant and $z. \hbar = z^d \hbar$ for any $z \in \cm$. We require isomorphisms between graded quantizations to be $\cm$-equivariant. The set of isomorphism classes of graded quantizations of $(X,\Omega)$ is denoted as $\quan(X, \Omega)^{gr}$ or $\quan(X)^{gr}$.
\end{defi}

We write
\begin{equation} \label{eq:per}
	\per(\OO_\hbar) =  \hbar^{-1}[\Omega] + \omega_1(\OO_\hbar) +  \hbar \omega_2(\OO_\hbar)+  \hbar^2 \omega_3 (\OO_\hbar)+ \cdots \in  \hbar^{-1}H^2_{DR}(X)[[\hbar]],
\end{equation}
where $\omega_k(\OO_\hbar) \in H^2_{DR}(X)$. Note that when $(X, \Omega)$ is graded, the de Rham class $[\Omega] \in H^2_{DR}(X)$ of the symplectic form $\Omega$ vanishes, thanks to the well-known Cartan homotopy formula. The following result is due to first named author (\cite[Section 2.2 and Prop. 2.3.2]{quant_iso}).

\begin{Prop}\label{prop:graded}
	Let $(X,\Omega)$ be a graded smooth symplectic variety.% Then
%	\begin{enumerate}
%		\item
		If $\OO_\hbar$ is a graded quatization of $(X, \Omega)$, then $\per(\OO_\hbar) \in H^2_{DR}(X, \C) \subset  \hbar^{-1}H^2_{DR}(X)[[\hbar]]$, i.e., $\omega_k(\OO_\hbar) = 0$ for all $k \geqslant 2$ in \eqref{eq:per}.
The converse is true if $X$ is admissible.
%		\item
%		If $\OO_\hbar$ is even, then $\per(\OO_\hbar) =  \hbar^{-1}[\Omega] = 0$.\footnote{Not sure if we need to define even quantizations. -- Shilin}
%		\item
%		If $X$ is, in addition, admissible, then a graded quantization $\OO_\hbar$ is even if and only if $\per(\OO_\hbar)=0$.  \footnote{Do we need the notion of even quantization? --Shilin}
%	\end{enumerate}
\end{Prop}

Then Theorem \ref{thm:quan_BK} and Proposition \ref{prop:graded} together imply the following result.

\begin{Cor}\label{cor:per_graded}
	Let $(X, \Omega)$ be a graded smooth symplectic variety. Suppose $H^1(X,\OO_X) = H^2(X,\OO_X)=0$. Then the period map $\per$ identifies the set $\quan(X, \Omega)^{gr}$ of isomorphism classes of graded quantizations of $(X,\Omega)$ with $H^2_{DR}(X) \subset  \hbar^{-1}H^2_{DR}(X)[[\hbar]]$.
\end{Cor}


\begin{Rem} \label{rmk:enhanced_period}
	In  \cite{Yu}, the second named author has defined an {\it enhanced (truncated) non-commutative period map} $\overline{\per}: \quan(X, \Omega) \to H^2_F(X)$ for any (not necessarily admissible) symplectic variety $X$ and proved that, for any quantization $\OO_\hbar$ of $(X, \Omega)$, the map $	c_1^{\operatorname{top}}: H^2_F(X) \to H^2_{DR}(X)$ sends $\overline{\per}(\OO_\hbar)$ to $ \omega_1(\OO_\hbar)$, where $\omega_1(\OO_\hbar)$ is as in \eqref{eq:per}. This is compatible with Corollary \ref{cor:per_graded}, since in the context there we have $H^2_F(X) \simeq H^2_{DR}(X)$. The enhanced period map will play a role in quantizations of Lagrangian subvarieties. See Remark \ref{rmk:enhanced_period_lag}.
\end{Rem}

\begin{comment}
\begin{defi}\label{defi:formal_quan_h}
	Let $(X,\Omega)$ be a graded symplectic variety so that the symplectic form $\Omega$ has weight $d$. A \emph{graded quantization over $\C[[\hbar]]$}, of $(X,\Omega)$ is the following datum
	\begin{itemize}
		\item
		a sheaf $\OO_\hbar$ of flat associative $\C[[\hbar]]$-algebras on $X$, complete and separated in the $\hbar$-adic topology, such that the commutator $[a, b] = a \cdot b - b \cdot a \in \hbar^d \OO_\hbar$ for any $a, b \in \OO_\hbar$,
		\item
		an isomorphism $\OO_\hbar / \hbar \OO_\hbar \xrightarrow{\sim} \OO_X$ of sheaves of commutative Poisson algebras, where the Poisson bracket on $ \OO_\hbar / \hbar \OO_\hbar$ is induced by $\hbar^{-d}[a,b]$, $a, b \in \OO_\hbar$, and
		\item
		a pro-rational  $\cm$-action on $\OO_\hbar$ by algebra automorphisms such that the isomorphism $\OO_\hbar / \hbar \OO_\hbar\xrightarrow{\sim} \OO_X$ is $\cm$-equivariant and $z. \hbar = z^d \hbar$ for any $z \in \cm$.
	\end{itemize}
    The set of ($\cm$-equivariant) isomorphism classes of graded quantizations over $\C[[\hbar]]$ of $(X, \Omega)$ is denoted as $\quan(X, \Omega)^{gr}$.
\end{defi}
\end{comment}

We finish this section with an example. Suppose $X$ be  a conical symplectic singularity
such that $\operatorname{codim}_X X^{sing}\geqslant 4$. Then
\begin{equation}\label{eq:orbit_coh_cohom_vanish}
H^i(X^{reg},\mathcal{O}_X)=0, \text{ for }i=1,2.
\end{equation}
Recall that in Section \ref{SS_quant_nilp} to $\lambda\in H^2(X^{reg},\C)$
we have assigned the filtered quantization $\A_\lambda$ of $\C[X]$. Let $\A_{\lambda,\hbar}$ denote
the completion of its modified Rees algebra. The microlocalization
$\A_{\lambda,\hbar}|_{\Orb}$ is the graded quantization with period $\lambda$,
see, e.g., \cite[Section 3.2]{orbit}.

\subsection{Quantization of Lagrangian subvarieties} \label{subsec:quan_lag}

\subsubsection{Definitions}

Suppose that $Y$ is a smooth closed Lagrangian subvariety of a symplectic variety $X$ with the embedding denoted by $\iota : Y \hookrightarrow X$. When $X$ is graded, we also assume $Y$ is $\cm$-stable. There is a natural notion of (graded) quantization of a ($\cm$-equivariant) vector bundle over $Y$ (\cite[Section 1]{BC}, \cite[Section 4.3]{LY}).

\begin{defi}
	Suppose $\OO_\hbar$ is a quantization of $(X,\Omega)$ and $\EE$ is a vector bundle of rank $r$  over $Y$. A {\it quantization} of $\EE$, or rather the sheaf $\iota_* \EE$ of $\OO_X$-modules, is an $\OO_\hbar$-module $\EE_\hbar$ flat over $\C[[\hbar]]$, complete and separated in the $\hbar$-adic topology, with an isomorphism $\EE_\hbar /\hbar \EE_\hbar \cong \iota_*\EE$ of $\OO_X$-modules.
	
	Now assume $(X,\Omega)$ is graded and $Y$ is a $\cm$-stable Lagrangian subvariety of $X$. Suppose $\OO_\hbar$ is a graded quantization of $X$ and $\EE$ is a $\cm$-equivariant ($\cm$-linearized) vector bundle of rank $r$ over $Y$. A {\it graded quantization} of $\EE$ is a quantization $\EE_\hbar$ of $\EE$ with a pro-rational $\cm$-action, such that the structure map $\OO_\hbar \otimes_{\C[[\hbar]]} \EE_\hbar \to \EE_\hbar$ is $\cm$-equivariant and the isomorphism $\EE_\hbar /\hbar \EE_\hbar \cong \iota_*\EE$ is also $\cm$-equivariant.
\end{defi}

\begin{comment}
	\subsubsection{Jet bundles}
	
	We want to define the jet bundle of $M_\hbar^0$.  First,
	consider the  sheaf of algebras $\hat{\A}_\hbar^{\Orb}\boxtimes \Str_{\Orb_K}$ on $\Orb\times \Orb_K$.
	Let $\mathcal{J}_\Delta$ denote the ideal sheaf of the diagonal in $\Str_{\Orb}\boxtimes
	\Str_{\Orb_K}$ and let $\mathcal{J}_{\Delta,\hbar}$ denote its preimage in
	$\hat{\A}_\hbar^{\Orb}\boxtimes \Str_{\Orb_K}$. Let $\jet_{\Orb_K} \hat{\A}_\hbar$
	denote the completion of  $\hat{\A}_\hbar^{\Orb}\boxtimes \Str_{\Orb_K}$ with respect to
	this ideal. It is nothing else but the specialization of $\jet \hat{\A}_\hbar^{\Orb}$
	to $\Orb_K$. It is a pro-vector bundle of algebras on $\Orb_K$ that comes with a flat
	connection.
	
	We can use a similar construction to define the jet bundle $\jet_{\Orb_K} M_\hbar^0$.
	This is a $\jet_{\Orb_K} \hat{\A}_\hbar$-module that is equipped with a compatible flat connection.
	Note that the sheaf $\jet_{\Orb_K} \hat{\A}_\hbar$ comes with an $\Str_{\Orb_K}$-linear
	involution induced by the anti-involution $\theta\boxtimes \operatorname{id}$
	on $\hat{\A}_\hbar^{\Orb}\boxtimes \Str_{\Orb_K}$. It will also be denoted by $\theta$. Using this
	involution, we can define the notion of a weakly HC $(\jet_{\Orb_K} \hat{\A}_\hbar,\theta)$-module
	similarly to what was done in \cite[Section 2.6]{HC_symp}. Namely,  these are finitely generated
	$\jet_{\Orb_K} \hat{\A}_\hbar$-modules $\mathfrak{M}^0_\hbar$ that come equipped with a compatible
	flat connection $\nabla$ and the Lie algebra action of $(\jet_{\Orb_K} \hat{\A}_\hbar)^{-\theta}$ subject
	to the following conditions:
	
	\begin{itemize}
		\item[(i)] The action map $(\jet_{\Orb_K} \hat{\A}_\hbar)^{-\theta}\otimes \mathfrak{M}^0_\hbar\rightarrow
		\mathfrak{M}^0_\hbar$ is flat. \footnote{Needs to be rewritten in terms of $\J\hbar$.  {\textcolor{red}{-- Shilin}}}
		\item[(ii)] $\mathfrak{M}^0_\hbar$ is complete and separated in the $\hbar$-adic topology.
		\item[(iii)] The successive quotients $\hbar^{k-1}\mathfrak{M}^0_\hbar/\hbar^k \mathfrak{M}^0_\hbar$
		are the jet bundles of vector bundles on $\Orb_K$.
	\end{itemize}
	
\end{comment}


\begin{comment}
	Moreover, the flat connections $\theta_\P$ of $\P$ and $\theta_{\P_\J}$ of $\P_\J$ are compatible in the sense that we have the commutative diagram
	\begin{equation}\label{diag:PPJ}
		\begin{tikzcd}
			\EE_{\P_\J} \arrow[r, hook] \arrow[Isom, "\theta_{\P_\J}"']{d}  & \iota^*\EE_{\P}  \arrow[Isom, d, "\theta_{\P}"']   \\
			(\derdj)_{\P_\J} \arrow[r, hook] & \iota^*\derd_{\P}
		\end{tikzcd}	
	\end{equation}
	where $\derd_{\P}$ is the Lie algebra bundle over $X$ assoicated to the torsor $\P$ and the $\autd$-module $\derd$ and $(\derdj)_{\P_\J}$ is the Lie algebra bundle over $Y$ assoicated to the torsor $\P_\J$ and the $\autdj$-module $\derdj$. The inclusion $(\derdj)_{\P_\J} \hookrightarrow \iota^*(\derdj)_{\P_\J}$ is induced by the inclusion $\derdj \hookrightarrow \derd$ of $\autdj$-modules.
	
	Moreover, $\P_\J$ is a transitive torsor since the conormal bundle to $Y$ is naturally identified with its tangent bundle due to the Lagrangian condition (\cite[Lemma 4.3]{BC}).
\end{comment}


\subsubsection{Picard algebroids associated to Lagrangian subvarieties} \label{subsec:lag_pic}

By \cite[\S\,4.1]{LY}, one can associate two related Picard algebroids over a Lagrangian subvariety $Y$ to a given quantization $\OO_\hbar$. Namely, let $\J_Y$ denote the ideal subsheaf of $\OO_X$ consisting of all
sections vanishing on $Y$. Set $\J_{Y,\hbar}$ and $\J_{Y,\hbar}'$ to be the preimage of the ideal $\J_Y$ and $\J_Y^2$ respectively under the projection $\OO_\hbar \twoheadrightarrow \OO_X$. We write $\J_{Y,\hbar}^2 = \J_{Y,\hbar} \cdot \J_{Y,\hbar} \subset \OO_\hbar$, which is a subsheaf of two-sided ideals of the algebra $\OO_\hbar$. Then $\hbar^{-1} \J_{Y,\hbar}^2 \subset \hbar^{-1}\J_{Y,\hbar}' \subset \hbar^{-1}\J_{Y,\hbar}$ are all subsheaves of Lie ideals in the sheaf $ \hbar^{-1}\OO_\hbar$ of Lie algebras. Consider the sheaf $\hbar^{-1} (\J_{Y,\hbar}/\J_{Y,\hbar}^2)$, which is a sheaf of Lie algebras supported on $Y$. The left multiplication of $\OO_\hbar$ on $\hbar^{-1}\J_{Y,\hbar}$ descends to a left $\OO_Y$-module structure on $\hbar^{-1} (\J_{Y,\hbar}/ \J_{Y,\hbar}^2)$. We have a short exact sequence of sheaves of Lie algebras and $\OO_Y$-modules
\begin{equation}\label{exsq:J}
	0 \to \hbar^{-1} (\J_{Y,\hbar}' / \J_{Y,\hbar}^2) \to \hbar^{-1} (\J_{Y,\hbar}/ \J_{Y,\hbar}^2) \to \hbar^{-1} (\J_{Y,\hbar} / \J_{Y,\hbar}') \to 0.
\end{equation}
By \cite[Lemma 5.4]{BGKP}, there are canonical isomorphisms
\begin{equation} \label{eq:J_isom}
	\hbar^{-1} (\J_{Y,\hbar} / \J_{Y,\hbar}') \simeq \T_Y, \quad \hbar^{-1} (\J_{Y,\hbar}' / \J_{Y,\hbar}^2) \simeq \OO_Y.
\end{equation}
Then  \eqref{exsq:J} and \eqref{eq:J_isom} together make $\hbar^{-1} (\J_{Y,\hbar}/ \J_{Y,\hbar}^2)$ into a Picard algebroid over $Y$, which is denoted as $\tatp$.

The second Picard algebroid arises from a different $\OO_Y$-module structure on $\hbar^{-1} (\J_{Y,\hbar}/ \J_{Y,\hbar}^2)$ defined as follows. There is a commutative non-assoicative \emph{symmetrized product} on $\OO_\hbar$ defined by $a \bullet b := \frac{1}{2}(a b + b a)$. The symmetrized product induces a second $\OO_\hbar$-module structure, which is also compatible with the short exact sequence \eqref{exsq:J} and isomorphisms \eqref{eq:J_isom},  and hence we get a second Picard algebroid structure on $\hbar^{-1} (\J_{Y,\hbar}/ \J_{Y,\hbar}^2)$ with the same Lie bracet. We denote the second Picard algebroid as $\tat$. By \cite[Section 1]{BC} and \cite[Section 4.1]{LY},  the identity map of $\hbar^{-1} (\J_{Y,\hbar}/ \J_{Y,\hbar}^2)$ induces a canonical isomorphism
\begin{equation}\label{eq:tatp2tat}
	\tatp \xrightarrow{\sim} \tat + \frac{1}{2}\TT(K_Y).
\end{equation}
of Picard algebroids, where $K_Y$ denotes the canonical line bundle of $Y$.

By \cite[\S\,5.3]{BGKP}, the Picard algebroid $\tat$ satisfies
\begin{equation}  \label{eq:c1_lag}
	c^{\operatorname{top}}_1(\tat)=\iota^*[\omega_1(\OO_\hbar)],
\end{equation}
where $\omega_1(\OO_\hbar)$ is as in \eqref{eq:per} and $\iota^*: H^2_{DR}(X) \to H^2_{DR}(Y)$ is the restriction map induced by the embedding $\iota: Y \hookrightarrow X$.

\begin{Rem} \label{rmk:enhanced_period_lag}
	It is shown in \cite[Proposition 6.4]{Yu} that $\iota^\circ\left(\overline{\per}(\OO_\hbar)\right) = \tat$, where $\overline{\per}: \quan(X, \Omega) \to \pic(X)$ is the enhanced period map mentioned in Remark \ref{rmk:enhanced_period}. This provides a refinement of the equality \eqref{eq:c1_lag}. The homogeneous case, which is the main focus of the current paper, is treated in Lemma \ref{Lem:quantum_comom_period} and Corollary \ref{Cor:equiv_condition}.
\end{Rem}

\begin{comment}
	
	To this end, we first consider an alternative definition of $\tat$ in terms of jet bundles. Let $\der(\OO_\hbar)_{\J_{Y,\hbar}}$ denote the sheaf of $\OO_Y$-linear derivations of $\OO_\hbar$ which preserves $\J_{Y,\hbar}$. Let $\J_{Y,\hbar}^\wedge$ be the completion of $\J_{Y,\hbar} \subset \OO_\hbar$ with respect to the $\I_t$-adic topology of $ \OO_\hbar$, which is an ideal subsheaf of the algebra $\OO^{\wedge}_t$. The natural flat connection on $\OO_\hbar^\wedge$ preserves the idea subsheaf $\J_{Y,\hbar}^\wedge$ (unlike the case of $\I_t^\wedge$) and the subsheaf of flat sections of $\J_{Y,\hbar}^\wedge$ is naturally identified with $\J_{Y,\hbar}$. Therefore $\hbar^{-1}\J_{Y,\hbar}^\wedge$ is a Lie subalgebroid of $\OO_\hbar^\wedge$, with the anchor map given by the restriction of $\tilde{\eta}_t^\wedge: \OO_\hbar^\wedge \to \T_X$ to $\hbar^{-1}\J_{Y,\hbar}^\wedge$. Let $\der^{\wedge}(\OO_\hbar)_{\J_{Y,\hbar}}$ denote the Lie subalgebroid of $\der^{\wedge}(\OO_\hbar)$ of derivations which preserve $\J_{Y,\hbar}^\wedge \subset \OO^{\wedge}_t$, then it is the jet bundle of $\der(\OO_\hbar)_{\J_{Y,\hbar}}$. Define the subsheaves of the tangent sheaf $\T_X$ of $X$
	\[ \T_X (\J_Y) = \{ \delta \in \T_X \mid \delta (\J_Y) \subset \J_Y  \}, \quad  \T_X (\J_Y^2 ) = \{ \delta \in \T_X \mid \delta (\J_Y) \subset \J_Y^2  \} \]
	then $\T_X (\J_Y)$ is a Lie subalgeroid of $\T_X$ and $\T_X (\mathscr{I}^2_Y)$ is a Lie  ideal of $\T_X (\J_Y)$, such that $\T_X (\J_Y) / \T_X (\mathscr{I}^2_Y) \simeq \iota_* \T_Y$. We have
	\[  \der^{\wedge}(\OO_\hbar)_{\J_{Y,\hbar}} =  (\eta_t^\wedge)^{-1} ( \T_X (\J_Y)), \quad (\tilde{\eta}_t^\wedge)^{-1} ( \T_X (\J_Y)) = \hbar^{-1}\OO_X + \hbar^{-1} \J_{Y,\hbar}^\wedge, \]
	(not a direct sum if $X$ is disconnected!).  We see that the restriction of $\der^{\wedge}(\OO_\hbar)_{\J_{Y,\hbar}}$, resp. $\hbar^{-1} \J_{Y,\hbar}^\wedge$, to $Y$ is naturally equipped with the structure of a Lie algebroid over $Y$ induced by the natural isomorphism
	\begin{equation} \label{eq:pic_res}
		\der^{\wedge}(\OO_\hbar)_{\J_{Y,\hbar}} \big|_Y \simeq \iota^\circ \der^{\wedge}(\OO_\hbar), \quad \text{resp.} \quad \hbar^{-1} \J_{Y,\hbar}^\wedge \big|_Y \simeq \iota^\circ ( \hbar^{-1}\OO_\hbar^\wedge  / \hbar^{-1} \OO_X ),
	\end{equation}
	of $\OO_Y$-modules, together with the surjective anchor maps given by the restriction of $\eta_t^\wedge$, resp. $\tilde{\eta}_t^\wedge$, to $Y$.
\end{comment}





\subsubsection{Existence and Uniqueness} \label{subsec:lag_vec}
Let $\EE$ be an algebraic vector bundle  over $Y$ and $\EE_\hbar$ be a quantization of $\EE$. In particular, $\EE_\hbar$ is a module over $\OO_\hbar$ and we have $\J_{Y,\hbar} \cdot \EE_\hbar \subset \hbar \EE_\hbar$. Since $\EE_\hbar$ is flat over $\C[[\hbar]]$, we have a Lie algebra action of $\hbar^{-1} \J_{Y,\hbar}$ on $\EE_\hbar$ given by $(\hbar^{-1} a). m = \hbar^{-1}(a \cdot m)$. This action descends to a Lie algebroid action of $\tatp$ on $\EE_\hbar / \hbar \EE_\hbar \simeq \EE$, which makes $\EE$ into a $\tatp$-module.


\begin{Rem} \label{rmk:pic_mod_t2}
	Note that the $\tatp$-module structure on $\EE$ is completely determined by the $\OO_\hbar / \hbar^2 \OO_\hbar$-module structure on $\EE_\hbar / \hbar^2 \EE_\hbar$. That is, in the definition of $\tatp$ from Section \ref{subsec:lag_pic}, we can replace $\J_{Y,\hbar}$ by $\wt{\J}_{Y,\hbar} := \J_{Y,\hbar} / \hbar^2 \OO_\hbar \subset \OO_\hbar / \hbar^2 \OO_\hbar$ and get a natural identification $\tatp \simeq \hbar^{-1} \wt{\J}_{Y,\hbar} / \wt{\J}^2_{Y,\hbar}$ of Picard algebroids. Apply the same construction as above to the $\OO_\hbar / \hbar^2 \OO_\hbar$-module $\EE_\hbar / \hbar^2 \EE_\hbar$ gives the same $\tatp$-module structure on $\EE$.
\end{Rem}
%\footnote{I.L.: Does this remark have some purpose? It's not referred to below and doesn't serve an %"educational purpose" like the previous one. \\  S.Y.: When I define the affine analogues of these Picard algebroid in Section \ref{subsec:strong_slice}, I use only the mod $t^2$ information about the quantization. I've added Lemma \ref{lem:Picard_loc} there, which claims that the localization of Picard Lie algebras give the Picard algebroids, and the reference to this remark above the lemma.}

Recall that we can decompose $\per(\OO_\hbar)$ as in \eqref{eq:per}, where $ \omega_k(\OO_\hbar) \in H^2_{DR}(X)$. The following theorem (\cite[Thm\,1.1]{BC}) gives a necessary and sufficient condition for $\EE$ to admit a quantization. The case when $\EE$ is a line bundle was proven in \cite[Theorem 1.4]{BGKP}.

\begin{Thm}\label{thm:lag_vec}
  Fix a quantization $\OO_\hbar$ of $(X,\Omega)$. Let $Y$ be a smooth Lagrangian subvariety $Y$ of $X$ and $\EE$ be a coherent $\tatp$-module over $Y$. Then $\EE$ admits a quantization $\EE_\hbar$ that lifts the $\tatp$-module structure on $\EE$, if and only if $\iota^*[\omega_k(\OO_\hbar)] = 0$ holds in $H^2_{DR}(Y)$, $\forall ~ k \geqslant 2$.
\end{Thm}


\begin{comment}
	\begin{Thm}\label{thm:lag_vec}
		\begin{enumerate}
			\item
			A vector bundle $E$ of rank $r$ over a smooth Lagrangian subvariety $Y$ admits a quantization if and only if the following conditions hold:
			\begin{enumerate}
				\item
				$\iota^* \omega_i(\OO_\hbar) = 0$ holds in $H^2_{DR}(Y)$, $\forall ~ i \geqslant 2$;
				\item
				the projectivization $\mathbb{\EE}$ admits a flat algebraic connection;
				\item \label{cond:lag_vec3}
				$\frac{1}{r} c_1(E) = \frac{1}{2} c_1(K_Y) + \tat$ holds in $\pic(Y)$.
			\end{enumerate}	
			\item
			If nonempty, the set $\qyr$  admits a free action of the group $\F((\OO_Y[[t]])^\times)$. The set of orbits for this action can be identified with the set of all isomorphism classes of $\PGLh$-bundles with a flat algebraic connection. 	
		\end{enumerate}
	\end{Thm}
\end{comment}

Note that, the quantization $\EE_\hbar$ in Theorem \ref{thm:lag_vec}, if exists, is usually far from being unique. The set of isomorphism classes of all quantizations of vector bundles with a fixed rank admit a free action of an often infinite-dimensional group. We refer the reader to \cite[Thm\,1.1]{BC} for more details.

From now on, we will restrict ourselves to the graded setting. One advantage of working under the graded assumption, besides many others, is that we have the uniqueness of a quantization under a certain mild condition, which we recall now. First of all, note that when $\EE_\hbar$ is a graded quantization, the induced $\tatp$-module structure on $\EE$ is compatible with the $\cm$-actions, i.e., $\EE$ is a graded $\tatp$-module in the sense of Definition \ref{defn:pic_module_action} . Modulo the Lie ideal $\OO_X$ of $\tatp$, the $\tatp$-module structure induces an algebraic flat connection on the projectivization $\mathbb{P}(\EE)$ of $\EE$, via which the differentiation of the $\cm$-action on $Y$ lifts to a Lie algebra $\C$-action on $\mathbb{P}(\EE)$ (see the discussion before \cite[Definition 4.3.2]{LY}). On the other hand, the $\cm$-action on $\EE$ descends to a $\cm$-action on $\mathbb{P}(\EE)$, which differentiates to a second $\C$-action. The following notion was introduced in \cite[Definition 4.3.2]{LY}.

\begin{defi} \label{defn:strong_cm}
	Let $\TT$ be a graded Picard algebroid and $\EE$ be a graded $\TT$-module over $Y$ that is $\OO_Y$-coherent. The $\cm$-action on $\EE$ is said to be \emph{strong} if the Lie algebra $\C$-action on the projectivization $\mathbb{P}(\EE)$, defined via the flat connection that is induced from the $\TT$-module structure, coincides with the $\C$-action induced by the $\cm$-action on $\EE$.
\end{defi}

\begin{comment}
The definition of strongness is motivated by the following lemma (\cite[Lemma 5.4.1]{LY}, which is a direct corollary of Cartan's homotopy formula and plays a crucial role in the uniqueness result.
Recall from Section \ref{subsec:pic_defn} that,
 for any $\OO_Y$-coherent weak $(\TT,\cm)$-module $\EE$,  there is a projective integrable connection on $\EE$, which induces a flat connection $\nabla^0$ on $\send^0(\EE) := \send_{\OO_Y}(\EE)/\OO_Y \cdot \Id_\EE$\footnote{I suggest this becomes a lemma: the section 4.1.1 has a lot of stuff in it. S.Y.: Since most parts of this section, including Lemma \ref{lem:strong} and Proposition \ref{prop:irred_strongness}, are not be used below and will be removed, I will not add any details.}. Then $\nabla^0$ is equivariant with respect to the conjugate $\cm$-action on $\send^0(E)$ and hence defines a $\cm$-action on the de Rham cohomology groups $H^\bullet_{\nabla^0}(Y, \send^0(E))$.

\begin{Lem}\label{lem:strong}
	For a strong $\cm$-action on a $\OO_Y$-coherent  $\TT$-module $\EE$, the induced $\cm$-action on $H^\bullet_{\nabla^0}(Y, \send^0(\EE))$ is trivial.
\end{Lem}\footnote{I.L.: There are no references to this lemma. Is it used below? S.Y.: No.}
\end{comment}

The strongness condition is vacuous when $\EE$ is a line bundle. More generally, we have the following sufficient condition to ensure the strongness condition. 
%First recall that by Section \ref{subsec:RH}, the analytification $\EE^{an}$ of an $\OO_Y$-coherent $\TT$-module $\EE$ determines a projective representation of the fundamental group $\pi_1(Y,y)$ on the geometric fiber $\EE_y$ at a base point $y \in Y$ via the (twisted) Riemann-Hilbert correspondence (even when $\EE$ itself is not regular). %If $Y^{an}$ is connected, then the irreducibility $\EE_y$ as a projective representation is independent of the choice of $y$.

\begin{Prop}[{\cite[Prop. 6.11]{Yu}}] \label{prop:irred_strongness}
	Let $\TT$ be a graded Picard algebroid over $Y$ and $\EE$ be an $\OO_Y$-coherent graded $\TT$-module over $Y$, whose analytification $\EE^{an}$ is irreducible as a $\TT^{an}$-module. Then the $\cm$-action on $\EE$ is strong.
\end{Prop}

\begin{comment}
\begin{proof}
	Observe that the irreducibility of $\EE^{an}$ implies that $Y^{an}$ is connected. As in Section \ref{subsec:pic_defn}, the $\TT$-module structure determines a projective connection on $\EE$, which gives a flat connection on the projectivization $\mathbb{P}(\EE^{an})$. Its monodromy gives an irreducible projective representation $\rho_y$ of the fundamental group $\pi_1(Y,y)$ on the geometric fiber $\EE_y$ at a base point $y \in Y$, which is the same as the representation given in Section \ref{subsec:RH} via the (twisted) Riemann-Hilbert correspondence.
	
	We have an infinitesimal Lie algebra $\C$-action ta shuon $\mathbb{P}(\EE^{an})$ induced by the projective connection and the $\cm$-action on $Y$. We restrict it to an action of the Lie subalgebra $i \R \subset \C$, which is the Lie algebra of the unit circle group $S^1 \subset \C$. We first claim that this $i \R$-action in fact integrates to an $S^1$-action \footnote{$S^1$-action? S.Y.: Corrected.} on $\mathbb{P}(\EE^{an})$. Note that the $\cm$-action on $Y$ gives a canonical homomorphism from $\Z \simeq \pi_1(S^1) = \pi_1(\cm)$ to $\pi_1(Y,y)$. Pick the generator of the group $\pi_1(S^1)$ given by $\gamma(\theta) = e^{i \theta} \in S^1$ and let $[\gamma]$ denote its image in $\pi_1(Y,y)$, which is a canonically defined central element of $\pi_1(Y,y)$ and is independent of the choice of $y$ (cf. Section \ref{subsec:RH}\footnote{Here it would be good to have a more detailed reference. E.g. put a remark in Sec 4.1.4.}). Then $\rho_y([\gamma])$ is a projective automorphism of $\mathbb{P}(\EE_y)$ and commutes with the action of $\pi_1(Y,y)$, i.e., $\rho_y([\gamma])$ is an automophism of the irreducible projective representation $\rho_y$. By Schur's Lemma, $\rho_y([\gamma])$ acts as the identity (modulo a scalar multiple). This implies that the monodromy of the flat connection on $\mathbb{P}(\EE^{an})$ along the loop $\theta \mapsto \gamma(\theta) \cdot y$ emanating from any point $y \in Y$ induced by the action of $S^1 \subset \cm$ is identity. Therefore we obtain a well-defined $S^1$-action on $\mathbb{P}(\EE^{an})$, denoted as $\beta$. Its differential is the Lie algebra $i\mathbb{R}$-action we have started with. The flat connection on $\mathbb{P}(\EE^{an})$ is preserved by $\beta$.
	
	The $\cm$-action on $\EE$ descends to one on $\mathbb{P}(\EE^{an})$. We restrict it to an $S^1$-action on $\EE$, denoted as $\alpha$. Both $\alpha$ and $\beta$ lift the $S^1$-action on $Y$ and preserve the projective connection. It follows that, for any $z \in S^1$, $\alpha(z)^{-1} \beta(z)$ is an automorphism of the bundle $\mathbb{P}(\EE^{an})$ and induces an automorphism of $\rho_y$ associated to any point $y$, therefore it must be identity again by the irreducibility of $\rho_y$. This implies that the differentials of the actions $\alpha$ and $\beta$ coincide, therefore we have the coincidence of infinitesimal $\C$-actions after complexifications.
\end{proof}
\end{comment}

%\begin{Cor}\label{Cor:strongness_sufficient}
%	If a weak $(\TT, \cm)$-module $\EE$ has regular singularities and is irreducible\footnote{S.Y.: Is this corollary about reducible $\EE$?} as an %algebraic $\TT$-module, then the $\cm$-action on $\EE$ is automatically strong. In particular, the strongness holds when
%\begin{enumerate}
%\item $\pi_1(Y)$ is finite, and $\EE$ has regular singularities.
%\item
%or an algebraic group $K$ acts on $Y$ with finitely many orbits, $\TT$ is strongly $K$-equivariant and $\EE$ is $(K, \kappa)$-equivariant for some %$\kappa \in (\kf^*)^K$.
%\end{enumerate}
%\end{Cor}
%\begin{proof}
%Taking flat sections gives a category equivalence between
%\begin{itemize}
%\item the category of $(\TT,\cm)$-modules,
%\item and the category of graded finite dimensional projective representations of $\pi_1(Y)$
%\end{itemize}
%(with Schur multiplier recovered from $\TT$, see Section \ref{subsec:RH}\footnote{Again, a more detailed ref would be helpful. S.Y.: Section %\ref{subsec:RH} only deals with the case when the twist is given by a multiple of the twist of some line bundle. }). It follows that the
%$\cm$-action on the irreducible summands\footnote{S.Y.: $\EE$ is irreducible, no?} of $\EE$ is strong. Hence the action on $\EE$ is strong.
%
%Case (1) trivially follows. Case (2) follows because the conditions of (1) are satisfied\footnote{S.Y.: Why does (2) imply $\pi_1(Y)$ is finite?}
%(the regularity follows as mentioned in the end of Section \ref{subsec:RH}).
%\end{proof}
%\footnote{Some comments on why this holds, e.g. by quoting results from the above. Then some intro %sentence for the theorem.}

\begin{comment}
	\begin{Lem}
		Let $\TT$ be a weakly $\cm$-equivariant Picard algebroid over $Y$. Let $\EE$ be an irreducible object in the abelian category $\coh(\mathscr{D}_\TT)^{gr}$ be a the category of $\OO_Y$-coherent weak $(\TT,\cm)$-modules over $Y$.Then the $\cm$-action on $\EE$ is automatically strong.
	\end{Lem}
\end{comment}

Now we get to the main result of this section.

\begin{Thm}[{\cite[Thm 4.3.3]{LY}}] \label{thm:quan_lag}
	Let $\EE$ be an $\OO_Y$-coherent graded $\tatp$-module over the Lagrangian subvariety $Y$.  Assume in addition that the $\cm$-action on $\EE$ is strong. Then a graded quantization $\EE_\hbar$ of $\EE$, if exists, is unique up to isomorphism, which gives rise to the given $\tatp$-module structure and the $\cm$-equivariant structure on $\EE$.
\end{Thm}

\begin{comment}
	\begin{Rem}
		The strongness condition is assumed in order to guarantee the uniqueness of the graded quantization (see \cite[Lemma 5.4.1]{LY} and \cite[Proposition 5.4.2]{LY}). This condition is not always satisfied. It can be shown that, if $E$ is irreducible as an analytic twisted local system over $Y$, any $\cm$-action on $E$ compatible with that on $Y$ is automatically strong. However, one can easily make an example not satisfying the strongness condition by taking two copies of such $E$ and twist the $\cm$-action on one of it by any nontrivial character of $\cm$ and then take the direct sum.
		
		However, in the $K$-homogenous case and other cases treated in this paper, the strongness condition is always satisfied. See Remark \ref{rem:strong_equiv}.
	\end{Rem}
	
	We have seen that $\quan(X)$ is a torsor over $H^2(\Omega_X^{\geqslant 1}) \cong \pic(X)$. On the other hand, the restriction map $\pic(X) \to \pic(Y)$ is a group homomorphism. Regard $\pic(Y)$ as a torsor over itself, then we have the following result. {\bf Note sure if we still need this.}
	
	
	\begin{Prop}\label{prop:atiyah}
		Under the isomorphism $\overline{\per}: \quan(X) \xrightarrow{\sim} \pic(X)$ in Proposition \ref{prop:period}, $\TT(\bullet,Y)$ can be identified with the restriction map $\iota^\circ: \pic(X) \to \pic(Y)$. In particular, the map $\TT(\bullet,Y): \quan(X) \to \pic(Y)$ respects the torsor structures, i.e.,
		\begin{equation}
			\TT([\OO_\hbar] + [\TT], Y) = \TT(\OO_\hbar,Y) + \iota^\circ(\TT),  	
		\end{equation}
		for any $\OO_\hbar \in \quan$ and $\TT \in \pic(X)$. The same is true for $\TT^+(\bullet,Y)$.
	\end{Prop}
	
	\begin{proof}
		The first claim is Proposition \ref{prop:picY=per}. The remaining statements follow from Proposition \ref{prop:period}.
	\end{proof}
	
\end{comment}



\subsection{Hamiltonian quantization}\label{SS_Ham_quant}
Let $G$ be an affine algebraic group. Suppose $G$ acts on the graded symplectic variety $X$ commuting with the $\cm$-action and preserving the symplectic form $\Omega$. Let $\alpha_G: G \times X \to X, (g, x) \mapsto g \cdot x$, denote the action map. Roughly speaking, a $G$-equivariant graded quantization of $(X, \Omega)$ is a graded quantization $\OO_\hbar$ of $(X, \Omega)$ equipped with a $G$-action by automorphisms of sheaves of $\C[[\hbar]]$-algebras that lifts the $G$-action on $X$, such that
\begin{enumerate}
	\item
	it commutes with the $\cm$-action on $\OO_\hbar$;
	\item
	the quotient homomorphism $\OO_\hbar \twoheadrightarrow \OO_\hbar / \hbar \OO_\hbar \simeq \OO_X$ is $G$-equivariant.
\end{enumerate}
In more detail, define the maps
\begin{align*}
	\pr_G: G \times X \to G, & \quad (g,x) \mapsto g \\
	\pr_X: G \times X \to X, & \quad (g,x) \mapsto x \\
	i: X \to G \times X, & \quad x \mapsto (e, x)
\end{align*}
Set
\begin{equation} \label{eq:quan_pullback}
	\alpha_G^* \OO_\hbar := \pr_G^{-1}  \OO_G[[\hbar]]  \widehat{\otimes}_{\C[[\hbar]]} \alpha_G^{-1} \OO_\hbar = \varprojlim_{k \to \infty} \pr_G^{-1}  \OO_G  \otimes_\C  \alpha_G^{-1} (\OO_\hbar / \hbar^k \OO_\hbar)
\end{equation}
and define $\pr_X^* \OO_\hbar$ similarly, so that both are sheaves of $\pr_G^{-1}  \OO_G[[\hbar]]$-algebras over $G \times X$. Here $\widehat{\otimes}$ denotes the completed tensor product. We can regard $G \times X$ as a variety over $G$ equipped with the relative symplectic form $\alpha_G^* \Omega = \pr_X^* \Omega$ and an action of $\cm$ on the second factor. The $G$-action on $X$ induces a $\cm$-equivariant isomorphism $\beta_0^G: \alpha_G^* \OO_X  \xrightarrow{\sim} \pr_X^* \OO_X$ of sheaves of Poisson algebras, satisfying the usual compatibility conditions. Then both $\alpha_G^*\OO_\hbar$ and $\pr_X^*\OO_\hbar$ can be regarded as graded quantizations of $G \times X$ over $G$. A {\it $G$-equivariant structure} on a graded quantization $\OO_\hbar$ consists of an isomorphism
\[ \beta_\hbar^G: \alpha_G^* \OO_\hbar \xrightarrow{\sim} \pr_X^* \OO_\hbar \]
of sheaves of $\pr_G^{-1}  \OO_G[[\hbar]]$-algebras, such that
\begin{itemize}
\item
it satisfies the usual associativity condition and  $i^* \beta_\hbar^G = \Id_{\OO_\hbar}$, \item $\beta_\hbar^G$ is $\cm$-equivariant (this is a more precise version of condition (1) above),
\item and $\beta_\hbar^G \mod \hbar$ equals $\beta_0^G: \alpha_G^* \OO_X  \xrightarrow{\sim} \pr_X^* \OO_X$ (this is a more precise version of condition (2)).
\end{itemize}

Given a $G$-equivariant graded quantization $(\OO_\hbar, \beta_\hbar^G)$, we can assign an action of the Lie algebra $\g$ on $\OO_\hbar$, i.e., a Lie algebra homomorphism $L_\hbar^\g : \g \to \der(\OO_\hbar)$, where $\der(\OO_\hbar) = \der_{\C[[\hbar]]}(\OO_\hbar)$ denotes the sheaf of $\C[[\hbar]]$-linear derivations of $\OO_\hbar$ (note that such derivations  are automatically continuous with respect to the $\hbar$-adic topology), as follows. Let $L_G : \g \to \Gamma(G, \T_G)$ denote the Lie algebra homomorphism generated by the left action of $G$ on itself. Then $L_G$ identifies $\g$ with the space of right invariant vector fields on $G$. Define
\begin{equation} \label{eq:inf_action_OO_h}
	L_\hbar^\g (\xi) s = i^* \langle (L_G (\xi) \boxtimes \Id) [\beta_\hbar^G (\alpha_G^*(s))] \rangle \in \send_{\C[[\hbar]]} (\OO_\hbar)
\end{equation}
for any $\xi \in \g$ and local section $s \in \OO_\hbar$, where $\alpha_G^*(s)$ stands for the local section $1 \otimes \alpha_G^{-1}(s)$ of the sheaf $ \pr_G^{-1}  \OO_G[[\hbar]]  \widehat{\otimes}_{\C[[\hbar]]} \alpha_G^{-1} \OO_\hbar = \alpha_G^* \OO_\hbar$. In other words, for any open subset $U \subset X$, $L_\hbar^\g(\xi)$ is the composite map
\begin{multline*}
	L_\hbar^\g (\xi) : \Gamma(U, \OO_\hbar) \xrightarrow{\alpha_G^*} \Gamma(\alpha_G^{-1}(U), \alpha_G^* \OO_\hbar) \xrightarrow{\beta_\hbar^G} \ \Gamma(\alpha_G^{-1}(U), \OO_G \widehat{\boxtimes} \OO_\hbar) \\ \xrightarrow{L_G(\xi) \boxtimes \Id} \Gamma(\alpha_G^{-1}(U), \OO_G \widehat{\boxtimes} \OO_\hbar) \xrightarrow{i^*} \Gamma(U, \OO_\hbar).
\end{multline*}
It is standard to check that $L_\hbar^\g (\xi) s$ lies in $\der(\OO_\hbar)$.

\begin{defi}\label{defi:Hamilt_action_quantization}
	Let $\OO_\hbar$ be a graded quantization of a graded symplectic variety of weight $d$. A \emph{(graded) quantum Hamiltonian $G$-action} on $\OO_\hbar$ consists of
	\begin{itemize}
		\item	
		a $G$-equivariant structure $\beta_\hbar^G$ on $\OO_\hbar$ in the above sense;
		\item
		a $G$-equivariant $\C$-linear map $\Phi_\hbar: \g \to \Gamma(X, \OO_\hbar)$	whose image lies in the degree $d$ part of $ \Gamma(X, \OO_\hbar)$,
	\end{itemize}
	such that the adjoint action of $\hbar^{-1} \Phi_\hbar (\xi)$ on $\OO_\hbar$, $u \mapsto \hbar^{-1} [\Phi_\hbar(\xi), u ]$, for any $\xi \in \g$, agrees with the infinitesimal action $L_\hbar^\g(\xi)$ on $\OO_\hbar$ induced by the $G$-equivariant structure $\beta_\hbar^G$ as above. The map $\Phi_\hbar$ is called a {\it quantum comoment map}.
\end{defi}

Note that the definition above implies that $ [\Phi_\hbar(\xi), \Phi_\hbar(\zeta)] = \hbar \Phi_\hbar ([\xi, \zeta]_\g)$, $\forall \, \xi, \zeta \in \g$. In other words, the composite map $\Phi_\hbar: \g \to \Gamma(X, \OO_\hbar) \xrightarrow{\hbar^{-1} \cdot} \hbar^{-1} \Gamma(X, \OO_\hbar)$ is a Lie algebra homomorphism, where the Lie bracket of $\hbar^{-1} \Gamma(X, \OO_\hbar)$ is given by taking commutators. Therefore $\Phi_\hbar$ induces a $\cm$-equivariant algebra homomorphism $\Phi_\hbar: \U_\hbar \g \to \Gamma(X, \OO_\hbar)$, where $\U_\hbar \g$ is the Rees algebra of  $U(\g)$
with respect to the PBW filtration.

For any Hamiltonian $G$-action $(\beta_\hbar^G, \Phi_\hbar)$ on a graded quantization $\OO_\hbar$, set $\phi$ to be the composition $\g \xrightarrow{\Phi_\hbar} \Gamma(X, \OO_\hbar) \to \Gamma(X, \OO_\hbar / \hbar \OO_\hbar) \simeq \C[X] $, then $\phi$ is dual to a classical moment map $X \to \g^*$, which is $G \times \cm$-equivariant when $\g^*$ is defined to be of weight $d$ and is equipped with the coadjoint action of $G$. We say that $\Phi_\hbar$ lifts $\phi$ and $(\OO_\hbar, \beta_\hbar^G, \Phi_\hbar)$ is a Hamiltonian quantization of the Hamiltonian $G$-space $(X, \Omega, \phi)$.

Now suppose $K$ is an algebraic subgroup of $G$ with Lie algebra $\kf$. Suppose further that  $Y$ is a closed Lagrangian subvariety in $X$ stable under the $K \times \cm$-action. In addition, we assume that the image of $\mathfrak{k}$ under $\phi$ lies in the ideal subsheaf $\J_Y \subset \OO_X$ of $Y$. Then for any $G$-Hamiltonian quantization $(\OO_\hbar,\Phi_\hbar)$ of $(X, \Omega)$, the restriction of the $G$-action to $K$ preserves the ideal subsheaf $\J_{Y,\hbar} \subset \OO_\hbar$ defined in Section \ref{subsec:lag_pic} and hence induces a weak $K$-action on the sheaf $\hbar^{-1}(\J_{Y,\hbar} / \J_{Y,\hbar}^2)$ over $Y$, which is compatible with both the two $\OO_Y$-module structures. Moreover, the restriction $\Phi_\hbar|_{\kf} : \kf \to \Gamma(X, \OO_\hbar)$ of the quantum comoment map has its images lying in $\Gamma(X, \J_{Y,\hbar})$. The composition of $\Phi_\hbar|_{\kf}$ with the quotient map $\J_{Y,\hbar} \twoheadrightarrow \J_{Y,\hbar} / \J_{Y,\hbar}^2$ induces strong $K$-actions on both $\tat$ and $\tatp$ which lift the weak $K$-actions and hence makes them into strongly $K$-equivariant graded Picard algebroids on $Y$. %The isomorphism \eqref{eq:tatp2tat} then also respects the strong $K$-actions, where the strong $K$-action on $\frac{1}{2} \TT(K_Y)$ is induced by the canonical $K$-action on the canonical bundle $K_Y$.
%We still denote it by  $\Phi_\hbar: \g \to \Gamma(Y, \TT^+(\OO_\hbar,Y))$


%Now suppose $E$ is a $K \times \C$-equivariant vector bundle over $Y$.
The following definition is the sheaf theoretic counterpart to Definition \ref{eq:HC_equiv}.

\begin{defi}\label{defi:equiv_quant_modules}
	With the assumptions and notations above, suppose $\EE$ is a $\cm$-equivariant vector bundle over $Y$. A {\it $(K,\kappa)$-equivariant graded Hamiltonian quantization} of $(Y,\EE)$ is a graded quantization $\EE_\hbar$ of $\EE$ equipped with a $\C[[\hbar]]$-linear $K$-action that lifts the $K$-action on $Y$, such that
	\begin{enumerate}
		\item the $K$-action commutes with the $\cm$-action on $\EE_\hbar$;
		\item the structure morphism $\OO_\hbar \otimes \EE_\hbar \to \EE_\hbar$ is $K$-equivariant;
		\item the differential of the $K$-action on $\EE_\hbar$ is given by
		\[ L_\hbar^\kf(\xi) m = \hbar^{-1} \Phi_\hbar(\xi) \cdot m - \langle\kappa,x\rangle m,\]
		for any $\xi \in \kf$ and $m \in \EE_\hbar$.
	\end{enumerate}
\end{defi}

Again we can unpackage the definition of a $K$-action on $\EE_\hbar$ as in the case of $G$-action on $\OO_\hbar$ as follows. Let $\alpha_K, \, \wt{\pr}_X : K \times X \to X$ denote the restriction of the morphisms $\alpha_G, \, \pr_X : G \times X \to X$ to $K \times X$. Then $\alpha_K^{-1}(Y) = \wt{\pr}_X^{-1}(Y) = K \times Y$ can be regarded as a closed relative Lagrangian subvariety of the relative symplectic variety $K \times X$ over $K$. The sheaves $\alpha_K^* \EE_\hbar$ and $\wt{\pr}_X^* \EE_\hbar$, defined similarly as in \eqref{eq:quan_pullback}, are naturally modules over the sheaves of algebras $\alpha_K^* \OO_\hbar$ and $\wt{\pr}_X^* \OO_\hbar$ respectively, therefore are quantizations of $\alpha_K^* \EE$ and $\wt{\pr}_X^* \EE$ respectively, thanks to the isomoprhism $\beta_\hbar^G |_{K \times X}: \alpha_K^* \OO_\hbar \xrightarrow{\sim} \wt{\pr}_X^* \OO_\hbar$. A $K$-equivariant structure on $\EE_\hbar$ then consists of a $\cm$-equvariant isomorphism $\beta^K_\hbar: \alpha_K^* \EE_\hbar \xrightarrow{\sim} \wt{\pr}_X^* \EE_\hbar$ of sheaves, such that it preserves the module structures and is compatible with $\beta_\hbar^G |_{K \times X}: \alpha_K^* \OO_\hbar \xrightarrow{\sim} \wt{\pr}_X^* \OO_\hbar$. We can define the differentiation of the $K$-action on $\EE_\hbar$ in a similar way as in \eqref{eq:inf_action_OO_h} to get a map $L_\hbar^\kf: \kf \to \send_{\C[[\hbar]]}(\EE_\hbar)$, such that
\[ L_\hbar^\kf (\xi) ( s \cdot m) = L_\hbar^\g(\xi) (s) \cdot m + s \cdot L_\hbar^\kf(\xi) ( m ), \quad \forall \, s \in \OO_\hbar, \, m \in \EE_\hbar. \]

It is clear that, if $\EE_\hbar$ is a $(K,\kappa)$-equivariant graded quantization of $\EE$, then such a structure induces an algebraic $K$-action on $\EE$ as an $\OO_X$-module, such that it commutes with the $\cm$-action on $\EE$. Moreover, the $K$-action on $\EE$ and the strong $K$-action on $\tatp$ mentioned above make $\EE$ into a $(K, \kappa)$-equivariant $\tatp$-module.



\section{W-algebras and restriction functors}
W-algebras are quantizations of Slodowy slices in $\g^*$. They were introduced
in \cite{Premet_slice}, alternative definitions and characterizations were given
in \cite{Wquant},\cite{HC} and other papers by the first named author.
In \cite[Section 6.1]{Wdim} the first named author constructed a restriction functor
for HC $(\U,\theta)$-modules that is going to play a central role in this paper.

The goal of this section is to recall these constructions. We will also relate the
restriction functor for Harish-Chandra modules with some constructions from
Section \ref{S_qclass_quant}. In particular, we will see that one can use the restriction
functor to recover some results from \cite{LY}.

Throughout the section $G$ is a semisimple algebraic group, $\g$ is its Lie algebra, and
$\sigma$ is an involution of $\g$. Let $\theta:=-\sigma$ be the corresponding anti-involution
of $\g$.

\subsection{Slodowy slices}\label{SS_Sl_sl}
Fix a nilpotent element $\chi\in \g^*$ with $\theta(\chi)=\chi$.
Let $\Orb$ denote its $G$-orbit.

For a $G$- and $\sigma$-invariant orthogonal form
$(\cdot,\cdot)$ on $\g$, we get an identification  $\g\cong \g^*$.
Let $e\in \g$ denote the image of $\chi$ under this identification.
We include $e$ into an $\slf_2$-triple $(e,h,f)$ such that
$\sigma(h)=h, \sigma(f)=-f$. Then we can consider the
Slodowy slice $e+\mathfrak{z}_{\g}(f)$. Let $S$ be its image in
$\g^*$, also to be referred to as the Slodowy slice. Note that
$S$ is a transverse slice to $\Orb$. Note also that it is $\theta$-stable.

Let us describe two reductive group actions on $S$.
Let $Q:=Z_G(e,h,f)$.
This is a maximal reductive subgroup of the stabilizer $G_\chi$. Note that
$Q$ acts on $S$. Also we have a $\C^\times$-action on $S$. Namely, let
$z\mapsto z^h$ be the one-parameter subgroup of $G$ corresponding to
the semisimple element $h$. Then $z \mapsto z^{-2} z^h$ stabilizes $\chi$
and preserves $S$. Moreover, the induced grading on $\C[S]$ is positive.
We note that $\C^\times$ commutes with $\sigma$, while $\sigma$ naturally
acts on $Q$ and for $s\in S, q\in Q$, we have $\sigma(qs)=\sigma(q)\sigma(s)$.

The algebra $\C[S]$ is Poisson with bracket of degree $-2$, see, e.g., \cite[Section 2]{GG}. The action of
$Q$ on $S$ is Hamiltonian.

\subsection{W-algebra}\label{SS_W_alg} The algebra $\C[S]$ admits a quantization, called a
finite W-algebra. We recall a construction following \cite[Sections 2.1,2.2]{W_prim}.

Recall that we write $\U$ for the universal enveloping algebra of $\g$.
Form the Rees algebra $\U_\hbar$ of $\U$. The element $\chi$ defines a
homomorphism $\C[\g^*]\rightarrow \C$. Consider its composition
with $\U_\hbar\twoheadrightarrow \C[\g^*]$. We consider the completion
$\U_\hbar^{\wedge_\chi}$ with respect to the kernel of this composition. The completion comes with an action of
\begin{equation}\label{eq:group_acting} \C^\times \times \left(Q\rtimes (\Z/2\Z)\right)
\end{equation}
Here the action of $\Z/2\Z$ is induced by
$\theta$ on $\g$ and $\hbar\mapsto\hbar$, this is a $\C[\hbar]$-linear
anti-involution. The group $Q$ acts by $\C[\hbar]$-linear automorphisms and
this action is Hamiltonian: the quantum comoment map, to be denoted by $\Phi_\U$, comes from $\mathfrak{q}\hookrightarrow
\U_\hbar$. The action of $\C^\times$ is given by $z.x=z^{-2}z^hx$, and it rescales $\hbar$ via $z.\hbar := z^2\hbar$.

Set $V:=T_\chi\Orb$. This is  a symplectic vector space.
The space again comes with a natural action of
(\ref{eq:group_acting}). Let $\omega$ denote the symplectic form on
$V^*$. We can form the homogeneous Weyl algebra
$$\Weyl_\hbar:=T(V^*)[\hbar]/(u\otimes v-v\otimes u-\hbar\omega(u,v)).$$
The group (\ref{eq:group_acting}) acts on $\Weyl_\hbar$:
\begin{itemize}
\item $Q$ acts by
$\C[\hbar]$-linear automorphisms coming from the $Q$-action on $V$,
\item the generator of $\Z/2\Z$ acts by a
$\C[\hbar]$-linear antiinvolution induced by its action on $V$,
\item and the $\C^\times$-action is induced from the action of $\C^\times$ on $\g$
that fixes $\chi$ (introduced in Section \ref{SS_Sl_sl}).
\end{itemize}
Note that the action of $Q$ is Hamiltonian, let us
write $\Phi_{\Weyl}$ for the corresponding quantum comoment map
$\q\rightarrow \Weyl_\hbar$.

Consider the completion $\Weyl_\hbar^{\wedge_0}$. There is an embedding
$\Weyl_\hbar^{\wedge_0}\hookrightarrow \U_\hbar^{\wedge_\chi}$ that is
equivariant with respect to (\ref{eq:group_acting}) and lifts the decomposition
$T_\chi\g^*\cong T_\chi S\oplus V$. Any two such embeddings are conjugate with
an appropriate automorphism of $\U^{\wedge_\chi}_\hbar$ that commutes with
(\ref{eq:group_acting}), see \cite[Section 2.1]{W_prim}.

We write $\Walg_\hbar'$ for the centralizer of $\Weyl_\hbar^{\wedge_0}$
in $\U_\hbar^{\wedge_\chi}$. This is a $\C[[\hbar]]$-algebra with an action of
(\ref{eq:group_acting}). We have the natural isomorphism, where $\widehat{\otimes}$
stands for the completed tensor product:
\begin{equation}\label{eq:decomp_U}
\Weyl_\hbar^{\wedge_0}\widehat{\otimes}_{\C[[\hbar]]}\Walg_\hbar'
\xrightarrow{\sim} \U_\hbar^{\wedge_\chi}.
\end{equation}
We note that $\Phi_\U(\xi)-\Phi_\Weyl(\xi)\in \Walg_\hbar'$ for all
$\xi\in \q$. Set $\Phi_\Walg:=\Phi_\U-\Phi_{\Weyl}$, this is
a quantum comoment map for the action of $Q$ on $\Walg'_\hbar$.


Let $\Walg_\hbar$ stand for the $\C^\times$-finite part of $\Walg_\hbar'$.
Since the grading on $\C[S]$ coming from the $\C^\times$-action is positive,
we see that $\Walg_\hbar'$ coincides with the completion of $\Walg_\hbar$
at $\chi$. Set $\Walg:=\Walg_\hbar/(\hbar-1)$. This is the finite W-algebra.
It comes with a filtration and also with a Hamiltonian action of $Q\rtimes (\Z/2\Z)$
preserving the filtration. It also comes with an automorphism $\zeta$ of order $2$,
it is induced by the action of $-1\in \C^\times$ on $\Walg_\hbar$.

\subsection{Restriction of Dixmier algebras}\label{SS_Dixmier_restriction}
The construction above can be adapted to construct the restriction functor between
categories of Harish-Chandra bimodules, \cite[Section 3]{HC}. We will need it in the case of Dixmier algebras. Recall that by a Dixmier algebra one means an associative algebra $\A$ that comes equipped
with a rational $G$-action by algebra automorphisms and satisfying the following two conditions
\begin{itemize}
\item
this action is Hamiltonian
\item
and that the quantum comoment map $\U\rightarrow \A$ turns $\A$ into a HC bimodule.
\end{itemize}
For example,
any filtered quantization of $\C[\Orb]$ is a Dixmier algebra, see the discussion in the end of
Section \ref{SS_quant_nilp}.

Let $\A$ be a Dixmier algebra.

\begin{Lem}\label{Lem:Dixmier_good_filtration}
There is a good algebra filtration on $\A$.
\end{Lem}

Note that this is claimed in the proof of \cite[Lemma 3.4.5]{HC}, but the proof is wrong.
We use this opportunity to correct the proof.

\begin{proof}
Take a $G$-stable subspace $V\subset \A$ that generates $\A$ as a left $\U$-module and contains
$1$. We can find a positive integer $k$ such that $V^2\subset \U_{\leqslant k}V$.
Set $\A_{\leqslant i}:=\U_{\leqslant i}1$ for $i<k$, and $\A_{\leqslant i}:=\U_{\leqslant i}1+\U_{\leqslant i-k}V$ otherwise. This filtration is $G$-stable, it is a left $\U$-module filtration, and it is good.
It remains to show that it is an algebra filtration. The inclusion $\A_{\leqslant i}\A_{\leqslant j}
\subset \A_{\leqslant i+j}$ is checked as follows:
\begin{itemize}
\item for $i<k$, this follows from the claim that we have a left $\U$-module filtration,
\item for $j<k$, this follows from the claim that we have a right $\U$-module filtration,
\item for $i,j\geqslant k$, we also need to use that $\U_{\leqslant i-k}V\U_{\leqslant j-k}V
\subset \U_{\leqslant i+j-2k}V^2\subset \U_{\leqslant i+j-k}V\subset \A_{\leqslant i+j}$.
\end{itemize}
\end{proof}

%We can pick a good algebra filtration on $\A$,
%see, for example, the proof of \cite[Proposition 3.4.5]{HC}, compatible with the PBW
%filtration on $\U$.


Let $\A_\hbar$ denote the Rees algebra of $\A$. We get
a graded algebra homomorphism $\U_\hbar\rightarrow \A_\hbar$.  Set
$$\A_{\hbar}^{\wedge_\chi}:=\U_\hbar^{\wedge_\chi}\otimes_{\U_\hbar}\A_\hbar.$$
This $\C[[\hbar]]$-module coincides with the completion of $\A_\hbar$ with respect to a suitable
topology, where the product is continuous. So it has a natural algebra structure making the map
$\U_\hbar^{\wedge_\chi}\rightarrow \A_\hbar^{\wedge_\chi}$ an algebra
homomorphism.

Let $\underline{\A}_\hbar'$ denote the centralizer of $\Weyl_\hbar^{\wedge_0}$
in $\A_\hbar^{\wedge_\chi}$, this is a subalgebra.  By \cite[Proposition 3.3.1]{HC},
we have
\begin{equation}\label{eq:decomp_A}
\Weyl_\hbar^{\wedge_0}\widehat{\otimes}_{\C[[\hbar]]}\underline{\A}_\hbar'
\xrightarrow{\sim} \A_\hbar^{\wedge_\chi}.
\end{equation}
Then we can pass from $\underline{\A}_\hbar'$ to $\underline{\A}_\hbar$
(the $\C^\times$-locally finite part). We set $\underline{\A}:=\underline{\A}_\hbar/(\hbar-1)\underline{\A}_\hbar$.
We call $\underline{\A}$ the {\it restriction} of $\A$ to $S$ -- this is motivated
by $\gr\underline{\A}$ being the pullback of the $\C[\g^*]$-algebra $\gr\A$
to $S\subset \g^*$.

Note that the action of $Q$ on $\underline{\A}_\hbar$ is Hamiltonian: if $\Phi_{\A}$
denotes the quantum comoment map for $\A_\hbar$, then
\begin{equation}\label{eq:q_comom_A}
\Phi_{\underline{\A}}:=\Phi_{\A}-\Phi_{\Weyl}
\end{equation}
is a quantum comoment map for $\underline{\A}_\hbar$.

Now suppose that $\A$ comes with a filtration preserving anti-involution
$\theta$ that is intertwined with the anti-involution $\theta$ on $\U$
by the quantum comoment map $\U\rightarrow\A$.
Then similarly to Section \ref{SS_W_alg}, $\underline{\A}$ inherits an
anti-involution, again denoted by $\theta$.

Again, as in the case of W-algebras, $\underline{\A}$ comes with an order two
automorphism $\zeta$.

\subsection{Restriction of HC modules}\label{SS:HC_mod_restr}
Now we recall a construction of the restriction functor for HC modules.
The construction was given in \cite[Section 6.1]{Wdim} for the case of
$\A=\U$. It extends to the general setting, where $\A$ is a Dixmier algebra, in a straightforward way.
For simplicity, assume that $K$ is connected. Let $\A$ be as in the previous section.
Assume that $\gr\A$ is commutative, so
it makes sense to speak about HC $(\A,\theta)$-modules (with $d=1$).

We consider the category $\wHC^{gr}(\A,\theta)^{K,\kappa}$. A functor to be constructed depends on a choice of an irreducible component in $\Orb\cap (\g/\kf)^*$.
The intersection  $\Orb\cap (\g/\kf)^*$ is a smooth Lagrangian subvariety in
the symplectic variety $\Orb$. Every irreducible component of this subvariety is
a single $K$-orbit. Pick an irreducible component $\Orb_K$. Let $\chi$ be a point of
this orbit.

We will need to rescale the gradings.
We pick an even number $d$ such that $\frac{1}{2}dh\in \kf$ gives rise to
a one-parameter subgroup $\gamma:\C^\times\rightarrow K$. For example, for $K\subset G$,
we can take $d=2$. Then we have an action of $\C^\times$ on any $M\in \wHC^{gr}(\A_\hbar,\theta)^{K,\kappa}$
given by $t.m=t^{di}\gamma(t)^{-1}m$ for $m\in M_i$. We also consider the similar $\C^\times$-actions
on $\U_\hbar,\A_\hbar,\Walg_\hbar$, and so on. Note that $\hbar$ has degree $d$. Take a primitive root of unity $\epsilon$ of order $d$. The action of $\epsilon$ on $\U_\hbar,\A_\hbar,\Walg_\hbar, M_\hbar$
is an order $d$ automorphism to be denoted by $\zeta$.

Consider the category
$\wHC^{gr}(\A_\hbar,\theta)^{K,\kappa}$, where now the grading on $\A_\hbar$ is rescaled.
 We are going to define a functor
$$\bullet_{\dagger,\chi}: \wHC^{gr}(\A_\hbar,\theta)^{K,\kappa}\rightarrow \wHC^{gr}(\underline{\A}_\hbar,
\theta).$$
Here
%$K_Q$ is the preimage of $Q$ in $K$, and
$\underline{\A}$ is the restriction of $\A$ to $S$ defined in the previous section.

Pick $M_\hbar\in \wHC^{gr}(\A_\hbar,\theta)^{K,\kappa}$.
Consider its completion
$$M_\hbar^{\wedge_\chi}:=\A_\hbar^{\wedge_\chi}\otimes_{\A_\hbar}M_\hbar.$$
This is an object $\wHC^{gr}(\A_\hbar^{\wedge_\chi},\theta)$, where now the $\C^\times$-action
comes from $t\mapsto t^{-d}\gamma(t)$.
We are going to produce an equivalence $$\wHC^{gr}(\A_\hbar^{\wedge_\chi},\theta)
\xrightarrow{\sim} \wHC^{gr}(\underline{\A}_\hbar^{\wedge_\chi},\theta).$$

For this consider the Lagrangian subspace $L\subset V^*$ given by $L:=(\kf.\chi)^*$
(where $L$ is embedded into $V^*$ via the decomposition $V=\kf.\chi\oplus \kf^\perp.\chi$).
Note that $L\subset \left(\U_{\hbar}^{\wedge_\chi}\right)^{-\theta}$, where $V^*$
is embedded into $\U_{\hbar}^{\wedge_\chi}$ via (\ref{eq:decomp_U}).
We write $M_\hbar^L$ for the annihilator of $L$ under the Lie algebra action.

\begin{Lem}\label{Lem:wHC_equivalence}
The assignment $M_\hbar\mapsto M_\hbar^L$ defines a category equivalence
$$\wHC^{gr}(\A_\hbar^{\wedge_\chi},\theta)
\rightarrow \wHC(\underline{\A}_\hbar^{\wedge_\chi},\theta).$$
A quasi-inverse functor is given by $N_\hbar\mapsto \C[[L,\hbar]]\widehat{\otimes}_{\C[[\hbar]]} N_\hbar$.
\end{Lem}
\begin{proof}
Recall the decomposition (\ref{eq:decomp_A}). Under this decomposition $\theta$ is the tensor product
of two involutions: the restriction of $\theta$ to $\Walg_\hbar'$ and the anti-involution on
$\Weyl_\hbar^{\wedge_0}$ that is $-1$ on $L$ and $1$ on its Lagrangian complement.
The claim is now standard, compare to \cite[Proposition 3.3.1]{HC} and Lemma \ref{Lem:wHC_formal_Weyl}.
\end{proof}

Taking the locally finite elements for $\C^\times$ gives rise to an equivalence
$$\wHC^{gr}(\underline{\A}_\hbar^{\wedge_\chi},\theta)\xrightarrow{\sim} \wHC^{gr}(\underline{\A}_\hbar,\theta).$$
The functor $\bullet_{\dagger,\chi}:\wHC^{gr}(\A_\hbar,\theta)^{K,\kappa}\rightarrow \wHC(\underline{\A}_\hbar,\theta)$ we need
is the composition
$$\wHC^{gr}(\A_\hbar,\theta)^{K,\kappa}\rightarrow \wHC^{gr}(\A_\hbar^{\wedge_\chi},\theta)\xrightarrow{\sim}
\wHC^{gr}(\underline{\A}^{\wedge_\chi}_\hbar,\theta)\xrightarrow{\sim} \wHC^{gr}(\underline{\A}_\hbar,\theta),$$
where the first functor is the completion functor. Note that $\bullet_{\dagger,\chi}$
is exact and  $\C[\hbar]$-linear by the construction. Similarly to
\cite[Section 3.4]{HC}, it descends to an exact functor
$$\HC(\A,\theta,\zeta)^{K,\kappa}\rightarrow \HC(\underline{\A},\theta,\zeta).$$

%We note that $\zeta$ acts on $\A$ by $1$. If $\zeta$ acts on $\underline{\A}$ by $1$
%(which will happen, for example, when $\underline{\A}=\C$), then
%$\HC(\underline{\A},\theta)= \HC_+(\underline{\A},\theta)\oplus \HC_-(\underline{\A},\theta)$,
%and $\bullet_{\dagger,\chi}$ restricts to
%$$\HC_+(\A,\theta)\rightarrow \HC_+(\underline{\A},\theta).$$

Let us examine what happens with the equivariance. Set
\begin{equation}\label{eq:K_Q_def}
K_Q:=K\cap Q.
\end{equation}
By the  construction of $\bullet_{\dagger,\chi}$, this functor lifts to a functor to the category
of weakly $K_Q$-equivariant modules. Now we proceed to the strong equivariance.

\begin{defi}\label{defi:half_character}
Let $\rho_L$ denote one half of the character of the
$K_Q$-action on $\Lambda^{top}L$, where $L(=(\kf.\chi)^*)$.
\end{defi}

Note that $\rho_L=-\rho_\omega$, where $\rho_\omega$ was defined in Section \ref{SS_intro_approach}.
Set
\begin{equation}\label{eq:shifted_character}
\kappa_Q:=\kappa|_{\kf_Q}-\rho_L = \kappa|_{\kf_Q} + \rho_\omega.
\end{equation}

\begin{Lem}\label{Lem:restriction_equivariance}
The functor $\bullet_{\dagger,\chi}$ sends
$\wHC^{gr}(\A_\hbar,\theta)^{K,\kappa}$ to
$\wHC^{gr}(\underline{\A}_\hbar,\theta)^{K_Q,\kappa_Q}$.
\end{Lem}
\begin{proof}
Pick $M_\hbar\in  \wHC(\A_\hbar,\theta)^{K,\kappa}$ and let $ N_\hbar\in \wHC(\underline{\A}_\hbar, \theta)$ denote its image. For $\xi\in \kf_Q$, we write
$\xi_M,\xi_N,\xi_L$ for the operators defined by $\xi$ on $M_\hbar^{\wedge_\chi},
N_\hbar^{\wedge_\chi}$ and $\C[[L,\hbar]]$ via the differentials of the $K_Q$-actions. From the $K_Q$-equivariant isomorphism
$M_{\hbar}^{\wedge_\chi}\cong \C[[L,\hbar]]\widehat{\otimes}_{\C[[\hbar]]}N_\hbar^{\wedge_\chi}$
we get
\begin{equation}\label{eq:vect_fields}
\xi_{M}=\xi_L+\xi_N.
\end{equation}
Note that $\Phi_{\Weyl}(\xi)$ acts on $\C[[L,\hbar]]$ as $\xi_L+\langle\rho_L,\xi\rangle$
so from (\ref{eq:vect_fields}) and (\ref{eq:q_comom_A}) we get the following equality of
operators on $M_\hbar^{\wedge_\chi}$ (via the Lie algebra actions)
$$\Phi_{\A}(\xi)-\langle\kappa,\xi\rangle=\Phi_\Weyl(\xi)-\langle\rho_L,\xi\rangle+
\Phi_{\underline{\A}}(\xi)-\langle\kappa-\rho_L,\xi\rangle.$$
The claim of the lemma follows.
\end{proof}

%We will then define the full subcategory $\HC(\underline{\A},K_Q,\kappa_Q)$ in the category
%of weakly equivariant $\underline{\A}$-modules and see that the essential image of
%$\bullet_{\dagger,\chi}$ lies in that subcategory.

\subsection{Properties of restriction functor}\label{SS_Prop_restr}
As mentioned in \cite[Section 6.1]{Wdim}, most of the properties of $\bullet_{\dagger,\chi}$
mentioned below are proved completely similarly to their counterparts for
Harish-Chandra bimodules from \cite[Sections 3.3,3.4,4]{HC}.

As we have mentioned, the functor $$\bullet_{\dagger,\chi}:
\HC(\A_\hbar,\theta,\zeta)\rightarrow\HC(\underline{\A}_\hbar,\theta,\zeta)$$ is exact. In fact,
up to an isomorphism, the functor only depends on $\Orb_K$, not on the choice of
$\chi$. An isomorphism is given by an element of $K$ that sends one point to the other.

For a closed $K$-stable subvariety $Y\subset (\g/\kf)^*$ we write
$\HC_Y(\A,\theta,\zeta)$ for the full subcategory of all modules whose associated
variety is contained in $Y$. By the construction, for $M\in \HC_Y(\A,\theta,\zeta)^{K,\kappa}$,
the module $M_{\dagger,\chi}$ is supported on $S\cap Y$, where $S$ is the Slodowy slice from  
Section \ref{SS_Sl_sl}. In particular, if
$\overline{\Orb}_K$ is an irreducible component in $Y$, then $\dim M_{\dagger,\chi}<\infty$.
And if $\Orb_K\not \subset Y$, we have $M_{\dagger,\chi}=\{0\}$.

Now suppose that (\ref{eq:codim_condition_K}) holds.

Let $\operatorname{V}(\A)$ denote the associated variety of $\A$. We further assume that
\begin{equation}\label{eq:assoc_var_condition}
\operatorname{V}(\A)=\overline{\Orb}.
\end{equation}

These conditions will be assumed until the end of the section.
In this case, it is known, \cite[Theorem 1.3]{Vogan}, that if $\overline{\Orb}_K$
is an irreducible component in the associated variety of an irreducible
HC module, then it coincides with the associated variety (recall that we have assumed that $K$ is connected). Consider
the Serre quotient category
\begin{equation}\label{eq:quotient_cat}
\HC_{\Orb_K}(\A,\theta,\zeta)^{K,\kappa}:=\HC_{\overline{\Orb}_K}(\A,\theta,\zeta)^{K,\kappa}/\HC_{\partial\Orb_K}(\A,\theta,\zeta)^{K,\kappa}.
\end{equation}
The functor $\bullet_{\dagger,\chi}$ factors through $\HC_{\Orb_K}(\A,\theta,\zeta)^{K,\kappa}$ and
maps to the category $(\underline{\A},\zeta)\operatorname{-mod}^{K_Q,\kappa_Q}_{fin}$ of
finite dimensional $(K_Q,\kappa_Q)$-equivariant $\Z/d\Z$-graded $\underline{\A}$-modules.
The functor has the
right adjoint
$$\bullet^{\dagger,\chi}:(\underline{\A},\zeta)\operatorname{-mod}^{K_Q,\kappa_Q}_{fin}
\rightarrow \HC_{\overline{\Orb}_K}(\A,\theta,\zeta)^{K,\kappa}.$$
The functor $\bullet^{\dagger,\chi}$  between the weakly HC categories is constructed as in the case of Harish-Chandra bimodules in
\cite[Section 3.3]{HC}, see \cite[Proposition 3.3.4]{HC} there and the preceding discussion.
Similarly to \cite[Section 3.4]{HC}, it carries to the HC categories.

Further  properties of $\bullet_{\dagger,\chi},\bullet^{\dagger,\chi}$
are summarized in the following proposition.

\begin{Prop}\label{Prop:adj_properties}
Suppose that (\ref{eq:codim_condition_K}) and (\ref{eq:assoc_var_condition}) hold. Then we have the following:
\begin{enumerate}
\item The kernel and the cokernel of the adjunction unit morphism
$M\rightarrow (M_{\dagger,\chi})^{\dagger,\chi}$ are supported on $\partial\Orb_K$.
In particular, $\bullet_{\dagger,\chi}: \HC_{\Orb_K}(\A,\theta,\zeta)^{K,\kappa}\rightarrow (\underline{\A},\zeta)\operatorname{-mod}_{fin}^{K_Q,\kappa_Q}$ is a full embedding. 
\item The image of $\HC_{\Orb_K}(\A,\theta,\zeta)^{K,\kappa}$ in $(\underline{\A},\zeta)\operatorname{-mod}_{fin}^{K_Q,\kappa_Q}$
is closed under taking subquotients.
\item If $\operatorname{codim}_{\overline{\Orb}_K}(\partial\Orb_K)>2$,
then $$\bullet_{\dagger,\chi}: \HC_{\Orb_K}(\A,\theta,\zeta)^{K,\kappa}\xrightarrow{\sim} (\underline{\A},\zeta)\operatorname{-mod}_{fin}^{K_Q,\kappa_Q}.$$
\end{enumerate}
\end{Prop}
\begin{proof}
(1) is proved in the same way as \cite[Proposition 3.4.1(5)]{HC}. (2) is proved as
\cite[Theorem 4.1.1]{HC}. (3) is proved as \cite[Theorem 4.4]{reg}.
\end{proof}

%\begin{Rem}\label{Rem:distinguished_filtration}
%In fact, as in \cite[Section 3.4]{HC}, $\bullet^{\dagger}$ comes from
%the functor $$\bullet^{\dagger,\chi}:
%\underline{\A}_\hbar\operatorname{-mod}^{K_Q,\kappa_Q})^{gr}_{fin}
%\rightarrow \wHC_{\overline{\Orb}_K}^{gr}(\A_\hbar,\theta)^{K,\kappa},$$
%where in the source we have graded modules finitely generated over $\C[\hbar]$.
%This shows that for $V\in \underline{\A}\operatorname{-mod}^{K_Q,\kappa_Q}_{fin}$,
%the HC module $V^{\dagger,\chi}$ comes with a distinguished filtration: we equip
%$V$ with the trivial filtration in degree $0$. Then %$V^{\dagger,\chi}=(R_\hbar(V))^{\dagger,\chi}/(\hbar-1)$ inherits the filtration.
%We always consider this filtration on $V^{\dagger,\chi}$.
%\end{Rem}

To finish this section, we want to understand the interaction between $\bullet_{\dagger,\chi}$
and the tensor products with HC bimodules. Here it will be convenient for us to work
with modules and bimodules over $\U$. Let $\B$ be a HC $\U$-bimodule and $M$ be a HC
$(\g,K)$-module. We have the restriction functor $\B\mapsto \B_\dagger: \HC^G(\U)
\rightarrow \HC^Q(\Walg)$, \cite[Section 3.4]{HC}.

\begin{Lem}\label{Lem:restriction_compatibilty}
We have a functorial $K_Q$-linear isomorphism of $\Walg$-modules $\B_\dagger\otimes_\Walg M_{\dagger,\chi}\xrightarrow{\sim}(\B\otimes_\U M)_{\dagger,\chi}$.
\end{Lem}
\begin{proof}
Pick good filtrations on $\B$ and $M$. We have the induced tensor product filtration on
$\B\otimes_{\U} M$, which is also good. The completion functor is tensor so we have a
natural isomorphism
\begin{equation}\label{eq:tensor_product_iso}\B_\hbar^{\wedge_\chi}\otimes_{\U_\hbar^{\wedge_\chi}}M_\hbar^{\wedge_\chi}
\xrightarrow{\sim} (\B_\hbar\otimes_{\U_\hbar}M_\hbar)^{\wedge_\chi}.
\end{equation}
Let
$\underline{\B}_\hbar'$ denote the centralizer of $\Weyl_\hbar^{\wedge_0}$
in $\B_\hbar^{\wedge_\chi}$. The left hand
side is
\begin{align*}
\left(\Weyl_\hbar^{\wedge_0}\widehat{\otimes}\underline{\B}_\hbar'\right)
\widehat{\otimes}_{\Weyl_\hbar^{\wedge_0}\widehat{\otimes}_{\C[[\hbar]]}
\Walg_\hbar^{\wedge_\chi}}(\C[[L,\hbar]]\widehat{\otimes}_{\C[[\hbar]]}N_\hbar^L)
\xrightarrow{\sim} \C[[L,\hbar]]\widehat{\otimes}_{\C[[\hbar]]}
(\underline{\B}_\hbar'\widehat{\otimes}_{\C[[\hbar]]}N_\hbar^L).\end{align*}
This isomorphism combined with (\ref{eq:tensor_product_iso}) and the constructions
of $\bullet_{\dagger,\chi}$ (Section \ref{SS:HC_mod_restr}) and
$\bullet_\dagger$ (\cite[Section 3.4]{HC}) yields
the claim of the lemma.
\end{proof}

\begin{Rem}\label{Rem:restriction_compatibility}
The lemma generalizes to the weakly HC (bi)modules. It also generalizes to Hom's instead of tensor
products (this will be needed in a subsequent paper by the first named author). Let $B_\hbar$
be a graded weakly HC $\U_\hbar$-bimodule, while $M_\hbar\in \wHC^{gr}(\U_\hbar,\theta)^{K,\kappa}$.
Then $\Hom_{\U_\hbar}(B_\hbar,M_\hbar)$ has a natural structure of an object in
$\wHC^{gr}(\U_\hbar,\theta)^{K,\kappa}$, and we have a natural isomorphism
$\Hom_{\U_\hbar}(B_\hbar,M_\hbar)_{\dagger,\chi}\xrightarrow{\sim}\Hom_{\Walg_\hbar}(B_{\hbar,\dagger},
M_{\hbar,\dagger,\chi})$. The proof repeats that of the lemma.
\end{Rem}

\subsection{The case of quantizations of nilpotent orbits}\label{subsec:W-algebra_nilpotent}
Now we would like to treat the case when the Dixmier algebra $\A$ in
Section \ref{SS_Dixmier_restriction} is a quantization of
$\C[\Orb]$. We assume that (\ref{eq:codim_condition}) and (\ref{eq:assoc_var_condition}) hold. Clearly, in this case $\underline{\A}=\C$.

Assume (\ref{eq:codim_condition}) holds.
%Note that this automatically
%implies (\ref{eq:codim_condition_K}).
In what follows we need to understand how the quantum comoment map
$\Phi_{\underline{\A}}:\q\rightarrow \underline{\A}$ behaves. Since $\underline{\A}=\C$,
the map $\Phi_{\underline{\A}}$ is just an element of $\q^{*Q}$.

Recall that the quantizations of $\C[\Orb]$ are classified by $H^2(\Orb,\C)$,
Section \ref{SS_quant_nilp}.
Note that $(\q^*)^Q$  is identified with $H^2(\Orb,\C)$ as a vector space, Lemma \ref{Lem:G_orbit_cohomology}.
On the lattice $\Hom(Q,\C^\times)\subset \q^{*Q}$ this identification sends
a character of $Q$ to the first Chern class of the corresponding $G$-equivariant line  bundle
on $\Orb$.

\begin{Lem}\label{Lem:quantum_comom_period}
Let $\lambda \in H^2(\Orb, \C)$ be the period of $\A$, see
Proposition \ref{prop:graded}. Then $\Phi_{\underline{\A}}=\lambda \in  \q^{*Q}$.
\end{Lem}
\begin{proof}
Let us write $\A_\lambda$ for the quantization $\A$ with period $\lambda$. Also recall the universal
quantization $\A_{\param}$, where $\param=H^2(\Orb,\C)$, Section \ref{SS_quant_nilp}.  We still have the quantum comoment map $\g\rightarrow \A_{\param}$,
which gives rise to a linear $Q$-equivariant map $\q\rightarrow \underline{\A}_{\param}=\C[\param]$. By degree reasons, the image of this map lies in $\q^{Q}\oplus \C$. The claim of the lemma
boils down to the claim that this map is the $Q$-equivariant  projection
$\q\twoheadrightarrow \q^{Q}$.

The first observation is that the $\q^Q$-part of $\q\rightarrow \q^Q\oplus \C$ is the projection.
To see this, one can consider the quasi-classical limit of the situation. Take a line
bundle $\mathcal{L}$ on $\Orb$, it corresponds to a character of $Q$, say $\mu$. Let $\A^\Orb_\lambda$ denote
the microlocalization of $\A_\lambda$ to $\Orb$.
We have $H^i(\Orb,\mathcal{O}_{\Orb})=0$ for $i=1,2$, (\ref{eq:orbit_coh_cohom_vanish}).
It follows that $\mathcal{L}$ deforms to a unique $G$-equivariant $\A^\Orb_{\lambda+\mu}$-$\A^\Orb_{\lambda}$-module, denote it by $\mathcal{L}'$, compare to
\cite[Proposition 5.2]{BPW} and its proof. Since $\operatorname{codim}_{\overline{\Orb}}\partial\Orb\geqslant 4$,  the global sections
$\Gamma(\mathcal{L}')$ is a HC $\A_{\lambda+\mu}$-$\A_\lambda$-bimodule whose microlocalization
back to $\Orb$ is $\mathcal{L}'$, compare to the proof of \cite[Theorem 4.4]{reg}.
Apply the restriction functor $\bullet_\dagger$ for Harish-Chandra bimodules, \cite[Section 3.4]{HC},
to $\Gamma(\mathcal{L}')$. We land in the category of $Q$-equivariant $\Walg$-bimodules. By \cite[Lemma 3.3.2]{HC}, we have that $\gr \Gamma(\mathcal{L}')_\dagger$
is the pullback of $\gr\mathcal{L}'$ to $S\cap \Orb$. So $\Gamma(\mathcal{L}')_\dagger$ is a one-dimensional
space and $Q$ acts via $\mu$. On the other hand, the adjoint action of $\mathfrak{q}\subset \Walg$
is by $\Phi_{\underline{\A}_{\lambda+\mu}}-\Phi_{\underline{\A}_\lambda}$. This difference is $\mu$. 
Equivalently, the $\q^Q$-part of the map $\q\rightarrow \q^Q\oplus \C$ is the projection.
%
%, compare
%to \cite{CWR}.

Finally, we need to show that the $\C$-part of the map $\q\rightarrow \q^Q\oplus \C$
is zero.  This is equivalent to the claim that for the canonical quantization $\A_0$, we have
$\Phi_{\underline{\A}_0}=0$. For this, consider the Rees algebra $R_\hbar(\A_0)$.
It comes with a parity anti-automorphism, denote it by $\varsigma$, see
\cite[Corollary 2.3.3]{quant_iso}. The anti-automorphism $\varsigma$ is
characterized by the property that it is the identity modulo $\hbar$ and sends
$\hbar$ to $-\hbar$. In this situation it makes sense to speak about
the symmetrized quantum comoment map to $R_\hbar(\A_0)$, the homogeneous
map with image in the $\varsigma$-invariants,  see \cite[Section 5.4]{quant_iso}.
Note that $\Phi_{\A_0}:\g
\rightarrow R_\hbar(\A_0)$ is symmetrized. We can assume that the image of the embedding
$V^*\hookrightarrow R_\hbar(\A_0)^\wedge$ is $\varsigma$-stable. So $\underline{\A}_{0,\hbar}$
inherits the anti-automorphism $\varsigma$. The quantum comoment map $\Phi_{\mathbb{A}}$
is symmetrized. From here and (\ref{eq:q_comom_A}) we see that
$\Phi_{\underline{\A}_0}:\q\rightarrow \underline{\A}_{0}=\C[\hbar]$
is also symmetrized. So,   $\Phi_{\underline{\A}_0}=0$.
\end{proof}

Here is an immediate corollary of Lemma \ref{Lem:quantum_comom_period}.

\begin{Cor}\label{Cor:equiv_condition} 
The category $\wHC^{gr}(\underline{\A}_\hbar,\theta)^{K_Q,\kappa_Q}$ is nonzero if and only if $\lambda|_{\kf_Q}-\kappa_Q$
integrates to a character of $K_Q^\circ$. This is the character by which $K_Q^\circ$
acts on every object in $\wHC^{gr}(\underline{\A}_\hbar,\theta)^{K_Q,\kappa_Q}$.
\end{Cor}


Now note that $(\underline{\A},\zeta)\operatorname{-mod}^{K_Q,\kappa_Q}$ is a semisimple category:
by Corollary \ref{Cor:equiv_condition} this is the category of $\Z/d\Z$-graded $K_Q$-modules, where
$\mathfrak{k}_Q$ acts by $\lambda|_{\kf_Q}-\kappa_Q$.
Recall that (\ref{eq:codim_condition}) implies (\ref{eq:codim_condition_K}) for each $\Orb_K$.
Proposition \ref{Prop:adj_properties} implies that $\HC_{\Orb_K}(\A,\theta,\zeta)^{K,\kappa}$
is the full subcategory that consists of direct sums of certain simple objects in
$(\underline{\A},\zeta)\operatorname{-mod}^{K_Q,\kappa_Q}$. So to determine $\HC_{\Orb_K}(\A,\theta,\zeta)^{K,\kappa}$ we just need to understand which simple
objects in $(\underline{\A},\zeta)\operatorname{-mod}^{K_Q,\kappa_Q}$ lie in the image of $\HC_{\Orb_K}(\A,\theta,\zeta)^{K,\kappa}$. The following lemma is
a straightforward corollary of Proposition \ref{Prop:adj_properties}.

\begin{Lem}\label{Lem:image_description}
Let $N$ be a simple object in $(\underline{\A},\zeta)\operatorname{-mod}^{K_Q,\kappa_Q}$. Then the following claims
are equivalent:
\begin{enumerate}
\item $N$ lies in the image of $\bullet_{\dagger,\chi}$.
\item $N^{\dagger,\chi}\neq 0$.
\item $(N^{\dagger,\chi})_{\dagger,\chi}\xrightarrow{\sim} N$.
\end{enumerate}
\end{Lem}

We finish this section with two remarks.

\begin{Rem}\label{Rem:removing_zeta}
Note that $\zeta$ acts trivially on both $\A$ and $\underline{\A}=\C$. It follows that $\HC(\A,\theta,\zeta)$
decomposes as the direct sum of $d$ copies of $\HC(\A,\theta):=\HC(\A,\theta,1)$ (where $\zeta$
acts on an object in $i$th copy by $\exp(2\pi\sqrt{-1}i/d)$). The same is true of $(\underline{\A},\zeta)\operatorname{-mod}^{K_Q,\kappa_Q}$.
Therefore, it makes sense to define $M_{\dagger,\chi}$ for $M\in \HC(\A,\theta,1)$
and $N^{\dagger,\chi}$ for $N\in \underline{\A}\operatorname{-mod}^{K_Q,\kappa_Q}$.
Below we also write $\operatorname{Rep}(K_Q,\lambda|_{\kf_Q}-\kappa_Q)$ for $\underline{\A}\operatorname{-mod}^{K_Q,\kappa_Q}$.
\end{Rem}

\begin{Rem}\label{Rem:distinguished_filtration}
In fact, as in \cite[Section 3.4]{HC}, $\bullet^{\dagger,\chi}$ comes from
the functor $$\bullet^{\dagger,\chi}:
\underline{\A}_\hbar\operatorname{-mod}^{K_Q,\kappa_Q,gr}_{fin}
\rightarrow \wHC_{\overline{\Orb}_K}^{gr}(\A_\hbar,\theta)^{K,\kappa},$$
where in the source we have graded modules finitely generated over $\C[\hbar]$.
This shows that for $N\in \underline{\A}\operatorname{-mod}^{K_Q,\kappa_Q}_{fin}$,
the HC module $N^{\dagger,\chi}$ comes with a distinguished filtration: we equip
$N$ with the trivial filtration in degree $0$. Then $N^{\dagger,\chi}=(R_\hbar(N))^{\dagger,\chi}/(\hbar-1)$ inherits the filtration.
We always consider this filtration on $N^{\dagger,\chi}$.
\end{Rem}

\subsection{Restriction functor vs quantizations}
The purpose of this section is relate the restriction functor in Sections \ref{SS:HC_mod_restr}-\ref{subsec:W-algebra_nilpotent} to
quantizations of vector bundles on Lagrangian subvarieties studied in Section
\ref{S_qclass_quant}. This will ultimately give us tools to examine the
equivalent conditions of Lemma \ref{Lem:image_description}.

In this section, $\Orb$ is a nilpotent orbit in $\g^*$ and $\A$ is a filtered quantization
of $\C[\Orb]$. We assume that $K$ is connected.
We do not put any conditions on the codimension of boundary at this point. Let
$\Orb_K$ denote a $K$-orbit in $\Orb\cap (\g/\kf)^*$. The goal of this section
is to produce an equivalence between $\wHC^{gr}(\underline{\A}_\hbar,\theta)^{K_Q,\kappa_Q}$
and a category of microlocal sheaves on $\Orb_K$. The functor $\bullet_{\dagger,\chi}$
will become the microlocalization functor, while  $\bullet^{\dagger,\chi}$ that exists
under condition (\ref{eq:codim_condition_K}) will
become the global sections functor.

To relate  $\wHC^{gr}(\underline{\A}_\hbar,\theta)^{K_Q,\kappa_Q}$ to the category of microlocal sheaves
we will use jet bundles.

\subsubsection{Microlocal version}
In this section we will write $\hat{\A}_\hbar$ for the $\hbar$-adic completion of the
Rees algebra $\A_\hbar$ of $\A$. Let $\hat{\A}_\hbar^{\Orb}$ denote the microlocalization of
$\hat{\A}_\hbar$ to $\Orb$ viewed as an open subvariety in $\operatorname{Spec}(\C[\Orb])$.
This is a sheaf in Zariski topology.


We can talk about  weakly HC $(\hat{\A}_\hbar^\Orb,\theta)$-modules, these are sheaves in the Zariski
topology on $\Orb$ defined similarly to Definition \ref{defi:wHC}.
Denote the corresponding category by $\wHC(\hat{\A}_\hbar^\Orb,\theta)$. Note that we have
the microlocalization functor from the category $\wHC_{\overline{\Orb}_K}(\hat{\A}_\hbar,\theta)$
to the subcategory $\wHC_{\Orb_K}(\hat{\A}^\Orb_\hbar,\theta)$ of all weakly HC $\hat{\A}^\Orb_\hbar$-
modules supported (set theoretically) on $\Orb_K$.

\begin{Lem}\label{Lem:global_sections}
If (\ref{eq:codim_condition_K}) holds, then we have the right
adjoint functor $\wHC_{\Orb_K}(\hat{\A}^\Orb_\hbar,\theta)\rightarrow
\wHC_{\overline{\Orb}_K}(\hat{\A}_\hbar,\theta)$ given by taking the global sections.
\end{Lem}
\begin{proof}
Let $M_\hbar^0\in \wHC_{\Orb_K}(\hat{\A}^\Orb_\hbar,\theta)$. We need to show that
\begin{itemize}
\item[(*)]
$\Gamma(M_\hbar^0)$ is a finitely generated $\hat{\A}_\hbar$-module.
\end{itemize}
Note that $\Gamma(M_\hbar^0)$ is automatically
complete and separated in the $\hbar$-adic topology because $M_\hbar^0$ is. So it is
enough to prove (*) when $M_\hbar^0$ is killed by $\hbar$.
Note that $M_\hbar^0/\hbar M_{\hbar}^0$
is supported on $\Orb_K$ scheme theoretically, this follows from (i)
of Definition \ref{defi:wHC}.
So $M_\hbar^0$
is a $K$-equivariant coherent sheaf on $\Orb_K$, hence a vector bundle. Since $\operatorname{codim}_{\overline{\Orb}_K}\partial\Orb_K\geqslant 2$, the global sections
$\Gamma(M^0_\hbar)$ is a finitely generated module over $\C[\overline{\Orb}_K]$. This finishes
the proof.
\end{proof}

Similarly, we can talk about the category $\wHC(\hat{\A}_\hbar^\Orb,\theta)^{K,\kappa}$.
Our goal is to  establish an equivalence $\wHC_{\Orb_K}(\hat{\A}_\hbar^{\Orb},\theta)^{K,\kappa}
\xrightarrow{\sim} \wHC(\underline{\hat{\A}}_\hbar,\theta)^{K_Q,\kappa_Q}$ that intertwines
$\bullet_{\dagger,\chi}$ with the microlocalization functor (and
$\bullet^{\dagger,\chi}$ with the global section functor).

To accomplish this, we will need (quantum) jet bundles.

\subsubsection{Jet bundles}
We can form the jet bundle of $\hat{\A}_\hbar^{\Orb}$ to be denoted
by $\jet \hat{\A}_\hbar^{\Orb}$. This is pro-vector bundle on $\Orb$ whose fiber at $x\in \Orb$
is the completion $\A_{\hbar}^{\wedge_x}$. It comes equipped with a flat connection to be
denoted by $\nabla$. We recover $\hat{\A}_{\hbar}^{\Orb}$ as the subalgebra of flat sections
$(\jet \A_\hbar^{\Orb})^\nabla$. Compare to \cite[Section 3.2]{BK},\cite[Section 3.2]{1dim_quant},\cite[Section 2.6]{HC_symp}.

Now let $M_\hbar^0\in \wHC_{\Orb_K}(\hat{\A}_{\hbar}^{\Orb},\theta)$.
%Also note that, for each $\chi\in \Orb_K$,
%the completed module $M_{\hbar}^{0\wedge_\chi}$ is a weakly HC $(\A_\hbar^{\wedge_\chi},\theta)$-module.
The algebra $\A_\hbar^{\wedge_\chi}$ is identified with the formal Weyl algebra of the symplectic vector space $T_\chi\Orb\cong T_\chi\Orb_K\oplus T^*_\chi\Orb_K$
so that the involution $\theta$ acts by $1$ on $T^*_\chi\Orb_K$
and by $-1$ on $T_\chi \Orb_K$. By Lemma
\ref{Lem:wHC_formal_Weyl} we have
\begin{equation}\label{eq:wHC_decomp}
M_\hbar^{0\wedge_\chi}\cong \C[[T_\chi\Orb_K,\hbar]]\otimes_{\C[[\hbar]]}N_\hbar.
\end{equation}

%Similarly to \cite[Proposition 3.3.1]{HC} we have the following %result.\footnote{Maybe it should be in an earlier section.}
%
%\begin{Lem}\label{Lem:wHC_Weyl_decomp}
%Let $N_\hbar$ denote the annihilator of $T_\chi\Orb$ in $M_\hbar^{0\wedge_\chi}$
%(for the Lie algebra action). Then $M_\hbar^{0\wedge_\chi}\cong %\C[[T_\chi\Orb,\hbar]]\otimes_{\C[[\hbar]]}N_\hbar$, where $\C[[T_\chi\Orb,\hbar]]$
%has a natural structure of a weakly HC module over the Weyl algebra.
%\end{Lem}
%It follows that $M_\hbar^{0\wedge_\chi}$ is identified
%with the direct sum of several copies of $\C[[T_\chi \Orb,\hbar]]$.

We want to define the jet bundle of $M_\hbar^0$.  First,
consider the  sheaf of algebras $\hat{\A}_\hbar^{\Orb}\widehat{\boxtimes} \Str_{\Orb_K}$ on $\Orb\times \Orb_K$ (the completed outer tensor product).
Let $\mathcal{J}_\Delta$ denote the ideal sheaf of the diagonal in $\Str_{\Orb}\boxtimes
\Str_{\Orb_K}$ and let $\mathcal{J}_{\Delta,\hbar}$ denote its preimage in
$\hat{\A}_\hbar^{\Orb}\widehat{\boxtimes} \Str_{\Orb_K}$. Let $\jet_{\Orb_K} \hat{\A}_\hbar$
denote the completion of  $\hat{\A}_\hbar^{\Orb}\widehat{\boxtimes} \Str_{\Orb_K}$ with respect to
this ideal. It is nothing else but the pullback of the quantum jet bundle $\jet \hat{\A}_\hbar^{\Orb}$
to $\Orb_K$. It is a $K$-equivariant pro-vector bundle of algebras on $\Orb_K$ that comes with a flat
connection. The fiber at $\chi$ is $\A_\hbar^{\wedge_\chi}$.

We can use a similar construction to define the jet bundle $\jet_{\Orb_K} M_\hbar^0$.
This is a $\jet_{\Orb_K} \hat{\A}_\hbar$-module that is equipped with a compatible flat connection.
Note that the sheaf $\jet_{\Orb_K} \hat{\A}_\hbar$ comes with an $\Str_{\Orb_K}$-linear
involution induced by the anti-involution $\theta\boxtimes \operatorname{id}$
on $\hat{\A}_\hbar^{\Orb}\boxtimes \Str_{\Orb_K}$. It will also be denoted by $\theta$. Using this
involution, we can define the notion of a weakly HC $(\jet_{\Orb_K} \hat{\A}_\hbar,\theta)$-module
similarly to what was done in \cite[Section 2.6]{HC_symp} for HC bimodules.

We now explain the details of the definition.
Define $\mathfrak{J}_\hbar\subset \jet_{\Orb_K} \hat{\A}_\hbar$ to be the preimage of the
ideal in $\jet_{\Orb_K} \mathcal{O}_\Orb$ generated by $(\jet_{\Orb_K} \mathcal{O}_\Orb)^{-\theta}$.

Consider locally finitely generated
$\jet_{\Orb_K} \hat{\A}_\hbar$-modules $\mathfrak{M}^0_\hbar$ that come equipped with a compatible
flat connection $\nabla$ and the Lie algebra action of $\mathfrak{J}_\hbar$ subject
to the following conditions:

\begin{itemize}
\item[(i)] The action map $\mathfrak{J}_\hbar\otimes \mathfrak{M}^0_\hbar\rightarrow
\mathfrak{M}^0_\hbar$ is flat (i.e. is compatible with the flat connection).
\item[(ii)] $\mathfrak{M}^0_\hbar$ is complete and separated in the $\hbar$-adic topology.
\item[(iii)] The successive quotients $\hbar^{k-1}\mathfrak{M}^0_\hbar/\hbar^k \mathfrak{M}^0_\hbar$
are the jet bundles of some vector bundles on $\Orb_K$.
\end{itemize}

Let $\wHC(\jet_{\Orb_K} \hat{\A}_\hbar,\theta)$ denote the category
of such modules $\mathfrak{M}^0_\hbar$. The morphisms are $\jet_{\Orb_K} \hat{\A}_\hbar$-
and $\mathfrak{J}_\hbar$-linear maps that respect the flat connections.

Then we have an exact functor
$$\jet: \wHC_{\Orb_K}(\hat{\A}_\hbar^{\Orb},\theta) \to \wHC(\jet_{\Orb_K} \hat{\A}_\hbar,\theta)$$
of taking the jet bundle.

\begin{Lem}\label{Lem:HC_jet}
The following claims are true:
\begin{enumerate}
\item
For $\mathfrak{M}^0_\hbar\in \wHC(\jet_{\Orb_K} \hat{\A}_\hbar,\theta)$, its submodule of flat sections, $(\mathfrak{M}^0_\hbar)^\nabla$, lies in
$\wHC_{\Orb_K}(\hat{\A}_\hbar^{\Orb},\theta)$.
\item  Moreover,
$$\bullet^\nabla: \wHC(\jet_{\Orb_K} \hat{\A}_\hbar,\theta)\rightarrow
\wHC_{\Orb_K}(\hat{\A}_\hbar^{\Orb},\theta) $$ is quasi-inverse to $\jet$.
\end{enumerate}
\end{Lem}
\begin{proof}
For (1) we just need to show that $(\mathfrak{M}^0_\hbar)^\nabla$ is locally finitely
generated. Thanks to conditions (ii), it is enough to do this when
$\mathfrak{M}^0_\hbar$ is annihilated by $\hbar$. In this case, by (iii), it is the jet bundle of a vector
bundle and the flat sections recover the vector bundle.

(2) is proved as \cite[Lemma 2.16]{HC_symp}.
\end{proof}
%with $\bullet^{\nabla}$ being right adjoint to $\jet$.

%The following claim is proved in exactly the same way as \cite[Lemma 3.2]{HC_symp}.
%
%\begin{Lem}\label{Lem:HC_jet}
%The functor $\jet$ is an equivalence (with quasi-inverse $\bullet^{\nabla}$).
%\end{Lem}

We can also define the graded analog $\wHC^{gr}(\jet_{\Orb_K} \hat{\A}_\hbar,\theta)$:
the objects $\mathfrak{M}^0_\hbar$ come with a pro-rational action of $\C^\times$ compatible
with the action on $\jet_{\Orb_K} \hat{\A}_\hbar$ and such that the connection is
$\C^\times$-invariant. A direct analog of Lemma \ref{Lem:HC_jet} holds.

\subsubsection{Equivariant version}\label{SSS_equiv_jets}
Note that the $G$-action on $\hat{\A}_\hbar^{\Orb}$ is Hamiltonian in the sense of
Definition \ref{defi:Hamilt_action_quantization}.
We now discuss the equivariance for weakly HC modules: we want to define a jet counterpart of $\wHC(\hat{\A}_\hbar^\Orb,\theta)^{K,\kappa}$, the category of $(K,\kappa)$-equivariant
objects in $\wHC(\hat{\A}_\hbar^\Orb,\theta)$, defined as in Definition \ref{defi:equiv_quant_modules}. Note that $K$ naturally acts on $\jet_{\Orb_K}\hat{\A}_\hbar^\Orb$.
For $\xi\in \kf$, we write $\xi_{\jet\A}$ for the induced derivation of
$\jet_{\Orb_K}\hat{\A}_\hbar$, $\xi_{\Orb_K}$ for the vector field on $\Orb_K$
induced by $\xi$, and $\xi_{\A}$ for the image of $\xi$ in $\A_\hbar$. Then we have
\begin{equation}\label{eq:deriv_equation}
\xi_{\jet\A}=\nabla_{\xi_{\Orb_K}}+\hbar^{-2}[\xi_{\A},\cdot].
\end{equation}
Compare to \cite[(3.1)]{1dim_quant}.

\begin{defi}\label{defi:wHC_jet_equiv}
By a {\it $(K,\kappa)$-equivariant weakly HC} $\jet_{\Orb_K}\hat{\A}_\hbar$-module
we mean a weakly $K$-equivariant object $\mathfrak{M}^0_\hbar\in \wHC(\jet_{\Orb_K}\hat{\A}_\hbar,\theta)$
satisfying the following condition:
\begin{equation}\label{eq:jet_equiv}
\xi_{\mathfrak{M}^0_\hbar}m=\nabla_{\xi_{\Orb_K}}m+\xi_{\A}.m-\langle\kappa,\xi\rangle m
\end{equation}
for all local sections $m$. The category of $(K,\kappa)$-equivariant weakly HC
$\jet_{\Orb_K}\hat{\A}_\hbar$-modules will be denoted by
$\wHC(\jet_{\Orb_K}\hat{\A}_\hbar,\Orb)^{K,\kappa}$.
\end{defi}

With this definition, the following is a straightforward consequence of
Lemma \ref{Lem:HC_jet}.

\begin{Lem}\label{Lem:HC_jet_equiv}
The functor $\jet$ is an equivalence
$$\wHC(\hat{\A}_\hbar^{\Orb},\theta)^{K,\kappa}\xrightarrow{\sim} \wHC(\jet_{\Orb_K}\hat{\A}_\hbar,\theta)^{K,\kappa}.$$
(with quasi-inverse $\bullet^{\nabla}$).
\end{Lem}

\subsubsection{Restriction to a point}\label{SSS_point_restriction}
Our goal here is to produce an equivalence
$$\wHC(\jet_{\Orb_K}\A_\hbar,\theta)^{K,\kappa}\xrightarrow{\sim} \wHC(\underline{\A}_\hbar,\theta)^{K_Q,\kappa_Q}.$$
First, we have the functor
\begin{equation}\label{eq:restr_to_pt}\wHC(\jet_{\Orb_K}\hat{\A}_\hbar,\Orb)^{K,\kappa}\rightarrow
\wHC(\A_\hbar^{\wedge_\chi},\theta)^{K_Q,\kappa|_{\kf_Q}}
\end{equation}
of restriction to $\chi$.

\begin{Lem}\label{Lem:restr_pt_equiv}
The functor (\ref{eq:restr_to_pt})  is a category equivalence.
\end{Lem}
\begin{proof}
We construct a quasi-inverse functor. Let $N_\hbar\in \wHC(\A_\hbar^{\wedge_\chi},\theta)^{K_Q,\kappa|_{\kf_Q}}$. The Lie algebra
$\mathfrak{k}_\chi$ acts on $N_\hbar$ via the Lie algebra action of $(\A_\hbar^{\wedge_\chi})^{-\theta}$
shifted by $\kappa$.
This action is compatible with that of $K_Q$ by the equivariance condition. So we get a pro-rational
action of $K_\chi$ on $N_\hbar$. So
we can form the corresponding homogeneous bundle $K\times^{K_\chi}N_\hbar$ on $\Orb_K$ with fiber $N_\hbar$ at $\chi$.
We recover the flat connection on this bundle
by (\ref{eq:jet_equiv}). We get an object of
$\wHC(\jet_{\Orb_K}\hat{\A}_\hbar,\theta)^{K,\kappa}$. The functors sending $\mathfrak{M}_\hbar^0$ to its fiber at $\chi$ and sending $N_\hbar$ to $K\times^{K_\chi}N_\hbar$ are quasi-inverse.
\end{proof}

Now recall, Section \ref{SS:HC_mod_restr}, that we have the equivalence
$$\wHC(\hat{\A}_\hbar^{\wedge_\chi},\theta)^{K_Q,\kappa|_{\kf_Q}}\xrightarrow{\sim}
\wHC(\underline{\A}_\hbar,\theta)^{K_Q,\kappa_Q}$$
given by $\C[[L,\hbar]]\otimes_{\C[[\hbar]]}\bullet$.
Composing this equivalence with
(\ref{eq:restr_to_pt}) we get the required equivalence
\begin{equation}\label{eq:interm_jet_equiv}
\wHC(\jet_{\Orb_K}\A_\hbar,\theta)^{K,\kappa}\xrightarrow{\sim}
\wHC(\underline{\A}_\hbar,\theta)^{K_Q,\kappa_Q}.
\end{equation}

\subsubsection{Conclusion}\label{SSS_conclusion_microloc}
Combining Lemma \ref{Lem:HC_jet_equiv} with (\ref{eq:interm_jet_equiv})  we conclude that we have an equivalence
\begin{equation}\label{eq:final_jet_equiv}
\wHC(\hat{\A}_\hbar^{\Orb},\theta)^{K,\kappa}\xrightarrow{\sim}
\wHC(\underline{\hat{\A}}_\hbar,\theta)^{K_Q,\kappa_Q}.
\end{equation}
By the construction of $\bullet_\dagger$ in Section \ref{SS:HC_mod_restr},  this equivalence intertwines the microlocalization functor
$\wHC(\hat{\A}_\hbar,\theta)^{K,\kappa}\rightarrow \wHC(\hat{\A}_\hbar^\Orb,\theta)^{K,\kappa}$
and the restriction functor $\bullet_{\dagger,\chi}:
\wHC(\A_\hbar,\theta)^{ K,\kappa}\rightarrow \wHC(\underline{\hat{\A}}_\hbar,\theta)^{K_Q,\kappa_Q}$.
It follows that the global section functor
$\wHC(\hat{\A}^\Orb_\hbar,\theta)^{K,\kappa}\rightarrow \wHC(\hat{\A}_\hbar, \theta)^{K,\kappa}$
(the right adjoint to the microlocalization functor)
gets intertwined with $\bullet^{\dagger,\chi}$ (the right adjoint of $\bullet_{\dagger,\chi}$).

The same works for the graded categories.

\subsubsection{Twisted local systems and quantizations}\label{SSS_tw_loc_quant}
Let $M$ be a $(K,\kappa)$-equivariant HC $(\A,\theta)$-module. Equip it with a good filtration
and form the completed Rees module $M_\hbar$. Consider the microlocalization $M_\hbar|_{\Orb}$.
The quotient $M_\hbar|_{\Orb}/ \hbar M_\hbar|_{\Orb}$ is a $K$-equivariant coherent sheaf on
$\Orb\cap (\g/\kf)^*$. Since the latter is a Lagrangian subvariety and the union of finitely
many $K$-orbits, $M_\hbar|_{\Orb}/ \hbar M_\hbar|_{\Orb}$ is a $K$-equivariant vector 
bundle on $\Orb\cap (\g/\kf)^*$, and $M_\hbar|_{\Orb}$ is its graded quantization. 
So $M_\hbar|_{\Orb}/(\hbar)$ is a $(K,\kappa)$-equivariant twisted local system on $\Orb^\theta$,
see the last paragraph in Section \ref{SS_Ham_quant}.

Let $V:=(M_\hbar/\hbar M_\hbar)_{\dagger,\chi}$. If the period of $\A$ is $\lambda$, then
$\kf_Q$ acts on $V$ by $\lambda|_{\kf_Q}-\kappa_Q$, see Corollary \ref{Cor:equiv_condition}.
This determines the twist of the restriction of $M_\hbar/\hbar M_\hbar$ to $\Orb_K$,
see Lemma \ref{Lem:homog_classif}. In particular, consider the situation when $\lambda=\kappa=0$.
Then this restriction has half-canonical twist. This gives an analog of \cite[Theorem 8.7]{Vogan}
in our setting.

%Section 3 and Section 4.1.

\section{Extension}\label{sec:extn}
Here we will give sufficient conditions for an object of
$\underline{\A}_\lambda\operatorname{-mod}^{K_Q,\kappa_Q}$
to lie in the image of $\bullet_{\dagger,\chi}$.  Our results
in this section are similar to what was obtained in \cite[Section 3]{HC_symp}.



\subsection{Extension to codimension $2$} \label{subsec:extn_codim2}
Recall that we assume that $\operatorname{codim}_{\overline{\Orb}_K}\partial\Orb_K\geqslant 2$.
Let $\Orb_K^2$ denote the union of $\Orb_K$ and all codimension $2$ $K$-orbits
in $\overline{\Orb}_K$. The closure $\overline{\Orb}_K$ is taken in $\operatorname{Spec}(\C[\Orb])$
and may be different from the closure in $\overline{\Orb}$.

Set $\partial \Orb_K^2:=\overline{\Orb}_K\setminus \Orb_K^2$, this is
closed subvariety in $\overline{\Orb}_K$ of codimension at least $3$.
Let $\iota$ denote the inclusion $\Orb_K\hookrightarrow \Orb_K^2$ and
$\iota'$ denote the inclusion $\Orb_K^2\hookrightarrow \overline{\Orb}_K$.

Recall that $\hat{\A}_\hbar$ stands for the $\hbar$-adic completion of
$\A_\hbar$, and  $\hat{\A}_\hbar^\Orb$ is the microlocalization of $\hat{\A}_\hbar$ to $\Orb$.
We write $\HC^{gr}(\hat{\A}_\hbar^{\Orb},\theta)$ for the full subcategory of
$\C[\hbar]$-flat objects in $\wHC^{gr}(\hat{\A}_\hbar^{\Orb},\theta)$.


Let $\mathcal{E}_\hbar$ be an object in $\HC^{gr}(\hat{\A}_\hbar^{\Orb},\theta)$. Assume that $\mathcal{E}_\hbar$ is supported on $\Orb_K$
so that $\mathcal{E}_\hbar/\hbar \mathcal{E}_\hbar$ is a twisted local system on $\Orb_K$, see
Section \ref{SSS_tw_loc_quant}.
We want to get a sufficient
condition for $(\iota'\circ\iota)_* \mathcal{E}_\hbar$ to be coherent.
See \cite[Section 2.5]{HC_symp} for a discussion of coherent modules and the pushforwards.


To start with note that if $(\iota'\circ\iota)_* \mathcal{E}_\hbar$  is coherent,
then so is $\iota_*\mathcal{E}_\hbar$. So, as the first step, we will find a necessary and sufficient condition for $\iota_*\mathcal{E}_\hbar$ to be coherent.

Take a $K$-orbit $\Orb_K'\subset \Orb^2_K\setminus \Orb_K$ and a point $\chi'\in \Orb_K'$. Let $\underline{\chi}$
denote its image in $\kf^\perp$. Consider an $\slf_2$-triple $(e',h',f')$
corresponding to $\underline{\chi}$ and the $\theta$-stable  Slodowy slice $\underline{S}\subset \g^*$
associated to this triple. Let $\check{X}$ denote a connected component of the preimage of $\underline{S}$
in $X=\operatorname{Spec}(\C[\Orb])$, this is a transverse slice to $\Orb_K'$ in $X$  and a 4-dimensional conical 
symplectic singularity. Also note that $\check{X}\subset X$ is $\theta$-stable.

Let $\Orb'=G \cdot \Orb'_K$ be the $G$-saturation of $\Orb'_K$ in $\g^*$.
Set $\check{V}:=T_{\chi'}\Orb'$, this is a symplectic vector space. Set
$\breve{X}:=\check{X}\times \check{V}$. Taking the product of contracting torus
actions on $\check{X}$ and $\check{V}$ we get a contracting action on $\breve{X}$,
such that the weight of the Poisson bracket is $-2$. Also $\check{X},\check{V}$ come with
Poisson anti-involutions induced from $X$. The resulting Poisson anti-involution on
$\breve{X}$ will be denoted by $\theta$.
We set $\check{X}^\times:=\check{X} \setminus \{0\}, \breve{X}^\times=\check{X}^\times\times \check{V}$.

Using the quantum slice construction, Section \ref{SS_Dixmier_restriction}, from $\A$ we get a quantization
of $\check{X}$ to be denoted by $\check{\A}$. Note that in Section \ref{SS_Dixmier_restriction},
the quantum slice algebra was denoted by $\underline{\A}$ but from now on we will use the notation
$\underline{\A}$ for the quantum slice associated to the open orbit $\Orb$. We let $\check{\A}_\hbar$
denote the Rees algebra of $\check{\A}$.

We form the quantization
$\breve{\A}_\hbar:=\check{\A}_\hbar\otimes_{\C[\hbar]}\Weyl_\hbar(\check{V})$ of $\breve{X}$.
Let $\breve{\A}_\hbar^\times$ (resp. $\check{\A}_\hbar^\times$) be the microlocalization of the $\hbar$-adic completion of
$\breve{\A}_\hbar$ (resp. $\check{\A}_\hbar$) to $\breve{X}^\times$ (resp. $\check{X}^\times$).


We write $X^{\wedge_{\chi'}}$ for the spectrum of the completion of
$\C[X]$ at $\chi'$. It coincides with $\breve{X}^{\wedge_0}$ and, by the construction
in Section \ref{SS_Dixmier_restriction}, we have
an isomorphism of quantizations $\A_\hbar^{\wedge_{\chi'}}\cong
\breve{\A}_\hbar^{\wedge_0}$.
Our first task is to define an object $\breve{\Ecal}_\hbar\in \HC^{gr}(\breve{\A}_\hbar,\theta)$ %\footnote{Again it seems to be $\wHC^{gr}$ instead of $\HC^{gr}$. -- Shilin;
%Well we only use this construction for flat modules, so it can stay as is -- I.L.}
whose restriction to $\breve{X}^\times\cap \breve{X}^{\wedge_0}$ coincides with
the restriction of $\Ecal_\hbar$ to $\Orb\cap X^{\wedge_{\chi'}}$.

 Take $n\in \Z_{>0}$.
Consider the quotient $\mathcal{E}_{\hbar,n}:=\mathcal{E}_\hbar/\hbar^n\mathcal{E}_\hbar$. This is
a weakly HC $(\hat{\A}_\hbar^\Orb,\theta)$-module. Since $\mathcal{E}_{\hbar,0}$ is a vector
bundle on $\Orb_K$ and $\operatorname{codim}_{\Orb_K^2}\partial \Orb_K\geqslant 2$,
the sheaf $\iota_* \mathcal{E}_{\hbar,0}$ is coherent. It follows that
$\iota_*\mathcal{E}_{\hbar,n}$ is coherent as well. Consider the completion
$(\iota_*\mathcal{E}_{\hbar,n})^{\wedge_{\chi'}} \in \wHC^{gr}(\A_\hbar^{\wedge_{\chi'}},\theta)$.
By Lemma  \ref{Lem:wHC_equivalence}, $(\iota_*\mathcal{E}_{\hbar,n})^{\wedge_{\chi'}}$
decomposes as
\begin{equation}\label{eq:decomposition11}[(\iota_*\mathcal{E}_{\hbar,n})^{\wedge_{\chi'}}]^{\check{L}}\widehat{\otimes}_{\C[[\hbar]]}
\C[[\check{L},\hbar]],
\end{equation}
where $\check{L}=\check{V}^{-\theta}$. This decomposition is $\C^\times$-equivariant
with respect to the induced $\C^\times$-action on $\check{L}$. We define
the weakly HC $(\breve{\A}_\hbar^{\wedge_0}, \theta)$-module $\breve{\Ecal}_{\hbar,n}'$ as (\ref{eq:decomposition11}), where we redefine
the $\C^\times$-action on $\C[[\check{L},\hbar]]$ to come from the dilation action.
We also define the weakly HC $(\check{\A}_\hbar^{\wedge_0}, \theta)$-module $\check{\Ecal}_{\hbar,n}'$ as the first factor in (\ref{eq:decomposition11}).
Let $\breve{\Ecal}_{\hbar,n},\check{\Ecal}_{\hbar,n}$ denote the $\C^\times$-finite parts
in $\breve{\Ecal}_{\hbar,n}', \check{\Ecal}_{\hbar,n}'$, so that they are weakly HC $(\breve{\A}_\hbar, \theta)$- and $(\check{\A}_\hbar, \theta)$-modules respectively. So we have
\begin{equation}\label{eq:decomposition115}
\breve{\Ecal}_{\hbar,n}=\check{\Ecal}_{\hbar,n}\otimes_{\C[\hbar]}\C[L,\hbar].
\end{equation}

Now we proceed to constructing a sheaf $\breve{\Ecal}^\times_\hbar$ of HC $(\breve{\A}_\hbar^\times, \theta)$-modules. Consider the restriction
$\breve{\Ecal}_{\hbar,n}^\times$ of $\breve{\Ecal}_{\hbar,n}$ to $\breve{X}^\times$.
It follows from the construction of $\breve{\Ecal}_{\hbar,n}$
that the kernel and the cokernel of
$$\breve{\Ecal}_{\hbar,n}/\hbar^{n-1}\breve{\Ecal}_{\hbar,n}
\rightarrow \breve{\Ecal}_{\hbar,n-1}$$
are supported on $\{0\}\times \check{L}$. Indeed, a similar claim holds for
the sheaves $\Ecal_{\hbar,n}$, hence for $(\iota_*\mathcal{E}_{\hbar,n})^{\wedge_{\chi'}}$,
hence for $\breve{\Ecal}_{\hbar,n}'$, hence for $\breve{\Ecal}_{\hbar,n}$.
Therefore
\begin{equation}\label{eq:sheaf_map}\breve{\Ecal}^\times_{\hbar,n}/\hbar^{n-1}\breve{\Ecal}^\times_{\hbar,n}
\xrightarrow{\sim} \breve{\Ecal}^\times_{\hbar,n-1}\end{equation}
We set
$$\breve{\Ecal}_\hbar^\times:=\varprojlim \breve{\Ecal}^\times_{\hbar,n}.$$
Similarly, we can define the sheaf $\check{\Ecal}^\times_\hbar$ of $(\check{\A}_\hbar^\times, \theta)$-modules. Observe that, by the construction, the sheaves $\breve{\Ecal}^\times_\hbar$
and $\check{\Ecal}^\times_\hbar$ are flat over $\C[[\hbar]]$ and hence HC.


Note that, thanks
to (\ref{eq:sheaf_map}) and its analog for $\check{\Ecal}_\hbar$,
$$\breve{\Ecal}^\times_{\hbar}/\hbar^n \breve{\Ecal}^\times_\hbar\xrightarrow{\sim}
\breve{\Ecal}^\times_{\hbar,n}, \quad \check{\Ecal}^\times_{\hbar}/\hbar^n \check{\Ecal}^\times_\hbar\xrightarrow{\sim}
\check{\Ecal}^\times_{\hbar,n}.$$
And since
$(\iota_*\mathcal{E}_{\hbar,n})^{\wedge_{\chi'}}\cong \breve{\Ecal}_{\hbar,n}^{\wedge_0}$,
we conclude that
\begin{itemize}
\item
the restrictions of $\Ecal_\hbar,\breve{\Ecal}^\times_\hbar$
to $X^{\wedge_{\chi'}}\cap \Orb\cong \breve{X}^{\wedge_0}\cap \breve{X}^\times$
are isomorphic.
\end{itemize}
%Let $X^{\wedge_{\chi'}}$ denote
%the spectrum of the completion $\C[\Orb]^{\wedge_{\chi'}}$. Note that
%the algebra  $\C[\Orb]^{\wedge_{\chi'}}$ splits into the completed tensor
%product of $\C[\check{X}]^{\wedge_0}$ and $\C[\check{V}]^{\wedge_0}$, where we write
%$V'$ for the tangent space $T_{\chi'}(G\chi')$. Similarly, the completion
%$\A_\hbar^{\wedge_{\chi'}}$ splits into the completed tensor
%product of $\check{\A}_\hbar^{\wedge_0}$ and the completed formal Weyl algebra
%of $\check{V}$. Consider the completion of $(\iota_*\mathcal{E}_{\hbar,n})^{\wedge_\chi}$.
%It also splits into the completed tensor product by Lemma
%\ref{Lem:wHC_equivalence}. Let $\check{E}'_{\hbar,n}$ denote
%the factor corresponding to $\check{\A}_\hbar^{\wedge_0}$. This is a weakly HC
%$(\check{\A}_\hbar^{\wedge_0},\theta)$-module. It comes with an Euler derivation $\mathsf{eu}$
%compatible with the $\C^\times$-action on $\check{\A}^{\wedge_0}_\hbar$.
%Let $\check{E}_{\hbar,n}$ denote the $\mathsf{eu}$-finite part of
%$\check{E}'_{\hbar,n}$ so that the latter is the completion of the
%former. It follows from the construction that
%the kernel and the cokernel of the natural map
%$$\check{E}_{\hbar,n}/\hbar^{n-1}\check{E}_{\hbar,n}
%\rightarrow \check{E}_{\hbar,n-1}$$
%are supported at $0\in \underline{X}.$ Let $\check{\mathcal{E}}_{\hbar,n}$
%denote the localization of $\check{E}_{\hbar,n}$ to $\check{X}\setminus \{0\}$.
%Then we have $\check{\mathcal{E}}_{\hbar,n}/\hbar^{n-1}\check{\mathcal{E}}_{\hbar,n}
%\xrightarrow{\sim} \check{\mathcal{E}}_{\hbar,n-1}$. Set
%$\check{\mathcal{E}}_\hbar:=\varprojlim \check{\mathcal{E}}_{\hbar,n}$.
%This is a HC bimodules over the microlocalization of $\check{\A}_\hbar$ to
%$\check{X}\setminus \{0\}$.

%\begin{Rem}
%Similarly, we can define the restriction of $\mathcal{E}_\hbar$ to $X^{\wedge_\chi}\cap \Orb$.
%Namely, we take the $\mathsf{eu}$-finite part in
%$(\iota_*\mathcal{E}_{\hbar,n})^{\wedge_\chi}$ and proceed similarly to what we have done above.
%We denote the resulting sheaf by $\mathcal{E}_\hbar|_{X^{\wedge_\chi}\cap \Orb}$.
%\end{Rem}

Here is a criterion for $\iota_*\mathcal{E}_\hbar$ to be coherent.

\begin{Lem}\label{Lem:coh_extension1}
The following two conditions are equivalent:
\begin{enumerate}
\item $\iota_*\mathcal{E}_\hbar$ is coherent.
\item For each codimension two orbit $\Orb'_K$ in $\overline{\Orb}_K$,
the restriction of $\Gamma(\check{\mathcal{E}}^\times_\hbar)$ (where $\check{\mathcal{E}}^\times_\hbar$
is constructed for that orbit) to $\check{X}^\times$ coincides with $\check{\Ecal}^\times_\hbar$.
\end{enumerate}
\end{Lem}
\begin{proof}
Consider the third condition:
\begin{itemize}
\item[(3)] For each codimension two orbit $\Orb'_K$ in $\overline{\Orb}_K$,
the restriction of $\Gamma(\mathcal{E}_\hbar|_{X^{\wedge_{\chi'}}\cap \Orb})$ to $X^{\wedge_{\chi'}}\cap\Orb$ coincides with $\mathcal{E}_\hbar|_{X^{\wedge_{\chi'}}\cap \Orb}$.
\end{itemize}
The equivalence of (1) and (3) is established in exactly the same way as in
the proof of \cite[Lemma 3.4]{HC_symp}. Now we explain why (2)$\Leftrightarrow$(3).

By the discussion before the lemma,
$$\mathcal{E}_\hbar|_{X^{\wedge_{\chi'}}\cap \Orb}\cong \breve{\Ecal}^\times_\hbar|_{\breve{X}^{\wedge_0}\cap\breve{X}^\times}.$$
Note that (3) is equivalent to the claim that the kernel and the cokernel
of
$$\Gamma(\breve{\Ecal}^\times_\hbar|_{\breve{X}^{\wedge_0}\cap\breve{X}^\times})/
\hbar^n \Gamma(\breve{\Ecal}^\times_\hbar|_{\breve{X}^{\wedge_0}\cap\breve{X}^\times})
\rightarrow \Gamma(\breve{\Ecal}^\times_{\hbar,n}|_{\breve{X}^{\wedge_0}\cap\breve{X}^\times})$$
are supported on $\{0\}\times L^{\wedge_0}$.

Now we are going to pass to the subspaces of $\C^\times$-finite elements.  Note that
$$\Gamma(\breve{\Ecal}^\times_\hbar|_{\breve{X}^{\wedge_0}\cap\breve{X}^\times})=
\varprojlim \Gamma(\breve{\Ecal}^\times_{\hbar,n}|_{\breve{X}^{\wedge_0}\cap\breve{X}^\times}).$$
Arguing as in Step 2 of the proof of \cite[Lemma 3.4]{HC_symp}, we see that
$\Gamma(\breve{\Ecal}^\times_{\hbar,n}|_{\breve{X}^{\wedge_0}\cap\breve{X}^\times})=
\Gamma(\breve{\Ecal}^\times_{\hbar,n})^{\wedge_0}$.
Therefore the subspace of $\C^\times$-finite elements in $\Gamma(\breve{\Ecal}^\times_{\hbar,n}|_{\breve{X}^{\wedge_0}\cap\breve{X}^\times})$
coincides with $\Gamma(\breve{\Ecal}^\times_{\hbar,n})$. Since the $\C^\times$-action on
$\breve{X}$ is contracting, we see that
\begin{equation}\label{eq:completion_coincide11}
\Gamma(\mathcal{E}_\hbar|_{X^{\wedge_{\chi'}}\cap \Orb})\cong \Gamma(\breve{\mathcal{E}}_\hbar)^{\wedge_0}.
\end{equation}

So (3) is equivalent to the claim that
the kernel and the cokernel of
$$\Gamma(\breve{\Ecal}^\times_\hbar)/
\hbar^n \Gamma(\breve{\Ecal}^\times_\hbar)
\rightarrow \Gamma(\breve{\Ecal}^\times_{\hbar,n}).$$
to be supported on $\{0\}\times L$.

Let us show that $\Gamma(\breve{\Ecal}^\times_{\hbar,n})=
\Gamma(\check{\Ecal}^\times_{\hbar,n})\otimes_{\C[\hbar]}\C[L,\hbar]$.
We cover $\check{X}^\times$ with principal open affine subsets
$\check{X}_{f_i}\subset \check{X}$. Then $\Gamma(\check{\Ecal}^\times_{\hbar,n})$
is the zeroth cohomology of the \v{C}ech complex for $\check{\Ecal}_{\hbar,n}$
associated with the cover $\check{X}^\times=\bigcup \check{X}_{f_i}$, its terms
are direct sums of localizations of $\check{\Ecal}_{\hbar,n}$. Similarly,
$\Gamma(\breve{\Ecal}^\times_{\hbar,n})$ is the zeroth cohomology of the \v{C}ech
complex associated to the cover $\breve{X}^\times=\bigcup \breve{X}_{f_i}$. This \v{C}ech
complex is obtained from the previous one by tensoring with $\C[L,\hbar]$ over $\C[\hbar]$.
Since $\C[L,\hbar]$ is flat over $\C[\hbar]$ our claim follows.

Since
$\Gamma(\breve{\Ecal}^\times_\hbar)=\varprojlim \Gamma(\breve{\Ecal}^\times_{\hbar,n})$, we see that
\begin{equation}\label{eq:decomposition12}
\Gamma(\breve{\Ecal}^\times_{\hbar})=
\Gamma(\check{\Ecal}^\times_\hbar)\widehat{\otimes}_{\C[[\hbar]]}\C[L][[\hbar]].
\end{equation}
It follows that (3) is equivalent to the claim that the kernels and the cokernels of
$$\Gamma(\check{\Ecal}^\times_\hbar)/
\hbar^n \Gamma(\check{\Ecal}^\times_\hbar)
\rightarrow \Gamma(\check{\Ecal}^\times_{\hbar,n})$$
are finite dimensional for all $n$. But this condition is equivalent to (2).

%Recall that $\check{E}'_{\hbar,n}$, by the construction,
%is the centralizer of $L$ in  $(\iota_*\mathcal{E}_{\hbar,n})^{\wedge_\chi}$,
%while $(\iota_*\mathcal{E}_{\hbar,n})^{\wedge_\chi}=\C[[L]]\widehat{\otimes}
%\check{E}'_{\hbar,n}$. We have  $(\iota_*\mathcal{E}_{\hbar,n})^{\wedge_\chi}\cong %\Gamma(\mathcal{E}_{\hbar,n}|_{X^{\wedge_\chi}\cap \Orb})$.
%From here we deduce that $\Gamma(\check{\mathcal{E}}_{\hbar,n})^{\wedge_0}=
%\check{E}'_{\hbar,n}$.
%So
%$$\Gamma(\check{\mathcal{E}}_{\hbar})^{\wedge_0}=\varprojlim
%\Gamma(\check{\mathcal{E}}_{\hbar,n})^{\wedge_0}= \varprojlim
%\check{E}'_{\hbar,n}.$$
%Also we have $$\Gamma(\mathcal{E}_\hbar|_{X^{\wedge_\chi}\cap \Orb})=
%\varprojlim(\iota_*\mathcal{E}_{\hbar,n}^{\wedge_{\chi'}}).$$
%Since $\iota_*\mathcal{E}_{\hbar,n})^{\wedge_{\chi'}}=\C[[L]]\otimes \check{E}'_{\hbar,n}$,
%we see that
%$$\Gamma(\check{\mathcal{E}}_{\hbar})^{\wedge_0}=
%\C[[L]]\widehat{\otimes} \Gamma(\check{\mathcal{E}}_{\hbar})^{\wedge_0}.$$
%Note that (2) is equivalent to the claim that the cokernel of
%$\Gamma(\check{\mathcal{E}}_{\hbar})^{\wedge_0}\rightarrow \check{E}_{\hbar,0}$
%is finite dimensional. Similarly, (3) is equivalent to the claim that
%the cokernel of  $\Gamma(\check{\mathcal{E}}_{\hbar})^{\wedge_0}
%\rightarrow \iota_*\mathcal{E}_{\hbar,0}^{\wedge_{\chi'}}$
%has finite length over $\C[[L]]$. An equivalence between (2) and (3)
%follows from the observation that
%$\iota_*\mathcal{E}_{\hbar,0}^{\wedge_{\chi'}}=\C[[L]]\widehat{\otimes} \check{E}_{\hbar,0}$.
%%The proof closely follows that of \cite[Lemma 3.4]{HC_symp}.
%%If in (2) we have used
%%$\mathcal{E}_\hbar|_{X^{\wedge_\chi}\cap \Orb}$ instead of $\check{\mathcal{E}}_\hbar$,
%%then the proof would be literally the same. Let us explain why we can use $\check{\mathcal{E}}_\hbar$
%%instead of $\mathcal{E}_\hbar|_{X^{\wedge_\chi}\cap \Orb}$. For (1)$\Rightarrow$(2),
%%this is easy.
%%because under the equivalence
%%between wHC modules for $\check{\A}^{\wedge_0}_\hbar$ and $\A_\hbar^{\wedge_{\chi'}}$
\end{proof}





\subsection{Extension outside of codimension $2$}
Now suppose that the equivalent conditions of Lemma
\ref{Lem:coh_extension1} hold. The following claim
is an analog of \cite[Lemma 3.5]{HC_symp}. Note that the conditions
in the statement are more restrictive: unlike in the setting of
\cite{HC_symp} not all orbits in $\overline{\Orb}_K\setminus \Orb_K^2$ have even
codimension so we can get an orbit of codimension $3$, while in the setting
of \cite{HC_symp} the minimal possible codimension is $4$.

\begin{Lem}\label{Lem:coh_extension2}
Suppose that the following assumption holds for all codimension $2$ orbits
in $\overline{\Orb}_K$ (and the corresponding sheaves $\check{\mathcal{E}}_\hbar$):
\begin{itemize}
\item[(*)] $\Gamma(\check{\mathcal{E}}_\hbar)/
\hbar \Gamma(\check{\mathcal{E}}_\hbar)\xrightarrow{\sim}\Gamma(\check{\mathcal{E}}_{\hbar}/\hbar \check{\mathcal{E}}_{\hbar})$.
%\item[(ii)] $\Gamma(\check{\mathcal{E}}_{\hbar}/\hbar \check{\mathcal{E}}_{\hbar})$ is a maximal
%Cohen-Macaulay module over $\C[\check{X}^\theta]$.
\end{itemize}
Then the microlocalization of $\Gamma(\mathcal{E}_\hbar)$ to $\Orb_K$
coincides with $\mathcal{E}_\hbar$.
\end{Lem}
\begin{proof}
Combining (\ref{eq:completion_coincide11}) and (\ref{eq:decomposition12}), we see that condition (*) is equivalent to the claim that
\begin{equation}\label{eq:quot_mod_h}
\Gamma(\mathcal{E}_\hbar|_{X^{\wedge_{\chi'}}\cap \Orb})/
\hbar \Gamma(\mathcal{E}_\hbar|_{X^{\wedge_{\chi'}}\cap \Orb})\xrightarrow{\sim}
\Gamma(\mathcal{E}_{\hbar,0}|_{X^{\wedge_{\chi'}}\cap \Orb}).
\end{equation}
In its turn, (\ref{eq:quot_mod_h}) for all codimension $2$ orbits imply that
$$(\iota_*\mathcal{E}_\hbar)/\hbar (\iota_*\mathcal{E}_\hbar)
\xrightarrow{\sim} \iota_*(\mathcal{E}_\hbar/\hbar \mathcal{E}_\hbar).$$

Note that $\check{\mathcal{E}}_{\hbar}/\hbar \check{\mathcal{E}}_{\hbar}$ is a vector bundle
on $\check{X}^\theta\setminus\{0\}$. From \cite[Corollary 3.12]{Burban_Drozd} we see that
$\Gamma(\check{\mathcal{E}}_{\hbar}/\hbar \check{\mathcal{E}}_{\hbar})$ is a maximal Cohen-Macaulay
module over $\C[\check{X}^\theta\setminus\{0\}]$, the normalization of $\C[\check{X}^\theta]$.
By \cite[Lemma 2.24]{Burban_Drozd},  $\Gamma(\check{\mathcal{E}}_{\hbar}/\hbar \check{\mathcal{E}}_{\hbar})$ is Cohen-Macaulay over $\C[\check{X}^\theta]$.

We conclude that $\iota_*(\mathcal{E}_\hbar/\hbar \mathcal{E}_\hbar)$
is a maximal Cohen-Macaulay module over $\Orb_K^2$. Since the complement to
$\Orb_K^2$ in $\overline{\Orb}_K$ has codimension $\geqslant 3$, we see that
$H^1(\Orb_K^2, \iota_*(\mathcal{E}_\hbar/\hbar \mathcal{E}_\hbar))$
is a finitely generated module over $\C[\overline{\Orb}_K]$. Now we can
finish the proof as in \cite[Lemma 3.5]{HC_symp}.
\end{proof}

\begin{comment}
\begin{Rem}\label{rem:sufficient_conditions}
Lemma \ref{Lem:coh_extension2} can be restated as follows.
Suppose that for
each codimension $2$ orbit in $\overline{\Orb}_K$ we have the following:
\begin{itemize}
\item We have $\Gamma(\check{\mathcal{E}}_\hbar)/\hbar
\Gamma(\check{\mathcal{E}}_\hbar)\xrightarrow{\sim} \Gamma(\check{\mathcal{E}})$.
\end{itemize}
Then $\Gamma(\mathcal{E}_\hbar)$ restricts to $\mathcal{E}_\hbar$.
\end{Rem}
\end{comment}
It turns out that (*) is not convenient to check, as it is hard to determine $\check{\mathcal{E}}^\times_\hbar$.
It is easy to determine $\check{\mathcal{E}}^\times:=\check{\mathcal{E}}^\times_\hbar/\hbar \check{\mathcal{E}}^\times_\hbar$: we will see in the next section that it is obtained by restricting the twisted local system $\mathcal{E}:=\mathcal{E}_\hbar/\hbar \mathcal{E}$ to the slice.
In the next section we will also see that $\check{\mathcal{E}}^\times$ admits a unique quantization to
an $\check{\A}_\hbar^\times$-module if we assume that the $\C^\times$-action on $\mathcal{E}$ is strong.


\subsection{Restrictions to slices}
We adopt the settings and notations of Section \ref{subsec:extn_codim2}.
Take an object $\mathcal{E}_\hbar \in \wHC^{gr}(\hat{\A}_\hbar^{\Orb},\theta)^{K,\kappa}$.  By Section \ref{subsec:lag_vec}, the sheaf $\EE = \EE_\hbar / \hbar \EE_\hbar$ is a coherent $(K,\kappa)$-equivariant $\TT^+(\hat{\A}_\hbar^\Orb, \Orb_K)$-module and hence a $K$-equivariant vector bundle over $\Orb_K$.
The following claim follows by combining Remark \ref{Rem:homog_analyt_classif} with Proposition \ref{prop:irred_strongness}. 

\begin{Lem}\label{Lem:irred_strongness_HC} 
	Suppose the geometric fiber of $\mathcal{E}_\chi$ at $\chi \in \Orb_K$ is an irreducible $K$-module.
	Then the induced $\cm$-action on $\EE$ is strong in the sense of Definition \ref{defn:strong_cm}.
\end{Lem}

Let $\check{Y} = \check{X}^\theta$ and $\check{Y}^\times = \check{Y} \setminus \{ 0 \} = (\check{X}^\times)^\theta$. In Section \ref{subsec:extn_codim2} we have produced a module $\check{\Ecal}^\times_\hbar$ over $\check{X}^\times$ from $\EE_\hbar$, which is supported on $\check{Y}^\times$, so that $\check{\mathcal{E}}^\times = \check{\mathcal{E}}_\hbar^\times/\hbar \check{\mathcal{E}}_\hbar^\times$ is a graded $\TT^+(\check{\A}_\hbar^\times, \check{Y}^\times)$-module,
where the action is restricted from the action $t.m:=t^{-2}\gamma'(t)m$ on $\mathcal{E}$, where $\gamma$ is the one-parameter subgroup in $G$ corresponding to $h'$.
%We want to know if the induced $\cm$-action\footnote{I.L.: perhaps, earlier I should recall what this %action is.} on $\check{\mathcal{E}}^\times$ is also strong, since if this is true, then %$\check{\mathcal{E}}^\times_\hbar$ can be uniquely recovered from
%$\check{\mathcal{E}}^\times := \check{\mathcal{E}}^\times_\hbar/\hbar \check{\mathcal{E}}^\times_\hbar$ %by Theorem \ref{thm:quan_lag}.

Here is the main result of this section.

\begin{Prop}\label{Prop:slice_strongness}
We have the following results:
	\begin{enumerate}
		\item
		The Picard algebroid  $\TT^+(\check{\A}_\hbar^\times, \check{Y}^\times)$ is naturally isomorphic to the pullback of the Picard algebroid $\TT^+(\hat{\A}_\hbar^{\Orb}, \Orb_K)$ via the natural closed embedding $\iota_{\check{Y}^\times} : \check{Y}^\times \hookrightarrow \Orb_K$. The similar result holds for $\TT(\check{\A}_\hbar^\times, \check{Y}^\times)$ and $\TT(\hat{\A}_\hbar^{\Orb}, \Orb_K)$.
		\item
		 Let $\mathcal{E}_\hbar\in\wHC^{gr}(\hat{\A}_\hbar^{\Orb},\theta)^{K,\kappa}$ and keep the notations in Section \ref{subsec:extn_codim2}. Then the $\TT^+(\check{\A}_\hbar^\times, \check{Y}^\times)$-module $\check{\mathcal{E}}^\times$ is naturally isomorphic to the pullback $\iota_{\check{Y}^\times}^* \EE$ of the $\TT^+(\hat{\A}_\hbar^{\Orb}, \Orb_K)$-module $\EE$ with respect to the isomorphism in (1).
		\item
		If the $\cm$-action on $\EE$ is strong, so is the induced $\cm$-action on $\check{\mathcal{E}}^\times$.
	\end{enumerate}	
\end{Prop}

To prove this proposition, we need to first understand the relation between the Picard algebroids $\TT^+(\hat{\A}_\hbar^\Orb, \Orb_K)$ and $\TT^+(\check{\A}_\hbar^\times, \check{Y}^\times)$ (as well as $\TT(\hat{\A}_\hbar^\Orb, \Orb_K)$ and $\TT(\check{\A}_\hbar^\times, \check{Y}^\times)$ ). We start with defining global analogues of them. The following definition is similar to Definition \ref{defn:pic} of Picard algebroids. Let $\der(?)$ denote the Lie algebra of derivations of an algebra ``?''.

\begin{defi}
	Suppose $\B$ is a commutative (Noetherian) $\C$-algebra. A  {\it (weak) Picard Lie algebra} over $\B$ is a quadruple $(\T, [-,-], \eta_\T, 1_{\T})$, where
	\begin{itemize}
		\item
		$\T$ is a finitely generated $\B$-module;
		\item
		$[-,-]$ is a $\C$-linear Lie bracket;
		\item
		$\eta_\T: \T \to \der(\B)$ is a Lie algebra morphism and a morphism of $\B$-modules;
		\item
		$1_\T$ is a central element in the Lie algebra $\T$ which lies in the kernel of $\eta_\T$,
	\end{itemize}
	such that the Lie bracket $[-,-]$, $\eta_\T$ and the $\B$-module structure of $\T$ are compatible with each other, in the sense that they satisfy a Leibniz rule similar to that in the definition of a Lie algebroid and $\eta_\T$ behaves like an anchor map. Note that we have a $\B$-linear morphism $\iota_\T: \B \to \T$, $b \mapsto b \cdot 1_\T$, and the composite map $\B \xrightarrow{\iota_\T} \T \xrightarrow{\eta_\T} \der(\B)$ is zero. This sequence, however, does not have to be exact.
	
	As in the case of Picard algebroids, one can similarly define the notion of a morphism $\tau: \T \to \F$ between two Picard Lie algebras $\T=(\T, [-,-]_\T, \eta_\T, 1_{\T})$ and $\F=(\F, [-,-]_\F, \eta_\F, 1_{\F})$ over a base algebra $\B$, by requiring $\tau$ to be a $\B$-linear Lie algebra homomorphism that intertwines the anchor maps $\eta_\T$ and $\eta_{\F}$, and $\tau(1_\T)=1_\F$. Unlike the case of Picard algebroids, however, a morphism of Picard Lie algebras needs not to be an isomorphism. We write $\mathrm{wPA}(\B)$ for the category of all (weak) Picard algebras over $\B$.
\end{defi}

Given a homomorphism $f: \B \to \mathcal{C}$ of commutative algebras over $\C$, we can also define the pullback Picard Lie algebra $f^\circ \T$ over $\mathcal{B}$ of $\T$ as the fiber product of the natural maps
  \[  \der(\mathcal{C}) \to \mathcal{C} \otimes_{\B} \der(\B) \quad \text{and} \quad 1 \otimes \eta_\T :  \mathcal{C} \otimes_{\B} \T \to \mathcal{C} \otimes_{\B} \der(\B).\]
The central element $1_{f^\circ \T}$ is taken to be $(0, 1_\T)$ in this fiber product. The anchor map of $f^\circ \T$ is just the projection map $f^\circ \T \to \der(\mathcal{C})$. The Lie bracket of $f^\circ \T$ can be defined in a similar way as how pullbacks of Picard algebroids are defined. We have the obvious functoriality: if $\tau: \T \to \F$ is a morphism of Picard algebras over $\B$, then it induces a morphism $f^\circ(\tau): f^\circ \T \to f^\circ \F$ of Picard algebras over $\mathcal{C}$ in the obvious way, such that  $\T \mapsto f^\circ \T$ defines a functor $f^\circ: \mathrm{wPA}(\B) \to \mathrm{wPA}(\mathcal{C})$; moreover,
 if $f: \B \to \mathcal{C}$ and $g: \mathcal{C} \to \Dcal$ are two homomoprhisms of commutative $\C$-algebras and $\T$ is a Picard algebra over $\B$, then we have a natural isomorphism $(g \circ f)^\circ  \simeq g^\circ \circ f^\circ : \mathrm{wPA}(\B) \to \mathrm{wPA}(\mathcal{\Dcal})$ of functors. We can hence take various localizations $\B \to S^{-1} \B$ of $\B$ and localize a Picard algebra $\T$ over $\B$ to a Lie algebroid over $\spec \B$ or any of its open subset.

One can also talk about modules over the Picard Lie algebra $(\T, [-,-], \eta_\T, 1_\T)$ similarly to modules over Picard algebroids. That is, a module over $\T$ is a $\B$-module $N$  equipped with a Lie algebra action of $\T$, which satisfies the Leibniz rule with respect to $\eta_\T$ and the $\B$-module structure such that $1_\T$ acts by the identity. Given a homomorphism $f: \B \to \mathcal{C}$ of commutative $\C$-algebras, for any $\T$-module $M$, the $\mathcal{C}$-module $f^*M := \mathcal{C} \otimes_{\B}  M$ is naturally a $f^\circ \T$-module. If $p: M \to N$ is a morphism of Picard algebras over $\B$, then it induces a morphism $f^*(p) : f^*M \to f^*N$ of modules over the Picard algebra $f^\circ \T$ over $\mathcal{C}$, so that $f^*(-)$ defines a functor from the category of modules over $\T$ to the category of modules over $f^\circ \T$.



\begin{comment}
We regard the tensor algebra $\Dcal_0 \otimes_\C \Dcal_0$ as a Possion algebra over $1  \otimes \Dcal_0$ by extending the Poisson bracket in the first factor to a $1 \otimes \Dcal_0$-linear Poisson bracket. Then taking the Possion bracket gives a $1 \otimes \Dcal_0$-linear Lie algebra morphism $\bar{\eta}: \Dcal_0 \otimes_\C \Dcal_0 \to \der(\Dcal_0) \otimes \Dcal_0$.

Set $\wt{\Dcal}_t^{(2)} := \wt{\Dcal}_t / t^2 \wt{\Dcal}_t $ and form the tensor algebra $\wt{\Dcal}_t^{(2)} \otimes \Dcal_0$. Then $\hbar^{-1} \wt{\Dcal}_t^{(2)} \otimes \Dcal_0$ is a Lie algebra with the Lie bracket $[-,-]_0$ given by taking commutators.  For any $f \in \wt{\Dcal}_t^{(2)} \otimes \Dcal_0$, let $\bar{f}$ denote its image under the projection $\wt{\Dcal}_t^{(2)} \otimes \Dcal_0 \twoheadrightarrow \Dcal_0 \otimes_\C \Dcal_0$ of algebras. Define a map $\epsilon: \der(\Dcal_0) \otimes \Dcal_0 \to \der(\Dcal_0)$ by $\epsilon(\xi \otimes g) = g \cdot \xi$ and a Lie algebra morphism $\eta: \hbar^{-1} \wt{\Dcal}_t^{(2)} \otimes \Dcal_0 \to \der(\Dcal_0)$ by $\eta(\hbar^{-1} f) = \hbar^{-1} \epsilon(\bar{\eta}(\bar{f}))$. We define a new Lie bracket on $\wt{\Dcal}_t^{(2)} \otimes \Dcal_0$   by
\begin{equation}\label{eq:Lie}
	[f, g] = [f, g]_0 + [1 \otimes \eta(f)] (g) - [1 \otimes \eta(g)] (f), \quad \forall \, f, g \in \wt{\Dcal}_t^{(2)} \otimes \Dcal_0,
\end{equation}
where $1 \otimes \xi$ for $\xi \in \der(\Dcal_0)$ means that $\xi$ acts on the second factor of $\wt{\Dcal}_t^{(2)} \otimes \Dcal_0$. The Lie bracket $[-,-]$ is compatible with the multiplication of of the algebra $\wt{\Dcal}_t^{(2)} \otimes \Dcal_0$, i.e., $[a, - ]$ acts as a derivation of $\wt{\Dcal}_t^{(2)} \otimes \Dcal_0$ for any $a \in \wt{\Dcal}_t^{(2)} \otimes \Dcal_0$. Therefore $[-,-]$ is not $1 \otimes \Dcal_0$-linear, but satisfies the Leibniz rule with respect to the $1 \otimes \Dcal_0$-module structure. This is an analog of the Lie bracket $[-,-]^t$ in \cite[Section 4.3]{Yu}.

Let $\I$ denote the kernel of the multiplication map $\Dcal_0 \otimes_\C \Dcal_0 \to \Dcal_0$ and let $\I_t$ denote the preimage of $\I$ under the projection $\wt{\Dcal}_t^{(2)} \otimes \Dcal_0 \twoheadrightarrow \Dcal_0 \otimes_\C \Dcal_0$ of algebras. Set
\[ \I_t \bullet \I_t := \operatorname{Span} \left\langle a \bullet b = \frac{1}{2}(ab + ba) \,|\, a, b \in \I_t \right\rangle \subset \I_t. \]
A similar argument as in the proof of \cite[Lemma 4.9]{Yu} shows that $\hbar^{-1} (1 \otimes \Dcal_0 +  \I_t \bullet \I_t )$ is Lie ideal under the Lie bracket $[-,-]$. Now set
\[ \TT_{\wt{\Dcal}_t} := \hbar^{-1} [ \wt{\Dcal}_t^{(2)} \otimes \Dcal_0 / (1 \otimes \Dcal_0 +  \I_t \bullet \I_t )],  \]\footnote{Again, I'm confused about using $\hbar$ instead of $t$. S.Y.: I've replaced $\TT_{\Dcal_\hbar}$ by $\TT_{\wt{\Dcal}_t}$.}
then $\TT_{\wt{\Dcal}_t}$ is naturally a finitely generated $\Dcal_0$-module. The Lie bracket $[-, -]$ in \eqref{eq:Lie} descends to a Lie bracket on $\TT_{\wt{\Dcal}_t}$, still denoted by $[-,-]$. The Lie algebra morphism $\eta$ descends to a Lie algebra morphism $\TT_{\wt{\Dcal}_t} \to \der(\Dcal_0)$, which is $\Dcal_0$-linear. There is a natural $\Dcal_0$-module morphism $\Dcal_0 \to \TT_{\wt{\Dcal}_t}$ induced by the inclusion $1 \otimes \Dcal_0 = \hbar^{-1} \cdot t (1 \otimes \Dcal_0 ) \hookrightarrow \hbar^{-1} \wt{\Dcal}_t^{(2)} \otimes \Dcal_0$, such that the image of $\Dcal_0$ lies in the kernel of $\eta$ and the image of $1_{\Dcal_0}$ is a central element in the Lie algebra $\TT_{\wt{\Dcal}_t}$, which is still denoted as $1_{\TT_{\wt{\Dcal}_t}}$ or simply $1$. Then the quadruple $(\TT_{\wt{\Dcal}_t}, [-,-], \eta, 1_{\TT_{\wt{\Dcal}_t}})$ forms a weak Picard Lie algebra over $\Dcal_0$. It is easy to check that this construction is compatible with localizations\footnote{I.L.: I don't quite understand what you mean here -- what you have is manifestly a sheaf.  S.Y. : I agree it is obvious. It is just one-sentence remark. }.






We can apply the construction above to the case when $\Dcal_\hbar$ is the quantization $\A_\hbar$ of $\C[\Orb]$ or the quantum slice $\check{\A}_\hbar$, with their canonical lattices. Localizing the Lie algebra $\TT_{\wt{\Dcal}_t}$ in question to the open subsets $\Orb$ or $\check{X}^\times$, we get the Picard algebroids $\overline{\per}(\hat{\A}_\hbar^{\Orb})$ or $\overline{\per}(\check{\A}_\hbar^\times)$. The same construction can also be carrried out when $\Dcal_\hbar$ is the completed algebra $\A_\hbar^{\wedge_{\chi'}} \cong \breve{\A}_\hbar^{\wedge_0}$ or $\check{\A}_\hbar^{\wedge_0}$ ($\Dcal_0$ is no longer finitely generated, but still Noetherian). The only modification is to take completed tensor products, e.g., $\wt{\Dcal}_t^{(2)} \widehat{\otimes} \Dcal_0 = (\wt{\Dcal}_t / t^2 \wt{\Dcal}_t)  \widehat{\otimes}  \Dcal_0$, etc. The Lie algebra $\TT_{\wt{\Dcal}_t}$ thus obtained is just the completion of the Picard Lie algebra assoicated to $\A_\hbar$ or $\check{\A}_\hbar$\footnote{Say something about the module of derivations in the completed case. e.g, it is finitely generated, the derivations are automatically continuous, etc. If we add a section about quantizaiton on `completed schemes', we can address these issues there. -- Shilin; I.L.: to add an explanation as a remark is a good idea, but perhaps we
	shouldn't add that section.}. Note that $\TT_{\wt{\Dcal}_t}$ inherits a natural $\cm$-action from that of $\wt{\Dcal}_t$ (here we take the diagonal $\cm$-action on the tensor product in the construction above). We remark that, when $\Dcal_\hbar = \check{\A}_\hbar^{\wedge_0}$, the subspace of $\cm$-finite vectors in $\TT_{ \check{\A}_\hbar^{\wedge_0}}$ is exactly $\TT_{\check{\A}_\hbar}$.


%If $\mathcal{C}_\hbar$ is another $\C[\hbar]$-algebra satisfying the same conditions of $\Dcal_\hbar$ and $f: \mathcal{\C}_\hbar \to \Dcal_\hbar$ is a homomorphism $\C[\hbar]$-algebras preserving the lattice structures, then $f$
\end{comment}

We adopt the setting and notations of Section \ref{SS_weakly_HC}. In particular, let $\Dcal_\hbar$ be an associate flat $\C[\hbar]$-algebra. Assume that $\Dcal_0 = \Dcal_\hbar / \hbar \Dcal_\hbar$ is a Noetherian commutative $\C$-algebra. Then the map $\Dcal_\hbar \times \Dcal_\hbar \to \Dcal_\hbar : (a,b) \mapsto \hbar^{-1}[a,b]$ descends to a Possion bracket on $\Dcal_0$. As in Section \ref{SS_weakly_HC}, suppose $\theta$ is a $\C[\hbar]$-linear anti-involution on $\Dcal_\hbar$. We are mostly interested in the case when $\Dcal_\hbar$ is the quantization $\A_\hbar$ of $\C[\Orb]$ or the quantum slice $\check{\A}_\hbar$. In what follows, we will define  in a similar manner the affine global analogues for $(\Dcal_\hbar, \theta)$ of the Picard algebroids $\TT^+(\hat{\A}_\hbar^\Orb, \Orb_K)$ and $\TT^+(\check{\A}_\hbar^\times, \check{Y}^\times)$ introduced in Section \ref{subsec:lag_pic}. According to \ref{rmk:pic_mod_t2}, those Picard algebroids can be defined using only the quantization mod $\hbar^2$. Therefore it is natural to set $\Dcal_\hbar^{(2)} := \Dcal_\hbar / \hbar^2 \Dcal_\hbar$ to replace $\Dcal_\hbar$.

 Let $\J$ be the ideal in $\Dcal_0$ generated by $\Dcal_0^{-\theta} $, which can also be regarded as a Lie algebra via commutator. Set $\Dcal'_0: = \Dcal_0 / \J$. Let $\bar{\J}_\hbar$ be the preimage of $\J$ in $\Dcal_\hbar^{(2)}$ under the quotient map $\Dcal_\hbar^{(2)} \to \Dcal_0$. Then the subspace $\bar{\J}_\hbar^2$ is a Lie ideal of $ \bar{\J}_\hbar$ and we can form the Lie algebra $\hbar^{-1} (\bar{\J}_\hbar / \bar{\J}_\hbar^2)$. Since $\hbar \bar{\J}_\hbar \subset \bar{\J}_\hbar^2$ by definition, the natural inclusion $\hbar \Dcal_\hbar^{(2)}  \subset  \bar{\J}_\hbar$ descends to a map
   \[  \iota_\T: \Dcal'_0 \simeq \Dcal_\hbar / \bar{\J}_\hbar \to \hbar^{-1} (\bar{\J}_\hbar / \bar{\J}_\hbar^2). \]
The action of $\hbar^{-1} \bar{\J}_\hbar$ on $\Dcal_\hbar^{(2)}$ by taking commutators preserves $\bar{\J}_\hbar$ and hence induces a Lie algebra morphism
   \[ \eta_\T: \hbar^{-1} (\bar{\J}_\hbar / \bar{\J}_\hbar^2) \to \der(\Dcal_\hbar^{(2)} / \bar{\J}_\hbar) \simeq \der(\Dcal'_0). \]
 As in Section \ref{subsec:lag_pic}, let $\T^+(\Dcal_\hbar, \theta)$ and $\T(\Dcal_\hbar, \theta)$ denote the same Lie algebra $\hbar^{-1} (\bar{\J}_\hbar / \bar{\J}_\hbar^2)$, but with the two $\Dcal'_0$-module structures induced by the left multiplication of $\Dcal_\hbar^{(2)}$ on $\bar{\J}_\hbar$ and the symmetrized product respectively. Clearly, the  maps $\eta_\T$ and $\iota_\T$ are morphisms of $\Dcal'_0$-modules with respect to both $\Dcal'_0$-module structures on $\hbar^{-1} (\bar{\J}_\hbar / \bar{\J}_\hbar^2)$. Moreover, $\eta_\T \circ \iota_\T = 0$. It is then routine to check that both $\T^+(\Dcal_\hbar, \theta)$ and $\T(\Dcal_\hbar, \theta)$ with the structure maps  $\eta_\T$ and $\iota_\T$ are weak Picard Lie algebras over $\Dcal'_0$. All the constructions are compatible with $\cm$-actions. Again our constructions are compatible with localizations. This together with Remark \ref{rmk:pic_mod_t2} immediately implies the following lemma.

 \begin{Lem}\label{lem:Picard_loc}
 	The localizations of $\T^+( \A_\hbar, \theta)$ and $\T( \A_\hbar, \theta)$ to $\Orb_K$ are naturally isomorphic to the Picard algebroids $\TT^+(\hat{\A}_\hbar^\Orb, \Orb_K)$ and $\TT(\hat{\A}_\hbar^\Orb, \Orb_K)$ respectively. Similarly the localizations of $\T^+( \check{\A}_\hbar, \theta)$ and $\T( \check{\A}_\hbar, \theta)$ to $\check{Y}^\times$ are naturally isomorphic to the Picard algebroids $\TT^+(\check{\A}_\hbar^\times, \check{Y}^\times)$ and $\TT(\check{\A}_\hbar^\times, \check{Y}^\times)$ respectively.
 \end{Lem}

As in the case of quantization of vector bundles over Lagrangian subvarieties (see Section \ref{subsec:lag_vec}), if $M_\hbar$ is a HC $(\Dcal_\hbar, \theta)$-module, then the $\Dcal'_0$-module $M_0 := M_\hbar / \hbar M_\hbar$ is naturally a module over $\T^+(\Dcal_\hbar, \theta)$. Again all the constructions are compatible with $\cm$-actions.

\begin{comment}
\begin{Lem} \label{lem:restriction_period}\footnote{S.Y.: this lemma is only needed for the $a_2$ singularity, which is however taken care of by other arguments. So we probably don't need the definition of $\TT_{\wt{\Dcal}_t}$.}
	The enhanced period $\overline{\per}(\check{\A}_t^\times)$ is naturally isomorphic to the pullback of $\overline{\per}(\hat{\A}_t^{\Orb})$ as a Picard algebroid via the embedding $\iota_{\check{X}}: \check{X} \hookrightarrow \Orb$. In particular, $\per(\check{\A}_t^\times) = \iota_{\check{X}}^* [\per(\hat{\A}_t^{\Orb})] \in H^2_{DR}(\check{X})$.
\end{Lem}

\begin{proof}
	We have a decomposition $\A_\hbar^{\wedge_{\chi'}} \simeq  \check{\A}_\hbar^{\wedge 0} \widehat{\otimes}_{\C[[\hbar]]}  \Weyl_\hbar(\check{V} )^{\wedge_0} $ which is compatible with the lattice structures\footnote{Say something more.  -- Shilin. I.L.: perhaps.}. Set $\Dcal_\hbar = \A_\hbar^{\wedge_{\chi'}}$ and $\mathcal{C}_\hbar = \check{\A}_\hbar^{\wedge 0}$, then $\Dcal_0 = \Dcal_\hbar/\hbar\Dcal \simeq \C[X^{\wedge_{\chi'}}]$ and $\mathcal{C}_0 = \mathcal{C}_\hbar / \hbar \mathcal{C}_\hbar \simeq \C[\check{X}^{\wedge_0}]$. Let $f: \C[X] \to \C[\check{X}]$ be the homomorphism of algebras corresponding the inclusion $\check{X} \hookrightarrow X$\footnote{I.L.: it's weird to use the homomorphism of noncompleted algebras when you only deal with completions. S.Y.: In the second paragraph I deal with $f$.}. It induces a morphism $\hat{f}: \Dcal_0 \to \mathcal{C}_0$ for the completed algebras.
	
	The inclusion $\mathcal{C}_\hbar \widehat{\otimes} 1 \subset \mathcal{C}_\hbar \widehat{\otimes}_{\C[[\hbar]]}  \Weyl_\hbar(\check{V} )^{\wedge_0}  \simeq \Dcal_\hbar $ induces a morphism $\TT_{\wt{\mathcal{C}}_t} \to \mathcal{C}_0 \widehat{\otimes}_{\Dcal_0} \TT_{\wt{\Dcal}_t}$. Its composition with the map $\mathcal{C}_0 \widehat{\otimes}_{\Dcal_0} \TT_{\wt{\Dcal}_t} \to \mathcal{C}_0 \widehat{\otimes}_{\Dcal_0} \der(\Dcal_0)$, induced by the anchor map of  $\TT_{\wt{\Dcal}_t}$, agrees with the composite map $\TT_{\wt{\mathcal{C}_t}} \to \der( \mathcal{C}_0) \to \mathcal{C}_0 \widehat{\otimes}_{\Dcal_0} \der(\Dcal_0)$.
	Hence we have a natural $\cm$-equivariant morphism $\TT_{\wt{\mathcal{C}_t}} \to \hat{f}^\circ \TT_{\wt{\Dcal}_t}$ of Picard Lie algebras over $\mathcal{C}_0$. Note that $\hat{f}^\circ \TT_{\wt{\Dcal}_t}$ is nothing else but the completion of $f^\circ \TT_{\tilde{\A}_t}$. Taking $\cm$-finite vectors gives a $\cm$-equivariant morphism $\TT_{\check{\A}_t} \to f^\circ \TT_{\tilde{\A}_t}$ of Picard Lie algebras over $\C[\check{X}]$, whose localization to $\check{X}^{\times}$ is a $\cm$-equivariant morphism $\overline{\per}(\check{\A}_t^\times) \to \iota^\circ  \overline{\per}(\hat{\A}_t^{\Orb})$ of Picard algebroids and must be an isomorphism.
\end{proof}
\end{comment}





\begin{proof}[Proof of Proposition \ref{Prop:slice_strongness}]
  %The proof of (1) is similar to that of Lemma \ref{lem:restriction_period}.
  We first prove part (1). We have the decomposition $\A_\hbar^{\wedge_{\chi'}} \simeq  \check{\A}_\hbar^{\wedge 0} \widehat{\otimes}_{\C[[\hbar]]}  \Weyl_\hbar(\check{V} )^{\wedge_0}$. Set $\Dcal_\hbar = \A_\hbar^{\wedge_{\chi'}}$ and $\mathcal{C}_\hbar = \check{\A}_\hbar^{\wedge 0}$, then $\Dcal_0 = \Dcal_\hbar/\hbar\Dcal \simeq \C[X^{\wedge_{\chi'}}]$ and $\mathcal{C}_0 = \mathcal{C}_\hbar / \hbar \mathcal{C}_\hbar \simeq \C[\check{X}^{\wedge_0}]$. Let $\J_\mathcal{C}$ and $\J_\Dcal$ be the ideals in $\mathcal{C}_0$ and $\Dcal_0$ generated by $\mathcal{C}_0^{-\theta}$ and $\Dcal_0^{-\theta}$ respectively. Let $\bar{\J}_{\mathcal{C},\hbar}$ and $\bar{\J}_{\Dcal,\hbar}$ be the preimage of $\J_\mathcal{C}$ and $\J_\Dcal$ in $\mathcal{C}_\hbar^{(2)}$ and $\Dcal_\hbar^{(2)}$, respectively. Set $\mathcal{C}'_0: = \mathcal{C}_0 / \J_\mathcal{C}$ and $\Dcal'_0: = \Dcal_0 / \J_\Dcal$. We have a natural morphism $\Dcal_0 \to \mathcal{C}_0$ intertwining $\theta$'s, which is the completion of the morphism $\C[X] \to \C[\check{X}]$ corresponding the inclusion $\check{X} \hookrightarrow X$. This induces a morphism $\hat{f}: \Dcal_0' \to \mathcal{C}_0'$. It is the completion of the morphism $f: \C[X] / \langle \C[X]^{-\theta} \rangle \to \C[\check{X}] / \langle \C[\check{X}]^{-\theta} \rangle$.

  The subspace $\bar{\J}_{\mathcal{C},\hbar} \widehat{\otimes} 1 \subset \mathcal{C}_\hbar^{(2)} \widehat{\otimes}_{\C[[\hbar]]}  \Weyl_\hbar(\check{V} )^{\wedge_0}  \simeq \Dcal_\hbar^{(2)}$ lies inside $\bar{\J}_{\Dcal,\hbar}$. The inclusion map $\bar{\J}_{\mathcal{C},\hbar} \widehat{\otimes} 1 \hookrightarrow \bar{\J}_{\Dcal,\hbar}$ induces a $\mathcal{C}'_0$-linear map $\T^+(\mathcal{C}_\hbar, \theta) \to \mathcal{C}'_0 \otimes_{\Dcal'_0} \T^+(\Dcal_\hbar, \theta) $. Take the composition of the map $\bar{\J}_{\mathcal{C},\hbar} \widehat{\otimes} 1 \hookrightarrow \bar{\J}_{\Dcal,\hbar}$ with the map $\mathcal{C}'_0 \otimes_{\Dcal'_0} \T^+(\Dcal_\hbar, \theta) \to \mathcal{C}'_0 \otimes_{\Dcal'_0} \der(\Dcal'_0)$ induced by the anchor map of $\T^+(\Dcal_\hbar, \theta)$. The resulting composite map $\bar{\J}_{\mathcal{C},\hbar}\rightarrow \mathcal{C}'_0 \otimes_{\Dcal'_0} \der(\Dcal'_0)$  coincides with the composite map 
  $$\bar{\J}_{\mathcal{C},\hbar}\twoheadrightarrow 
  \bar{\J}_{\mathcal{C},\hbar}/\bar{\J}_{\mathcal{C},\hbar}^2\xrightarrow{\sim}\T^+(\mathcal{C}_\hbar, \theta) \to \der(\mathcal{C}'_0) \to \mathcal{C}'_0 \otimes_{\Dcal'_0} \der(\Dcal'_0).$$
  This implies that we have a natural $\cm$-equivariant morphism $\T^+(\mathcal{C}_\hbar, \theta) \to \hat{f}^\circ [\T^+(\Dcal_\hbar, \theta)]$ of Picard Lie algebras over $\mathcal{C}'_0$. Note that $\hat{f}^\circ [\T^+(\Dcal_\hbar, \theta)]$ is nothing else but the completion of $f^\circ [\T^+(\A_\hbar, \theta)]$. %\footnote{S.Y.: Elaborate? We need some statements like completion of $\T^+(\A_\hbar, \theta)$ equals $\T^+(\Dcal_\hbar, \theta)$ and taking completions intertwines $f^\circ$ and $\hat{f}^\circ$. This should follow from functoriality once we assume all of our algebras are Noetherian and modules are finitely generated. But first of all, we have to show that the completion of the $\J$-ideal of $\A_\hbar$ is the same as the $\J$-ideal of $\Dcal_\hbar$, or even more elementary, the completion of the ideal $\langle \C[X] \rangle^{-\theta} \subset \C[X]$ is the same as the ideal $\langle \C[X^{\wedge_{\chi'}}] \rangle^{-\theta}$ of $\Dcal_0 = \Dcal_\hbar/\hbar\Dcal \simeq \C[X^{\wedge_{\chi'}}]$. This is not so hard but also not entirely obvious.}.  
  Taking the $\cm$-finite vectors gives a $\cm$-equivariant morphism $\T^+(\check{\A}_\hbar, \theta) \to f^\circ [\T^+(\A_\hbar, \theta)]$ of Picard Lie algebras over $\C[\check{Y}]$, whose localization to $\check{Y}^{\times}$ is a $\cm$-equivariant morphism $\TT^+(\check{\A}_\hbar^\times, \check{Y}^\times) \to \iota_{\check{Y}^\times}^\circ [ \TT^+(\hat{\A}_\hbar^{\Orb}, \Orb_K)]$ of Picard algebroids by Lemma \ref{lem:Picard_loc} and hence must be an isomorphism. The same argument implies the second statement of part (1).

Now we prove part (2). Keeping the notations from Section \ref{subsec:extn_codim2}, we have a natural isomorphism $\check{\Ecal}_{\hbar,1}' \xrightarrow{\sim}  \mathcal{C}'_0 \otimes_{\Dcal'_0} (\iota_*\mathcal{E}_{\hbar,1})^{\wedge_{\chi'}}$ induced by the decomposition \eqref{eq:decomposition11} for $n=1$ and the inclusion $\check{\Ecal}_{\hbar,1}' \simeq \check{\Ecal}_{\hbar,1}' \widehat{\otimes} 1 \subset (\iota_*\mathcal{E}_{\hbar,1})^{\wedge_{\chi'}}$. By \eqref{eq:decomposition11} for $n=2$, we see that the restriction of the Lie algebra action of $\bar{\J}_{\Dcal, \hbar}$ on $(\iota_*\mathcal{E}_{\hbar,2})^{\wedge_{\chi'}}$ to an action of $\bar{\J}_{\mathcal{C},\hbar} \widehat{\otimes} 1$ preserves the subspace $\check{\Ecal}_{\hbar,2}' \widehat{\otimes} 1$. This implies that the isomorphism $\check{\Ecal}_{\hbar,1}' \xrightarrow{\sim}  \mathcal{C}'_0 \otimes_{\Dcal'_0} (\iota_*\mathcal{E}_{\hbar,1})^{\wedge_{\chi'}} = (\iota_{\check{Y}}^*\iota_*\mathcal{E})^{\wedge_{0}}$ intertwines the $\T^+(\mathcal{C}_\hbar, \theta)$-module structure on the domain and the $\hat{f}^\circ [\T^+(\Dcal_\hbar, \theta)]$-module structure on the codomain, via the morphism $\T^+(\mathcal{C}_\hbar, \theta) \to \hat{f}^\circ [\T^+(\Dcal_\hbar, \theta)]$ of Picard Lie algebras. Taking $\cm$-finite parts and localizing to $\check{Y}^\times$, we get (2).

Finally we prove part (3). Recall that to construct  $\check{\mathcal{E}}^\times_\hbar$ and $\check{\mathcal{E}}^\times$ together with their $\cm$-actions, we first take the action $\alpha_{h'}$ of the one-parameter subgroup of $K$ generated by the semisimple element $h'$ of the $\slf_2$-triple $(e',h',f')$ corresponding to $\chi'$, which commutes with the original $\cm$-action on $\EE_\hbar$, then twist the $\cm$-action on $\EE_\hbar$ by the action $\alpha_{h'}$ to make it compatible with the Kazhdan action. Since $\EE_\hbar$ is a $(K, \kappa)$-equivariant graded Hamiltonian quantization of $(\Orb_K, \EE)$, the sheaf $\EE$ is a $(K, \kappa)$-equivariant $\tatp$-module. This implies that the differential of the $K$-action on the projectivization $\mathbb{P}(\EE)$ of $\EE$ coincides with the $\kf$-action induced by the projective flat connection and the $\kf$-action on $\Orb_K$. The Kazhdan $\cm$-action is obtained by twisting the orginal strong $\cm$-action by $\alpha_{h'}$. Now the action $\alpha_{h'}$ is the restriction of the $K$-action on $\EE$ to a torus subgroup, hence also satisfies the strongness condition. The action $\alpha_{h'}$  commutes with the orginal $\cm$-action, so the twisted $\cm$-action on $\EE$ is still strong.
Then in the construction of quantum slice and $\check{\mathcal{E}}^\times_\hbar$, we restrict the twisted action on $\EE$ to $\check{\mathcal{E}}^\times$ and the natural isomorphisms in part (1) and (2) are $\cm$-equivariant with respect to the twisted $\cm$-actions. Therefore the resulting $\cm$-action on $\check{\mathcal{E}}^\times$ is also strong.
\end{proof}

%\begin{Rem}\footnote{What is the purpose of this remark?  S.Y.: Just talk to myself, to make sure there is no contradiction with Lemma \ref{Lem:restriction_equivariance}. Can be removed.}
%	If we restrict the strong $K$-action on $\TT^+(\hat{\A}_t^{\Orb}, \Orb_K)$ to a strong $K_Q'$-action, this action restricts to a strong $K_Q'$-action on $\TT^+(\check{\A}_t^\times, \check{Y}^\times)$ and $\check{\mathcal{E}}^\times$ becomes a $(K, 0)$-equivariant module over $\TT^+(\check{\A}_t^\times, \check{Y}^\times)$. However, this strong $K_Q'$-action on $\TT^+(\check{\A}_t^\times, \check{Y}^\times)$ is not compatible with quantum comoment map for the quantum slice $\check{\A}_t$. On the other hand, $\check{\mathcal{E}}^\times$ is a $(K, -\rho_L)$-equivariant HC module from the quantum slice construction. This is because the isomorphism $\TT(\check{\A}_t^\times, \check{Y}^\times) \simeq \iota_{\check{Y}^\times}^\circ \TT(\hat{\A}_t^{\Orb}, \Orb_K)$ is compatible with the strong $K_Q'$-actions. Therefore there is an isomophism of strongly $K$-equivariant Picard algebroids between $\TT^+(\check{\A}_t^\times, \check{Y}^\times)$ and the sum of the Picard algebroid $\TT^+(\check{\A}_t^{\Orb}, \Orb_K)$ and $\frac{1}{2} \TT(\wedge^{\operatorname{top}} N_{\check{Y}^\times / \Orb_K})$, where $N_{\check{Y}^\times / \Orb_K}$ is the normal bundle of $\check{Y}^\times$ inside $\Orb_K$ and $\wedge^{\operatorname{top}}$ stands for taking top degree wedge product. One can show that the contracting Kazhdan action identifies the vector bundle $\check{Y}^\times$ with the trivial vector bundle $L \times \check{Y}^\times$ in a $K_Q'$-equivariant way, where the latter is endowed with the diagonal $K_Q'$-action. Since $N_{\check{Y}^\times / \Orb_K}$ is trivial, $\TT^+(\check{\A}_t^\times, \check{Y}^\times)$ is isomorphic to $\TT^+(\hat{\A}_t^{\Orb}, \Orb_K)$ as as Picard algebroids, and this is the isomorphism in part (1) of Lemma \ref{lem:restriction_strong}. However, for the $\kf_Q'$-actions, $\frac{1}{2} \TT(\wedge^{\operatorname{top}} N_{\check{Y}^\times / \Orb_K})$ contribute the term $\rho_L$.
%\end{Rem}
%
%\begin{proof}[Proof of Proposition \ref{Prop:slice_strongness}]
%	The proposition now follows by combining (3) of Lemma \ref{lem:restriction_strong}
%	with  Lemma \ref{Lem:irred_strongness_HC}.
%\end{proof}


\subsection{Description of the image}\label{SS_image_description}
The following proposition gives a description of the image of $\bullet_{\dagger,\chi}$
that will be used to characterize the image of
$\overline{\HC}(\A_\lambda,\theta)^{K,\kappa}$ in $\operatorname{Rep}(K_Q,\lambda|_{\kf_Q}-\kappa_Q)$.

Let $\Orb_K^{i},i=1,\ldots,\ell,$ be all
codim $2$ orbits in $\overline{\Orb}_K$, $\check{X}^i$ be the corresponding
4-dimensional slices in $X$ (at points $\chi^i\in \Orb_K^i$), and $\theta^i$ be induced Poisson anti-involutions of $X^i$. We write $\check{X}^{i\times}$ for $\check{X}^i\setminus \{0\}$.
Set $\check{\A}^i_\lambda:=(\A_\lambda)_{\dagger,\chi^i}$, these are quantizations of $\C[\check{X}^i]$.
For $V\in \operatorname{Rep}(K_Q,\lambda|_{\kf_Q}-\kappa_Q)$, we write $\mathcal{E}$ for the corresponding
twisted local system on $\Orb_K$. We write $\check{\mathcal{E}}^i$ for the restriction of $\mathcal{E}$ to $\check{X}^{i\times}$. We note that $\mathcal{E}$ is an equivariant twisted system on an orbit,
so it is regular, see the last paragraph of Section \ref{subsec:RH}. By Proposition \ref{Prop:regular_local_systems}, $\check{\mathcal{E}}^i$ is regular as well.

\begin{Prop}\label{Lem:image_description11}
Let $V$ be an irreducible object in $\operatorname{Rep}(K_Q,\lambda|_{\kf_Q}-\kappa_Q)$. The following claims are true:
\begin{enumerate}
\item If $V\in \operatorname{Im}\bullet_{\dagger,\chi}$, then there are
$M^i\in \HC(\check{\A}^i_\lambda,\theta^i)$ with good filtrations such that
$\gr M^i|_{\check{X}_i^\times}\cong \check{\mathcal{E}}^i$, a $\C^\times$-equivariant isomorphism of twisted
local systems.
\item Conversely, assume that $V$ satisfies the following condition:
there are $M^i\in \HC(\check{\A}^i_\lambda,\theta^i)$ with good filtrations such that
$\gr M^i\cong \Gamma(\check{\mathcal{E}}^i)$, an isomorphism of graded weakly HC modules.
Then $V\in \operatorname{Im}\bullet_{\dagger,\chi}$.
\end{enumerate}
\end{Prop}
\begin{proof}
Set $M=V^{\dagger,\chi}$, see Section \ref{SS_Prop_restr}.
This module comes with the natural filtration, Remark \ref{Rem:distinguished_filtration}.

To prove (1), observe that $M_{\dagger,\chi}=V$, Lemma \ref{Lem:image_description}. Then
$(\gr M)_{\dagger,\chi}=\gr (M_{\dagger,\chi})=V$.
We have $\gr M|_{\Orb_K}=\mathcal{E}$ by the description of $\bullet_{\dagger,\chi}$
as the microlocalization functor, Section \ref{SSS_conclusion_microloc}.
It follows that $\gr (M_{\dagger,\chi^i})=(\gr M)_{\dagger,\chi^i}$, so we can
take $M^i:=M_{\dagger,\chi^i}$ in (1). This finishes the proof of (1).

We proceed to proving (2). Let $\mathcal{E}_\hbar$ be the object
in $\HC(\hat{\A}_\hbar^{\Orb},\theta)^{K,\kappa}$ corresponding to the
$\hbar$-adic completion of $R_\hbar(V)$ (with trivial filtration) under
equivalence (\ref{eq:final_jet_equiv}).
We need to show $M_{\dagger,\chi}\cong V$. Thanks to Section
\ref{SSS_conclusion_microloc}, this is equivalent to the claim
that the microlocalization of $\Gamma(\mathcal{E}_\hbar)$ to
$\Orb_K$ is $\mathcal{E}_\hbar$.

We are now going to show that
the assumption in (2) implies that (*) of Lemma \ref{Lem:coh_extension2}
holds for all $i$. Let $M^i_\hbar$ denote the $\hbar$-adic completion of $R_\hbar(M^i)$. We have $M^i_\hbar/\hbar M^i_\hbar\xrightarrow{\sim} \Gamma(\check{X}^{i\times}, \check{\mathcal{E}}^i)$. Let $\check{\mathcal{E}}^i_\hbar$
denote the microlocalization of $M^i_\hbar$ to $\check{X}^{i\times}$.
By Lemma  \ref{Lem:irred_strongness_HC}, the $\C^\times$-action on $\mathcal{E}$ is strong.
By Proposition \ref{Prop:slice_strongness}, the $\C^\times$-action on $\check{\mathcal{E}}^i$ is strong as well.
So $\check{\mathcal{E}}^i_\hbar$ is the unique quantization
of $\check{\mathcal{E}}^i$ by Theorem \ref{thm:quan_lag}.
We have natural homomorphisms $M^i_\hbar\rightarrow \Gamma(\check{\mathcal{E}}^i_\hbar)$.
The composition of $M^i_\hbar/\hbar M^i_\hbar\rightarrow \Gamma(\check{\mathcal{E}}^i_\hbar)/\hbar \Gamma(\check{\mathcal{E}}^i_\hbar)$ with the inclusion $\Gamma(\check{\mathcal{E}}^i_\hbar)/\hbar \Gamma(\check{\mathcal{E}}^i_\hbar)
\hookrightarrow \Gamma(\check{\mathcal{E}}^i)$ is an isomorphism. It follows that
$M^i_\hbar\xrightarrow{\sim} \Gamma(\check{\mathcal{E}}^i_\hbar)$ and (*) of Lemma
\ref{Lem:coh_extension2} holds for all $i$.

Now Lemma \ref{Lem:coh_extension2} implies that
the microlocalization of $\Gamma(\mathcal{E}_\hbar)$ coincides with $\mathcal{E}_\hbar$.
This finishes the proof.
\end{proof}
\begin{comment}
\begin{Lem}\label{Lem:filtration_independence}
Suppose that we are have a HC $(\check{\A},\theta)$-module $M$ with a good filtration such that $\gr M\cong \Gamma(\check{X}^\times, \check{\mathcal{E}})$, an isomorphism of weakly HC $(\check{\A}_\hbar,\theta)$-modules. Then (*)
of Lemma \ref{Lem:coh_extension2} holds.
\end{Lem}
\begin{proof}
Let $M_\hbar$ denote the $\hbar$-adic completion of $R_\hbar(M)$. We have $M_\hbar/\hbar M_\hbar\xrightarrow{\sim} \Gamma(\check{X}^\times, \check{\mathcal{E}})$. Let $\check{\mathcal{E}}_\hbar$
denote the microlocalization of $M_\hbar$ to $\check{X}^\times$. Then it is the unique quantization
of $\check{\mathcal{E}}$ by Proposition \ref{Prop:slice_strongness} combined with Theorem \ref{thm:quan_lag}. We have natural homomorphisms $M_\hbar\rightarrow \Gamma(\check{\mathcal{E}}_\hbar)$.
The composition of $M_\hbar/\hbar M_\hbar\rightarrow \Gamma(\check{\mathcal{E}}_\hbar)/\hbar \Gamma(\check{\mathcal{E}}_\hbar)$ with the inclusion $\Gamma(\check{\mathcal{E}}_\hbar)/\hbar \Gamma(\check{\mathcal{E}}_\hbar)
\hookrightarrow \Gamma(\check{\mathcal{E}}^\times)$ is the isomorphism. It follows that
$M_\hbar\xrightarrow{\sim} \Gamma(\check{\mathcal{E}}_\hbar)$ and (*) of Lemma
\ref{Lem:coh_extension2} holds.
\end{proof}
\end{comment}


\section{Harish-Chandra modules in dimension $4$}
\subsection{Goal}\label{SS_goals}
%Current section 7 with minimum of the necessary information on the nilpotent orbits
%(from Section 6) plus Section 2.6.
We start by introducing some notation. Let $X$ be a conical symplectic singularity,
 and $\theta$ its Poisson anti-involution.
Assume that $\operatorname{codim}_X X^{sing}\geqslant 4$. In this case, as was explained in
Section \ref{SS_quant_nilp}, we can consider
the universal filtered quantization $\A_{\param}$ of $\C[X]$.
The anti-involution $\theta$ naturally lifts to $\A_{\param}$, denote the
lift by the same letter. Recall the category $\HC(\A_\lambda,\theta):=\HC(\A_\lambda,\theta,1)$ and its Serre quotient $\overline{\HC}(\A_\lambda,\theta)$. See Section \ref{SS_quant_nilp}.


In the present section we describe the categories $\overline{\HC}(\A_\lambda,\theta,\zeta)$, where
\begin{itemize}
\item
$\A_\lambda$ is a quantization of a suitable four dimensional conical symplectic singularity $X$
with $X^{sing}=\{0\}$,
\item and $\theta$ is as above.
\end{itemize}

More precisely, we care about the following situation.
Let $\Orb$ be a nilpotent orbit satisfying (\ref{eq:codim_condition}). For $X$ we take
a slice to a codimension $4$ orbit in $\operatorname{Spec}(\C[\Orb])$. This is a conical
symplectic singularity. The list of possible options for $X$ is known in  all
cases, see \cite{KP1} for $\g=\mathfrak{sl}_n$, \cite{KP2} for the other classical Lie algebras,
and \cite{FJLS1} and \cite{FJLS2} for the exceptional Lie algebras.
Note that these papers deal with slices to open orbits in $\partial\Orb$,
where $\Orb$ is viewed inside of $\g^*$. Our slices will be obtained for those by passing to the normalization and taking a connected component.

%For $(\A_\lambda,\theta)$ we take the quantization of $X$ and its
%anti-involution induced by those for $\operatorname{Spec}(\C[\Orb])$.
Here is the list of possible varieties $X$ that can appear, we will elaborate on
the orbits $\Orb$ in Section \ref{S_classif_thm}.

1) Type $a_2$: the variety $X$ is the closure of the minimal nilpotent orbit in
$\mathfrak{sl}_3$.

2) Type $c_2$: the variety $X$ is the closure of the minimal nilpotent orbit
in $\mathfrak{sp}_4$. Equivalently, $X=\C^4/\{\pm 1\}$.

3) $X=\C^4/\mathbb{Z}_3$. Here we write $\mathbb{Z}_3$ for the cyclic group with
3 elements thought of as the group of roots of unity of order $3$.
The action on $\C^4$ is as follows: an element $\zeta\in \mathbb{Z}_3$
acts by $\operatorname{diag}(\zeta,\zeta,\zeta^{-1},\zeta^{-1})$.

4) Type $a_2/S_2$ (see \cite[Prop. 3.7]{FJLS2}): let $\tilde{X}$ be such as in 1). Consider the involution
of $\tilde{X}$ obtained by restricting the outer involution of $\mathfrak{sl}_3$,
it is Poisson. This involution gives rise to an $S_2$-action. We take $X=\tilde{X}/S_2$.

5) Type $\chi$. Let $\Orb$ be the principal orbit in $\mathfrak{sl}_5$. Let $\tilde{\Orb}$
be its universal cover (with Galois group $\mathbb{Z}_5$). Consider the variety
$\Spec(\C[\tilde{\Orb}])$. It is smooth in codimension $2$. For $X$ we take
the slice to a codimension $4$ orbit in $\Spec(\C[\tilde{\Orb}])$. It is studied in details in \cite{BBFJLS1}.

We note that cases 3)-5) only appear for exceptional $\g$.

It will be convenient for us to use the following terminology.

\begin{defi}\label{defi:loc_sys_quantized}
Let $X$ be a conical symplectic singularity with $\operatorname{codim}_X X^{sing}\geqslant 4$, $\theta$ its graded Poisson anti-involution,
$\A$ a filtered quantization of $\C[X]$, and $\mathcal{E}$ a graded regular twisted local system on $X^{reg}\cap X^\theta$.
\begin{itemize}
\item[(i)]
We say that $\mathcal{E}$ is {\it $\A$-quantizable} if there is
a HC $(\A,\theta)$-module $M$ with a $\C^\times$-equivariant isomorphism $\gr M|_{X^{reg}}\cong \mathcal{E}$
of twisted local systems on $X^{reg}\cap X^\theta$.
\item[(ii)] Moreover, if in this case we additionally have
$\gr M\xrightarrow{\sim} \Gamma(\mathcal{E})$ (for some HC $(\A,\theta)$-module $M$ as in (i)), then we say that $\mathcal{E}$ is {\it strongly $\A$-quantizable}.
\end{itemize}
\end{defi}

In what follows, we will always be in the following situation:
\begin{itemize}
\item The category of regular twisted local systems on $X^\theta\cap X^{reg}$ is semisimple with finitely many simples,
\item and each simple admits a grading.
\end{itemize}

In this case, thanks to Lemma \ref{Lem:graded_loc_sys}, it is enough to check (i) and (ii) only
for irreducible regular twisted local systems.


\begin{defi}\label{defi:obstructive_slice}
Let $X$ be one of the five 4-dimensional conical symplectic singularities mentioned above
and $\theta$ be a graded Poisson anti-involution of $\C[X]$. We say that $(X,\theta)$
is {\it unobstructive} if for every quantization $\A$, every graded regular twisted local system whose
$\C^\times$-action (with twist recovered from the microlocalization of $\A$ to $\Orb$, compare to Theorem 
\ref{thm:lag_vec}) is strong on $X^\theta\setminus \{0\}$ is strongly $\A$-quantizable.
\end{defi}

The latter definition is motivated by Proposition \ref{Lem:image_description11}: for an unobstructive slice the conditions of this proposition are always satisfied, so it does not obstruct extending
a HC module from $X^{reg}$ to $X$.

\subsection{Slices}
As was mentioned in the previous section, we have a description of the (formal) slices to all codimension 4 $G$-orbits, $\Orb'$, in
$X:=\operatorname{Spec}(\C[\Orb])$. However, we need more: we need to know the Poisson structure, the subalgebra of $\C^\times$-finite elements,
and the Poisson anti-involution $\theta$. In this and the next section, we will see that these structures
are recovered essentially uniquely.

Here we are going to prove that
\begin{itemize}
\item[(*)]
the slice to $\Orb'$ in $X$ at a $\theta$-stable
point $\chi'$ is well-defined as a formal Poisson scheme with an anti-Poisson involution
and an action of $\C^\times$ that rescales the bracket.
\end{itemize}

First, we explain what we mean by a formal slice. The completion $\C[X]^{\wedge_{\chi'}}$ decomposes as
\begin{equation}\label{eq:decomp_101}\C[X]^{\wedge_{\chi'}}=\C[[x_1,y_1,\ldots,x_k,y_k]]\widehat{\otimes}
A\end{equation}
where $2k=\dim \Orb'$, the factor $\C[[x_1,\ldots,y_k]]$ is the algebra of functions on the formal
symplectic polydisc with Darboux basis $x_1,\ldots,y_k$, and $A$ is a complete local ring that coincides
with the centralizer of $\C[[x_1,\ldots,y_k]]$ in $\C[X]^{\wedge_{\chi'}}$. See, for example, \cite[Theorem 2.3]{Kaledin_symplectic}, for a more general result. A quotient $A$ of $\C[X]^{\wedge_{\chi'}}$ that factors through (\ref{eq:decomp_101}) will be called a {\it transverse slice} in $X^{\wedge_{\chi'}}$. Note that if a reductive group $\tilde{Q}$ acts on
$\C[X]^{\wedge_{\chi'}}$ by algebra automorphisms rescaling the bracket,
then one can find $x_1,\ldots,y_k$
such that $\operatorname{Span}_{\C}(x_1,\ldots,y_k)$ is $\tilde{Q}$-stable.
This gives a $\tilde{Q}$-action on $A$.

In particular, taking the formal neighborhood at $\chi'$ of an actual slice in $X$ we get a formal slice in the sense of the 
previous paragraph. So, the following lemma implies (*). 
%\footnote{It is worthwhile make a one sentence remark that, as a consequence of Lemma 7.1, if there exists an actual $\theta$-stable and %$\cm$-stable slice such that the restriction of the $\cm$-action on the slice is contracting, then all such slices are conjugate by Hamiltonian %automorphisms. -- Shilin;
%I.L.: well, the lemma is about the formal setting. So perhaps, the remark should appear somewhere else.
%5/24: and your claim is actually false -- there's a problem with an automorphism, it doesn't extend to the ambient variety. S.Y.: Are you saying %that we cannot make the automorphism $\cm$-equivariant?}

\begin{Lem}\label{Lem:slice_well_defined}
Suppose $\tilde{Q}$ is a reductive group acting on $\C[X]^{\wedge_{\chi'}}$ by algebra automorphism
and rescaling the Poisson bracket such that the decomposition is $\tilde{Q}$-stable. Then every transverse slice to $\operatorname{Spec}(\C[[x_1,\ldots,y_k]])$ in $X^{\wedge_{\chi'}}$ is conjugate
to the standard slice $x_1=\ldots=y_k=0$ via a  Hamiltonian
automorphism. If the transverse slice is $\tilde{Q}$-stable, we can pick this
Hamiltonian automorphism to be  $\tilde{Q}$-equivariant.
\end{Lem}
\begin{proof}
By the implicit function theorem (for formal completions), every transverse slice
%\footnote{Say something about what slice means in this context, so that this is true. -- Shilin}
in $X^{\wedge_{\chi'}}$ is defined by the ideal
$$(x_i-F_i, y_i-G_i| i=1,\ldots,k)$$
for suitable elements $F_i,G_i$ that lie in the ideal $I\subset \C[X]^{\wedge_{\chi'}}$ generated by the maximal ideal of a fixed transverse slice $A$
and the homogeneous quadratic polynomials in $x_i,y_i$. We can modify the elements $F_i,G_i$
without changing the ideal and assume that they actually lie in the maximal ideal of $A$.

Now consider the element $f:=\sum_{i=1}^k (-y_i F_i+x_i G_i)$. The Hamiltonian automorphism
$\exp(\{f,\cdot\})$ sends $x_i$ to $x_i-F_i$ and $y_i$ to $y_i-G_i$ for all $i$. The condition
that the slice is $\tilde{Q}$-stable is equivalent to the map $\operatorname{Span}_\C(x_1,\ldots,y_k)
\rightarrow \operatorname{Span}_\C(F_1,\ldots,G_k)$ being $\tilde{Q}$-equivariant. From here it is
easy to see that $f$ transforms under $\tilde{Q}$ with the character opposite of that of the bracket.
Hence $\exp(\{f,\cdot\})$ is $\tilde{Q}$-equivariant.
\end{proof}

\subsection{Anti-involutions}
Let $A'$ be the completion at $0$ of the algebra of functions on a conical symplectic
singularity $X$. Let $\param$ be a vector space and $A'_\param$ be a $\C^\times$-equivariant Poisson deformation of $A'$ over $\C[[\param]]$ with $\param^*$ in degree $d$, where $-d$
is the degree of the Poisson bracket on $A$. The algebras $A',A'_\param$ come with a tautological $\C^\times$-action denoted by $\gamma$.
Let $\A_{\param,\hbar}'$ be a graded formal quantization of $A'_\param$
with $\hbar$ of degree $d$. Assume that
\begin{itemize}
\item[(i)] 
$A'$ has no non-zero elements of degrees $1,\ldots,d-1$.
\item[(ii)] The degree $d$ part $A'_d$ is reductive as a  Lie algebra w.r.t $\{\cdot,\cdot\}$, whose action via $\{\cdot,\cdot\}$
on $A'$ integrates to a (faithful) pro-rational action of a connected reductive algebraic group, $\underline{G}$.
\item[(iii)] Moreover, for any nontrivial one-parameter subgroup $\gamma_0:\C^\times\rightarrow \underline{G}$,
there is $i\in \Z$ such that the dimensions of the eigenspaces in $A'$ corresponding to the character $t\mapsto t^i$
are different for $\gamma$ and $\gamma\gamma_0$. 
\end{itemize}

Suppose that $\gamma_1$ is another contracting $\C^\times$-action on $\A'_{\param,\hbar}$
such that the degree of the Poisson bracket on $A'_\param$ is
$-kd$ and the degrees of $\param^*$ and $\hbar$ are $kd$.
Let $\theta$ be a $\C[[\hbar]]$-linear anti-involution of $\A'_{\param,\hbar}$
that commutes with $\gamma_1$. Our goal is to get a more explicit description of the involution $\theta$.

Note that $\theta$ preserves $\C[[\param,\hbar]]$
because the latter subalgebra is the center of $\A'_{\param,\hbar}$. Since $\theta$
commutes with $\gamma_1$, it preserves $\param^*\oplus \C\hbar$.
%We can modify the embedding
%$\param^*\hookrightarrow \A_{\param,\hbar}'$ by adding a multiple of $\hbar^d$ and assume that
%$\theta$ preserves $\param^*$.
%We can shift $\theta$
%and assume that $\theta$ preserves $\param^*$ (but it may act nontrivially on that space).

The main result of this section is as follows.

\begin{Lem}\label{Lem:action_conjugation}
Suppose that (i)-(iii) are satisfied. 
Then there is an automorphism of $\A'_{\param,\hbar}$ of the form $\exp(\hbar^{-1}[F,\cdot])$
such that after conjugating $\gamma$ with this automorphism we can achieve
that $\gamma_1$ and $\theta$ commute with $\gamma$.
\end{Lem}
\begin{proof}
We consider the group $\mathcal{G}$ of all $\C$-linear automorphisms
of $\A'_{\param,\hbar}$ that
\begin{itemize}
\item  are automorphisms or anti-automorphisms of the $\C$-algebra $\A'_{\param,\hbar}$,
\item and act on $\param^*\oplus \C\hbar$ either by rescaling if they are automorphisms or by rescaling composed with $\theta$
if they are anti-automorphisms.
\end{itemize}
Note that $\mathcal{G}$ contains a normal subgroup of $\C[[\param,\hbar]]$-linear algebra
automorphisms, denoted by $\mathcal{G}_0$. It follows from \cite[Proposition 2.14]{orbit}
that all Poisson derivations of $A'$ are inner (note that \cite[Proposition 2.14]{orbit} deals with $\C[X]$ but 
the Lie algebra of Poisson derivations of $A'$ is the completion at $0$ of that for $\C[X]$ so the result for 
$A'$ follows). Therefore
all $\C[[\param,\hbar]]$-linear derivations of $\A'_{\param,\hbar}$ are of the form $\hbar^{-1}[F,\cdot]$
for $F \in \A_{\param,\hbar}'$. By the condition that $A'_i=0$ for $0<i<d$, we see that
every derivation preserves $\prod_{i\geqslant j} (\A'_{\param,\hbar})_i$ for all $j$. It follows that
$\mathcal{G}_0$ decomposes as $\mathcal{G}_0^\gamma\ltimes U$, where
we write $\mathcal{G}_0^\gamma$ for the subgroup of $\mathcal{G}_0$ consisting of all
elements commuting with $\gamma$, and $U$ is a pro-unipotent pro-algebraic group.
Note that $\mathcal{G}_0^\gamma$ is an algebraic subgroup whose connected component is $\underline{G}$.

Let $\mathcal{G}^0$ denote the normal subgroup of $\mathcal{G}$ consisting of all automorphisms. This is an index $2$ subgroup. 

Observe that $\mathcal{G}^0=Z_{\mathcal{G}^0}(\gamma)\ltimes U$. Indeed, $\mathcal{G}^0=\gamma(\C^\times)\mathcal{G}_0$ and we can use the decomposition
$\mathcal{G}_0=\mathcal{G}_0^\gamma\ltimes U$ from the previous paragraph.

The group $Z_{\mathcal{G}^0}(\gamma)$ is algebraic.
It follows that $\mathcal{G}^0$ contains the unique maximal
normal pro-unipotent subgroup, denote it by $\tilde{U}$. We have $\mathcal{G}^0=
H^0\ltimes \tilde{U}$, where $H^0$ is a maximal reductive subgroup of $Z_{\mathcal{G}^0}(\gamma)$.
We can include $H^0$ as an index $2$ normal subgroup to a maximal reductive subgroup  $H\subset \mathcal{G}$
(this is because $H^0$ and $\theta H^0\theta^{-1}$ are conjugate by an element of $\tilde{U}$ and so 
$N_{\mathcal{G}}(H^0)\twoheadrightarrow \mathcal{G}/\mathcal{G}^0$).
Every reductive subgroup of $\mathcal{G}$ is conjugate to a subgroup in $H$
by an element of $\tilde{U}$. Note that every element of $\tilde{U}$ is of the form 
$\exp(\hbar^{-1}[F,\cdot])$.
In particular, we can conjugate $\gamma_1$ and $\theta$ by an automorphism of this form and assume that 
$\langle\theta\rangle\times\gamma_1(\C^\times)\subset H$.

It remains to show that $\gamma(\C^\times)$ is central in $H$. It is central in $H^0$ by construction.
In particular, $\gamma$ commutes with $\theta \gamma \theta^{-1}$. Set $\gamma_0:=\theta\gamma \theta^{-1}\gamma^{-1}$.
This is a one-parameter subgroup in $H$ that acts trivially on $\param^*\oplus \C$. So the image is in 
$\underline{G}$. The dimensions of eigenspaces in $A'$ with the same character of $\gamma$ and $\theta\gamma \theta^{-1}$ on $A'$ are the same. 
Thanks to (iii), $\gamma_0$ is trivial. So $\theta$ and $\gamma$  commute finishing the proof.  
%The centralizer $\mathcal{G}^{\gamma_1}$ is an algebraic subgroup
%that contains $\gamma(\C^\times)$ and $\theta$
%
%Hence $\mathcal{G}=H\ltimes \tilde{U}$, where $H$ is a reductive group containing $H^0$.
%Since $\gamma_1$ and $\theta$ commute, we see that the conjugation by $\theta$ acts trivially on $H^0/(H^0\cap \mathcal{G}_0)$
%(note that the image of $\gamma_1(\C^\times)$ in $H^0/(H^0\cap \mathcal{G}_0)$\footnote{S.Y.: why is $\gamma_1(\C^\times)$ in $H^0$? We didn't %assume $\gamma$ commutes with $\gamma_1$. Maybe this is typo and you meant the image of $\gamma(\C^\times)$ in $H^0/(H^0\cap \mathcal{G}_0)$?}
%is the connected component of $1$ in $H^0/(H^0\cap \mathcal{G}_0)$,
%and the component group of $H^0/(H^0\cap \mathcal{G}_0)$ is
%$\Z/2\Z$).
%From here we see that $\gamma(\C^\times)$ is central in $H$. Also it is a standard fact that any reductive
%subgroup in $\mathcal{G}$ can be conjugated to a subgroup of $H$ by an element of $\tilde{U}$. Such an
%element is a converging series of the form $\exp(d)$, where $d$ is a $\C[[\param,\hbar]]$-linear derivation.
%Such a derivation has the form $\hbar^{-d}[F,\cdot]$, this has been discussed in the
%first paragraph of the proof.
%We apply this to the reductive subgroup $\langle\theta\rangle\times \operatorname{Im}\gamma_1$
%to finish the proof of the lemma.
\end{proof}

Note that if $\gamma$ and $\gamma_1$ commute, then the subalgebras of locally finite elements in
$\A'_{\param,\hbar}$ under $\gamma$ and $\gamma_1$ coincide. This is because the eigenspaces for
both $\gamma$ and $\gamma_1$ in $\A'_{\param,\hbar}$ are finite dimensional.
%Indeed, we have that $\gamma_1\gamma^{-k}$ \footnote{Maybe $\gamma^{-k}\gamma_1$ instead?  Also the %condition that $\gamma$ and $\gamma_1$ commutes should directly imply that the subalgebras of finite %vectors coincide, no?  -- Shilin}
%is a one-parameter subgroup in the centralizer of both $\gamma$ and $\gamma_1$. It acts locally
%finitely on the subalgebras of locally finite elements in both $\gamma$ and $\gamma_1$. Hence
%the coincidence of the subalgebras.

Now we are going to deduce two corollaries.

\begin{Cor}\label{Cor:a2_c2_involution}
Suppose that $X$ is of type $a_2$ or $c_2$.
Let $\Orb$ be the open orbit in $X$,
$\param=H^2(\Orb,\C)$, and $\A'_{\param,\hbar}$ is the completion at $0$
of the (modified) Rees algebra of the universal filtered quantization $\A_{\param}$ of $\C[X]$.
Then the subalgebra of locally
finite (w.r.t. $\gamma_1$) elements in $\A'_{\param,\hbar}$ is stable under $G=\operatorname{SL}_3$
or $\operatorname{Sp}_4$ and $\theta$ induces an involution on $\g$. This involution preserves
the minimal orbit $\Orb_{min}$ in $\g$
and the action on $\param$ is induced from the action on $\Orb$.
\end{Cor}
Note that the conclusion of the lemma identifies $\theta$ completely.
\begin{proof}
We apply Lemma \ref{Lem:action_conjugation} to the action $\gamma$ coming from the dilation
action on the Lie algebra $\g$ of $G$. Here $d=1$ so  (i) is satisfied automatically. 
The Lie algebra $A'_1$ is nothing else but $\g$, so (ii) holds. And for a non-trivial one-parameter 
subgroup $\gamma_0:\C^\times\rightarrow \underline{G}=\operatorname{Ad}(\g)$, there is a vector 
$x\in \g$ with $\gamma_0(t)x=t^j x$ with $j<0$. So (iii) is satisfied (take $i=1+j$). 

Since $\theta$ commutes with $\gamma$, it preserves the degree
$1$ component in $\A_{\param,\hbar}'$. That component is $\g\oplus \param^*\oplus \C\hbar$. It follows that $\theta$
preserves $\g$ and induces a Lie algebra anti-involution on $\g$.
The claim about the action of $\theta$ on $\param$ follows because
$\A_{\param}$ is the universal filtered quantization of $\C[X]$
so that the action of $\theta$ on $\param$ must be the natural one by universality
(see \cite[Proposition 3.5]{orbit} for the description of what it means to be universal). 
%\footnote{S.Y.: Is it completely trivial? It seems that you are using the properties of the extended Namikawa-Weyl group and the fact that this group is trivial for $X$. I.L.: Yes, this is from the definition. S.Y.: I am not sure. If we compose $\theta$ with the one in the statement of the corollary, then it is an automorphism that is identity on $\Orb$, but can act a piori non-trivially on $\param$. Then one can argue that this automorphism is an element from the Namikawa-Weyl group, which is trivial in our case. I.L.: I stand by what I said: if you have the nontrivial Weyl group, then this deformation is not universal. S.Y.: I see $\A_{\param}$ is mentioned at the very beginning of the section. I've edited the statement slightly to remind the reader.}
\end{proof}



%\begin{Cor}\label{Cor:a_2_deformed}
%Suppose that $\operatorname{Spec}(A')$ is of type $a_2$ and the deformation $A'_\param$
%is the completion at $0$ of the universal graded deformation Poisson of the type $a_2$
%singularity. Suppose $$
%\end{Cor}

\begin{Cor}\label{Cor:tau_a_2S_2_involution}
Suppose that $X$ is of type $\tau$ or $a_2/S_2$. Then for $\gamma$
we can take
\begin{itemize}
\item for $\tau$, the action induced from the dilation action on $\C^4$ (with $d=2$),
\item for $a_2/S_2$, the action induced from the dilation action on $a_2$ (with $d=1$).
\end{itemize}
In particular, in both these cases, the subalgebra of locally finite elements with respect
to $\gamma_1$ is the standard one.
\end{Cor}
\begin{proof}
We need to check the condition (i)-(iii). For $a_2/S_2$, this is very similar to  
Corollary \ref{Cor:a2_c2_involution}: here $d=1$ and $A'_1=\mathfrak{so}_3$. The details are left to the reader. 

Let us consider the case of $\tau$. Here $d=2$. (i) follows because the group $\Z/3\Z$ acts on $\C^4$ without nonzero fixed points.  
The Lie algebra $A'_2$ is $\mathfrak{gl}_2$. The $\operatorname{GL}_2$ acts on $X=(T^*\C^2)/\Z_3$
via its natural action on $\C^2$ and $\underline{G}$ is the quotient of $\operatorname{GL}_2$ by the central $\Z_3$. 
So (ii) follows. Now we check (iii). If $\gamma_0$ is not central in $\underline{G}$, we can argue as 
in the proof of Corollary \ref{Cor:a2_c2_involution}. So suppose $\gamma_0$ is central. 
Then $\gamma\gamma_0$ acts on $A'_2$ with eigenvalue $t\mapsto t^2$. It follows that all other eigenvalues
of $\gamma\gamma_0$ in $A'$ are of the form $t\mapsto t^i$ for $i\geqslant 3$. In particular this forces the eigenvalues 
of $\gamma_0$ in $S^3(\C^2), S^3(\C^{2*})$ to be trivial, and hence $\gamma_0$ to be trivial. This checks (iii).
%In the case when the singularity is $\tau$, we need to show that $A'_1=\{0\}$.
%This follows .
\end{proof}

We will describe possible anti-involutions $\theta$ in these two cases later.

Starting the next section we proceed to classifying Harish-Chandra modules with full support
for the four out of five singularities of interest: $a_2,c_2,\C^4/\mathbb{Z}_3, a_2/S_2$.
We will see that the varieties of types $c_2,\C^4/\mathbb{Z}_3$ are unobstructive for all
choices of involutions. The variety of type $a_2$ is unobstructive for one of the two possible
choices of an involution.

\subsection{Type $a_2$}\label{SS_type_a2}
%For each irreducible component $Y^\times$ of $X^\theta\setminus \{0\}$ and a quantization parameter
%$\lambda$, we aim to describe the twisted local systems
%of the form $(\gr M)|_{Y^\times}$ for $M\in \HC_Y(\A_\lambda,\theta)$ (where $Y$ is the closure of %$Y^\times$; we note that for the types $a_2,c_2$ this reduces to describing the image of the functor %$\bullet_{\dagger,\chi}$). Further, for each $\mathcal{L}$ of the form $(\gr M)|_{Y^\times}$ (and given %$\lambda$)
%we determine when there is $\tilde{M}\in \HC_Y(\A_\lambda,\theta)$ with a good filtration such that
%$\gr\tilde{M}=\Gamma(Y^\times,\mathcal{L})$. Since the latter is a Cohen-Macaulay module over
%$\C[Y]$, see the proof of Lemma \ref{Lem:coh_extension2}, we refer to the existence of $\tilde{M}$
%and a good filtration as the Cohen-Macaulay property. A motivation for these considerations is
%Proposition \ref{Lem:image_description11}.
%In each case we examine the Cohen-Macaulay property: given an irreducible component $Y^\times$ of
%$X^\theta\setminus \{0\}$, a quantization parameter $\lambda$, a twist $\kappa$ a twisted local system %$\mathcal{L}$
The first case to consider is when $X$ is the closure of the minimal orbit (Jordan type $(2,1)$)
in $\g:=\mathfrak{sl}_3$. We have $\param=\C$. The quantizations of $\C[X]$ have the form
$\mathscr{D}^{\lambda}(\mathbb{P}^2)$, the algebra of $\lambda$-twisted
differential operators on $\mathbb{P}^2$. The period of such a quantization is $\lambda-\frac{1}{2}c_1(K_{\mathbb{P}^2})$. We consider anti-involutions $\theta$ of the universal quantization $\A_\param$ coming from $\mathfrak{sl}_3$.

By Corollary \ref{Cor:a2_c2_involution}, $\theta=-\sigma$ for an involution $\sigma$ of $\g$.
Up to conjugation, the Lie algebra $\mathfrak{sl}_3$ has two involutions
$\sigma$. One has $\kf=\mathfrak{gl}_2$ and the other has $\kf=\mathfrak{so}_3$. When $\kf=\mathfrak{gl}_2$,
the involution $\sigma$ is inner, and $\sigma$ acts trivially
on $\param$. When $\kf=\mathfrak{so}_3$, the involution $\sigma$ is outer, and
$\sigma$ act on $\param$ by $-1$. Note that the  action of $-1$ on $T^*\mathbb{P}^2$
is contained in the dilation action of $\C^\times$, a connected group. So the actions of
$\sigma$ and $\theta$ on the $H^2(T^*\mathbb{P}^2,\C)$ coincide.

Let us write $\mathscr{D}^\lambda_{\Pb^2}$ for the sheaf of $\lambda$-twisted
differential operators on $\Pb^2$. We have the global section
functor $\mathscr{D}^\lambda_{\Pb^2}\operatorname{-mod} \rightarrow \mathscr{D}^\lambda(\Pb^2)\operatorname{-mod}$,
where the source category is that of coherent $\mathscr{D}^\lambda_{\Pb^2}$-modules and the target is the category
of finitely generated $mathscr{D}^\lambda(\Pb^2)$-modules.

We claim that
the functor $\Gamma$ induces an equivalence between the
quotients of $\mathscr{D}^\lambda_{\Pb^2}\operatorname{-mod}$ and $\mathscr{D}^\lambda(\Pb^2)\operatorname{-mod}$
by the Serre subcategories of $\mathcal{O}$-coherent twisted $\mathscr{D}$-modules and
finite dimensional $\mathscr{D}^\lambda(\Pb^2)$-modules, respectively. Indeed, let $\operatorname{Loc}: \mathscr{D}^\lambda_{\Pb^2}\operatorname{-mod} \to \mathscr{D}^\lambda(\Pb^2)\operatorname{-mod}$
denote localization functor, which is the left adjoint of $\Gamma$ (see \cite[Section 3.3]{BB}). The kernels and cokernels of the adjunction unit and 
counit morphisms for the adjoint pair $(\operatorname{Loc},\Gamma)$ are killed by the microlocalization
to $\Orb=\overline{\Orb}\setminus \{0\}$ because the resolution morphism $T^*\Pb^2\rightarrow \overline{\Orb}$
is an isomorphism over $\Orb$. It follows that they are supported at $\{0\}$, which implies our claim.  

The group $K$ acts on $\Pb^2$ with finitely many orbits. Thanks to this we can classify
irreducible $(K,\kappa)$-equivariant $\mathscr{D}^\lambda_{\Pb^2}$-modules, see Lemma \ref{Lem:fin_many_orbit_classif}. To apply the lemma, we need the following notation.
Let $v\in \C^3$ be a nonzero vector and $[v]$ be the corresponding line
in $\Pb^2$. Let $P_{[v]}$ denote the stabilizer of $[v]$ in $\operatorname{SL}_3$. We identify the character
lattice $\mathfrak{X}(P_{[v]})$ with $\Z$, with $1$ corresponding to the character of
$P_{[v]}$ in $\C v$ (and so the $G$-equivariant line bundle $\mathcal{O}(1)$ corresponds to $-1$). So $\lambda$ can be viewed as 
a complex number. This also gives an identification of 
$H^2(T^*\mathbb{P}^2,\C)$ with $\C$ (with $c_1(\mathcal{O}(1))=-1$ and $c_1(K_{\mathbb{P}^2})=3$).

%\footnote{S.Y.: This is where Lemma ({Lem:tw equi loc sys}) was in v15. I don't know why it was gone.}

%Shall we add the following lemma back?
\begin{comment}
	Consider the orbit $K/K_{[v]}$, where $K_{[v]}=K\cap P_{[v]}$.
	The following lemma is standard.
	
	\begin{Lem}\label{Lem:tw_equi_loc_sys}
		The following claims are true:
		\begin{enumerate}
			\item If there is a $(K,\kappa)$-equivariant $\mathcal{O}$-coherent
			$D^\lambda_{K/K_{[v]}}$-module, then $\lambda|_{\kf_{[v]}}-\kappa|_{\kf_{[v]}}$
			comes from a character of $K_{[v]}$.
			\item If this is the case, then the irreducible $(K,\kappa)$-equivariant
			$\mathcal{O}$-coherent $D^\lambda_{K/K_{[v]}}$-modules are in bijection with
			the irreducible representations of $K_{[v]}/K_{[v]}^\circ$.
		\end{enumerate}
	\end{Lem}
\end{comment}


\subsubsection{The case of $\kf=\mathfrak{gl}_2$}\label{SS_HC_sl3_gl2}
Here $\theta$ comes from the inner involution of $\mathfrak{sl}_3$. We identify $\kf^{*K}$ with $\C$ so that the element $1\in \C$
corresponds to $-\operatorname{tr}\in \kf^{*K}$. So we can view $\kappa$ as a complex number. 

Our goal is to establish the following claim.

\begin{Lem}\label{Lem:sl3_gl2_unobstructive}
Let $X$ be the closure of the minimal orbit in $\mathfrak{sl}_3$ and $\theta:=-\sigma$,
where $\sigma$ is an inner involution of $\mathfrak{sl}_3$. Then $(X,\theta)$ is unobstructive.
\end{Lem}
\begin{proof}
{\it Step 1}.
In this case we have three $K$-orbits on $\mathbb{P}^2$. Let us write points of $\Pb^2$
as $[a:b:c]$. We embed $\operatorname{GL}_2$ into
$\operatorname{SL}_3$ as $x\mapsto \operatorname{diag}(\det(x)^{-1},x)$.
The representatives of the three orbits are $[1:1:0],[0:1:0]$ and
$[1:0:0]$. 

Let us see what the classification of irreducible twisted equivariant twisted local systems (see Lemma
\ref{Lem:homog_classif}) on the three orbits looks like:

1) {\it The case of $[1:1:0]$}. The group $K_{[v]}$ consists of all matrices of the form
$$\begin{pmatrix}x&0&0\\0&x&y\\0&0&x^{-2}\end{pmatrix},$$
and the orbit $K \cdot [v]$ is open in $\Pb^2$. We identify $\mathfrak{X}(K_{[v]})$ with $\Z$, where $1$
corresponds to the character that sends the above matrix to $x$. %In particular, we can view $\kappa$ as a complex number. 
The restriction map $\mathfrak{X}(P_{[v]})\rightarrow \mathfrak{X}(K_{[v]})$ sends
$1$ to $1$.
So, by Lemma \ref{Lem:homog_classif}, the existence of a $(K,\kappa)$-equivariant $\mathcal{O}$-coherent
$\mathscr{D}^\lambda_{K/K_{[v]}}$-module is equivalent to $\lambda-\kappa\in \Z$.
If this is the case we have a unique twisted
equivariant irreducible $\mathcal{O}$-coherent module because $K_{[v]}$ is connected.

2) {\it The case of $[0:1:0]$}. The group $K_{[v]}$ consists of all matrices of the form
$$\begin{pmatrix}x_1&0&0\\0&x_2&y\\0&0&x_1^{-1}x_2^{-1}\end{pmatrix},$$
and therefore the orbit $K \cdot [v]$ is closed and isomorphic to $\Pb^1$. The restriction of $1\in \mathfrak{X}(P_{[v]})$ sends the matrix above to $x_2$,
while the restriction of $\det^{-1}$ from $K$ to $K_{[v]}$ sends that matrix
to $x_1$. We conclude that the existence of a $(K,\kappa)$-equivariant
$\mathcal{O}$-coherent $\mathscr{D}^\lambda_{K/K_{[v]}}$-module is equivalent to
$\lambda,\kappa\in \Z$. Again, an irreducible module is unique.

3) {\it The case of $[1:0:0]$}. The group $K_{[v]}$ is $K$ itself and the orbit $K \cdot [v]$ only consists of the point $[1:0:0]$. The restriction
map $\mathfrak{X}(P_{[v]})\rightarrow \mathfrak{X}(K)$ sends $1$ to $\det^{-1}$.
The conclusions regarding the existence of a twisted equivariant coherent
module and the classification of irreducibles are as in 1).

The above case by case considerations show that in all cases we have
\begin{itemize}
\item three irreducible
$(K,\kappa)$-equivariant $\mathscr{D}^\lambda_{\Pb^2}$-modules when both $\kappa,\lambda$
are integral,
\item two irreducible $(K,\kappa)$-equivariant $\mathscr{D}^\lambda_{\Pb^2}$-modules when $\lambda-\kappa$
is integral but $\kappa$ is not integral,
\item no irreducible $(K,\kappa)$-equivariant $\mathscr{D}^\lambda_{\Pb^2}$-modules when $\lambda-\kappa$
is not integral.
\end{itemize}
We only have one $\mathcal{O}$-coherent $(K,\kappa)$-equivariant $\mathscr{D}^\lambda_{\mathbb{P}^2}$-module
when $\kappa$ and $\lambda$ are both integral. We conclude that we have two non $\mathcal{O}$-coherent
$(K,\kappa)$-equivariant $\mathscr{D}^{\lambda}_{\mathbb{P}^2}$-modules if $\lambda-\kappa$ is an integer
and none otherwise.

{\it Step 2}.
Now let us examine $\Orb\cap \kf^\perp$. We identify $\g$ with $\g^*$ using the trace pairing. The space $\kf^\perp$ consists of the nonzero
matrices of the following form:
$$\begin{pmatrix}0&?&?\\
?&0&0\\?&0&0\end{pmatrix}.$$
Recall that $\Orb$ consists of the nilpotent rank one matrices. So $\Orb\cap \kf^\perp$
consists of two components
\[
  \Orb^1_K = \left\{
     \begin{pmatrix}
       0&?&?\\
       0&0&0\\
       0&0&0
    \end{pmatrix}
  \right\} \cong \C^{2*}\setminus \{0\}
  \quad
   \Orb^2_K = \left\{
  \begin{pmatrix}
  0&0&0\\
  ?&0&0\\
  ?&0&0
  \end{pmatrix}
  \right\} \cong \C^{2}\setminus \{0\},
\]
with the usual linear actions of $\operatorname{GL}_2$ on $\C^2$ and its dual $\C^{2*}$.
We claim that the condition $\lambda-\kappa\in \Z$ is equivalent to the condition of
Corollary \ref{Cor:equiv_condition}.
Take $\chi\in \Orb^2_K$ to be the unit matrix $E_{21}$. The corresponding group $K_Q$
is $\{\operatorname{diag}(x,x,x^{-2})| x\in \C^\times\}$.
The element $\rho_L$ from Definition \ref{defi:half_character}
is easily seen to lie in $\Z+1/2$. And since the period is $\lambda+3/2$, the condition
$\lambda-\kappa\in \Z$ is indeed equivalent to the condition in Corollary \ref{Cor:equiv_condition}.
The case of $\Orb^1_K$ is similar. 

Thanks to Corollary \ref{Cor:equiv_condition}, if $\HC_{\Orb^i_K}(\mathscr{D}^\lambda(\Pb^2))^{K,\kappa}\neq \{0\}$
then $\lambda-\kappa\in \Z$.
Our case by case analysis tells us for each of the orbits $\Orb_K=\Orb^i_K,i=1,2,$ in $\Orb\cap \kf^\perp$ we have
\begin{align*}%\HC_{\Orb_K}(D^\lambda(\Pb^2))\neq 0\Leftrightarrow \lambda-\kappa\in \Z
%\Rightarrow
\HC_{\Orb_K}(\mathscr{D}^\lambda(\Pb^2))\xrightarrow{\sim}\operatorname{Rep}(K_Q, \lambda|_{\kf_Q}- \kappa_Q).
\end{align*}

Note that $\lambda - \kappa \in \Z$ is the condition for $\kappa$ to have
a strong $(K,\kappa)$-action on $\mathcal{E}$, this also follows from
Lemma \ref{Lem:homog_classif}.
So every weakly $\C^\times$-equivariant twisted local system on $X^\theta\setminus \{0\}$
is $\A$-quantizable for all quantizations $\A$ of $\C[X]$.

{\it Step 3}. The goal of this step is to show that every weakly $\C^\times$-equivariant twisted local
system on $X^\theta\setminus \{0\}$ is strongly $\A$-quantizable for any quantization $\A$. This will finish the proof. Note that the category of weakly $\C^\times$-equivariant twisted local systems on
$X^\theta\setminus \{0\}$ is semisimple so we can restrict to irreducible objects $\mathcal{E}$ there.

Consider the closed $K$-orbit of the single point $[1:0:0] \in \Pb^2$ as in 3). The unique $(K,\kappa)$-equivariant irreducible $\mathscr{D}^\lambda_{\Pb^2}$-module supported on this point has a natural filtration such that the associated graded is a line bundle over the characteristic variety of this $\mathscr{D}^\lambda_{\Pb^2}$-module, which is the cotangent vector space of $[1:0:0]$ in $\Pb^2$. Note that this cotangent vector space is mapped isomorphically to the closure $\overline{\Orb}^1_K \subset \overline{\Orb}\cap \kf^\perp$ via the resolution map $T^*\Pb^2 \to \overline{\Orb}$. By identifying this $\mathscr{D}^\lambda_{\mathbb{P}^2}$-module with its global sections, we see that the corresponding HC $(\A,\theta)$-module $M$ has a filtration whose associated graded module $\gr M$ is a line bundle over $\overline{\Orb}^1_K \cong \C^{2*}$. Since $\overline{\Orb}^1_K$ is normal and the natural map $\gr M \to \Gamma((\gr M)|_{\Orb^1_K})$ is an isomorphism, we conclude that $\Gamma((\gr M)|_{\Orb^1_K})$ is Cohen-Macaulay and coincides
with the associated graded of $\Gamma(M|_{\Orb^1_K})$ for any $\lambda \in \C$ satisfying $\lambda - \kappa \in \Z$.
%\footnote{There's a, likely, expositional issue here and below. Lemma 5.2 refers to a very particular %filtration, uniquely determined, essentially. We don't check that the filtration we use is
%that distinguished filtration. Most likely, this is not needed, while we have the coincidence for
%some filtration, we have it for the distinguished filtration. This is clear when our HC modules are %classical -- $a_2,c_2$ types because the $K$-types coincides. But this still requires an argument %somewhere.   \textcolor{red}{Answer} (Shilin): for the $\mathscr{D}$-modules supported on the closed $K$-orbit, the %two filtrations always coincide. }

To analyze the HC-module supported on $\overline{\Orb}^2_K$, we take $\varepsilon: g \to (g^t)^{-1}$ to be the outer automorphism of $\SL_3$ which preserves $K=\operatorname{GL}_2$. We also denote the induced automorphism of the Lie algebra $\slf_3$ by $\varepsilon$, so that $\varepsilon$ commutes with $\theta$ and interchanges the two $\operatorname{GL}_2$-orbits in $\Orb\cap \kf^\perp$. Since  $\A_\mu = \mathscr{D}^{\mu + \frac{3}{2}}(\mathbb{P}^2)$, the involution $\varepsilon$ induces an isomorphism $\A_\mu \cong \A_{-\mu}$ of algebras ($\A_0 = \mathscr{D}^{ \frac{3}{2}}(\mathbb{P}^2)$ is the canonical quantization of $\overline{\Orb}$). Twisting the $\A_\mu$-module structures and the $K$-actions by $\varepsilon$ gives equivalences of categories $\HC(\A_\mu,\theta)^{K,\kappa} \cong \HC(\A_{-\mu},\theta)^{K,-\kappa}$ and $\HC_{\overline{\Orb}^1_K}(\A_\mu,\theta)^{K,\kappa} \cong \HC_{\overline{\Orb}^2_K}(\A_{-\mu},\theta)^{K,-\kappa}$ since $\varepsilon$ interchanges $\Orb^1_K$ and $\Orb^2_K$. Therefore we end up with the same conclusion for HC-modules supported on $\overline{\Orb}^2_K$ as for those  HC-modules supported on $\overline{\Orb}^1_K$  above.
\end{proof}

\subsubsection{The case of $\kf=\mathfrak{so}_3$}\label{SS:sl3_so3}
Here $\theta=-\sigma$ for an outer involution $\sigma$ of $\mathfrak{sl}_3$.

\begin{Lem}\label{Lem:sl3_so3_obstructive}
The following claims are true:
\begin{enumerate}
\item The variety $X^\theta\setminus \{0\}$ is $(\C^2\setminus \{0\})/\Z_4$ (where a generator of $\Z_4$ acts on $\C^2$ by $\sqrt{-1}$). In particular, we consider the ordinary local systems on
    $X^\theta\setminus \{0\}$. Their category is $\operatorname{Rep}(\Z_4)$.
\item The two $\operatorname{SO}_3$-equivariant (corresponding to representation of $\Z_4$, where
$-1$ acts by $1$) local systems on $X^\theta\setminus \{0\}$
are strongly $\A$-quantizable for any quantization $\A$.
\item If $\A=\mathscr{D}^\lambda(\mathbb{P}^2)$  and $\lambda\not\in \Z+\frac{1}{2}$, then any $\A$-quantizable
local system is $\operatorname{SO}_3$-equivariant.
\item If $\A=\mathscr{D}^\lambda(\mathbb{P}^2)$  and $\lambda\in \Z+\frac{1}{2}$, then only one of the two non-$\operatorname{SO}_3$-equivariant
irreducible local systems is $\A$-quantizable. Which one, depends on the parity of $\lambda$. The $\A$-quantizable irreducible local system is strongly $\A$-quantizable if and only if
$\lambda\in \{1/2,3/2,5/2\}$.
\end{enumerate}
\end{Lem}
\begin{Rem}
	The case when $\lambda=3/2$ in part (4) is well-known, see \cite{RawnsleySternberg} and \cite{Torasso}. It was also pointed out by Vogan in \cite[Example 12.4]{Vogan}. 
%\footnote{S.Y.: Probably this should be mentioned earlier in the Introduction. Also the whole lemma is known to Vogan (not surprisingly) and was %communicated to me long time ago. I don't quite remember the details. The solution may involved coherent continuation or cohomological induction. %But we need to compute the $K$-types of the represnetations, so maybe our way is more convenient.} 
\end{Rem}

\begin{proof}
{\it Step 1}.
Let $K=\operatorname{SL}_2$ mapping to $\operatorname{SL}_3$ with image $\operatorname{SO}_3$.
Corollary \ref{Cor:a2_c2_involution} shows that the HC $(\A_\lambda,\theta)$-modules are
just usual HC $(\A_\lambda,K)$-modules.

In this case we have two $K$-orbits in $\mathbb{P}^2$ that correspond to the non-degenerate and
isotropic loci with respect to the form. Let  the $K$-invariant quadratic form
on $\C^3$ be given by $q(x_1,x_2,x_3)=x_1^2+x_2^2+x_3^2$.  Then the representatives
are $[0:1:0],[1:\sqrt{-1}:0]$. Let $\tau_n$ denote the unique irreducible representation of $\SL_2$ of dimension $n$.

1) The case of $[v]=[0:1:0]$. Then $K_{[v]}$ is  the normalizer $N_K(T_K)$ of a maximal torus $T_K$ in $K$ and the $K$-orbit $Q_o := K \cdot [v]$ is open in $\Pb^2$. The Lie algebra $\kf_{[v]} = \operatorname{Lie}(T_K)$ has no nonzero
$K_{[v]}$-invariant characters. Since $N_K(T_K) / T_K \simeq \Z_2$, we use Lemma \ref{Lem:homog_classif} to conclude that for all values of $\lambda$, there are two irreducible $K$-equivariant $\mathcal{O}$-coherent
$\mathscr{D}^\lambda_{K/K_{[v]}}$-modules $\phi_\chi$, where $\chi \in  \Z_2$ is regarded as a character of $K_{[v]}$. Then $\phi_\chi$ are actually $\SO(3)$-equivariant. Let $i: Q_o \hookrightarrow \Pb^2$ denote the inclusion map, then $\I(\lambda, Q_o, \phi_\chi): = i_+ \phi_\chi = i_* \phi_\chi$ are $\mathscr{D}^\lambda_{\Pb^2}$-modules. Let $a\in \{0,1\}$ be such that $\chi$ sends $1+2\Z \in \Z/2\Z$ to $(-1)^a$,
we write $\chi_a$ for $\chi$ to indicate the dependence. Set $M_{2a} = \Gamma(\Pb^2, \I(\lambda, Q_o, \phi_\chi))$. An explicit calculation using Frobenius reciprocity shows that there is an isomorphism
  \begin{equation}\label{eq:K_Q_open}
  	 M_{2a} \big|_K \simeq \operatorname{Ind}_{K_{[v]}}^K \chi_a = \bigoplus_{n \equiv 2a+1 \, (\operatorname{mod} 4)} \tau_n,
\end{equation}
of $K$-representations.

2) The case of $[v]=[1:\sqrt{-1}:0]$. Here the stabilizer is the Borel subgroup of
$\operatorname{SL}_2$,  therefore the $K$-orbit $Q_c: = K \cdot [v]$ is closed and isomorphic to $\Pb^1$. The restriction map $\mathfrak{X}(P_{[v]})
\rightarrow \mathfrak{X}(K_{[v]})$ sends $1$ to $2$. So if $\lambda$
is an integer, then we have a unique irreducible HC module and it is $\operatorname{SO}_3$-equivariant.
When $\lambda\in \frac{1}{2}+\Z$, then we still get a unique irreducible but it is only
$\operatorname{SL}_2$-equivariant -- the action of the center is nontrivial.

{\it Step 2}. Let us construct the HC module corresponding to the closed orbit explicitly.
Let $j: Q_c \hookrightarrow \Pb^2$ denote the inclusion map. Let $\mathscr{D}^\lambda_{Q_c} $ denote the sheaf of twisted differential operators over $Q_c$ corresponding to the restriction of $\lambda$ to
the stabilizer of the $K$-orbit $Q_c$, see the discussion before Lemma \ref{Lem:fin_many_orbit_classif}.
%The sheaf of twisted differential operators over $Q_c$ induced by $j$ is \[ D^\lambda_{Q_c} := \{ A \in \Dscr_\lambda ~|~ [A, \mathscr{I}_{Q_c}] \subset \mathscr{I}_{Q_c} \cdot \Dscr_\lambda \} / \mathscr{I}_{Q_c} \cdot \Dscr_\lambda, \]where $\mathscr{I}_{Q_c}$ is the ideal sheaf of $Q_c$.
Define the $\mathscr{D}^\lambda_{\Pb^2}$ - $\mathscr{D}^\lambda_{Q_c}$-bimodule
\[ \mathscr{D}^\lambda_{\Pb^2 \leftarrow Q_c} := j^{-1} \mathscr{D}^\lambda_{\Pb^2} \otimes_{j^{-1} \OO_{\Pb^2}} \omega_{Q_c/\Pb^2},  \]
where $\omega_{Q_c/\Pb^2} = j^* \omega_{\Pb^2}^{-1} \otimes_{\OO_{Q_c}} \omega_{Q_c} = \OO_{\Pb^1}(4)$ is the relative canonical bundle of $Q_c \simeq \Pb^1$ in $\Pb^2$, which is also the normal bundle $N_{Q_c}$ of $Q_c$. Each $\lambda \in \frac{1}{2} \Z$ determines a  $K$-equivariant line bundle $L_\lambda = \OO_{\Pb^2}(-2 \lambda)$ on $Q_c$, which is a $(\mathscr{D}^\lambda_{Q_c}, K)$-module. We push $L_\lambda$ forward to $\Pb^2$ via $j$ to get a $\mathscr{D}^\lambda_{\Pb^2}$-module
\[  \mathcal{I}(\lambda, Q_c, L_\lambda) = j_+(L_\lambda) := j_* (\mathscr{D}^\lambda_{\Pb^2 \leftarrow Q_c} \otimes_{\mathscr{D}^\lambda_{Q_c}} L_\lambda ).  \]
We set $M(\lambda, Q_c, L_\lambda) = \Gamma(\Pb^2, \I(\lambda, Q_c, L_\lambda))$. There is a natural $K$-stable ascending filtration $\{ \mathcal{F}_k \I \}_{k \geqslant 0}$ of coherent $\OO_{Q_c}$-submodules in $\mathcal{I}(\lambda, Q_c, L_\lambda)$, where $\mathcal{F}_k \I$ is the subsheaf of local sections of degree of transversal derivatives
less or equal than $k$, so that $\mathcal{F}_0 \I = \omega_{Q_c/\Pb^2} \otimes_{\OO_{Q_c}} L_\lambda$ and  the associated graded sheaf is
  \[ \gr \I(\lambda, Q_c, L_\lambda) = \bigoplus_{k \geqslant 0} S^k N_{Q_c} \otimes  \omega_{Q_c/\Pb^2} \otimes  L_\lambda \simeq \bigoplus_{k \geqslant 0} \OO_{\Pb^1}(-2\lambda + 4k + 4). \]
When $\lambda \in \frac{1}{2} \Z$ and $\lambda \leqslant \frac{5}{2}$, all the line bundle summands in the last sum have vanishing higher cohomology. Therefore the higher cohomology of $\mathcal{F}_k \I$ also vanish and $\{ \mathcal{F}_k \I \}_{k \geqslant 0}$ induces an ascending $K$-stable good filtration of $M(\lambda, Q_c, L_\lambda)$, such that there is an isomorphism
  \begin{equation}\label{eq:K_Q_closed}
	\gr M(\lambda, Q_c, L_\lambda) \simeq \bigoplus_{k \geqslant 0} \Gamma(\Pb^1, \OO_{\Pb^1}(-2\lambda + 4k + 4)) = \bigoplus_{k \geqslant 0} \tau_{-2\lambda + 4k + 5}.
\end{equation}
of $K$-representations.

{\it Step 3}. Let us conclude the classification of irreducible HC modules.
When $\lambda\not\in \frac{1}{2}\Z$,
we get two irreducible $\SO_3$-equivariant $\mathscr{D}^\lambda_{\mathbb{P}^2}$-modules, coming from case 1) above. When $\lambda$ is integral, one of the two $\operatorname{SO}_3$-equivariant $\mathscr{D}^\lambda_{\mathbb{P}^2}$-modules $\I(\lambda, Q_o, \phi_\chi)$ from case 1) is irreducible and the other one has a $\OO_{\mathbb{P}^2}$-coherent submodule whose quotient is the $\operatorname{SO}_3$-equivariant $\mathscr{D}^\lambda_{\mathbb{P}^2}$-module $\I(\lambda, Q_c, L_\lambda)$ from case 2). Therefore we end up with two irreducible $\SO(3)$-equivariant $\mathscr{D}^\lambda_{\mathbb{P}^2}$-modules that are not $\mathcal{O}$-coherent.
When $\lambda\in \frac{1}{2}+\Z$ we get two $\SO(3)$-equivariant modules and one which is only $\SL_2$-equivariant.

\begin{comment}
	When $\lambda$ is integral, one of the three $\operatorname{SO}_3$-equivariant $\mathscr{D}^\lambda_{\mathbb{P}^2}$-modules is $\mathcal{O}$-coherent
	(this is one of the two  from case 1)). When we mod it out we end up
	with two $\SO(3)$-equivariant modules. When $\lambda\not\in \frac{1}{2}\Z$,
	we still get two $\SO(3)$-equivariant modules, coming from case 1 above.
	And when $\lambda\in \frac{1}{2}+\Z$
	we get two $\SO(3)$-equivariant modules and one which is only $\SL(2)$-equivariant.
\end{comment}

{\it Step 4}.
Now we examine $\Orb\cap \kf^\perp$.  It consists of all symmetric matrices of rank $1$. These matrices form one $K$-orbit $\Orb_K$. The stabilizer in $\SO(3)$ is the semidirect product of the subgroup of strictly upper triangular matrices (where we now realize $\operatorname{O}(3)$ as the subgroup of matrices stabilizing the quadratic form $x\mapsto x_2^2+x_1x_3$) and the subgroup $\{\operatorname{diag}(\epsilon,1,\epsilon)| \epsilon^2=1\}$. So the orbit
is $(\C^2\setminus \{0\})/\mathbb{Z}_4$. This proves (1) of the lemma.

One conclusion is that if we take $K=\operatorname{SO}_3$, the functor  $\bullet_{\dagger}$ associated
to the open $K$-orbit  is an equivalence. For $K=\operatorname{SL}_2$ it is not: in the interesting case
of $\lambda\in \frac{1}{2}+\Z$, we have exactly one object in the target (identified with $\Z/4\Z\operatorname{-mod}$) which is not in the image.

This proves part (3) and the part of (4) about quantizable twisted local systems.

{\it Step 5.} Let us prove the claims about strongly quantizable local systems ((2) and the remaining part of (4)).
We can choose $e \in \Orb_K$ so that
  \[  K_e = \left\{
     \begin{pmatrix}
  	   t & x\\
  	   0 &t^{-1}\\
      \end{pmatrix}
      \,\Bigg| \, t^4 = 1
  \right\}
  \]
We only consider those representations of $K_e$ that are trivial on the identity component. There are four irreducible ones given by
  \[  \chi_k
    \begin{pmatrix}
  	  t & x\\
  	  0 &t^{-1}\\
    \end{pmatrix}
    = t^k, \quad k = 0,1,2,3. \]
Let $L_k$ denote the corresponding $K$-equivariant line bundles over $\Orb_K$. Then we have isomorphisms of $K \times \cm$-representations
  \begin{equation} \label{eq:K_Lk}
  	 \Gamma(\Orb_K, L_k) = \operatorname{Ind}_{K_e}^K \chi_k = \bigoplus_{n \equiv k+1 \, (\operatorname{mod} 4)} \tau_n.
\end{equation}

We write $\mu$ for $\lambda-3/2$ and $\A_\mu$ for the quantization with period $\mu$, so that
$\A_\mu=\mathscr{D}^{\lambda}(\mathbb{P}^2)$.
The outer automorphism $\varepsilon: g \to (g^t)^{-1}$ of $\SL_3$ fixes $\SO(3)$ and its action on $\slf_3$ coincides with $-\theta$. Moreover, it preserves $\Orb_K$ and maps each line bundle  $L_k$ to itself. Again twisting the $\A_\mu$-module structures by $\varepsilon$ gives equivalences of categories $\HC(\A_\mu,\theta)^K \cong \HC(\A_{-\mu},\theta)^K$ and $\HC_{\overline{\Orb}_K}(\A_\mu,\theta)^K \cong \HC_{\overline{\Orb}_K}(\A_{-\mu},\theta)^K$. Therefore we can focus on the case when $ \mu = \lambda- \frac{3}{2} \leqslant 0$.

When $\lambda \notin \frac{1}{2} \Z$, the global section modules $M_{2a}$ from case 1) are irreducible as $(\A_\mu, \theta)$-modules.  By the last part of the proof of Lemma \ref{Lem:HC_bijection2}, there exists a $K$-stable good filtration on each $M_{2a}$ such that the natural map $\gr M_{2a} \to \Gamma(\gr M_{2a}|_{\Orb_K})$ is injective. By the K-type decompositions \eqref{eq:K_Q_open} and \eqref{eq:K_Lk}, we see that this map is actually an isomoprhism and $\gr M_{2a}|_{\Orb_K} \simeq L_{2a}$, $a = 1,2$. 
  %\footnote{This special filtration should concides with the one given by orders of poles along $Q_c$. Also I suspect the quantizaiton of $L_0$ and $L_2$ is the pushforward of the two $K$-equivariant $\mathcal{O}$-coherent $\mathscr{D}^\lambda_{K/K_{[v]}}$-modules in case 1) via $i_{!}$. -- Shilin; I.L.: any particular relevance of this footnote?}

When $\lambda \in \frac{1}{2} + \Z$ and $\lambda \leqslant \frac{3}{2}$, consider the irreducible $(\A_\mu, \theta)$-module $M(\lambda, Q_c, L_\lambda)$ from case 2) and denote it by $M_k$, where $k \in \{1, 3\}$ and $k \equiv 2 \lambda \, (\operatorname{mod} 4)$. Recall that $M_k$ has a natural $K$-stable filtration. By \eqref{eq:K_Q_closed} and \eqref{eq:K_Lk}, we have $\Gamma(\gr M_k |_{\Orb_K}) \simeq L_{k}$, $k \in \{1, 3\}$. We have $\gr M_k \simeq \Gamma(\Orb_K, L_{k})$ only when $\lambda = \frac{3}{2}$ and $\lambda = \frac{1}{2}$.

This finishes the proof.
%The case when $\lambda < $
\end{proof}

\subsection{Type $c_2$}\label{SS_c2}
Here we consider the algebra $\A=D(\C^2)^{\{\pm 1\}}$,
the only quantization of $\overline{\Orb}$, where
$\Orb$ is the minimal nilpotent orbit in $\mathfrak{sp}_4$.
We claim  that $\A$ is a simple algebra. Here is a more general result.

\begin{Lem}\label{Lem:simplicity}
Suppose that $\tilde{\A}$ is a simple algebra and a  filtered quantization of $\C[\tilde{X}]$, where
$\tilde{X}$ is a conical symplectic singularity. Further, let $\Gamma$ be a finite
group acting on $\tilde{\A}$ by filtered algebra automorphisms such that the induced
action on $\tilde{X}$ is faithful. Then the algebras $\tilde{\A}\#\Gamma$ and $\tilde{\A}^\Gamma$
are Morita equivalent and both are simple.
\end{Lem}
\begin{proof}
The proof is standard. For $\gamma\in \Gamma$,  let $\tilde{\A}_\gamma$ denote the bimodule obtained
from the regular bimodule by twisting with $\gamma$ on the right. This $\tilde{\A}$-bimodule is simple,
and $\tilde{\A}_\gamma\not\cong \tilde{\A}_{\gamma'}$ for $\gamma\neq \gamma'$, the latter can be seen by comparing the associated varieties of these bimodules (that are the graphs of $\gamma,\gamma'$ in $X\times X$). We conclude that the only $\Gamma$- and $\tilde{\A}$- sub-bimodules in $\tilde{\A}\#\Gamma$ are
$\{0\}$ and $\tilde{\A}\#\Gamma$. Therefore, $\tilde{\A}\#\Gamma$ is a simple algebra. The remaining claims follow from here.
\end{proof}
%$\Weyl(V)^\Gamma$
%is simple for any finite dimensional symplectic vector space $V$
%and any finite group $\Gamma$ of symplectomorphisms.
%Indeed, $W(V)$ is simple, hence so is $W(V)\#\Gamma$, and hence
%W(V)^\Gamma$ is Morita equivalent to $W(V)\#\Gamma$, and, therefore is also simple.
In this section we will apply this lemma to $\tilde{X}=\C^4$, its only quantization $\A$, the Weyl algebra,
and $\Gamma=\{\pm 1\}$.

The algebra $\mathfrak{sp}_4$ has two involutions up to conjugation:
for one we have $\kf=\mathfrak{sl}_2\times \mathfrak{sl}_2$
and for the other $\kf=\mathfrak{gl}_2$. For the former,
$\kf^\perp\cap \Orb=\varnothing$, so we only consider the case
$\kf=\mathfrak{gl}_2$.

\begin{Lem}\label{Lem:c2_unobstructive}
Let $X:=\C^4/\{\pm 1\}$ and $\theta$ be the graded Poisson anti-involution corresponding to the
symmetric subalgebra $\mathfrak{gl}_2$. The pair $(X,\theta)$ is unobstructive.
\end{Lem}

We will prove this lemma after some preparation.

\subsubsection{Lifts of HC modules}
This section contains a general result used to classify irreducible $(\A,\theta)$-modules
in this and the subsequent two cases.

Consider the following situation. Let $\tilde{X}$ be a conical symplectic singularity
such that $\operatorname{codim}_{\tilde{X}} \tilde{X}^{sing}\geqslant 4$.
%For example,
%we can take a $G$-equivariant cover $\widetilde{\Orb}$ of $\Orb$ such that $\Orb$ satisfies
%(\ref{eq:codim_condition}) and set $\widetilde{X}:=\operatorname{Spec}(\C[\widetilde{\Orb}])$.
%This variety is indeed a conical symplectic singularity, see \cite[Lemma 2.5]{HC_symp}.


Let $\tilde{\A}$ be a filtered
quantization of $\tilde{X}$ equipped with an anti-involution $\theta$ and a compatible $\Z/d\Z$-grading, let
$\zeta$ denote the corresponding order $d$ automorphism of $\tilde{\A}$. Let $\Gamma$ be
a group acting on $\tilde{\A}$ by filtration preserving automorphisms so that
\begin{enumerate}
\item  the action of $\Gamma$ commutes with $\theta$,
\item
and  the induced action of $\Gamma$ on $\tilde{X}^{reg}$
is free.
\end{enumerate}

Let $\tilde{Y}$ be an irreducible component of $\tilde{X}^{\theta}$.

Set $X:=\tilde{X}/\Gamma, \A:=\tilde{\A}^\Gamma$. We have the induced anti-involution
of $\A$ to be also denoted by $\theta$. Let $Y$ be the image of $\tilde{Y}$ in $X$,
this is an irreducible component of $X^\theta$.

It makes sense to speak about $\Gamma$-equivariant $\Z/d\Z$-graded HC $(\tilde{\A},\theta)$-modules.
Denote their category by $\HC^\Gamma(\tilde{\A},\theta,\zeta)$. We can consider the subcategory
$\HC^\Gamma_{\Gamma\tilde{Y}}(\tilde{\A},\theta,\zeta)$ of all modules supported on $\Gamma\tilde{Y}$
and its quotient $\overline{\HC}^\Gamma_{\Gamma\tilde{Y}}(\tilde{\A},\theta,\zeta)$ of all modules with full support.

Taking the $\Gamma$-invariants, we get a functor $\HC^\Gamma(\tilde{\A},\theta,\zeta)\rightarrow
\HC(\A,\theta,\zeta)$ that restricts to $\HC^\Gamma_{\Gamma\tilde{Y}}(\tilde{\A},\theta,\zeta)\rightarrow
\HC_{Y}(\A,\theta,\zeta)$ and then descends to
\begin{equation}\label{eq:equiv_HC_invariants}
\overline{\HC}^\Gamma_{\Gamma\tilde{Y}}(\tilde{\A},\theta,\zeta)\rightarrow
\overline{\HC}_{Y}(\A,\theta,\zeta).
\end{equation}


The main result of this section is as follows.

\begin{Lem}\label{Lem:lift}
The functor (\ref{eq:equiv_HC_invariants}) is an equivalence.
%Let $M$ be a HC $(\A,\theta)$-module supported on $Y$. Then there is a HC $(\tilde{\A},\theta)$-module %$\tilde{M}$ supported on $\Gamma \tilde{Y}$
%whose microlocalization to $\tilde{X}^{reg}$ coincides with that of
%$\tilde{\A}\otimes_{\A}M$.\footnote{I.L.: I should write a better version: one with a category %equivalence. This will simplify the exposition in the subsequent sections.}
\end{Lem}
\begin{proof}
Let us write $\pi$ for (\ref{eq:equiv_HC_invariants}). Note that $\pi$ is exact and, since
the action of $\Gamma$ on $X^{reg}$ is free, faithful. We claim that $\tilde{\A}\otimes_{\A}\bullet$
is right inverse of $\pi$, this will finish the proof. What needs to be proved is that for
$M\in \HC_Y(\A,\theta,\zeta)$, there is an object $\tilde{M}\in \HC^\Gamma_{\Gamma\widetilde{Y}}(\tilde{\A},\theta,\zeta)$
whose microlocalization to $X^{reg}$ coincides with that of $\tilde{\A}\otimes_{\A}M$.

Clearly, $\tilde{\A}\otimes_\A M$ is supported on $\Gamma\tilde{Y}$. We will work with the
modified Rees algebras and modules, $M_\hbar,\A_\hbar, \tilde{\A}_\hbar$. Note that for
$a\in \A_\hbar^{-\theta},b\in \tilde{\A}_\hbar, m\in M$ we have
$$a(b\otimes m)=[a,b]\otimes m+ b\otimes am=\hbar [(\hbar^{-1}[a,b])\otimes m+ b\otimes (\hbar^{-1}(am))].$$
It follows that for $b'\in  \A_\hbar^{-\theta}\tilde{\A}_\hbar$ (equivalently -- due to $[\tilde{\A}_\hbar,\tilde{\A}_\hbar]\subset \hbar \tilde{\A}_\hbar$-- for $b'\in  \tilde{\A}_\hbar\A_\hbar^{-\theta}$) and $m'\in \tilde{\A}_\hbar\otimes_{\A_\hbar}M_\hbar$ we have
$b'm\in \hbar (\tilde{\A}_\hbar\otimes_{\A_\hbar}M_\hbar)$.
We write $(\tilde{\A}_\hbar\otimes_{\A_\hbar}M_\hbar)^{reg}$ for  the microlocalization of the $\tilde{\A}_\hbar\otimes_{\A_\hbar}M_\hbar$ to $\tilde{X}^{reg}$.
Set $\tilde{M}_\hbar:=\Gamma((\tilde{\A}_\hbar\otimes_{\A_\hbar}M_\hbar)^{reg})$.
%Arguing as in the proof of Lemma \ref{Lem:HC_bijection2}, we see that this module is finitely generated.

We observe that $\tilde{M}_\hbar$ is a finitely generated $\tilde{\A}_\hbar$-module and the kernel and
cokernel of the natural homomorphism $\tilde{\A}_\hbar\otimes_{\A_\hbar}M_\hbar
\rightarrow \tilde{M}_\hbar$ are supported on $\tilde{X}^{sing}$. This is because $\operatorname{codim}_{\tilde{Y}}(\tilde{Y}\setminus \tilde{X}^{reg})\geqslant 2$, compare
to the proof of Lemma \ref{Lem:HC_bijection2}. So, it remains to check that
$\tilde{M}_\hbar$ is a HC $(\tilde{\A}_\hbar,\theta)$-module. This will follow if we show that
$\tilde{\A}_\hbar^{-\theta}\tilde{M}_\hbar\subset \hbar \tilde{M}_\hbar$.

Note that the
$\tilde{\A}_\hbar$-module $\tilde{M}_\hbar$ is supported on $\Gamma\tilde{Y}$. Its
microlocalization to $\tilde{X}^{reg}$ coincides with $(\tilde{\A}_\hbar\otimes_{\A_\hbar}M_\hbar)^{reg}$.
By the construction, we have $b' \tilde{m}\in \hbar \tilde{M}_\hbar$ for all $b'\in \tilde{\A}_\hbar\A_\hbar^{-\theta}$
 and $\tilde{m}\in \tilde{M}_\hbar$. Fix a $\C^\times$- and  $\theta$-stable affine subset $U\subset X^{reg}$
 and its preimage $\tilde{U}\subset \tilde{X}$. We have $b'\tilde{m}\in \hbar^d\tilde{M}_\hbar(\tilde{U})$ for all $b'\in \tilde{\A}_\hbar(\tilde{U})\A_\hbar(U)^{-\theta}$
and $\tilde{m}\in \tilde{M}_\hbar(\tilde{U})$. 

We claim that $b'\tilde{m}\in \hbar\tilde{M}_\hbar(\tilde{U})$
for all $b'\in \tilde{\A}_\hbar(\tilde{U})^{-\theta}$. This will follow once we show that
$\tilde{\A}_\hbar(\tilde{U})\A_\hbar(U)^{-\theta}\supset\tilde{\A}_\hbar(\tilde{U})^{-\theta}$.
The latter equality will follow from its classical analog:
\begin{equation}\label{eq:classic_analog}\C[\tilde{U}]^\theta
\C[U]^{-\theta}=\C[\tilde{U}]^{-\theta}.
\end{equation}
 And that equality will follow if we prove its analog
for the completions at an arbitrary point $u\in Y\cap U$. We write $\hat{A}$ for
$\C[U]^{\wedge_u}$. Then $$\C[U]^{-\theta,\wedge_u}=\hat{A}^{-\theta},
\C[\tilde{U}]^{\theta,\wedge_u}=(\hat{A}^\theta)^{\oplus \Gamma},
\C[\tilde{U}]^{-\theta,\wedge_u}=(\hat{A}^{-\theta})^{\oplus \Gamma}.$$
The completed analog of (\ref{eq:classic_analog}) follows.

So we see that $b'\tilde{m}\in \hbar\tilde{M}_\hbar(U)$ for all $b'\in \tilde{\A}_\hbar(U)^{-\theta}$
and $\tilde{m}\in \tilde{M}_\hbar(U)$. It follows that $\tilde{\A}_\hbar^{-\theta} \tilde{M}_\hbar\subset \hbar \tilde{M}_\hbar$. This finishes the proof.
\end{proof}

\subsubsection{Proof of Lemma \ref{Lem:c2_unobstructive}}
\begin{proof}
{\it Step 1}.
Consider the algebra $\tilde{\A}=D(\C^2)$ so that $\A$ is the invariants in $\tilde{\A}$
under the action of $\{\pm 1\}$. We will view $\tilde{\A}$ as the Weyl algebra
of the symplectic vector space $V:=\C^4$. The anti-involution $\theta$ of $\A$
comes from an anti-symplectic involution of $V$. To give such an involution is the
same thing as to decompose $V$ into the direct sum of two Lagrangian subspaces,
$L_+$ and $L_-$, the involution is $1$ on $L_+$ and $-1$ on $L_-$. In particular,
there are two irreducible components of $X^\theta$, these are $L_+/\{\pm 1\}$
and $L_-/\{\pm 1\}$.

{\it Step 2}. We claim that for each of $L_{\pm}$, there are two
non-isomorphic simple HC $(\A,\theta)$-modules supported at $L_\pm/\{\pm 1\}$ and the category
$\HC_{L_{\pm}/\{\pm 1\}}(\A,\theta)$ is semisimple.

To prove this claim observe first that
the involution $\theta$ admits two lifts to $\tilde{\A}$ that differ by
$-1$. The fixed point set in $V$ for one of them is $L_+$ and for the
other it is $L_-$. Both are $\{\pm 1\}$-stable.

The algebra $\A$ is simple by Lemma \ref{Lem:simplicity}. So $\bar{\HC}(\A)=\HC(\A)$
thanks to Lemma \ref{Lem:lag_antiinvolution}.
We take $Y=L_+$.
Thanks to Lemma \ref{Lem:lift}, $\HC_{L_{+}/\{\pm 1\}}(\A,\theta)\xrightarrow{\sim}
\HC^{\pm 1}_{L_+}(\tilde{\A},\theta)$. Analogously to Lemma \ref{Lem:wHC_formal_Weyl},
every HC $(\tilde{\A},\theta)$-module supported on $L_+$ is the direct sum of several copies of
$\C[L_+]$. So taking the annihilator of $L_+\subset \tilde{\A}$ defines an equivalence
$\HC^{\pm 1}_{L_+}(\tilde{\A},\theta)\xrightarrow{\sim} \operatorname{Rep}(\{\pm 1\})$.
This proves the claim in the beginning of the step.
\begin{comment}

Let us write $\Gamma$ for the group $\{\pm 1\}$ that acts on $\tilde{\A}$.

The algebra $\A$ is simple, see the beginning of Section
\ref{SS_c2}. So $\bar{\HC}(\A)=\HC(\A)$ and the semisimplicity
claim follows because $\bar{\HC}(\A)=\HC_{L_{+}/\{\pm 1\}}(\A,\theta)\oplus
\HC_{L_{-}/\{\pm 1\}}(\A,\theta)$,  and each $\HC_{L_{+}/\{\pm 1\}}(\A,\theta)$
is semisimple as $\bullet_{\dagger,\chi}$
is a full embedding, see Section \ref{SS_Prop_restr}.

The involution $\theta$ admits two lifts to $\tilde{\A}$ that differ by
$-1$. The fixed point set in $V$ for one of them is $L_+$ and for the
other it is $L_-$. Let us deal with simples in
$\HC_{L_{+}/\{\pm 1\}}(\A,\theta)$. Then we take the lift of $\theta$ whose fixed
point locus is $L_+$, denote it also by $\theta$. Let $M$ be such a simple object.
Consider $\tilde{\A}\otimes_\A M$. By Lemma
\ref{Lem:lift}, there is an object $\tilde{M}\in
\HC^\Gamma(\tilde{\A},\theta)$ whose microlocalization to $\C^4\setminus \{0\}$
coincides with that of $\tilde{\A}\otimes_\A M$. Since $\tilde{\A}$ is simple, we have a $\Gamma$-equivariant $\tilde{\A}$-module homomorphism
$\tilde{M}\cong \tilde{\A}\otimes_\A M$.
%So $\tilde{\A}\otimes_\A M$
%is the sum of several copies of $\tilde{\A}/\tilde{\A}L_+$.

 So we get
a functor $\HC(\A,\theta)\rightarrow \HC^\Gamma(\tilde{\A},\theta)$.
It has right adjoint (and left inverse), $\tilde{M}\rightarrow \tilde{M}^\Gamma$.
This functor is inverse thanks to a Morita equivalence between
$\tilde{\A}\#\Gamma$ and $\A=\tilde{\A}^\Gamma$.

On the other hand, we have a category equivalence $\operatorname{Rep}(\Gamma)
\xrightarrow{\sim} \HC^\Gamma(\tilde{\A},\theta), U\mapsto U\otimes (\tilde{\A}/\tilde{\A}L_+)$,
where we consider $\tilde{\A}/\tilde{\A}L_-$ with its natural equivariant structure. This finishes the
proof of the lemma.
\end{comment}

{\it Step 3}. In Step 2 we have seen that the simple objects in
$\HC(\A,\theta,\zeta)$ are precisely the $\{\pm 1\}$-isotypic components in $\tilde{\A}/\tilde{\A}L_-$.
So they have good filtrations whose associated graded modules are the isotypic components
of $\C[L_{\pm}]$. In particular, both associated graded modules are Cohen-Macaulay.
\end{proof}

\subsection{The case of $X=\C^4/\mathbb{Z}_3$}\label{SS_tau_check}
Set $\Gamma=\mathbb{Z}_3,\tilde{\A}=D(\C^2)$.
The variety $X$ is $\mathbb{Q}$-factorial terminal, and the only filtered quantization of $X$ is $\A=\tilde{\A}^\Gamma$, see, e.g., \cite[Section 3.6]{orbit}.

The goal of this section is to prove the following result.

\begin{Lem}\label{Lem:tau_unobstructive}
Let $X:=\C^4/\mathbb{Z}_3$. For any graded Poisson anti-involution $\theta$ of $\C[X]$,
the pair $(X,\theta)$ is unobstructive.
\end{Lem}
\begin{proof}
{\it Step 1}.
Here we are going to classify the anti-symplectic involutions of $X$ that commute
with the standard contracting $\C^\times$-action, see Corollary \ref{Cor:tau_a_2S_2_involution}. The classification will be up to an action of $\operatorname{GL}_2$.
Note that we can lift the involution from $X$ to an anti-symplectic transformation of
$\C^4$ that commutes with $\C^\times$, equivalently, linear. This lift, together with
the $\Gamma$-action gives a linear action on $\C^4$ of a group $\hat{\Gamma}$ of order $6$ that contains
$\Gamma$ as a normal subgroup. So we need to consider two cases:

{\it i}) $\hat{\Gamma}=\Z_3\times \Z_2$. In this case, the lift of $\theta$ is unique, denote it also
by $\theta$. Let $L_\pm\subset \C^4$ denote the eigenspaces for $\theta$ with eigenvalues
$1,-1$. They are $\Gamma$-stable. There are three choices for the pair $(L_+,L_-)$
(up to $\operatorname{GL}_2$-conjugacy). The first two are when $L_+,L_-$ are the
eigenspaces for $\Gamma$. The remaining choice is when $L_+$ intersects both of these
eigenspaces by 1-dimensional spaces. In all of these cases, $X^\theta=\C^2/\Gamma$
(where the actions of $\Gamma$ are by $\zeta\mapsto \operatorname{diag}(\zeta,\zeta)$
in the first case and $\zeta\mapsto \operatorname{diag}(\zeta,\zeta^{-1})$).


{\it ii}) $\hat{\Gamma}=S_3$. In this case, a lift of $\theta$ permutes the eigenspaces for $\Gamma$
and this determines a lift uniquely up to $\operatorname{GL}_2$-conjugation. The preimage of
$X^\theta$ in $\C^4$ is the union of three Lagrangian subspaces permuted simply transitively
by $\Gamma$. The regular locus of $X^\theta$ is $\C^2\setminus \{0\}$.

{\it Step 2}. We now show that all (twisted) local systems on $X^\theta\setminus \{0\}$
are strongly quantizable.
We start with case {\it i}). This is very similar to the case of $c_2$ considered above
and is based on Lemma \ref{Lem:lift}.
For each of the three involutions, the category of HC modules is semisimple and has
three simple objects. All of these simples admit a good filtration whose associated graded
is Cohen-Macaulay.

Consider case {\it ii}). Since $X^\theta\cap X^{reg}=\C^2\setminus \{0\}$, there is no twist.
The category of regular local systems is semisimple with one simple object, the structure
sheaf of $X^\theta\cap X^{reg}$. Thanks to the discussion after Definition \ref{defi:loc_sys_quantized}, it is enough to show that the structure sheaf is strongly quantizable.
%It follows that the category $\HC(\A,\theta)$ is semisimple and the number of simple objects
%is equal to $0$ or $1$. On the other hand, we can produce an object in $\HC(\A,\theta)$.
Fix a lift of $\theta$ and consider the corresponding subspaces $L_\pm$.
Then $\tilde{\A}/\tilde{\A}L_-$ is an irreducible HC $(\A,\theta)$-module
(the irreducibility follows because the associated variety of this module is irreducible,
the multiplicity is $1$, and the algebra $\A$ is simple).
Its usual filtration is good, and the associated graded is Cohen-Macaulay.
\end{proof}

\subsection{The case of $a_2/S_2$}\label{SS_a2S_2}
Here we write $\tilde{X}$ for the closure of the  minimal nilpotent orbit in
$\mathfrak{sl}_3$ and $\Gamma$ for the group (of two elements) generated by the
outer involution of $\mathfrak{sl}_3$ so that $X=\tilde{X}/\Gamma$. Note
that this variety is $\mathbb{Q}$-factorial terminal. Hence it has a unique quantization, $\A=\tilde{\A}^\Gamma$, where
$\tilde{\A}$ is the canonical quantization of $\tilde{X}$, i.e.,
$\tilde{\A}=\Gamma(\mathscr{D}^{3/2}_{\mathbb{P}^2})$.

To state the main result we need to understand the symmetry of the situation. Note that $\operatorname{SO}_3$ acts on $X$ and $\A$. The actions are Hamiltonian. We consider the grading
on $\C[X]$ with $d=1$, it is restricted from the grading on $\C[X]$.  So $\theta$ preserves $\mathfrak{so}_3$ and therefore normalizes $\operatorname{SO}_3$
acting on $\C[X]$. It follows, in particular, that $X^\theta$ is stable under the action of a maximal
torus of $\operatorname{SO}_3$ whose Lie algebra is fixed by $\theta$, denote it by $T_0$.

\begin{Lem}\label{Lem:a2S2_strongly_quantizable}%\footnote{For the main theorem of the next section, the strong equivariance will need to be checked explicitly -not pleasant. Or cheat using Atlas to check the algebraic cases: one will just need to compare the number of HC modules to the number of twisted local systems.}
Let $\chi$ be a character of $T_0$.
Any graded $(T_0,\chi)$-equivariant twisted local system on $X^\theta\setminus \{0\}$
is strongly quantizable.
\end{Lem}

In the process of proof we will fully classify HC $(\A,\theta)$-modules.

\subsubsection{Classification of anti-involutions}\label{SSS_a2S2_involutions}
Here we are going to classify Poisson anti-involutions of $X$ that commute
with the obvious contracting $\C^\times$-action. Again, the anti-Poisson
involution lifts to $\tilde{X}$ giving an action of an order $4$ group
$\tilde{\Gamma}$. Note that the generator of $\Gamma$ is not a square of
any Poisson anti-automorphism of $X$. Indeed, such an anti-automorphism
must come from an anti-isomorphism of $\mathfrak{sl}_3$ hence its square is an inner
automorphism.

It follows that
$\tilde{\Gamma}\cong \Z_2\times \Z_2$, i.e. $\theta$ admits two different
lifts to Poisson anti-involutions of $\tilde{X}$. We need to consider
two cases:
\\

\noindent {\it Case 1}. None of the lifts comes from the anti-involution $x\mapsto -x$
of $\mathfrak{sl}_3$.
Since the generator
of $\Gamma$ is an outer involution we recover both  Poisson anti-involutions
of $\tilde{X}$ considered in Section \ref{SS_type_a2}
as possible lifts of $\theta$. Denote them by $\theta_{in}$
(the inner lift) and $\theta_{out}$ (the outer lift). So $\theta_{in}$
is $-\operatorname{Ad}(g)$ for a non-trivial involution $g\in \operatorname{SO}_3$.
And $\theta_{out}$ is given by $x\mapsto \operatorname{Ad}(g)(x^t)$.

Now we describe the fixed point locus $X^\theta$. It is the union of
$\tilde{X}^{\theta_{in}}/\Gamma$ and $\tilde{X}^{\theta_{out}}/\Gamma$.
Recall that $\tilde{X}^{\theta_{in}}$ consists of two irreducible components.
These components are permuted by $\Gamma$. So $\tilde{X}^{\theta_{in}}/\Gamma$
is an irreducible variety whose smooth locus is $\C^2\setminus \{0\}$.

Now we consider the $\tilde{X}^{\theta_{out}}/\Gamma$. In analysing its
structure it will be convenient for us to make a modification. Namely,
we can assume that $\theta_{out}$ sends $x$ to $x^t$ and $\Gamma$
acts by $x\mapsto -\operatorname{Ad}(g)(x^t)$, where $g$ is an involution in $\operatorname{SO}_3$.
This choice is conjugate to the initial one. So $\tilde{X}^{\theta_{out}}$
consists of symmetric nilpotent rank one matrices. Such matrices
have the form
$$\begin{pmatrix} x_1^2& x_1x_2& x_1x_3\\
x_2x_1& x_2^2& x_2x_3\\
x_3x_1& x_3x_2& x_3^2\end{pmatrix},$$
where the numbers $x_1,x_2,x_3$ satisfy $x_1^2+x_2^2+x_3^2=0$ and are determined uniquely up to simultaneous multiplication by $-1$.  We can assume that $g=\operatorname{diag}(1,-1,-1)$.
Then the action of the generator of $\Gamma$ on $\tilde{X}^{\theta_{out}}$ is induced by $(x_1,x_2,x_3)
\mapsto (\sqrt{-1}x_1,-\sqrt{-1}x_2,-\sqrt{-1}x_3)$. It follows that the smooth part of
$\tilde{X}^{\theta_{out}}/\Gamma$ is $(\C^2\setminus \{0\})/\Z_8$.
\\

\noindent {\it Case 2}. One of the lifts of $\theta$ is $x\mapsto -x$. The only lift
that has fixed points in $\tilde{X}$ is $x\mapsto x^t$. So $X^\theta$
is the quotient of $\{x\in \tilde{X}| x=x^t\}$ by $\{\pm 1\}$. The smooth
locus is still $(\C^2\setminus \{0\})/\Z_8$.

\subsubsection{Classification of HC modules}
Since the algebra $\tilde{\A}$ is simple, so is $\A=\tilde{\A}^\Gamma$, Lemma \ref{Lem:simplicity}.

We concentrate on case 1.
We need to treat two (very different) cases. First, consider the HC $(\A,\theta)$-modules supported
on $\tilde{X}^{\theta_{in}}/\Gamma$. This case is treated as case {\it ii}) in
Section \ref{SS_tau_check}. We get a single irreducible HC module whose associated graded
is $\C[\C^2]$.

Now we proceed to HC modules supported on $\tilde{X}^{\theta_{out}}/\Gamma$.
By Lemma \ref{Lem:lift}, taking the $\Gamma$-invariants defines an equivalence
$$\HC^\Gamma(\tilde{\A},\theta_{out}) \xrightarrow{\sim}  \HC_{\tilde{X}^{\theta_{out}}/\Gamma}(\A,\theta).$$

Now we need to understand the structure of
$\HC^\Gamma(\tilde{\A},\theta_{out})$. First, we claim that $\Gamma$ acts trivially
on the set of irreducibles in $\HC(\tilde{\A},\theta_{out})$. Indeed, an irreducible
object is completely determined by the irreducible local system on $(\tilde{X}^{\theta_{out}})^{reg}
=(\C^2\setminus \{0\})/\Z_4$. The action of $\Gamma$ on the irreducible HC modules coincides
with the action on the set of isomorphism classes of irreducible local systems, which is trivial because the nontrivial element in
$\Gamma$ acts by $-1$ on $(\C^2\setminus \{0\})/\Z_4$.

We conclude that there are six irreducibles in $\HC_{\tilde{X}^{\theta_{out}}/\Gamma}(\A,\theta)$.
Each is a $\Gamma$-isotypic component in one of the three irreducible objects in $\HC(\tilde{\A},\theta_{out})$.
In particular, each admits a good filtration whose associated graded is Cohen-Macaulay.

In case 2, the only component of $X^\theta$ is handled completely analogously to $\tilde{X}^{\theta_{out}}/\Gamma$ in case 1.

\begin{proof}[Proof of Lemma \ref{Lem:a2S2_strongly_quantizable}]
By the analysis in this section, all local systems on $[\tilde{X}^{\theta_{in}}/\Gamma]\setminus \{0\}=\C^2 \setminus \{0\}$ are strongly quantizable. So we need to concentrate on the local systems
on $[\tilde{X}^{\theta_{out}}/\Gamma]\setminus \{0\}=(\C^2\setminus \{0\})/\Z_8$. The action of $T_0$
on $\C^2/\Z_8$ is induced from the action of the maximal torus of $\operatorname{SO}_3$ on $\C^2/\{\pm 1\}$. So, under the equivalence of the category of local systems with $\operatorname{Rep}(\Z_8)$,
the $T_0$-equivariant ones correspond to the representations of $\Z_8$, where $\Z_2\subset \Z_8$
acts trivially. The analysis above in this section shows that all such local systems are
strongly quantizable.
\end{proof}

%For each $M_k$, take some $K$-invariant good filtration. As in the last part of the proof of Lemma \ref{Lem:HC_bijection2}, let $M_\hbar^{k, loc}$ be the microlocalization of the Rees modules of $M_k$. Let $\tilde{M}^k_\hbar:=\Gamma(M_\hbar^{k,loc}|_{\Orb})^{fin}$ and let $\tilde{M}^k$ be the specialization of $\tilde{M}^k_\hbar$ at $\hbar = 1$. Note that $M_0$ and $M_2$ are irreducible unless $\lambda \in \Z$. When any of them is reducible, it has a unique proper submodule of infinite dimension whose quotient module is of finite dimension. The image of the natural map $M \to \tilde{M}$ cannot be supported on $\partial \Orb_K$, therefore the map $M \to \tilde{M}$ must be injective.

\section{Classification theorems}\label{S_classif_thm}
\subsection{Main results}\label{SS_main_results}
In this section $G$ is a  simple algebraic group. Let $\sigma$ be an involution of
$G$ and $K$ a connected algebraic group equipped with a homomorphism $K\rightarrow G$ that induces
an isomorphism $\kf\xrightarrow{\sim} \g^\sigma$. We fix a nilpotent orbit $\Orb\subset \g^*$
with codimension of the boundary at least $4$. Further, fix a $K$-orbit $\Orb_K\subset \Orb\cap
\kf^\perp$. Fix a character $\kappa\in \kf^{*K}$ and $\lambda\in H^2(\Orb,\C)$.

Set $X:=\operatorname{Spec}(\C[\Orb])$. The anti-involution $\theta=-\sigma$ of $\g$ induces a
Poisson anti-involution on $X$ also to be denoted by $\theta$.
Let $\Orb'_K$ be a $K$-orbit of codimension $2$ in the
closure of $\Orb_K$ in $X$. For $\chi'\in \Orb'_K$, we have an induced Poisson anti-involution
on a transverse slice $\check{X}$ to $\chi'$ in $X$.

%Consider the following condition:
%
%\begin{itemize}
%\item[($\clubsuit$)] $\check{X}$ is not of type $a_2$ or $a_2/S_2$.
%Furthermore, if $\check{X}$ is of type $a_2$,
%then $\theta$ comes from the inner involution of $\mathfrak{sl}_3$.\footnote{The second case requires %more care as there's some interplay between the twist of the action and equivariance.}
%\end{itemize}
%Current sections 8-10. Plus a discussion of nonzero parameters for type A, and that example
%in $E_6$.

Pick $\chi\in \Orb_K$ and let $K_Q,\kappa_Q$ have the same meaning as in Section \ref{SS:HC_mod_restr}. Recall the restriction functor $\bullet_{\dagger,\chi}: \HC_{\Orb_K}(\A_\lambda,\theta)^{K,\kappa}
\rightarrow \underline{\A}_\lambda\operatorname{-mod}^{K_Q,\kappa_Q}=\operatorname{Rep}(K_Q,\lambda|_{\kf_Q}-\kappa_Q)$. Here the target category
is the category of finite dimensional rational representations of $K_Q$, where $\kf_Q$ acts by
$\lambda|_{\kf_Q}-\kappa_Q$.
The following is the first of our main results.

\begin{Thm}\label{Thm:main1}
Suppose one of the following conditions holds:
\begin{itemize}
\item[(i)] $K\hookrightarrow G$.
\item[(ii)] There are no slices of type $\chi$. And for all slices of type $a_2$, the anti-involution of this slice comes from the inner involution of $\mathfrak{sl}_3$.
\end{itemize}
Then the restriction functor $\HC_{\Orb_K}(\A_\lambda,\theta)^{K,\kappa}\rightarrow
\operatorname{Rep}(K_Q,\lambda|_{\kf_Q}-\kappa_Q)$ is an equivalence.
\end{Thm}



\begin{Cor}\label{Cor:c_2_tau_slices}
Suppose $\g$ is of types $B,C,D$, or $\Orb$ is one of the following orbits in the exceptional types:
\begin{itemize}
\item $(E_6,2A_2+A_1), (E_8, D_4(a_1)+A_1), (E_8, A_3+A_2+A_1), (E_8,A_3+2A_1)$.
\end{itemize}
Then $\HC_{\Orb_K}(\A_\lambda,\theta)^{K,\kappa}\xrightarrow{\sim}
\operatorname{Rep}(K_Q,\lambda|_{\kf_Q}-\kappa_Q)$.
\end{Cor}
\begin{proof}
In the cases listed all slices $\check{X}$ are $\C^4/\{\pm 1\}$ or $\C^4/\mathbb{Z}_3$ or $\C^4$.
For the classical Lie algebras this follows from \cite{KP2}, while for the exceptional
groups this follows from \cite{FJLS1}. We are done by (ii) of Theorem \ref{Thm:main1}.
\end{proof}

The following corollary will be proved in subsequent sections.

\begin{Cor}\label{Cor:type_A_inner}
Suppose that $\g\cong \mathfrak{sl}_n$, and $\kf\neq \mathfrak{so}_n$. Then
$\HC_{\Orb_K}(\A_\lambda,\theta)^{K,\kappa}\xrightarrow{\sim}
\operatorname{Rep}(K_Q,\lambda|_{\kf_Q}-\kappa_Q)$.
\end{Cor}

Note that Theorem \ref{Thm:main1} and Corollaries \ref{Cor:c_2_tau_slices} and
\ref{Cor:type_A_inner} imply Theorem \ref{Thm:omnibus} from Introduction.

The situation when neither of the conditions of Theorem \ref{Thm:main1} hold
is more complicated. In the classical types it only occurs when
$(\g,K)=(\mathfrak{sl}_n, \operatorname{Spin}_n)$ (with $n\geqslant 3$).
We will state a classification result in this case in a later section.
We will also analyze some cases in the exceptional types.

We will adopt the following terminology of \cite{Noel1, Noel2} in this section. 

\begin{defi}\label{defn:sc-admissible}
	If $\Orb_K$ is part of at least one admissible $K$-orbit datum (in the sense of Section \ref{SS_intro_approach}) with respect to the $K$ of an adjoint group $\GR$ (in terms of the complex groups this means that $G$ is of adjoint type and $K=(G^\sigma)^\circ$), we simply say that it is \emph{admissible}. If the same condition holds for the $K$ of some finite covering group of the adjoint group but not the adjoint group itself, then we say that $\Orb_K$ is \emph{sc-admissible}. If $\Orb_K$ is not admissible under the simply connected cover, then we say that it is \emph{not admissible}.
\end{defi}


\subsection{Proof of Theorem \ref{Thm:main1}}
In this section we prove Theorem \ref{Thm:main1}.

We use the notation of Section \ref{SS_image_description}, in particular,
$\Orb_K^i, i=1,\ldots,\ell$, are all codimension $2$ orbits in $\overline{\Orb}_K$.
%Recall that the image is closed under taking direct summands, Proposition
%\ref{Prop:adj_properties}. Let $\Orb_K^{i},i=1,\ldots,\ell,$ be all
%codim $2$ orbits in $\overline{\Orb}_K$, $\check{X}^i$ be the corresponding
%4-dimensional slices in $X$ (at points $\chi^i\in \Orb_K^i$), and $\theta_i$ be induced Poisson %anti-involutions of $X^i$. We write $\check{X}^{i\times}$ for $\check{X}^i\setminus \{0\}$.
%Set $\check{\A}^i_\lambda:=(\A_\lambda)_{\dagger,\chi^i}$, these are quantizations of $\C[\check{X}^i]$.

Pick an irreducible representation $V\in \operatorname{Rep}(K_Q,\lambda|_{\kf_Q}-\kappa_Q)$ and let
$\mathcal{E}:=\mathcal{E}_V$ be the corresponding $(K,\kappa)$-equivariant $\lambda$-twisted local system
on $\Orb_K$. Then it makes sense to speak about the restrictions $\check{\mathcal{E}}^i$
of $\mathcal{E}$ to $\check{X}^{i\times}$.

%The following claim is a consequence of Proposition \ref{Lem:image_description11}.
%
%\begin{Cor}\label{Cor:club_suit}
%Suppose condition ($\clubsuit$) from Section \ref{SS_main_results} holds for all $i=1,\ldots,\ell$.
%Then the restriction functor $\bullet_{\dagger,\chi}$ is an equivalence.
%\end{Cor}
%\begin{proof}
%We just need to check that (2) of Proposition \ref{Lem:image_description11} holds. If $\check{X}^i$
%is $c_2$ (resp., $\C^4/\Z_3$; $a_2$ with inner involution),
%this was checked in the ``Cohen-Macaulay property'' part of Section \ref{SS_c2} (resp.,
%\ref{SS_tau_check}, \ref{SS_HC_sl3_gl2}). The only remaining case is when $\check{X}=\C^4$.
%But here $\check{X}^\theta=\C^2$, we have only one (twisted) local system -- the trivial one,
%and the condition in (2) of Proposition \ref{Lem:image_description11} is obviously satisfied.
%\end{proof}

\subsubsection{Three lemmas}
In the proof of the theorem we will need three lemmas.

%It remains to consider the case $K\hookrightarrow G$.
Assume for a moment that $\check{X}^i$ is of type $a_2$. Then the connected component of the  group of graded Poisson automorphisms of $\check{X}^i$ is $\operatorname{PGL}_3$ (indeed, the Lie algebra of this group is $\mathfrak{sl}_3\subset \C[\check{X}^i]$).
And if $\check{X}^i$ is of type $a_2/S_2$, then $\operatorname{SO}_3$ acts. Let
$Q^i$ denote the reductive part of $Z_G(\chi^i)$, this reductive group acts on
$\check{X}^i$ (from the realization of $\check{X}^i$ as a Slodowy slice). As before, we can assume that $Q^i$ is $\sigma$-stable and $\check{X}^i$ is $\theta$-stable.
Let $K^i_Q$ denote the preimage of $Q^i\subset G$ in $K$.

\begin{Lem}\label{Lem:Qi_action}
In the $a_2$ case, the action of $(Q^i)^\circ$ on $\check{X}^i$ factors through $(Q^i)^\circ\twoheadrightarrow
\operatorname{PGL}_3$.
%In the $a_2/S_2$ case,   $(Q^i)^\circ\xrightarrow{\sim}
%\operatorname{SO}_3$ 
%(in the $a_2/S_2$-case)\footnote{S.Y.: Now that we have stated Lemma \ref{Lem:character_integrality} for any covering group, do we still this part? %I.L.: yes -- without it you don't know what the group in that lemma is. Maybe it is the maximal torus in $\operatorname{SL}_2, in which 
%case the lemma doesn't give anything$}.
\end{Lem} 
\begin{proof}
The algebra $\g$ is either type $A$ or exceptional. We divide the proof into two cases.

{\it Case 1}.
Let $\g=\mathfrak{sl}_n$. We can replace $G$ with any connected reductive group such that
$[\g,\g]=\mathfrak{sl}_n$. So we take $G=\operatorname{GL}_n$ instead.

Let $\Orb$ be a coadjoint $G$-orbit with Jordan type $\tau=(\tau_1,\ldots,\tau_k)$ (we append $\tau$
with zeros as necessary). It is easy to
show, for example, using results of \cite{KP1}, that the condition $\operatorname{codim}_{\overline{\Orb}}\overline{\Orb}\setminus \Orb\geqslant 4$ is equivalent to $\tau_i-\tau_{i+1}\leqslant 1$ for all $i$. The codimension $4$ orbits in $\overline{\Orb}$
are in bijection with the indices $i$ such that $\tau_i=\tau_{i+1}+1=\tau_{i+2}+2$.
The corresponding codimension $4$-orbit $\Orb^i$ has Jordan type $$\tau^i=(\tau_1,\ldots,\tau_{i-1},\tau_{i+1},\tau_{i+1},\tau_{i+1},\tau_{i+3},\ldots).$$
Pick $e^i\in \Orb^i$. The reductive part $Q^i$ of the centralizer of $e^i$ is
$\prod_j \operatorname{GL}_{d_j}$, where $j$ is a part that occurs in $\tau^i$
and $d_j$ is its multiplicity. So we can write $Q^i=\operatorname{GL}_3\times \underline{Q}^i$,
where the first factor corresponds to the part $\tau_{i+1}$ and the second factor is the product of
the $\operatorname{GL}$-factors corresponding to the remaining parts.

Now let $e_0\in \mathfrak{gl}_3\subset \mathfrak{q}^i$ be an element from the minimal nilpotent orbit
$\Orb_{min}$.  We claim that $e^i+e_0\in \Orb$. For this it is enough to consider the similar question
when $e^i$ corresponds to the partition $(\tau_{i+1},\tau_{i+1},\tau_{i+1})$. If $\tau_{i+1}=1$, there
is nothing to prove. If $\tau_{i+1}>1$, then it is easy to see that $e^i+e_0$ has the same rank as
$e^i$. Also $(e^i+e_0)^{k}=0$ if and only if $k>\tau_{i+1}$. Since the orbit of $e^i+e_0$ has
$e^i$ in its closure, we see that $e^i+e_0\in \Orb$.

The map
$x\mapsto e^i+x$ embeds $\overline{\Orb}_{min}$ into $(e^i+\mathfrak{z}_{\g}(f^i))\cap \overline{\Orb}$
as a closed subvariety. As the latter is 4-dimensional and irreducible (in fact, normal), we get an
isomorphism. It is clear from the construction that $Q^i$ acts on the slice $\overline{\Orb}_{min}$
via the projection $Q^i\twoheadrightarrow \operatorname{PGL}_3$.

{\it Case 2}. Now suppose $\g$ is exceptional and $\check{X}^i$ is of type $a_2$. One can extract
the information about the pairs $(\Orb,\Orb^i)$ from \cite{FJLS1} and the information about $\mathfrak{q}$
and $\mathfrak{q}^i$, for example, from the tables in \cite[Section 13.1]{Carter}. %The following table contains the result (we only list orbits $\Orb$ whose boundary has codimension at least $4$). We also indicate the type of the slice in $\overline{\Orb}$ (that may fail to be normal).

Suppose, for the sake of contradiction, that the $\slf_3$-factor in $\mathfrak{q}^i$
(or both of such factors in the case when $\Orb$ is $A_2+2A_1$ in $E_6$) act trivially
on $\check{X}^i$. Then we have $\mathfrak{q}^i\hookrightarrow \mathfrak{q}$.
This is clearly impossible in all cases but when $\Orb$ is $A_4+A_1$ in type $E_8$.
In that case, suppose $\mathfrak{q}^i\subset \mathfrak{q}$. The centralizer of a Cartan subalgebra
in $\mathfrak{q}^i$ is a minimal Levi containing $e^i$, i.e., of type $D_4+A_2$. Then we can choose $e$
in the centralizer of a semisimple element in that Levi. This is impossible, because the minimal
Levi containing $e$ is of type $A_4+A_1$ and there's no inclusion of $\mathfrak{sl}_5\times \mathfrak{sl}_2$
into $\mathfrak{so}_8\times \mathfrak{sl}_3$. This contradiction finishes the proof (this conclusion can be also 
extracted from the treatment of these orbits in \cite{FJLS1}).  
%\footnote{S.Y.: The way \cite{FJLS1} find these $a_2$ already proves that the $\slf_3$-factor in $\mathfrak{q}^i$ acts non-trivially on the slice %(see Prop. 3.3 in \cite{FJLS1} and the proof in Section 4.5 there is case-by-case). Also for the $2 a_2$ singularity in $E_6$, both %$\slf_3$-factors act non-trivially on the slice and the component group $Q' / (Q')^\circ \simeq \Z_2$ acts transitively on the irreducbile %components of the slice. See Proposition 3.3, Remark 3.5 and Proposition 2.1 of \cite{FJLS1}.
%I.L.: perhaps, but having a short self-contained proof has its virtues.}
%
%{\it Case 3}. There is only one case of the $a_2/S_2$ slice: $\Orb$ is of type $A_4+A_1$, while
%$\Orb^i$ is of type $A_3+A_2+A_1$ in $E_7$.
%%One can show that the $(K^i)$ is isomorphic to $\operatorname{SO}_3$, \cite[Table 4, Row %20]{Alexeevskii},
%%in the case when the group $G$ is adjoint. We want to show that when $G$ is simply connected, then the %connected component of the reductive %centralizer is $\operatorname{SO}_3$.
%According to the table on page 403 of \cite{Carter}, $\mathfrak{q}^i\cong \mathfrak{sl}_2$.
%The claim that $Q_i^\circ\subset G$ is isomorphic to $\operatorname{SO}_3$ is equivalent to
%the condition that $h$\footnote{S.Y.: What's $h$? The semisimple element in the $\mathfrak{sl}_2$-triple for $\Orb^i$?} is conjugate to twice the %element in the coweight lattice of $G$
%(the co-root lattice for the simply connected group), this is because there is a representation of
%$G$, where $Q_i^\circ$ acts faithfully.
%Inspecting the tables in \cite[Section 8.4]{CM},
%we see that the corresponding semisimple element in the $\mathfrak{sl}_2$-triple is conjugate to $2\varpi_4$, where $\varpi_4$ is the fundamental %co-weight corresponding to the trivalent vertex.
%The co-weight $\varpi_4$ is in the co-root lattice and we are done.
%%, so $2\varpi_4$ acts
%%with even eigenvalues on every representation of the Lie algebra $E_7$.
%%We conclude that the image of $\operatorname{SL}_2$ corresponding to this  $\mathfrak{sl}_2$-triple is %$\operatorname{SO}_3$ in the simply %connected group as well.
\end{proof}

The second lemma concerns the case of $a_2/S_2$. According to \cite{FJLS1}, the singularity of type $a_2/S_2$ occurs only for the orbit $\Orb$ of type $A_4+A_1$ in $E_7$  along the orbit $\Orb^i$ of type $A_3 + A_2 + A_1$ in its boundary. Suppose $G$ is a  simple algebraic group of type $E_6$. Let $\sigma$ be an involution of $G$ and $K$ a connected algebraic group equipped with a homomorphism $K\rightarrow G$ that induces
an isomorphism $\kf\xrightarrow{\sim} \g^\sigma$.  Let $i$ be such that $\check{X}^i$ has this type.
%We write $\chi^i$ for a point $\chi'$ in the corresponding codimension $2$ $K$-orbit.
%Let $Q^i$ denote the stabilizer of the corresponding $\slf_2$-triple, and $K^i_Q:=K\cap Q^i$.

\begin{Lem}\label{Lem:character_integrality}
The following claims hold:
\begin{enumerate}
\item 
Suppose $\dim \kf^i_Q=1$.  Then the element $\rho_{L^i}$ for the orbit $\Orb^i_K$ (see Definition \ref{defi:half_character})
restricts to character of the 1-dimensional torus $T_0\subset \operatorname{SO}_3$ with Lie algebra $\kf^i_Q$. 
\item Suppose $\dim \kf^i_Q\neq 1$. Then the action of $(K^i_Q)^\circ$ on $\check{X}^i$ factors through $\operatorname{SO}_3$. 
\end{enumerate}
%\footnote{S.Y.: Are we taking $K \subset G$ here? This lemma even works for non-linear covers of $E_{7(7)}$ since $(K^i_Q)^\circ$ are independent of $K$ by computations. S.Y.: Done.}
\end{Lem}

\begin{proof}
	According to the classification of nilpotent orbits for real forms of $E_7$ and their closure relations (see \cite{Dk3, Dk4, Dk5}), if $\Orb$ and $\Orb^i$ intersect $\kf^\perp$, the only possible real form of $E_7$ is $E_{7(7)}$. In this case the closure $\overline{\Orb}_K$ intersect $A_3 + A_2 + A_1$ at four $K$-orbits, labelled by 39, 40, 41 and 42.
The Lie algebra $\mathfrak{q}^i$ is $\mathfrak{sl}_2$ (see, e.g., tables in \cite[Section 13.1]{Carter}), while the Lie algebras $\kf^i_Q$ for $39$ and $40$ are $\slf_2$ and 1-dimensional for $41$ and $42$ by \cite[Table XI, pg. 513]{Dk0}. Since the $G$-orbit $A_3 + A_2 +A_1$ is even, by \cite[Chapter 2]{Schwartz} (see also \cite[Corollary 7.28]{Vogan}), all the four $K$-orbits in $A_3 + A_2 +A_1$ are admissible in the sense of Definition \ref{defn:sc-admissible}. 

Note that $\mathfrak{q}$ is a two-dimensional abelian Lie algebra, \cite[Section 13.1]{Carter}. It follows that the action of $(Q^i)^\circ$
on $\check{X}_i$ is nontrivial. Since $\Orb^i$ is an even orbit, this action factors through $\operatorname{SO}_3$.
In particular,  the conclusion of (1) holds for the $K$-orbits $\#41$ and $42$, while the conclusion of (2) holds for 
the $K$-orbits $\#39$ and $40$.
\end{proof}

The third lemma concerns type $\chi$ (the case (5) from Section \ref{SS_goals}).

\begin{Lem}\label{Lem:type_chi}
Suppose $\Orb$ is the only orbit $A_4+A_3$ in $E_8$ whose closure has type $\chi$ singularity.
The conclusion of (i) of Theorem \ref{Thm:main1} is true in this case.
\end{Lem}

\begin{proof}
	The orbit $\Orb$ of type $A_4+A_3$ in $E_8$ is rigid by \cite[Table 3]{GraafElashvili}  and $\dim_\C \Orb=200$ by \cite[pg. 133]{CM}. 
According to \cite[Appendix D]{LMM}, the ideal $\J$, the kernel of the quantum comoment map for the unique quantization of $\C[\Orb]$
is the maximal ideal with infinitesimal character $\rho/5$, where $\rho$ is the half sum of positive roots of $E_8$. Inside $\Orb$, there is only one $K$-orbit $\Orb_K$ for the split real form $E_{8(8)}$, which is labeled as $\# 57$ in \cite[Table XIV]{Dk0}, and no $K$-orbit for the other real form $E_{8(-24)}$ by \cite[Table XV]{Dk0}. 
Since the adjoint group $G$ of type $E_8$ is simply connected, $K=G^\sigma$ is semisimple and connected by \cite[Table 3]{realgrp}, and corresponds to the adjoint $E_{8(8)}$. Hence $\kappa=0$. By \cite[Table 11, pg. 273]{King}, the $K$-component group of orbit $\# 57$ is trivial, so there is at most one $K$-admissible line bundle over $\Orb_K$.  The Atlas command  \verb|all_parameters_gamma| lists the Langlands parameters of all irreducible HC $(\g,K)$-modules of a given infinitesimal character, and the command \verb|GK_dim| computes the Gelfand-Kirillov dimension of the representation corresponding to a given parameter. We find that there is only one irreducible HC $(\g,K)$-module infinitesimal character $\lambda = \frac{1}{5} \rho$ of Gelfand-Kirillov dimension $100$. Therefore there is the unique irreducible object in $\operatorname{Rep}(K_Q,\lambda|_{\kf_Q}-\kappa_Q)$ and it is in the image of the restriction functor $\bullet_{\dagger,\chi}$, so the claim holds.
\end{proof}


\subsubsection{Proof of Theorem \ref{Thm:main1}} \label{SS:proof_main1}
Let us first deduce the theorem from the lemmas.

\begin{proof}
If (ii) holds, then all pairs $(\check{X}^i,\theta^i)$ except $a_2/S_2$ are unobstructive, by Lemma
\ref{Lem:sl3_gl2_unobstructive}, Lemma \ref{Lem:c2_unobstructive} and Lemma \ref{Lem:tau_unobstructive}.
The claim of the theorem follows from Proposition \ref{Lem:image_description11}.
%Recall that, thanks to Corollary \ref{Cor:club_suit} we only need to consider the case when %$K\hookrightarrow G$.
\begin{comment}
It remains to handle the extension to slices of types $a_2, a_2/S_2$.

We start with the case of $a_2$. First, consider the case when the anti-involution on $\check{X}^i$
comes from the inner involution of $\slf_3$. Note that, by the paragraph after
Lemma \ref{Lem:twisted_local}, $\mathcal{E}$ can be viewed a $(K,\kappa)$-equivariant weakly HC $(\hat{\A}^\Orb_{\param,\hbar},\theta)$-module,
where $\hbar$ acts by $0$ and the Lie algebra action of $(\param^*)^{-\theta}$ is by the projection of
$\lambda$ to $(\param^*)^{-\theta}$\footnote{It doesn't seem that the formula for the twist in that lemma is used anywhere. It's just that a wHC module on the smooth part is a twisted local system and vice versa. Perhaps, a more basic -- and shorter lemma -- will do.}. Consider its pushforward $\underline{M}$. Then consider $(M_0)_{\dagger,\chi^i}$. It is a weakly HC $(\A^i_{\param^i,\hbar},\theta^i)$-module, where
$\param^i\cong \C$, with zero action of $\hbar$ and the Lie algebra action of
$(\param^{i*})^{-\theta}=\param^{i*}$ factoring through a point. So its restriction to
$(\check{X}^i)^\theta=(\C^2\setminus \{0\})\sqcup (\C^{2*}\setminus \{0\})$ is a twisted local system. And its twisted equivariant for the action of $K_Q^i$ with respect to some character of $\kf_Q^i$.
Thanks to Lemma \ref{Lem:Qi_action},
$Q^i$ acts on $\check{X}^i=\overline{\Orb}_{min}$ via an epimorphism to $\operatorname{PGL}_3$.
So $K_Q^i$ must act on $(\check{X}^i)^\theta$ via an epimorphism to $\operatorname{GL}_2$
(note that the latter acts faithfully). We conclude that the restriction of $\check{\mathcal{E}}^i$
is a twisted $\operatorname{GL}_2$-equivariant twisted local system on $(\check{X}^{i})^\theta\setminus \{0\}$\footnote{This is a general fact and should be treated as such, perhaps in the current section 5.
And 4.6 (current 4.5) should spell a connection between wHC modules w. $\hbar=0$ and twisted local systems more explicitly. Overall, there seem to be a lot of claims in the paper that are not used. This needs to be fixed}. Thanks to Section \ref{SS_HC_sl3_gl2}\footnote{Perhaps, each of these sections should have a lemma with the key result}, there is $M^i\in \HC(\check{\A}^i_\lambda, \theta)$ such that
$\gr M=\Gamma((\check{X}^{i})^\theta\setminus \{0\}, \check{\mathcal{E}}^i)$.
\end{comment}

So we need to prove the theorem under assumption (i) or if the singularity $a_2/S_2$ is present. We can assume that the singularity $\chi$
does not occur, thanks to Lemma \ref{Lem:type_chi}. The only two singularities we need to consider
are $a_2$ and $a_2/S_2$.

Consider the case when $\check{X}^i$ is of type $a_2$, and the involution in question is outer (and $K\subset G$).
We note that $\mathcal{E}$ is weakly $K$-equivariant, hence $\check{\mathcal{E}}^i$ is weakly $K_{Q}^i=Q^i\cap K$-equivariant. Recall
that $(Q^i)^\circ$ acts on $\check{X}^i$ via $(Q^i)^\circ\twoheadrightarrow
\operatorname{PGL}_3$, Lemma \ref{Lem:Qi_action}. The corresponding epimorphism $\mathfrak{q}^i\twoheadrightarrow \mathfrak{sl}_3$
is equivariant with respect to the action of the reductive group $\langle\sigma\rangle\ltimes (Q^i)^\circ$. So we can find a $\langle\sigma\rangle\ltimes (Q^i)^\circ$-equivariant
section $\mathfrak{sl}_3\rightarrow \mathfrak{q}^i$ (where $\sigma$ acts as an outer involution on the source) that lifts uniquely to an algebraic group homomorphism $\operatorname{SL}_3\rightarrow (Q^i)^\circ$ whose composition with $(Q^i)^\circ\twoheadrightarrow \operatorname{PGL}_3$ is the natural epimorphism 
$\operatorname{SL}_3\twoheadrightarrow \operatorname{PGL}_3$.   Passing to the invariants of the involutions, we get homomorphisms $\operatorname{SO}_3\rightarrow K_{Q}^i\twoheadrightarrow \operatorname{SO}_3$ whose composition is the identity. The action of $(K_{Q}^i)^\circ$ on $\check{X}^i$
factors through the epimorphism $K_{Q}^i\twoheadrightarrow \operatorname{SO}_3$. It follows that $\check{\mathcal{E}}$ is
weakly $\operatorname{SO}_3$-invariant. Thanks to Lemma \ref{Lem:sl3_so3_obstructive},
$\check{\mathcal{E}}$ is strongly $\check{\A}^i_\lambda$-quantizable. Applying Proposition
\ref{Lem:image_description11} finishes the proof in the $a_2$-case.

Now we handle the case of $a_2/S_2$. Recall, Section \ref{SS_a2S_2}, that there may be two components of $(\check{X}^i)^\theta$, denote them by $Y_{in}$ and $Y_{out}$. If the former is present, then its nonzero locus is $(\C^2\setminus \{0\})$ and there is a HC $(\check{A}^i_\lambda,\theta^i)$-module $M_{in}$ with a good filtration such that $\gr M_{in}=\Gamma(Y_{in}^\times, \check{\mathcal{E}}^i)$. So we just need to show that
there is a HC module $M_{out}$ with $\gr M_{out}=\Gamma(Y_{out}^\times, \check{\mathcal{E}}^i)$.

The action of $\operatorname{SO}_3$ on
$\check{X}^i$. In Case 2 of Section \ref{SSS_a2S2_involutions}, the anti-involution $\theta$
of $\check{X}^i$ is $\operatorname{SO}_3$
is stable. This case is handled as the $a_2$-case above. Suppose now that we are in Case 1 of
Section \ref{SSS_a2S2_involutions}. Then $\theta$ is equivariant with respect to the maximal torus $(K^i_Q)^\circ$ of $\operatorname{SO}_3$.

As was mentioned before Lemma \ref{Lem:character_integrality}, the singularity of type $a_2/S_2$ occurs only for the orbit of type $A_4+A_1$ in $E_7$. This orbit is birationally rigid in the sense of \cite[Definition 1.2]{orbit}, see, e.g., \cite[Proposition 3.1]{Fu}.
So $\lambda=0$.
As noted in the proof of Lemma  \ref{Lem:character_integrality}, we only need to consider the real form $E_{7(7)}$. In this case, the group $K$ is semisimple (see \cite[Table 3]{realgrp}).
So $\kappa=0$. It follows that $\check{\mathcal{E}}^i$
is $(K^i_Q,-\rho^i_{L})$-equivariant. By Lemma \ref{Lem:character_integrality},
$\rho^i_L$ comes from a character of $(K^i_Q)^\circ$. By Lemma \ref{Lem:a2S2_strongly_quantizable},
$\check{\mathcal{E}}^i|_{Y_{out}^\times}$ is strongly quantizable. 
%\footnote{S.Y.: I suppose we can cite Lemma \ref{Lem:character_integrality} at the very beginning of the proof and claim that $a_2/S_2$ is %unobstructive in either case (i) or (ii).}

Now the claim that $\bullet_{\dagger,\chi}$ is essentially surjective under the assumptions of the theorem
follows from Proposition \ref{Lem:image_description11}.
\end{proof}





\subsection{Restriction of twisted local systems to slices}\label{SS_loc_sys_slice_restriction}
As we have seen in Section \ref{SS:sl3_so3}, when the conditions of Theorem \ref{Thm:main1} are not satisfied, the functor $\bullet_{\dagger,\chi}$ may fail to be essentially surjective. In this case,
to apply Proposition \ref{Lem:image_description11}, we need to be able to compute the restriction
of the $(K,\kappa)$-equivariant twisted local system $\mathcal{E}$ to $(\check{X}^{i\times})^\theta$.
This is what the present section is concerned with.

Recall, Proposition \ref{Prop:regular_local_systems},
that the category of regular twisted local systems on $\Orb_K$
is equivalent to the category of projective representations of $\pi_1(\Orb_K)$ (with Schur multiplier
determined by the twist). All $(K,\kappa)$-equivariant local systems on $\Orb_K$ are regular, and the category of projective representations is the category $\operatorname{Rep}(K_Q,\lambda|_{\kf_Q}-\kappa_Q)$.
So the Schur multiplier for the projective representation of $K_Q/K_Q^\circ$ (that is a quotient
of $\pi_1(\Orb_K)$) is read off from the character $\lambda|_{\kf_Q}-\kappa_Q$.

Now pick an irreducible component $\check{Y}$ of $(\check{X}^{i})^\theta$ and write
$\check{Y}^\times$ for $Y\setminus \{0\}$. The inclusion $\check{Y}^\times\hookrightarrow \Orb_K$
gives rise to a group homomorphism $\pi_1(Y^\times)\rightarrow \pi_1(\Orb_K)$ and hence
$\pi_1(Y^\times)\rightarrow K_Q/K_Q^\circ$.

The following lemma is a direct consequence of Proposition \ref{Prop:regular_local_systems}.

\begin{Lem}\label{Lem:local_sys_restr1}
Let $\mathcal{E}$ correspond to a projective representation $V$ of $K_Q/K_Q^\circ$.
Then the restriction of $\mathcal{E}$ to $\check{Y}^\times$ corresponds to the pullback of
$V$ to $\pi_1(\check{Y}^\times)$.
\end{Lem}

We will apply this construction in the case when $\check{X}^i$ is of type $a_2$
and its anti-involution $\theta$ corresponds to the outer involution of $\mathfrak{sl}_3$.
By (1) of Lemma \ref{Lem:sl3_so3_obstructive}, we have $\pi_1(Y^\times)\cong \Z_4$.
So what we need to determine is a homomorphism $\Z_4\rightarrow K_Q/K_Q^\circ$.

Here is how we are going to compute this homomorphism.
Choose a base point $\chi$ used
to compute $\pi_1(\Orb_K)$ on the image of $Y^\times$ in $\Orb_K$.
This choice gives rise to an embedding $(K_{Q}^i)_\chi\rightarrow K_\chi$.
%and hence to $(Q^i)_x^\sigma\rightarrow Q^\sigma$.
Recall, Lemma \ref{Lem:Qi_action}, that $(Q^i)^\circ$ acts on $\check{X}^i$ via an epimorphism to $\operatorname{PGL}_3$.
The action of $K_Q^i$ on $\check{X}^i$ factors through that of $Q^i$.
Hence $K_{Q}^{i,\circ}$ acts on $\check{X}^i$ via an epimorphism onto $\operatorname{SO}_3$. There is the unique homomorphism
$\operatorname{SL}_2 \rightarrow K_{Q}^{i,\circ} $ such that the composition
$\operatorname{SL}_2 \to K_{Q}^{i,\circ} \twoheadrightarrow \operatorname{SO}_3$
is the projection $\operatorname{SL}_2\twoheadrightarrow \operatorname{SO}_3$.
The homomorphism $\Z_4\rightarrow K_Q/K_Q^\circ$ we need is the composition 
\begin{equation} \label{eq:hom_Z4}
	\Z_4\hookrightarrow (\SL_2)_\chi \rightarrow (K_{Q}^{i,\circ})_\chi\hookrightarrow K_Q\twoheadrightarrow K_Q/K_Q^\circ.
\end{equation}
We note that the first map is an identification of $\Z_4$ with a maximal reductive subgroup of $(\SL_2)_\chi$
and so is not canonically defined, but its composition with the projection of 
$(\SL_2)_\chi$ to its component group is defined canonically. 

%The homomorphism we need is the composition
%\begin{equation}\label{eq:Z4_KQ_comp_homom}
%Z_4\hookrightarrow (\SL_2)_x\rightarrow
%\end{equation}

\subsection{Type A} \label{SS:type_A}
Throughout this section, $G=\operatorname{SL}_n, \g=\mathfrak{sl}_n$ (unless otherwise
noted), $\Orb$
is an orbit in $\g$ with codimension of the boundary at least $4$, $\tau$
is the corresponding partition. The goal of this section is to complete the classification of irreducible
$(K,\kappa)$-equivariant HC $(\A_\lambda,\theta)$-modules with GK dimension equal to $\frac{1}{2}\dim \Orb$.


Let $I$ be the set of all indices $i$
such that $\tau_{i}=\tau_{i+1}+1=\tau_{i+2}+2$. As mentioned in Case 1
of the proof of Lemma \ref{Lem:Qi_action},
the set $I$ parameterizes codimension 4 orbits in $\overline{\Orb}$.
For $i\in I$, we write $\Orb^i$ for the corresponding orbit, it corresponds
to the partition $\tau^i$ with $\tau^i_i=\tau^i_{i+2}=\tau_i-1$ and $\tau^i_j=\tau_j$
for $j\neq i,i+2$.

We note that, by \cite{KP_normal}, $\overline{\Orb}$ is normal, so $X=\overline{\Orb}$.

\begin{Rem}\label{Rem:outer_involutions}
Now we discuss involutions $\sigma$ of $\mathfrak{sl}_n$ that are outer automorphisms. They correspond
to orthogonal or, for even $n$, symplectic forms on $\C^n$. We note that each irreducible representation
of $\slf_2$ carries a unique (up to proportionality) bilinear form, where $e,f$ are self-adjoint.
This form is symmetric. It follows that for an $\slf_2$-triple $(e,h,f)$ with $\sigma(e)=-e, \sigma(f)=-f$, the multiplicity space of each irreducible representation carries a nondegenerate form
of the same type as the form on $\C^n$. In particular, if $\kf=\mathfrak{sp}_n$, then each multiplicity
space is even dimensional. It follows that there are no codimension $4$ orbits in $\overline{\Orb}$
in this case.
\end{Rem}

\subsubsection{Quantizations and their slices}\label{SS_type_A_quant}
Let us discuss quantization parameters for the quantizations of $\C[\Orb]$.
Let $P$ be the parabolic subgroup in
$G$ corresponding to the partition $\tau^t$ so that $T^*(G/P)$ is a symplectic resolution of
$\overline{\Orb}$. Then the space of quantization
parameters is $H^2(G/P,\C)=\operatorname{Pic}(G/P)\otimes_{\mathbb{Z}} \C$,
\cite[Section 3.2]{orbit}.
This space is described as follows. To each column
of $\tau$ we assign a complex number, we get $\tau_1$ numbers $\lambda_1,\ldots,
\lambda_{\tau_1}$. Then $\operatorname{Pic}(G/P)\otimes_{\mathbb{Z}} \C$
is the space $\C^{\tau_1}$ modulo the line given by $\lambda_1=\lambda_2=\ldots=\lambda_{\tau_1}$.

We can identify
$H^2(G/P,\C)$ with $(\mathfrak{l}/[\mathfrak{l},\mathfrak{l}])^*$.
The subalgebra $\mathfrak{l}$ is the algebra of block diagonal matrices
with blocks of sizes $\tau^t_1,\ldots, \tau^t_{\tau_1}$.
For $\lambda_i$ we take the coefficient in $\lambda$ of the trace of its $i$th block.

The quantization of $X$ corresponding to $\lambda\in (\mathfrak{l}/[\mathfrak{l},\mathfrak{l}])^*$
is $\mathscr{D}^{\lambda+\rho_{G/P}}(G/P)$, where $\rho_{G/P}$ is one half of the character of $\mathfrak{l}$
in $\Lambda^{top}(\g/\mathfrak{p})$.

Now take the codimension $4$ $G$-orbit in $X$ corresponding
to index $i$.
Consider the slice algebra $\check{\A}^i_\lambda$. It is a quantization of $\overline{\Orb}_{min}$,
where $\Orb_{min}$ is the minimal nilpotent orbit in $\mathfrak{sl}_3$, so is of the form
$\mathscr{D}^{\lambda^i+3/2}(\mathbb{P}^2)$ for $\lambda^i\in \C$.

\begin{Lem}\label{Lem:param_corresp_A}
We have $\lambda^i=\lambda_{i}-\lambda_{i+1}$.
\end{Lem}
\begin{proof}
By \cite[Lemma 3.3]{Perv}, the period of the slice quantization $\check{\A}^i_\lambda$
is the pullback of the period of $\A_\lambda$ under the inclusion map $\mathbb{P}^2
\hookrightarrow T^*(G/P)$. So we just need to compute the pullback.

Consider the basis in $\C^n$ labelled by the boxes of the Young diagram $\tau$.
Let $e$ be the element of $\Orb$ that acts by sending the basis element labelled by a
box to the one labelled
by the box directly to the right (and to zero if there's no such).
Let $e^i$ be the nilpotent element defined similarly
to $e$ but from $\tau^i$.

The preimage of $\check{X}^i(\cong \overline{\Orb}_{min})$ in $T^*(G/P)$ is $T^*\mathbb{P}^2$.
The zero fiber $\mathbb{P}^2$ is the parabolic Springer fiber consisting of all
partial flags $\{0\}=V_{\tau_1}\subset V_{\tau_1-1}\subset\cdots \subset\cdots V_{1}=\C^n$ with
$\dim V_i/V_{i+1}=\tau^t_i$ and $e^i V_i\subset V_{i+1}$ for all $i$. These conditions
fix all subspaces in the flag but $V_{i+1}$, while the condition on $V_{i+1}$ is that
it lies between $\operatorname{Im}(e^i)^{\tau_{i+1}-1}$ and $\operatorname{Im}(e^i)^{\tau_{i+1}}$.
The subspace       $V_j/\operatorname{Im}(e^i)^{\tau_{i+1}-1}$ is one-dimensional,
while the space $\operatorname{Im}(e^i)^{\tau_{i+1}-1}/\operatorname{Im}(e^i)^{\tau_{i+1}}$
is three-dimensional. This is how the parabolic Springer fiber is identified with
$\mathbb{P}^2$. So the pullback map $\operatorname{Pic}(G/P)\rightarrow
\operatorname{Pic}(\mathbb{P}^2)$ sends $(\lambda_1,\ldots,\lambda_{\tau_1})$
to the class of $\mathcal{L}^{\otimes \lambda_{i+1}}\otimes (\Lambda^2[\mathcal{V}/\mathcal{L}])^{\otimes \lambda_{i}}$. Here $\mathcal{L}(\cong \mathcal{O}(-1))$ is the tautological line bundle
and $\mathcal{V}$ is the trivial rank 3 vector bundle. Under the identification of the target lattice
with $\Z$ that we use, the class of that line bundle is
$\lambda_{i}-\lambda_{i+1}$.
\end{proof}

\subsubsection{The case of inner involution}\label{SS:inner}
\begin{proof}[Proof of Corollary \ref{Cor:type_A_inner}]
%\footnote{Details with the classification of K-orbits? I.L.: perhaps -- will you write this? Also, is %this true already on the level of $G$-orbits: if $\Orb'\subset \overline{\Orb}$
%both intersect $\mathfrak{k}^\perp$, then codimension is at least $6$ if it's bigger than $2$? In this %case we just need to know the relevant orbits, instead of the classification of the $G$-orbits.}
By Remark \ref{Rem:outer_involutions}, we only need to consider the case when
$\mathfrak{k}=(\mathfrak{gl}_k\times \mathfrak{gl}_{n-k})\cap
\mathfrak{sl}_n$. Here the nilpotent orbits of $K$ in $\mathfrak{k}^\perp$ are parameterized
by ``$ab$-diagrams'' due to \cite[Section 4]{KP_normal}, also see \cite[Section 2.2]{Ohta1}. These are Young diagrams with an $a$ or a $b$ in each box so that there are $k$ $a$'s and $n-k$ $b$'s, and the labels $a$ and $b$ alternate in each row.
By \cite[Theorem 3]{Ohta2}, the closure relation is generated by pairs $(\tau>\tau')$, where $\tau'$ is obtained from
$\tau$ by moving one box up (we draw Young diagrams so that the number of boxes decreases
bottom to top) preserving the $ab$-label.

Let us describe $\mathfrak{k}^i_Q$ for a nilpotent $K$-orbit labelled by a given $ab$-diagram
of shape $\tau^i$.
We have $\mathfrak{k}_Q^i=(\prod_{j=1}^{\tau_1}\mathfrak{k}^i_{Q,j})\cap \slf_n$
with $\mathfrak{k}^i_{Q,j}\subset \mathfrak{gl}_{n_j}$, where $n_j$ is the number of rows of length $j$ in $\tau^i$, determined as follows.  Fix $j$ and consider the rows in $\tau^i$ of length $j$.
Suppose that the first column of these rows contains $n_j^a$ $a$'s and
$n_j^b$ $b$'s.   We have $\mathfrak{k}^i_{Q,j}=\mathfrak{gl}_{n_j^a}\times
\mathfrak{gl}_{n_j^b}$.

Now suppose that $j=\tau_{i+1}$ so that $n_j=3$. By Case 1 in the proof of Lemma 
\ref{Lem:Qi_action}, we have the inclusion $\operatorname{SL}_3\hookrightarrow Q^i$
and the projection $(Q^i)^\circ\twoheadrightarrow \operatorname{PGL}_3$ such that the composition
$\operatorname{SL}_3\rightarrow \operatorname{PGL}_3$ is surjective and $Q^i$
acts on $\check{X}^i$ via its projection to $\operatorname{PGL}_3$. So $\mathfrak{k}^i_Q$
acts on $\check{X}^i$ via the image of $\operatorname{GL}_2$ in $\operatorname{PGL}_3$.
This excludes the case when $\theta^i$ comes from the outer involution.
Now we can use Theorem \ref{Thm:main1} to finish the proof of the corollary.
%Let us describe the embedding of $Q_{K,j}$ into $\operatorname{GL}_k\times \operatorname{GL}_{n-k}$.
%The factor $\operatorname{GL}_{n_i^+}$ embeds into both $\operatorname{GL}_k$ and
%$\operatorname{GL}_{n-k}$ via the sum of several copies of the trivial representation
%and of the tautological representation. The multiplicity of the tautological representation
%for the embedding into $\operatorname{GL}_k$ equals to the number of pluses in
%the rows of length $i$ whose first box is a plus. This equals $\lceil i/2\rceil$.
%For the embedding into $\operatorname{GL}_{n-k}$ the multiplicity similarly equals
%$\lfloor i/2\rfloor$. And for the factor $\operatorname{GL}_{n_i^-}$ the role of
%pluses and minuses gets swapped. In particular, this shows how to restrict
%characters from $K$ to $Q_K$.
%
%Each codimension $4$ orbit gives a slice algebra, say $\underline{\A}$
%and a character $\underline{\kappa}$ obtained from $\kappa$ by a shift
%by $\rho_{K,Q}$.
%Using the discussion of the previous paragraph together with Lemma
%\ref{Lem:param_corresp_A} we can check the conditions in the end of Section
%\ref{SS_HC_sl3_gl2} for each codimension 2 $K$-orbit\footnote{This should be elaborated}.
%Assume, from now on, that these conditions are satisfied for a codimension $2$
%$K$-orbit $\Orb_K$. Let $Y_1=\overline{\Orb}_K$.
\end{proof}

Until the end of the section we concentrate on the case when $\tilde{K}=\operatorname{Spin}_n$ (and $K=\operatorname{PSO}_n$).

\subsubsection{The group $\tilde{K}_Q/\tilde{K}_Q^\circ$}
Let $\Orb_K$ be a $K$-orbit in $\kf^\perp\cap \Orb$.
We want to determine the reductive part (of the centralizer in $\tilde{K}$ of a point in $\Orb_K$) $\tilde{K}_Q$ and, especially, its component group.

For $\bar{K}:=\operatorname{SO}_n$, this is  easy. Let $\bar{K}_Q$ denote the image of $\tilde{K}_Q$ in $\bar{K}$. For $i\in \{1,\ldots,\tau_1\}$, let $m_i$
denote the multiplicity of $i$ in $\tau$. Consider the group $\prod_{i=1}^{\tau_1}
\operatorname{O}_{m_i}$, it is embedded into $\operatorname{O}_n$ via
$(g_i)\mapsto \operatorname{diag}(g_i^{\oplus i})$. This is the reductive part of
the centralizer in $\operatorname{O}_n$. The element  $\operatorname{diag}(g_i^{\oplus i})$
belongs to $\operatorname{SO}_n$ if and only if $\prod_{i\text{ odd}}\det(g_i)=1$.
In particular, $\bar{K}_Q/\bar{K}_Q^\circ\cong (\Z/2\Z)^{\oplus (\tau_1-1)}$ (we cannot have only even
parts in $\tau$ because of the inequality $\tau_i-\tau_{i+1}\leqslant 1$). We also note that outside of the trivial case of $\tau=(1^2)$, every character
of $\bar{K}_Q$ factors through $\bar{K}_Q/\bar{K}_Q^\circ$. If there is a codimension $4$ orbit in
$\overline{\Orb}$, then $(\mathfrak{k}_Q^*)^{\tilde{K}_Q}=\{0\}$.
Hence $\operatorname{Rep}(\tilde{K}_Q,\lambda|_{\kf_Q}-\kappa_Q)=\operatorname{Rep}(\tilde{K}_Q/\tilde{K}_Q^\circ)$.


Recall that $\tilde{K}=\operatorname{Spin}_n$. The group $\tilde{K}_Q/\tilde{K}_Q^\circ$ either coincides with
$\bar{K}_Q/\bar{K}_Q^\circ\cong(\Z/2\Z)^{\oplus \tau_1-1}$ or is a central extension of the latter group by
$\Z/2\Z$. Our first task is to determine which option holds.

\begin{Lem}\label{Lem:comp_group_Spin1}
If $m_i>1$ for some odd $i$, then $\tilde{K}_Q/\tilde{K}_Q^\circ\cong (\Z/2\Z)^{\oplus \tau_1-1}$.
Otherwise, $\tilde{K}_Q/\tilde{K}_Q^\circ$ is a central extension of $(\Z/2\Z)^{\oplus \tau_1-1}$
by $\Z/2\Z$.
\end{Lem}
\begin{proof}
The group $\tilde{K}$ is a central extension of $\bar{K}$ by $\Z/2\Z$. Such an extension is parameterized by a class in $H^2(\bar{K},\Z/2\Z)$,
denote it by $s$. Let $\iota$ denote the inclusion $\bar{K}_Q\hookrightarrow \bar{K}$.
The group $\tilde{K}_Q$ is the central extension of $\bar{K}_Q$ by $\Z/2\Z$, which corresponds
to the class $\iota^*(s)\in H^2(\bar{K}_Q,\Z/2\Z)$. Let $\iota^\circ$ be the composition
of this inclusion with $\bar{K}_Q^\circ\hookrightarrow \bar{K}_Q$. We have
$\tilde{K}_Q/\tilde{K}_Q^\circ\xrightarrow{\sim} \bar{K}_Q/\bar{K}_Q^\circ$ if and only if
$\iota^{\circ*}(s)\neq 0$, otherwise $\tilde{K}_Q/\tilde{K}_Q^\circ$ is a central extension of $(\Z/2\Z)^{\oplus \tau_1-1}$
by $\Z/2\Z$.

Let $s_{m_i}$ denote the class of  the Spin extension of $\operatorname{SO}_{m_i}$. We claim
that $\iota^{\circ*}(s)=\sum_{j \text{ odd}} s_j$. Indeed, the pullback of $s$ under the
obvious inclusion of $\operatorname{SO}_{m_i}$ is $s_{m_i}$. So $\iota^{\circ*}(s)=\sum is_{m_i}$,
which implies our claim because $H^2(?,\Z/2\Z)$ is 2-torsion.

%The first option in the lemma corresponds to the case when the central extension of
%$\bar{K}_Q^\circ$ by $\Z/2\Z$ does not split, and the second option corresponds to
%the split case.
We have $s_{m_i}=0$ if and only if $m_i=1$.
 Now if we have
two connected groups $H_1,H_2$ and a central extension
$$1\rightarrow \Z/2\Z\rightarrow H\rightarrow H_1\times H_2\rightarrow 1,$$
then this extension splits if and only if the induced extensions of both $H_1$ and $H_2$
split. So $\iota^{\circ*}(s)=0$ if and only if
$is_{m_i}=0$ for all $i$. We arrive at the conclusion of the lemma.
\end{proof}

Now we describe $\tilde{K}_Q/\tilde{K}_Q^\circ$ in the case when this group is different from
$\bar{K}_Q/\bar{K}_Q^\circ$. Note that $\bar{K}_Q/\bar{K}_Q^\circ$ embeds
into the subgroup $\operatorname{SO}_{1^n}$ of $\operatorname{SO}_n$ consisting
of all elements that are diagonalizable in a given orthogonal basis. Namely,
we view a generator of $\operatorname{O}_i/\operatorname{SO}_i$ as an
element $\operatorname{diag}(-1,1,\ldots,1)\in \operatorname{O}(i)$. Taking the direct
sum over all $i$ we get the required embedding.
We also have a similarly defined subgroup $\operatorname{O}_{1^n}\subset \operatorname{O}_n$.

Now we define the ``unit Pin group'' $\operatorname{Pin}_{1^n}$. By definition,
it has generators $e_1,\ldots,e_n,-1$, where $-1$ is central and the relations
are $e_i^2=1, e_ie_j=-e_je_i$ for $i\neq j$. So, in particular, every permutation
of $\{1,\ldots,n\}$ gives an automorphism of $\operatorname{Pin}_{1^n}$.

Inside, we have the unit Spin group
$\operatorname{Spin}_{1^n}$ generated by $-1$ and the elements $e_ie_j, i\neq j$.
Then the preimage of $\operatorname{SO}_{1^n}$ in $\operatorname{Spin}_n$
is $\operatorname{Spin}_{1^n}$.
%We note the following relations in
%$\operatorname{Spin}_{1^n}$:
%\begin{equation}\label{eq:Spin_relations}
%(e_ie_j)^2=-1,\quad (e_ie_j)(e_je_k)=e_ie_k,\quad (e_ie_j)(e_ke_\ell)=(e_ke_\ell)(e_ie_j).
%\end{equation}
%The second relation is when  $i,j,k$ are pairwise distinct and the third
%relation is when all $i,j,k,\ell$ are pairwise distinct.

Now we describe the image of $\tilde{K}_Q/\tilde{K}_Q^\circ$ in $\operatorname{Spin}_{1^n}$.
For $i=1,\ldots,\tau_1$, we define $E_i\in \operatorname{Pin}_{1^n}$
as the product  $e_{j+1}\ldots e_{j+i}$, where $j=i(i-1)/2$. So
$E_1=e_1, E_2=e_2e_3$, etc. Then, with a suitable choice of an orthogonal
basis in $\C^n$,  $\tilde{K}_Q/\tilde{K}_Q^\circ$ embeds as
the intersection of the subgroup of $\operatorname{Pin}_{1^n}$
generated by $E_1,\ldots, E_{\tau_1}$, to be denoted by $\Gamma$, with $\operatorname{Spin}_{1^n}$.
Note that the following relations among the $E_i$'s hold:
\begin{equation}\label{eq:E_relations}
\begin{split}
& E_i E_j=(-1)^{ij}E_j E_i,\\
& E_i^2=1, \text{ if } i\equiv 0,1\text{ mod }4,\\
& E_i^2=-1, \text{ else}.
\end{split}
\end{equation}
The elements of the form $E_iE_j$, where $i,j$ are both odd and $i+j$ is divisible by
$4$ are central in $\Gamma$. The square of such an element is equal to $1$. The elements $E_k$
with even $k$ are central as well.
%It follows that $\tilde{K}_Q/\tilde{K}_Q^\circ$ splits
%into the direct product of the following two factors
Note that
\begin{enumerate}
\item[(i)]  The subgroup, $\Gamma_0$, generated
by the elements $E_{4i}$ and $E_{4i+1}E_{4i+3}$ is a  central subgroup
isomorphic to $(\Z/2\Z)^{q}$, where $q=\lfloor \tau_1/4\rfloor+\lfloor (\tau_1-3)/4\rfloor$.
\item[(ii)] The subgroup, $\Gamma_1$, generated by the elements $E_{4i+1}$ is naturally
identified with $\operatorname{Pin}_{(1^r)}$, where $r=\lfloor (\tau_1-1)/4\rfloor$.
\item[(iii)] The subgroup, $\Gamma_2$, generated by the elements $E_{4i+2}$ is
isomorphic to the quotient of $(\Z/4\Z)^{t}$ by $(\Z/2\Z)^{t-1}$ generated
by the differences of the elements of order $2$ in each factor $\Z/4\Z$. Here $t=\lfloor (\tau_1-2)/4\rfloor$.
\end{enumerate}
Note that both (ii) and (iii) have distinguished central elements of order $2$.
Then $\Gamma$ is the quotient of  the product of (i), (ii) and (iii)
by the difference of these elements. 
%\footnote{S.Y.: Does this mean $\Gamma$ is this quotient? Also we don't really need this description of $\Gamma$ below, since you have removed the %part about genuine repns of $\Gamma$. Or maybe in revised version later, we can count the number of quantizable local systems. It would be a good %summary of the somewhat complication computations.}
Define the element $-1$ in the quotient to be the image of either
of the two elements.

Then $\tilde{K}_Q/\tilde{K}_Q^\circ$ is identified with the intersection of $\Gamma$ with $\operatorname{Spin}_{1^n}$.
%The inter intersection will be  denoted by $\underline{\Gamma}_1$, it is generated
%by pairwise products of the generators.


\begin{comment}
Our next task is to describe the irreducible representations of $\tilde{K}_Q/\tilde{K}_Q^\circ$\footnote{Strictly speaking, we don't need this part, can be removed}.
The first task here is to understand the irreducible representations of
$\operatorname{Spin}_{1^n}$. First, there are $2^n$ one-dimensional representations,
the elements $-1$ acts trivially there. Now we study the irreducible
representations, where $-1$ acts by $-1$. We call them {\it genuine}.

\begin{Lem}\label{Lem:Pin_irreps}
If $n$ is odd, then $\operatorname{Spin}_{1^n}$ has a unique genuine irreducible representation
of dimension $2^{(n-1)/2}$. If $n$ is even, then $\operatorname{Spin}_{1^n}$ has two non-isomorphic
genuine irreducible representations, of dimensions $2^{(n-2)/2}$.
\end{Lem}
\begin{proof}
The proof is by induction on $n$. The case of $n=1$ is clear, here $\operatorname{Spin}_1=\Z/2\Z$.
We have a natural embedding $\operatorname{Spin}_{1^{n-1}}\hookrightarrow \operatorname{Spin}_{1^n}$
as the subgroup generated by the pairwise products of the elements $e_1,\ldots,e_{n-1}$. The image is normal and has index $2$.

Suppose $n$ is even. Then the restriction of each genuine
irreducible representation to $\operatorname{Spin}_{(1^{n-1})}$
must be irreducible for the dimension reasons. And we must have two non-isomorphic representations
of $\operatorname{Spin}_{(1^n)}$. The conjugation with the element $e_{n-1}e_n$ swaps these two representations.

Now suppose that $n$ is odd.  The restriction of an irreducible representation to
$\operatorname{Spin}_{(1^{n-1})}$ must contain both irreducible representations of
$\operatorname{Spin}_{(1^n)}$ with the same multiplicity. By dimension reasons,
we have a unique genuine irreducible representation.
\end{proof}

Using this description it is easy to describe the irreducible representations of
$\tilde{K}_Q/\tilde{K}_Q^\circ$. We will only care about genuine ones, those which do not
factor through $\bar{K}_Q/\bar{K}_Q^\circ$. The contribution of the factor $\Gamma_0$
is easy to understand, we ignore it below. There are two options for
a genuine irreducible representation of $\Gamma_1\Gamma_2$:
\begin{enumerate}
\item It is genuine as a representation of $\Gamma_1$ and factors through $\Gamma_2/\{\pm 1\}$.
Such a representation is the tensor product of an irreducible genuine $\Gamma_1$-representation
described by Lemma \ref{Lem:Pin_irreps} and a character of $\Gamma_2$ vanishing on $-1$.
\item It is genuine as a representation of $\Gamma_2$ and factors through
$\Gamma_1/\{\pm 1\}$. It is 1-dimensional (the group $\Gamma_2$ is abelian).
\end{enumerate}
\end{comment}

\subsubsection{Homomorphism $\Z_4\rightarrow \tilde{K}_Q/\tilde{K}_Q^\circ$}\label{SSS_Z4_homom_typeA}
As we have seen in Section \ref{SS_loc_sys_slice_restriction}, the image of $\bullet_{\dagger,\chi}$
is controlled by homomorphisms $\Z_4\rightarrow \tilde{K}_Q/\tilde{K}_Q^\circ$, one for each $\Orb_K^i$.
The goal of this section is to compute these
homomorphisms. We have one homomorphism $\Z_4\rightarrow \tilde{K}_Q/\tilde{K}_Q^\circ$ for each codimension $2$ $K$-orbit in
$\overline{\Orb}_K$, i.e.,
for each part of $\tau$ that occurs with multiplicity $1$ and is different from
$\tau_1$.

Note that the map $\Orb_K\mapsto \Orb$
defines a bijection between the $\operatorname{O}_n$-orbits in $\kf^\perp$
and the $\operatorname{SL}_n$-orbits in $\g^*$. The $\operatorname{O}_n$-orbit
with Jordan type $\tau$ splits into the union of two $\operatorname{SO}_n$-orbit if
and only if all parts in $\tau$ are even. The partition $\tau$ of $\Orb$ never satisfies this condition.
And the partition  $\tau^i$ of a codimension $4$ orbit $\Orb^i\subset \overline{\Orb}$
only satisfies this condition when $\tau=(3,2,1)$ (and $\tau^i=(2,2,2)$).
In this case $\Orb_K^i$ will denote any of the two orbits.
%We will not consider this case.\footnote{Unless Shilin wants so...  S.Y.: Maybe the outer automorphism in %$\operatorname{O}_n / \operatorname{SO}_n$ might help.}
%%\footnote{This requires a check: the claim would be true for $\operatorname{O}_n$-orbits, but not %necessarily $\operatorname{SO}_n$-orbits}.

Our first step is to understand the embedding $\underline{\tilde{K}}^i\hookrightarrow \tilde{K}_Q$, where we write
$\underline{\tilde{K}}^i$ for the  reductive part of the stabilizer $(\tilde{K}^{i,\circ}_Q)_\chi$ of the point $\chi \in \check{Y}^{i\times}$.  %It is more convenient to consider these subgroups in $\operatorname{GL}_n$ instead of $\operatorname{SL}_n$:
%we then get $\tilde{K}_Q=\prod_{i=1}^{\tau_1}\operatorname{O}_i$.
We start with constructing a point in $\check{Y}^{i\times}$, the locus of $\theta$-stable points in
$\check{X}^{i\times}$. We present $\chi$
as a Jordan matrix with blocks of sizes $(\tau_1,\ldots,\tau_k)$ coming in this order.
It is symmetric with respect to the following form. For each
$\ell\in \{1,\ldots,\tau_1\}$, consider the linear span $V_\ell$
of all Jordan vectors that belong to blocks of size $\ell$
so that $\C^n=\bigoplus_\ell V_\ell$ (as a based
space). On each
$V_\ell$ we consider the form given by the unit anti-diagonal matrix
so that every Jordan matrix with blocks of the same size is symmetric. Then
for the form on $\C^n$ we take the direct sum of the forms on $V_\ell$'s.

Set $\ell=\tau_{i+1}$. Set $\chi=\chi^i+e_0$, where $e_0$ is the matrix
$\begin{pmatrix}0&0&1\\0&0&0\\0&0&0\end{pmatrix}\in \mathfrak{gl}_3$ embedded
into $\operatorname{End}(V_{\ell}^{\oplus 3})$ via $x\mapsto x\otimes \operatorname{id}_{V_{\ell}}$,
compare to Case 1 of the proof of Lemma \ref{Lem:Qi_action}.
%For the point $\chi$ we are going to construct only the $V_\ell$-component,
%that is, the only component different from that in $\chi^i$, compare to Case 1 in
%Section \ref{SSS_proof_lem_a2}.
%Namely, we can decompose $V_\ell$ into a $\chi^i$-stable
%direct sum $V_\ell^1\oplus V_\ell^2\oplus V_\ell^3$, where each $V_\ell^i$ is $\C^\ell$ and the %restriction
%of $\chi^i$ to $V_\ell^j, j=1,2,3,$ is a single Jordan block.
%Let $I_{13}$ be the map $V_\ell\rightarrow V_\ell$
%that sends $V_\ell^1,V_\ell^2$ to zero and sends $V_\ell^3$ to $V_\ell^1$ by the identity map
%under the identifications $V_\ell^1,V_\ell^3\cong \C^\ell$. \footnote{S.Y.: You want to say $\chi = \chi' %+ I_{13}$?}For example, for $\ell=2$ we get
%the following blocks of $\chi',\chi$ corresponding to $V_\ell$:
%
%$$\chi'|_{V_2}=\begin{pmatrix}0&1&0&0&0&0\\
%0&0&0&0&0&0\\ 0&0&0&1&0&0\\ 0&0&0&0&0&0\\
%0&0&0&0&0&1\\ 0&0&0&0&0&0\end{pmatrix},
%\chi|_{V_2}=\begin{pmatrix}0&1&0&0&1&0\\
%0&0&0&0&0&1\\ 0&0&0&1&0&0\\ 0&0&0&0&0&0\\
%0&0&0&0&0&1\\ 0&0&0&0&0&0\end{pmatrix}.$$
As was mentioned in that proof,
the Jordan type of $\chi$ is $\tau$, and $\chi\in \check{X}^i$.
And $\chi$ is symmetric with respect to our form, hence $\theta$-stable.
So we are done.

Now we can describe the homomorphism $\underline{\tilde{K}}^i\rightarrow \tilde{K}_Q$. For this, we will
describe the corresponding embedding $\underline{Q}^i\hookrightarrow Q$, where
$\underline{Q}^i$ is the reductive part of the stabilizer of $\chi$ in $Q^i$.
It is convenient to replace $\operatorname{SL}_n$ with $\operatorname{GL}_n$
(to get to the original setting we just restrict the embedding to the subgroups of matrices with determinant $1$).


We have $Q=\prod_{i}\operatorname{GL}_{m_i}\times \GL_{m_{\ell-1}}\times \GL_1
\times \GL_{m_{\ell+1}}$, where in the first product $i$ runs over
$\{1,\ldots,\tau_1\}\setminus \{\ell-1,\ell,\ell+1\}$. The involution preserves
each factor $\GL_{m_i}$,  and the restriction comes from the orthogonal form on $\C^{m_i}$ from Remark \ref{Rem:outer_involutions}. Next, we have
$Q^i=\prod_{i}\operatorname{GL}_{m_i}\times \GL_{m_{\ell-1}-1}\times \GL_3
\times \GL_{m_{\ell+1}-1}$, again with an involution that preserves each factor.
Then $\underline{Q}^i=\prod_{i}\operatorname{GL}_{m_i}\times \GL_{m_{\ell-1}-1}\times \GL_1^2
\times \GL_{m_{\ell+1}-1}$. Here $\GL_1^2$ is embedded into $\GL_3$ via
$(t_1,t_2)\mapsto \operatorname{diag}(t_1,t_2,t_1)$. On the common factor $\prod_{i}\operatorname{GL}_{m_i}$
the embedding $\underline{Q}^i\hookrightarrow Q$ is the identity. Note also that the
embedding $\underline{Q}^i\hookrightarrow Q$
\begin{itemize}
\item[(i)] is compatible with the involution,
\item[(ii)] and intertwines the restrictions of the tautological representation of
$\operatorname{GL}_n$.
\end{itemize}

Thanks to (ii), $\operatorname{GL}_{m_{\ell-1}-1}$ embeds into $\operatorname{GL}_{m_{\ell-1}}$
as the stabilizer of a vector and a complimentary subspace. Thanks to (i), this vector
is non-isotropic. This determines the embedding up to an orthogonal automorphism.
The same consideration applies to $\operatorname{GL}_{m_{\ell+1}-1}
\hookrightarrow\operatorname{GL}_{m_{\ell+1}}$. Similar considerations show that, in the case when $\ell>1$,
the embedding of $\operatorname{GL}_1^2$ into $\GL_{m_{\ell-1}}\times \GL_1
\times \GL_{m_{\ell+1}}$ is as follows:
\begin{equation}\label{eq:embed_GL}
	(t_1,t_2)\mapsto (\operatorname{diag}(t_1,1,\ldots,1), t_2, \operatorname{diag}(t_1,1,\ldots,1)),
\end{equation}
where the first entries in the first and the third components correspond to the non-isotropic
vectors in $\C^{m_{\ell-1}},\C^{m_{\ell+1}}$, when we write diagonal matrices, we use an orthogonal basis in each $\C^{m_i}$. When $\ell=1$, we omit the first component in
the image.

Once the embedding $\underline{Q}^i\hookrightarrow Q$ is known, we get the embedding
$\underline{\bar{K}}^i\hookrightarrow \bar{K}_Q$ (where $\bar{\bullet}$ indicates that we are considering the subgroups in $\operatorname{SO}_n$ (and not in 
$\operatorname{Spin}_n$) by passing to the fixed points under the involution.
And $\underline{\tilde{K}}^i$ is obtained from $\underline{\bar{K}}^i$ by taking the central extension pulled back from $\tilde{K}_Q\twoheadrightarrow \bar{K}_Q$.

Now we proceed to describing the homomorphism $\Z_4\rightarrow \tilde{K}_Q/\tilde{K}_Q^\circ$
associated to $\ell$. Recall, Section \ref{SS_loc_sys_slice_restriction}, that it is computed in the following steps:

\begin{itemize}
\item[(I)] The choice of $\ell<\tau_1$ such that $\ell$ occurs in $\tau$
with multiplicity $1$ (equivalently, of a codimension $2$ $K$-orbit)
corresponds to a direct summand $\mathfrak{so}_3$
of $\kf^i_Q$. We consider the corresponding homomorphism $\operatorname{SL}_2\rightarrow \tilde{K}^i_Q$.
\item[(II)] We pass to the homomorphism $\Z_4\rightarrow \underline{\tilde{K}}^i$ between the reductive
parts of the stabilizers.
\item[(III)] Then we compose $\Z_4\rightarrow \underline{\tilde{K}}^i$ with the homomorphism
$\underline{\tilde{K}}^i\rightarrow \tilde{K}_Q/\tilde{K}_Q^\circ$ to get the required homomorphism
$\Z_4\rightarrow \tilde{K}_Q/\tilde{K}_Q^\circ$.
\end{itemize}

The cases of even $\ell$ and odd $\ell$ behave differently. We start with even $\ell$.
As we have seen in the proof of Lemma \ref{Lem:comp_group_Spin1}, in this case, the pullback of the central extension
from $\operatorname{SO}_n$ to $\operatorname{SO}_3$ is trivial. Equivalently,
the homomorphism $\operatorname{SL}_2\rightarrow \tilde{K}^i_Q$ factors through $\operatorname{SO}_3$.
The resulting homomorphism $\Z_4\rightarrow \tilde{K}_Q/\tilde{K}_Q^\circ$ is not injective and
this is the only thing that we need to know about it: by Lemma \ref{Lem:sl3_so3_obstructive}
the quantization of any $K$-equivariant local system on $\Orb_K$ extends to $\Orb_K^i$.

Now we consider the case when $\ell$ is odd. Here we do get an injective homomorphism
$\operatorname{SL}_2\rightarrow \tilde{K}^i_Q$. The stabilizer in $\operatorname{SO}_3$ of the
point $\chi$ as above is generated by $\operatorname{diag}(-1,1,-1)$. By \eqref{eq:embed_GL}, this element goes to
$E_{\ell-1}E_{\ell+1}\in \tilde{K}_Q/\tilde{K}_Q^\circ$ (up to $\pm 1$) when $\ell>1$ and to
$E_2$ when $\ell=1$.  In particular, the images of the homomorphisms
$\Z_4\rightarrow \tilde{K}_Q/\tilde{K}_Q^\circ$ all land in the abelian normal subgroup
$\Gamma_0\Gamma_2$.


\subsubsection{The main result}
Note that each $E_i$ is actually defined up to a sign. We can change the signs any way we want.
In particular, we arrive at the following characterization of the irreducible representations
of $\tilde{K}_Q/\tilde{K}_Q^\circ$ that lie in the image of $\bullet_{\dagger,\chi}$. Let us start with the special case of canonical quantizations.

\begin{Prop}\label{Prop:spin_HC_classif_canon}
Assume the quantization parameter $\lambda$ is equal to $0$.
For a suitable choice of the generators $E_i$ for even $i$, the following two conditions are
equivalent:
\begin{enumerate}
\item a representation
$V$ of $\tilde{K}_Q/\tilde{K}_Q^\circ$ lies in the image of $\bullet_{\dagger,\chi}$,
\item
for all odd $\ell<\tau_1$, the element $E_{\ell-1}E_{\ell+1}\in \tilde{K}_Q/\tilde{K}_Q^\circ$
(or $E_2$ for $\ell=1$) acts on $V$ without eigenvalue $-\sqrt{-1}$.
\end{enumerate}
\end{Prop}
\begin{proof}
By Lemma \ref{Lem:sl3_so3_obstructive}, all irreducible $\operatorname{SL}_2$-equivariant
local systems $\check{\mathcal{E}}^i$ but one (where the action of $\operatorname{SL}_2$ is faithful)
there is a HC $(\check{\A}^i_\lambda,\operatorname{SL}_2)$-module $M^i$ with
$\gr M^i=\Gamma(\check{\mathcal{E}}^i)$. The remaining local system does not correspond to any HC module.

By changing the generator of $\Z_4$, we can assume that the
remaining local system corresponds to the eigenvalue $-\sqrt{-1}$. We can further assume that
the generator maps to $E_{i-1}E_{i+1}$ (or $E_2$ for $i=1$) -- here we may change their signs.  Now the claim of this proposition
follows from Proposition \ref{Lem:image_description11}.
\end{proof}

We can also describe the image of $\bullet_{\dagger,\chi}$ for a general quantization parameter.
Let $\lambda^i$ be the quantization parameter for $\check{\A}^i_\lambda$, by Lemma
\ref{Lem:param_corresp_A}, we have $\lambda^i=\lambda_i-\lambda_{i+1}$ in the notation
of Section \ref{SS_type_A_quant}.

\begin{Thm}\label{Thm:spin_HC_classif_general}
Fix a choice of the generators $E_i$ for even  $i$ as in Proposition \ref{Prop:spin_HC_classif_canon}.
Then $V\in \operatorname{Im}\bullet_{\dagger,\chi}$ if and only if for all odd $i$ we have the following:
\begin{itemize}
\item if $\lambda^i\not\in \Z$, then the corresponding generator ($E_{i-1}E_{i+1}$ for $i>1$, or
$E_2$ for $i=1$) acts on $V$ with eigenvalues $\pm 1$,
\item if  $\lambda^i\in \Z$, then the corresponding generator acts on $V$ without eigenvalue
$\sqrt{-1}^{2\lambda^i-1}$.
\end{itemize}
\end{Thm}
\begin{proof}
First, we are going to reduce to the case when all integer $\lambda^i$ (for odd $i$) are zero.
The proof of Lemma \ref{Lem:param_corresp_A} shows that there is a line bundle $\mathcal{L}$
on $T^*(G/P)$ such that the restriction of $\mathcal{L}$ to the parabolic Springer fiber $\mathbb{P}^2$
of $\chi^i$ is $\mathcal{O}(-\lambda^i)$ if $\lambda^i$ is an integer. Let $\tilde{\lambda}$ denote the parameter obtained from $\lambda$ by replacing the integral entries by $0$.

For every line bundle
$\mathcal{L}'$ on $G/P$ and every quantization parameter $\lambda'$, consider
$\mathcal{L}'\otimes D^{\lambda'}_{G/P}$, this is a $D^{\lambda'+c_1(\mathcal{L}')}_{G/P}$-$D^{\lambda'}_{G/P}$-bimodule. Let $\A_{\lambda'}(\mathcal{L}')$
denote the global sections of this bimodule. Note that since $\mathcal{L}'\otimes D^{\lambda'}_{G/P}$ is an invertible
bimodule, and $T^*(G/P)\rightarrow X$ is an isomorphism over $\Orb$, the functor of tensoring
with $\A_{\lambda'}(\mathcal{L}')$ is an equivalence $\overline{\HC}(\A_{\lambda'},\theta)\xrightarrow{\sim} \overline{\HC}(\A_{\lambda'+c_1(\mathcal{L}')},\theta)$.

Let $\pi$ denote the projection $T^*(G/P)\twoheadrightarrow G/P$
and $\check{\mathcal{L}}'^i$ denote the restriction of $\mathcal{L}'$
to the copy of $\mathbb{P}^2\subset G/P$ corresponding to $i$.
Suppose $\mathcal{L}'$ is sufficiently ample, in particular,
$H^1(T^*(G/P),\pi^*\mathcal{L}')=0$. Then by \cite[Proposition 3.8]{Perv},
we have $\A_{\lambda'}(\mathcal{L}')_{\dagger,\chi^i}=\check{\A}^i_\lambda(\check{\mathcal{L}}'^i)$,
here we use the restriction functor for HC $\U$-bimodules.
Thanks to Lemma \ref{Lem:restriction_compatibilty}, for any $M\in \HC(\A_{\lambda},\theta)$,
we have a $\U$-linear isomorphism
\begin{equation}\label{eq:restr_iso}
\left(\A_{\lambda'}(\mathcal{L}')\otimes_{\U}M\right)_{\dagger,\chi^i}\cong
\check{\A}^i_\lambda(\check{\mathcal{L}}'^i)\otimes_{\U_{\dagger,\chi^i}}M_{\dagger,\chi^i}.
\end{equation}
Note that the action of $\U$ on $M$ factors through $\A_{\lambda}$, hence the action of
$\U_{\dagger,\chi^i}$ on $M_{\dagger,\chi^i}$ factors through $\check{\A}^i_\lambda$.
The latter is a simple algebra because its period is integral (and hence the twist for
differential operators on $\mathbb{P}^2$ is in $\frac{1}{2}+\Z$). So in the right hand side
of (\ref{eq:restr_iso}) we have $\check{\A}^i_\lambda(\check{\mathcal{L}}'^i)\otimes_{\check{\A}^i_\lambda}M_{\dagger,\chi^i}$.

Let $\psi$ denote the nontrivial automorphism of $\Z_4$ (so that $\psi^2=1$)
and let $\zeta_i$ denote the homomorphism $\Z_4\rightarrow \tilde{K}_Q/\tilde{K}_Q^\circ$
corresponding to the index $i$.
If $\mathcal{L}'^i=\mathcal{O}(d_i)$, then, on the level of representations of $\Z_4$, tensoring with $\check{\A}^i_\lambda(\check{\mathcal{L}}'^i)$ results in twisting with $\psi^{d_i}$. We can choose $\lambda'\in (\mathfrak{l}/[\mathfrak{l},\mathfrak{l}])^*$ such that $\lambda'-\lambda, \lambda'-\tilde{\lambda}$
are sufficiently dominant. We get equivalences
$$\overline{\HC}(\A_{\lambda},\theta)\xrightarrow{\sim} \overline{\HC}(\A_{\lambda'},\theta)
\xleftarrow{\sim} \overline{\HC}(\A_{\tilde{\lambda}},\theta).$$
The resulting self-equivalence $\eta,\tilde{\eta}$ of $\operatorname{Rep}(\tilde{K}_Q/\tilde{K}_Q^\circ)$ satisfy the following properties
for each $i$:
\begin{equation}\label{eq:twists_reps}
\zeta_i\circ\eta=\psi_*^{\lambda_i'-\lambda_i}\circ \zeta_i,
\zeta_i\circ\tilde{\eta}=\psi_*^{\lambda_i'-\tilde{\lambda}_i}\circ \zeta_i.
\end{equation}
Here the subscript $*$ indicates the twist by the corresponding automorphism.
%composing them with the pullbacks under the homomorphism $\Z_4\rightarrow %\tilde{K}_Q/\tilde{K}_Q^\circ$\footnote{Introduce notation like $\psi_i$} is the twist with the nontrivial automorphism
%of $\Z_4$\footnote{Introduce the notation and write the commutative diagram} $\lambda_i'-\lambda_i$
%and $\lambda'_i-\tilde{\lambda}_i$ times.
Equations (\ref{eq:twists_reps}) reduce the claim of the theorem to the case of $\lambda=\tilde{\lambda}$.
This is what we assume until the rest of the proof.

Recall, Lemma \ref{Lem:sl3_so3_obstructive}, that if $\lambda_i$ is not an integer, then the image of $\overline{\HC}(\check{\A}^i_\lambda,\theta)$  in $\operatorname{Rep}(\Z_4)$ is $\operatorname{Rep}(\Z_2)$. Moreover, every $\operatorname{SO}_3$-equivariant local system
is strongly $\check{\A}^i_\lambda$-quantizable, by Lemma \ref{Lem:sl3_so3_obstructive}. We now apply Proposition \ref{Lem:image_description11} to finish the proof of the theorem.
\end{proof}

\subsection{Exceptional types}\label{SS_exceptional}

Throughout this section, $\g$ is a simple Lie algebra of exceptional type, $G$ is its adjoint group (since we do not require $K$
to be a subgroup of $G$, this is a harmless assumption), $\Orb$ is a nilpotent $G$-orbit in $\g$ whose closure contains an orbit $\Orb'$ of codimension $4$, such that $\Orb'$ is open in $\overline{\Orb} \backslash \Orb$ and the singularity of $\overline{\Orb}$ along $\Orb'$ is of type $a_2$, $a_2 / S_2$ or $2a_2$. Note that in all the cases discussed here, $\Orb$ is always birationally rigid in the sense of \cite[Definition 1.2]{orbit}, except for the two orbits of type $A_2 + A_1$ and $A_2+2A_1$ in $\mathfrak{e}_6$ (for the latter one see Section \ref{SS:orbit_a} below).

\begin{Lem}\label{lem:only_a2}
	Under the assumption above, let $X = \spec \C[\Orb]$. Then $X$ contains a unique codimension $4$ symplectic leaf $X'$ with $a_2$ singularity and has no codimension $2$ leaves. Moreover, in this case $X'$ is the only  codimension $4$ leaf, except when $\g = \mathfrak{e}_7$ and $\Orb$ is of type $A_4+A_1$. For this orbit, there is another codimension $4$ leaf with $a_2/S_2$-singularity.
\end{Lem}

\begin{proof}
	The lemma follows by inspecting the closure diagram of nilpotent orbits of Lie algebra of exceptional types and the generic singularities of their closures in \cite{FJLS1}. The only case that requires some care is the orbit $\Orb$ of type $A_2 + A_1$ in $\mathfrak{e}_6$. There is only one orbit $\Orb'$ in $\overline{\Orb}$ with codimension $4$, which is of type $A_2$, and the singularity here is $2a_2$. Let $f: X=\spec \C[\Orb] \to \overline{\Orb}$ be the normalization morphism and $e$ be any point in $\Orb'$. By the discussions in \cite{FJLS1} above Proposition 2.1 there, the action of the component group $A(e')$ (with respect to the adjoint group $G$) of $\Orb'$ at a point $e' \in \Orb'$ on the connected components of $f^{-1}(e')$, i.e., the branches of $\overline{\Orb}$ at $e'$, is transitive. 
	Now $A(e) \simeq S_2$ by the table on page 402 of \cite{Carter}. Therefore the preimage of the orbit $\Orb'$ under the normalization map $f$ is a connected double cover $\tilde{\Orb}'$ of $\Orb'$. This is a symplectic leaf $X'$ in $X$, along which the signularity is $a_2$.
	%The Springer representation attached to $A_2+A_1$ is $64'_p$. The table at the right bottom of \cite[Pg. 603]{BS}.
	 %and its slice $S'$ in $\overline{\Orb}$. By \cite[Remark 3.5]{FJLS1}, each irreducible component of $S'$ corresponds to the minimal orbit closure $\check{X}^i \simeq a_2$ in a unique summand $\slf_3^i \simeq \slf_3$ of $\mathfrak{q}' = \slf_3^1 \oplus \slf_3^2$ and the other summand $\slf_3^{2-i}$ acts trivially on $\check{X}^i$. The two varieties $\check{X}^i $ meet at the origin. Let $X = \spec \C[\Orb]$ denote the normal closure of $\Orb$. The preimage of the orbit $\Orb'$ of type $A_2$ under the normalization map $X \to \overline{\Orb}$ is a connected double cover $\tilde{\Orb}'$ of $\Orb'$ by Lemma \ref{lem:A2+A1_normal} below. The component group of $\Orb'$ is $\Z_2$. The preimage of $S'$ in $X$ can be regarded as the disjoint union of $\check{X}^1$ and $\check{X}^2$ and the singularity of $X$ along $\tilde{\Orb}'$ is isomorphic to a single copy of $\Orb_{min}^{\slf_3}$. Therefore in all cases, $X = \spec \C[\Orb]$ can contain at most one codimension $4$ symplectic leaf with $a_2$ singularity. 
\end{proof}


%Throughout the section, $\g$ is a complex simple Lie algebra of exceptional type, $\g_\R$ is a (non-compact) real form of $\g$ and $G_\R = \operatorname{Int}(\g_\R)$ denotes the real adjoint group associated to $\g_\R$. Fix a maximal compact subgroup $\KR$ of $\GR$ and let $K$ denote the complexification of $\KR$. The list of all connected real simple exceptional groups of adjoint type and their $\KR$ can be extracted from \cite[Table 3]{realgrp}. 
%In the table, $F_4$, $E_6$ and $E_7$ stand for the simply connected compact Lie groups with Lie algebras $\mathfrak{f}_4$, $\mathfrak{e}_6$ and $\mathfrak{e}_7$ respectively.  We also use $\operatorname{Ss}(4n) := \Spin(4n) / \Z_2 $ to denote the semi-spin group (which is not isomorphic to $\SO(4n)$). The notation $K_1 \times K_2 / \Z_n$ is used, e.g., $(\Sp(3) \times \SU(2)) / \Z_2$, $(\operatorname{U}(1) \times E_6) / \Z_3$, $(\Sp(3) \times \SU(2)) / \Z_2$, etc. This means that the center of each $K_i$ contains a canonical copy of $\Z_n$ and we consider the diagonal embedding of $\Z_n$ into $K_1 \times K_2$\footnote{I.L.: I'm not happy with this description. For $n>2$, the group $\Z_n$ has nontrivial automorphisms. So even if the copy of $\Z_n$ is canonical, the diagonal embedding is not. \textcolor{red}{S.Y.} Me neither. It is not clear from the reference. I don't know a better reference. But note that we don't really need the description of $\KR$ below. All we need is the fundamental group and the Lie algebra $\kf_\R$.}.


Now assume $K$ is the fixed subgroup of $G$ with respect to some involution $\sigma$. All such symmetric pairs $(G,K)$ of noncompact type correspond bijectively to connected real simple exceptional groups of adjoint type, whose information can be extracted from \cite[Table 3]{realgrp}. Note that $K$ has finite fundamental group except for the symmetric pairs of type $E_{6(-14)}$ and $E_{7(-25)}$. Also note that there is no relevant $K$-orbit of $E_{7(-25)}$. For $E_{6(-14)}$, $K=(\operatorname{U}(1)  \times \Spin(10)) / \Z_4$, $\pi_1(K) \simeq \Z$. If $\pi_1(K)$ is finite, let $\tilde{K}$ denote the simply connected cover of $K$. Otherwise $\tilde{K}$ can mean any finite connected cover of $K$. 

Let $X = \spec \C[\Orb]$ and $X'$ be the unique codimension $4$ symplectic leaf along which the singularity of $X$ is of type $a_2$ as in Lemma \ref{lem:only_a2}. Let $\Orb_K$ be a nilpotent $K$-orbit in $\Orb \cap \kf^\perp$ and $\overline{\Orb}_K$ be its closure in $X$. Let $\Orb'_K$ be a $K$-orbit contained in $\overline{\Orb}_K \cap X'$. Fix a point $\chi$ in $\Orb_K$. Let $Q, K_Q, \tilde{K}_Q$ denote the reductive parts of the centralizers of $\chi$ in $G,K,\tilde{K}$,
respectively. Let $Z_K = K_Q / K_Q^\circ$ and $\tilde{Z}_K = \tilde{K}_Q / \tilde{K}_Q^\circ$ denote the component groups. Let $Q'$, $K'_Q$ and $\tilde{K}'_Q$ be the groups defined similarly for a fixed element $\chi' \in \Orb'_K$ as for $\chi \in \Orb_K$. As in Section \ref{SS_main_results}, we take a $\theta$-stable slice $\check{X}$ in $X$ that is transversal to $X'$ at $\chi'$, then $\check{X}$ is isomorphic to the minimal nilpotent orbit in $\mathfrak{sl}_3$. Then the group $Q'$ acts on $\check{X}$, $K'_Q$ and $\tilde{K}'_Q$ acts on $\check{X}^\theta$.


%In Tables \ref{Table:a2}, \ref{Table:2a2} and \ref{Table:a2/S_2}, we list, for various $\GR$, all cases in which the singularity of $\Orb$ along $\Orb'$ is of type $a_2$, $2a_2$ and $a_2/S_2$. We also record all associated groups and Lie algebras and indicate whether $\Orb$ is admissible or sc-admissible in the column ``Adm".
%Here is the meaning of each column. Let $\Orb_K$ denote a $K$-orbit included in $\Orb$ and $\Orb'_K$ denote a $K$-orbit in the intersection $\overline{\Orb}_K \cap \Orb'$, where $\overline{/\Orb}_K$ is the closure of $\Orb_K$ in $\kf^\perp$. \footnote{Where is the closure taken? Also do you mean codimension $2$ orbits? As written, the intersection is a single orbit. And the sentence is not correct grammatically . S. Y.: corrected.}. The column $\kf_Q'$ gives the isomorphism type of the Lie algebra of the stabilizer $K_{Q'}$ in $K$ of a normal $\slf_2$-triple of $\Orb'_K$\footnote{I.L.: no such column at least in Table 3. And one reads $\kf_Q'$ from $\q'_\R$},

%Let $\Orb_\R$ be the real form of $\Orb$ corresponding to $\Orb_K$ under the The column $\q'_\R$ gives the Lie algebra of the corresponding real form $Q'_\R$ of $Q'$ under the Sekiguchi correspondence, so that the isomorphism type of the Lie algebra $\kf_Q'$  of the stabilizer $K_{Q'}$ in $K$ of a normal $\slf_2$-triple of $\Orb'_K$. Then $K'_Q$ is the complexification of a maximal compact subgroup of $Q'_\R$.



By \cite[Proposition 3.3]{FJLS1} or Lemma \ref{Lem:Qi_action}, the action of $Q'$ on $\check{X}$ has a dense orbit. Therefore the restriction of $\theta$ on $\check{X}$ can be read off from the restriction of $\sigma$ on $\q'$, which can be found in \cite{Dk0} (in fact it is the corresponding anti-holomorphic involution that is given there). We assume that $\chi$ lies in $\check{X} \cap \Orb_K$ and the induced involution on $\check{X}$ is outer.  In all cases, we find that $K_Q'^\circ \simeq \SO_3$. Let $\underline{K}$ ($\underline{\tilde{K}}$ resp.) denote the reductive centralizer (=the component group of the centralizer) of $\chi$ for the action of $K_Q'^\circ$ ($\tilde{K}'^\circ_Q$ resp.) on $\check{X}\cap \Orb_K$. Then $\underline{K} \simeq \Z_2$. Analogous to \eqref{eq:hom_Z4}, we have the composite map
  \[  \Z_4\hookrightarrow (\SL_2)_\chi \rightarrow (\tilde{K}'^\circ_{Q})_\chi \hookrightarrow \tilde{K}_Q\twoheadrightarrow \tilde{K}_Q / \tilde{K}_Q^\circ = \tilde{Z}_K.\]
  %\[\Z_4 \simeq \pi_0[(\SL_2)_\chi] \twoheadrightarrow \pi_0[(\tilde{K}_{Q}'^\circ)_\chi] = \underline{\tilde{K}} \to \tilde{K}_Q / \tilde{K}_Q^\circ = \tilde{Z}_K \]
and a similarly defined map $\Z_4 \to Z_K$ for $K$. As we have seen in Section \ref{SS_loc_sys_slice_restriction}, the image of $\bullet_{\dagger,\chi}$ is controlled by homomorphism $\Z_4 \rightarrow \tilde{Z}_K$ in \eqref{eq:hom_Z4}. If this map is not injective, then similarly to the discussion in the end of Section \ref{SSS_Z4_homom_typeA}, the quantization of every irreducible $(\tilde{K},\kappa)$-equivariant
suitably  twisted local system on $\Orb_K$ extends to $\Orb'_K$. In order to determine whether the map $\Z_4 \to \tilde{Z}_K$ is injective, we need to compute $\tilde{Z}_K$ and $\underline{\tilde{K}}$. We have a natural surjective map $\underline{\tilde{K}} \to \underline{K}$, whose kernel is naturally identified with the kernel of the covering map $\tilde{K}'^\circ_Q \twoheadrightarrow K_Q'^\circ$. Therefore we will have to compute the covering morphism $\tilde{K}'^\circ_Q \twoheadrightarrow K_Q'^\circ$ to determine $\underline{\tilde{K}}$. In most cases, knowing the covering map $\tilde{K}_Q^\circ \twoheadrightarrow K_Q^\circ$, the group $\underline{\tilde{K}}$ and the map $\underline{\tilde{K}} \to \tilde{Z}_K$ could help us determine $\tilde{Z}_K$, but not in all cases. %\footnote{It is possible to generalize the methods in \cite{King} to determine $\tilde{Z}_K$ in all cases, but this is not our main purpose here. S.Y.: Not sure about the first half of the sentence. Maybe ATLAS can do this in the future. This footnote should be removed. } 
Sometimes we also consider certain double covering groups $\bar{K}$ of $K$, if they exist. We can similarly define the groups $\bar{K}$, $\bar{Z}_K$, $\bar{K}_Q$, $\bar{K}'_Q$, and $\underline{\bar{K}}$.


%We will mainly apply Lemma \ref{Lem:sl3_gl2_unobstructive}, Lemma \ref{Lem:sl3_so3_obstructive} and Lemma \ref{Lem:a2S2_strongly_quantizable} \footnote{I.L.: bad idea... And I think tables need to be placed after they are first mentioned. S.Y.: I've added references to the lemmas below}.


%According to \cite[Table 6]{GraafElashvili}, the orbit $A_2 + A_1$ is birationally\footnote{S.Y.: Need some Springer theory.} induced from the minimal nilpotent orbit $\Orb_{min}^{\mathfrak{so}_{10}}$ of a Levi subgroup of type $D_5$. Hence the quantization parameter space of $X$ is $H^2(\Orb,\C) \simeq \C$. 

We will list in Theorem \ref{thm:exceptional} below all the relevant $K$-orbits as above by cases and make conclusions about them. We follow the labeling of $K$-orbits in  \cite{Dk1, Dk6, Dk5, Dk2, Dk3, Dk4, Dk7, Dk8}, where the closure relations of the orbits can also be found. The papers \cite{Noel1, Noel2} also use the same labeling. 

\begin{Thm}\label{thm:exceptional}
The list of (sc-)admissible $K$-orbits $\Orb_K$ as above, together with their component groups, can be divided into the following three cases:  
\begin{enumerate}
\item[Case 1:]
the involution on the $a_2$-slice is inner or $\overline{\Orb}_K$ does not intersect the leaf $X'$ with the $a_2$-singularity,
\begin{itemize}
	\item 
	   $E_{6(2)}$: $\#12$, $\#13$ and $\#14$ in $A_2 + 2A_1$, $\tilde{Z}_K = Z_K = 1$;
	\item 
	    $E_{6(-14)}$: $\# 7$ in $A_2+A_1$, $Z_K= 1$; $\# 8$ in $A_2+A_1$: $Z_K= 1$ ($K$ is not semisimple);
	    %, $\tilde{Z}_K = \ker(\tilde{K} \to K)$. might not be true
	\item 
	    $E_{7(-5)}$: $\#27$ in $A_4 + A_1$, $Z_K = \Z_2$, $\tilde{Z}_K \simeq \Z_4$ or $\Z_2 \times \Z_2$; 
	\item 
	   	$E_{8(8)}$: $\#42$ in $A_4+2A_1$, $\tilde{Z}_K=Z_K =1$;  $\#43$ in $A_4+A_1$, $\tilde{Z}_K=Z_K \simeq \Z_2$;
	 \item 
	    $E_{8(-24)}$: $\#26$ in $A_4+A_1$, $\tilde{Z}_K=Z_K =1$.
\end{itemize}
\item[Case 2:]
the involution on the $a_2$-slice is outer and $\tilde{Z}_K = Z_K$,
\begin{itemize}
	\item 
	  $E_{6(6)}$: $\# 8$ in $A_2+A_1$, $\tilde{Z}_K \simeq \Z_2$;
	\item 
	  $E_{8(8)}$:  $\# 38$ in $A_4+A_1$, $\tilde{Z}_K \simeq \Z_4$;
\end{itemize}
\item[Case 3:]
the involution on the $a_2$-slice is outer and $\tilde{Z}_K \neq Z_K$,
\begin{enumerate}[(a)]
	\item 
	$E_{6(6)}$: $\# 10$ in $A_2+2A_1$, $Z_K \simeq \Z_2$, $\tilde{Z}_K \simeq \Z_4$;
	\item 
	$E_{7(7)}$: $\# 50$ in $A_4+A_1$, $Z_K = \Z_4$, $\bar{Z}_K \simeq \Z_4 \times \Z_2$, $\tilde{Z}_K \simeq \Z_8 \times \Z_2$ or $ \Z_4 \times \Z_4$.
	\item 
	$E_{8(8)}$: $\#44$ in $A_4 + 2A_1$, $Z_K = (\Z_2)^2$, $\tilde{Z}_K$ is a central extension of $Z_K$ by $\Z_2$ and can only be $\Z_2 \times \Z_4$, the dihedral group of order $8$ or the quaternion group.
\end{enumerate}
\end{enumerate}
Then
  \begin{enumerate}[(i)]
  	 \item 
  	   For Case 1 and 2, the functor \[ \bullet_{\dagger,\chi} : \HC_{\Orb_K}(\A_\lambda,\theta)^{\tilde{K},\kappa}\xrightarrow{\sim} \operatorname{Rep}(\tilde{K}_Q,\kappa_Q)\] 
  	   is an equivalence for any quantization parameter $\lambda$.
  	 \item 
  	   For Case 3, we have $\underline{\tilde{K}} \simeq \Z_4$ and the map $\Z_4 \to \tilde{Z}_K$
  	   is injective. Moreover, $K$ is simple and $\rho_\omega=0$, so $\kappa =0$ and $\kappa_Q=0$ (for $\lambda=0$). These imply that the functor
  	   \begin{equation} \label{eq:HC_bad_exc_orbits}
  	      \HC_{\Orb_K}(\A_0,\theta)^{\tilde{K},\kappa} = \HC_{\Orb_K}(\A_0,\theta)^{\tilde{K},0} \to \operatorname{Rep}(\tilde{K}_Q,\kappa_Q) = \operatorname{Rep}(\tilde{K}_Q)
        \end{equation} 
  	   is not essentially surjective. This means that there always exists at least one irreducible $\tilde{K}$-equivariant local system (of rank $1$) over $\Orb_K$, which does not quantize to a HC $(\A_0, \tilde{K})$-module.
  \end{enumerate}
\end{Thm}

\begin{proof}
	The lists are obtained by case-by-case inspections and computations. 
	%The information about admissibility  of $\Orb_K$ is given in \cite{Noel1, Noel2}. In particular, we know for each $\Orb_K$ whether it is admissible, sc-admissble or not admissbile in the sense of Definition \ref{defn:sc-admissible}. 	
	%If a set $\{\vartheta^{\chi}_1, \ldots, \vartheta^{\chi}_m\}$ of basis vectors of $\tf_\chi$ are chosen, we write $d \delta_\chi (z) = \sum_j  b_j z_j $, $b_j \in \Z$, for the value of $d \delta_\chi$ at the element $\sum_j z_j \vartheta^{\chi}_j \in \tf_\chi$, where $z_j \in \C$\footnote{S.Y.: it is probably not necessary to include $d \delta_\chi$, unless when $d \delta_\chi = 0$. When $d \delta_\chi$ is nonzero, it is not trivial to compute the twisted local systems explicitly and it is not our purpose.}.
	Except for the two orbits of $E_{6(-14)}$ in Case 1, $K$ is always semisimple. Assume $K$ is semisimple. By the discussions above the theorem, we need to compute the covering morphisms $\tilde{K}_Q^\circ \twoheadrightarrow K_Q^\circ$ and $\tilde{K}'^\circ_Q \twoheadrightarrow K_Q'^\circ$. 
	For this purpose, we choose a Cartan subalgebra $\tf_\chi$ of the Lie algebra $\kf_Q$ with corresponding Cartan subgroup $T_\chi$ of $K_Q^\circ$, together with basis vectors $\{\vartheta^{\chi}_i \}$ of $\tf_\chi$ as follows. For each $K$-orbit $\Orb_K$, No{\"e}l gives in \cite{Noel1, Noel2} an explicit nilpotent element $\chi \in \Orb_K$ and $\tf_\chi$ spanned by certain basis vectors, each of which is an integral linear combination of coroots $\alpha^\vee_i$ of $\g$ (which No{\"e}l denotes by $H_i$). Let $\omega^\vee_1, \ldots, \omega^\vee_l$  be the fundamental coweights corresponding to the simple roots $\Delta = \{ \alpha_1, \ldots, \alpha_l \}$ of $\g$, where we follow the labeling of simple roots in \cite{Bourbaki}. In all cases, $\vartheta^{\chi}_i$ is chosen to be of the form $\sum_{k=1}^{l} a_k \omega^\vee_{k}$, where $a_i$ are coprime integers. Moreover, $\{2\pi i \vartheta^{\chi}_i \}$ forms a basis of the lattice $\ker(\exp(\cdot): \tf_\chi \to T_\chi)$, i.e., the coweight lattice of $K_Q^\circ$\footnote{No{\"e}l's basis vectors of $\tf_\chi$ are integral linear combinations of the coroots of $\g$. He mistakenly takes this basis as the coweight lattice of $K_Q^\circ$ and uses it to decide admissibility of the orbit. This then means that his $G$ is not adjoint but simply connected. Therefore what he has really computed is not about admissibility for adjoint real groups, but for the group $G(\R)$ of $\R$-points of the simply connected simple complex groups. His conclusions about sc-admissiblity are still correct. For adjoint groups, we have to redo the calculations using our basis. Fortunately, the conclusions about admissibility of the orbits in the list of Theorem \ref{thm:exceptional} are not affected by this discrepancy, even though in some cases his lattice is strictly smaller than ours. However, there are $K$-orbits, for which $\operatorname{codim}_{\overline{\Orb}}\partial\Orb\geqslant 6$, that are not admissible from our computations, but are marked as admissible by No{\"e}l.}. Then all $\exp(2\pi i \vartheta^{\chi}_i)$ lie in $\ker(\tilde{K} \twoheadrightarrow K) = \pi_1(K)$ and generate the subgroup $\ker(\tilde{K}_Q^\circ \twoheadrightarrow K_Q^\circ)$.
	Let $\{ \varpi_1, \ldots \varpi_n \}$ denote the fundamental weights of $\kf$. Evaluation of $\{ \varpi_1, \ldots \varpi_n \}$ at each basis vector $\vartheta^{\chi}_i$ tells us about the order of the element $\exp(2\pi i \vartheta^{\chi}_i)$ in $\tilde{K}$. Indeed, let $m_i$ be the least positive integer such that all $\varpi_k, 1 \leq k \leq n$, have integral values at $m_i \vartheta^{\chi}_i$. Then $m_i$ is the order of $\exp(2\pi i \vartheta^{\chi}_i)$ in $\tilde{K}$.  In all cases except for $E_{7(-5)}$, $\pi_1(K)$ is a finite cyclic group, therefore $\exp(2\pi i \vartheta^{\chi}_i)$ is uniquely determined by its order and hence we know $\ker(\tilde{K}_Q^\circ \twoheadrightarrow K_Q^\circ)$ as a subgroup of $\pi_1(K)$. In the same way we can compute $\ker(\tilde{K}'^\circ_Q \twoheadrightarrow K_Q'^\circ)$. For $E_{7(-5)}$, we have $K \simeq (\operatorname{Ss}(12) \times \SU(2))/\Z_2$ and $\pi_1(K) \simeq \Z_2 \times \Z_2$. For the orbit $\#27$ in $E_{7(-5)}$, $\tilde{K}_Q^\circ$ is a $2$-fold cover of $K_Q^\circ$ by computation above.  In all cases, however, our computations already tell whether $\tilde{Z}_K = Z_K$ or not. 
	
	We first prove (i). Lemma \ref{lem:only_a2} implies that, in both Case 1 and 2, $X=\spec \C[\Orb]$ contains only one codimenison $4$ leaf, along which the singularity is of type $a_2$. Therefore the orbits in Case 1 satisfy (ii) of Theorem \ref{Thm:main1} and we are done.  For Case 2,  both orbits have $\kf_Q=\mathfrak{so}(3)$ according to \cite{Dk0.5}, which are semisimple. Therefore by Section \ref{SSS_homog_case}, the only strongly $K$-equivariant Picard algebroid over $\Orb_K$ is the trivial one (up to isomorphism). Together with the fact that $\tilde{Z}_K = Z_K$, this implies that $\tilde{K}$-equivariant admissible vector bundles are in fact $K$-equivariant and are nothing else but (untwisted) $K$-equivariant local systems (of rank $1$), for any quantization parameter $\lambda$ (recall that $A_2+A_1$ in $\mathfrak{e}_6$ is not birationally rigid). Therefore (i) of Theorem \ref{Thm:main1} applies.
	%Lemma \ref{Lem:sl3_so3_obstructive} implies that the map $\Z_4 \to Z_K$ is not injective and neither is the map $\Z_4 \to \tilde{Z}_K \simeq Z_K$. Now apply Proposition \ref{Lem:image_description11} and Lemma \ref{Lem:sl3_so3_obstructive} to get the conclusion.
	
	% or $\overline{\Orb}_K$ does not intersect with the $2a_2$-slices ($\#27$ of  $E_{7(-5)}$ in $A_4 + A_1$). 
	
	For (ii), we will give a detailed computation of the component group $\tilde{Z}_K$ for the orbit (a) in Section \ref{SS:orbit_a} and a less detailed one for the orbit (b) in Section \ref{SS:orbit_b}, as demonstrations of the general strategy outlined above. The analysis of orbit (c) is similar and we omit it. These computations confirm $\underline{\tilde{K}} \simeq \Z_4$ and the injectivity of the map $\Z_4 \to \tilde{Z}_K$. From \cite{Noel1, Noel2}, we also know the Lie algebra character $d \delta_\chi := (\omega_{\Orb_K})_\chi \in (\kf_Q^*)^{K_Q}$ (as in Section \ref{SS_intro_approach}) for $\Orb_K$, so that $\rho_\omega = \frac{1}{2} d \delta_\chi$ is also known. Then we can again apply Proposition \ref{Lem:image_description11} and Lemma \ref{Lem:sl3_so3_obstructive}.
\end{proof}


\begin{comment}

 A case-by-case computation shows the following uniform result.

\begin{Prop}\label{prop:Z4_exceptional}
	 Assume $\Orb_K$ is (sc-)admissible and the induced involution on an $a_2$-slice in $\spec \C[\Orb]$ is outer and $\tilde{Z}_K \neq Z_K$. Then $\underline{\tilde{K}} \simeq \Z_4$ and the composition
	   \[ \Z_4 \simeq \underline{\tilde{K}} \hookrightarrow \underline{\tilde{K}} \hookrightarrow \tilde{K}_Q \twoheadrightarrow \tilde{Z}_K \]
	 is injective.
\end{Prop}

Note that for all the $3$ orbits above, $d\delta_\chi = 0$ and the groups $K$ are simple. These observations yield the following result.

\begin{Cor}
	For all the three orbits satisfying Proposition \ref{prop:Z4_exceptional}, there exists some $\lambda$ such that the functor \[ \HC_{\Orb_K}(\A_\lambda,\theta)^{\tilde{K},\kappa} = \HC_{\Orb_K}(\A_\lambda,\theta)^{\tilde{K},0} \to \operatorname{Rep}(\tilde{K}_Q,\kappa_Q)\]  is not essentially surjective. This means that there always exists at least one $\tilde{K}$-admissible vector bundle over $\Orb_K$, which does not quantize to a Harish-Chandra $(\g, \tilde{K})$-module.
\end{Cor}
\end{comment}


	
	
\subsubsection{The $K$-orbit $\# 10$ of $E_{6(6)}$ in $\Orb = A_2 + 2A_1$}\label{SS:orbit_a} \, \\

Now let $(G, K)$ be the symmetric pair of type $E_{6(6)}$ where $K = \Sp(4, \C)/\Z_2$. Then $\tilde{K} = \Sp(4, \C)$. Let $\alpha_i$, $1 \leqslant i \leqslant 36$, be the positive roots of $\g = \mathfrak{e}_6$ as in \cite{Dk2} and choose $\Delta = \{ \alpha_1, \alpha_2, \ldots, \alpha_6 \}$ to be the set of simple roots. All $36$ positive roots written as integral linear combinations of the simple roots $\alpha_i$, $1 \leqslant i \leqslant 6$, can be found in \cite[Table 13]{Dk2}. A negative root $-\alpha_i$ will be also written as $\alpha_{-i}$. For each root $\alpha$, let $\g^\alpha$ denote the root space of $\alpha$. Choose root vectors $X_i \in \g^{\alpha_i}$ for $\pm i \in \{ 1, \ldots, 36\}$ so that together with the coroots $\alpha^\vee_i$, $i= 1, \dots, 6$, they form a Chevalley basis of $\g$. Set
\[ \beta_1 = \alpha_2, \quad \beta_2 = \alpha_4, \quad \beta_3 = \frac{\alpha_3 + \alpha_5}{2}, \quad \beta_4 = \frac{\alpha_1 + \alpha_6}{2} , \quad \beta_0 = \beta_1 + 2 \beta_2 + 3 \beta_3 + 2 \beta_4. \]
Then $\Delta_\kf = \{-\beta_0, \beta_4, \beta_3, \beta_2 \}$ forms a set of simple roots for $\kf = \mathfrak{sp}(4,\C)$. The fundamental weights of $\kf$ are $(\varpi_1, \ldots, \varpi_4)$ given by
\[
\begin{pmatrix}
	\varpi_1 \\ \varpi_2 \\ \varpi_3 \\ \varpi_4
\end{pmatrix}
=
\begin{pmatrix}
	1  &  1  &  1  &  \frac{1}{2}   \\
	1  &  2  &  2 &   1  \\
	1  &  2  &  3  &   \frac{3}{2}  \\
	1  &  2  &  3  &   2 \\
\end{pmatrix}
\begin{pmatrix}
	-\beta_0 \\ \beta_4 \\ \beta_3 \\ \beta_2
\end{pmatrix}.
\]


Let $\Orb_K$ be the $K$-orbit $\# 10$ in $\Orb=A_2+2A_1$. We have $Z_K = \Z_2$ and $d \delta_\chi = 0$, so the orbit is admissible for any finite cover of $K$, so in particular $K$-admissible. Using the explicit nilpotent element $\chi$ in \cite[Table 4]{Noel2}, we can take the Cartan subalgebra of $\kf_Q$ to be $\tf_\chi =  \C \vartheta^{\chi}$ where \footnote{In this case, No{\"e}l's basis vector in terms of coroots on the right hand side coincides with ours in terms of fundamental coweights. Note that there is a typo in No{\"e}l's expression for $\chi$: the last root vector $X_{\alpha_1 - \alpha_2 - \alpha_3 - 2 \alpha_4 - \alpha_5 - \alpha_6}$ should be $X_{-\alpha_1 - \alpha_2 - \alpha_3 - 2 \alpha_4 - \alpha_5 - \alpha_6}$ instead. The expression in the former subscript is not even a root since the coefficients have different signs. The latter root is  noncompact imaginary by \cite[Table 3]{Noel2}.}
\[ \vartheta^{\chi} = \omega^\vee_1 -2 \omega^\vee_4 + \omega^\vee_6 = \alpha^\vee_1 -2 \alpha^\vee_4 + \alpha^\vee_6.\]
Evaluating $(\varpi_1, \ldots, \varpi_4)$ at the basis vector $\vartheta^{\chi}$ of $\tf_\chi$ gives $(0, 1, 0, -2)$. This implies that $\exp(2\pi i \vartheta^{\chi})$ is the identity in $\tilde{K}$. This means that the nontrivial element in the center of $\tilde{K}$ does not lie in $\tilde{K}_Q^\circ$, hence  $\tilde{K}_Q^\circ=K_Q^\circ$ and $\tilde{Z}_K$ is a central extension of $Z_K \simeq \Z_2$ by $\Z_2$. We will show below that $\tilde{Z}_K \simeq \Z_4$.

The closure $\overline{\Orb}$ is normal in codimension $4$ and hence we can identify the symplectic leaf $X'$ in $X$ with the orbit $\Orb'$ of type $A_2 + A_1$. The closure $\overline{\Orb}_K$ intersect with $\Orb'$ of type $A_2 + A_1$ at the $K$-orbit $\Orb'_K  = \# 8$ and the involution on the corresponding slice is outer. As in \cite[Table 4]{Noel2}, choose a nilpotent element $\chi'$ in $\Orb'_K$ so that the Cartan subalgebra of $\kf'_Q$ can be taken as $\tf_{\chi'} = \C \vartheta^{\chi'}$ where \footnote{Again here No{\"e}l's basis vector coincides with ours.}
\[ \vartheta^{\chi'} = \omega^\vee_3 - \omega^\vee_4 + \omega^\vee_5 = \alpha^\vee_1 + \alpha^\vee_2 + 2\alpha^\vee_3 + 2\alpha^\vee_4 + 2\alpha^\vee_5+ \alpha^\vee_6.\]
Evaluating $(\varpi_1, \ldots, \varpi_4)$ at the basis vector $\vartheta^{\chi}$ of $\tf_{\chi'}$ gives $(-1/2, 0, 1/2, 0)$. This means that $\exp(2\pi i \vartheta^{\chi'})$ is the unique non-trivial element of order $2$ in the center of $\tilde{K}$ %Let $T'=\exp(\tf_{\chi'})$ denote the maximal torus in $K_Q'^\circ$ whose Lie algebra is $\tf_{\chi'}$. Then the computation above implies that the lift $\tilde{T}'$ of $T'$ in $\tilde{K}'_Q$ is a double cover of $T'$ such that the kernel of $\tilde{T}' \to T'$ is $\{ 1, \exp(2\pi i H')\} \simeq \Z_2$.
and $\tilde{K}_Q'^\circ$ is a double over of $K_Q'^\circ \simeq \SO_3$, hence $\tilde{K}_Q'^\circ \simeq \SL_2$. We see that $\underline{\tilde{K}} \simeq \Z_4$ with $2 \underline{\tilde{K}} = 2 \Z_4$ being the center of $\tilde{K}$. Since the center $2 \underline{\tilde{K}}$ is not in the kernel of the map $\underline{\tilde{K}} \to \tilde{Z}_K$ by the discussion above, we conclude that $\underline{\tilde{K}} \to \tilde{Z}_K$ is injective. This also implies that the map $\underline{K} \to Z_K$ is also injective. %\footnote{S.Y.: This can also be deduced by looking at the reductive part $\underline{G}$ of the stabilizier for the action of $Q'$ on $\check{X}$ and its inclusion into $Q$. $\underline{G}$ is a two dimensional torus and embeds as a maximal torus of $Q = Q^\circ \simeq \GL(2)$. However, it does not seem to be enough for the claim about the map $\underline{\tilde{K}} \to  \tilde{Z}_K$ for the simply connected group.}
We then conclude that $\tilde{Z}_K \simeq \Z_4$ and the natural map $\underline{\tilde{K}} \to \tilde{Z}_K$ is an isomorphism. The kernel of the natural map $\tilde{Z}_K \to Z_K$ equals $\ker(\tilde{K} \twoheadrightarrow K) \simeq \Z_2$. Therefore we only get three irreducbile HC $(\A_0, \tilde{K})$-modules. This finishes the proof of part (ii) of Theorem \ref{thm:exceptional} in this case.

 As mentioned at the beginning of this section, $A_2 + 2A_1$ in $\mathfrak{e}_6$ is not birationally rigid. Let us explain the classification of irreducible HC modules with full support over $\Orb_K$ for any quantization parameter $\lambda$ in this case. The component group
 for the orbit $A_2+2A_1$ is trivial, \cite[Section 8.4]{CM}. According to \cite[Table 6]{GraafElashvili}, the orbit $A_2 + 2A_1$ is (automatically birationally) induced from the zero nilpotent orbit of a Levi subgroup $L \subset G$ of type $A_4 + A_1$. Choose a parabolic subgroup $P$ in $G$ with Levi factor $L$ and form the generalized Springer map $T^*(G/P) \to \overline{\Orb}$. One can show that the parabolic Springer fiber over any point $\chi' \in \Orb'$ is identified with $\mathbb{P}^2$ and the pullback map $\operatorname{Pic}(G/P)\rightarrow
 \operatorname{Pic}(\mathbb{P}^2)$ is an isomorphism of free abelian groups of rank $1$. Using Lemma \ref{Lem:sl3_so3_obstructive}, we can also apply the same argument in the proof of Theorem \ref{Thm:spin_HC_classif_general} and get similar classification as in the case of type $A$ groups. We leave the details to interested readers. Note that in this case, $\kf_Q = \mathfrak{so}(2) \subset \gl(2) = \q$, hence the natural striction map $(\q^*)^Q \to (\mathfrak{k}_Q^*)^{K_Q}$ is zero. Therefore restriction of the twist $\lambda$ to $\Orb_K$ is always trivial by Corollary \ref{Cor:equiv_condition}. Moreover $d \delta_\chi = 0$, therefore $\kappa_Q=0$ and $\operatorname{Rep}(\tilde{K}_Q, \kappa_Q)$ is equivalent to $\operatorname{Rep}(\tilde{Z}_K) \simeq \operatorname{Rep}(\Z_4)$. 




\subsubsection{$K$-orbit $\# 50$ of $E_{7(7)}$ in $\Orb=A_4+A_1$}\label{SS:orbit_b} \, \\

Let $(G,K)$ be the symmetric pair of type $E_{7(7)}$ with $K = \SL(8) / \Z_4$, so $\tilde{K}=\SL(8)$. We have a unique double cover $\bar{K} := \tilde{K} / 2 \Z_4 = \SL(8) / 2\Z_4$ of $K$. Let $\Orb_K$ be the $K$-orbit $\# 50$ in $\Orb=A_4+A_1$. We have $Z_K \simeq \Z_4$. In this case $\kf_Q=0$, so $d\delta_\chi = 0$ and the $K$-orbit $\Orb_K$ is admissible.  

Let $X=\spec\C[\Orb]$, then $X$ is normal in codimenision $4$ and it has two codimension $4$ leaves, identified with the orbit $A_4$ with $a_2$-singualrity and the orbit $A_3 + A_2 + A_1$ with $a_2 /S_2$-singularity, respectively. By Lemma \ref{Lem:character_integrality} and the same argument in the proof of Theorem \ref{Thm:main1} (Section  \ref{SS:proof_main1}), the $a_2/S_2$-singularity poses no obstruction to the quantization of (twisted) local systems over $\Orb_K$. Therefore we only need to consider the $a_2$-singularity along the orbit $A_4$.

The closure $\overline{\Orb}_K$ intersects with $A_4$ at the $K$-orbit $\Orb'_K  = \# 43$ and the involution on the corresponding slice comes from an outer involution of $\slf_3$. Using the explicit representative element in $\Orb'_K$ from \cite[Table 3]{Dk5}, we can compute that $\tilde{K}'^\circ_Q$ is a double cover of $K_Q'^\circ$, while $\bar{K}_Q'^\circ \simeq K_Q'^\circ$, in a similar way as in Section \ref{SS:orbit_a}. We see that $\underline{K} \simeq \underline{\bar{K}} \simeq \Z_2$ and $\underline{\tilde{K}} \simeq \Z_4$, so that $2 \underline{\tilde{K}}$ is identified the unique subgroup of order $2$ of the center $\Z_4$ of $\tilde{K}$.
%Let $\bar{D}$ and $\tilde{D}$ denote the preimage of $\underline{K}$ in $\bar{K}$ and $\tilde{K}$ respectively. Then $\bar{D} = \underline{\bar{K}} \times (\Z_4/\overline{2}\Z_4) = \Z_2 \times (\Z_4/\overline{2}\Z_4)$, where the second component is the center of $\bar{G}_\R$. Let $\tilde{D}$ denote the preimage of $\bar{D}$ in $\underline{\tilde{K}}$, then $\tilde{D}$ is isomorphic to $\Z_4$. Also the center $\Z_4$ of $\tilde{G}_\R$ is contained in $\underline{\tilde{K}}$. The two copies of $\Z_4$ commute with each other and intersect at a subgroup of order $2$. Then the only possibility is $\underline{\tilde{K}} \simeq \Z_4 \times \Z_2$.
%\footnote{S.Y.: Maybe say that $\underline{K}$ is in fact $\underline{K}_\R \cap K_Q'^\circ_\R$, so the two $\Z_4$ commutes with}



Since $\kf_Q=0$, it follows that $K_Q$, $\bar{K}_Q$ and $\tilde{K}_Q$ are finite groups and coincide with their component groups $Z_K$, $\bar{Z}_K$ and $\tilde{Z}_K$. Therefore the maps from $\underline{K}$, $\underline{\bar{K}}$ and $\underline{\tilde{K}}$ respectively to $K_Q$, $\bar{K}_Q$ and $\tilde{K}_Q$ respectively are all injective, and so we can regard the former as subgroups of the latter. The group $\tilde{Z}_K$ is a central extension of the cyclic group $Z_K \simeq \Z_4$ by the center $\Z_4$ of $\tilde{K}$, hence $\tilde{Z}_K$ has to be abelian and so is $\bar{Z}_K$.
Pick a generator $a$ of $Z_K = \Z_4$ and a preimage $\bar{a}$ of $a$ under the projection map $\bar{Z}_K \twoheadrightarrow Z_K$. If $\bar{a}$ is of order $8$, then $\bar{Z}_K \simeq \Z_8$. But then $\underline{\bar{K}} \subset \bar{Z}_K$ is the unique subgroup of $\bar{Z}_K$ of order $2$, which is also the kernel of the projection $\bar{Z}_K \twoheadrightarrow Z_K$. This contradicts with the fact that $\underline{K}$ is the isomorphic image of $\underline{\bar{K}}$ under $\bar{Z}_K \twoheadrightarrow Z_K$. Therefore we conclude that $\bar{a}$ is of order $4$ and generates a cyclic subgroup of $\bar{Z}_K$ of order $4$, so that $\bar{Z}_K \simeq \Z_4 \times \Z_2$, where the second factor $\Z_2 \simeq \Z_4 / 2 \Z_4$ is the center of $\bar{K}$.

Let $\bar{G}$ be the simply connected complex simple group of type $E_7$. Then $\bar{K}$ is the fixed locus of an involution of $\bar{G}$ and hence satisfies the condition (i) of Theorem \ref{Thm:main1}. We conclude that there are totally $4$ distinct irreducible admissible $K$-equivariant line bundles on $\Orb_K$ and give rise to $4$ irreducible HC $(\g,K)$-modules supported on $\Orb_K$. There are $8$ irreducible admissible $\bar{K}$-equivariant line bundles on  $\Orb_K$, $4$ of which are $K$-equivariant ones mentioned before. The other $4$ genuine ones quantize to $4$ distinct genuine irreducible HC $(\g, \bar{K})$-modules. This coincides with the outcome of Atlas.

%We now analyze the case of the pair $(\g,\tilde{K})$. The closure $\overline{\Orb}_K$ intersect $A_3 + A_2 + A_1$ at four $K$-orbits, labelled by 39, 40, 41 and 42. Let $\Orb_K^{i}$, $i=39, 40, 41, 42$, denote these orbits, and write the relevant groups as $K_Q^i$ and $\tilde{K}_Q^i$ etc. Using the explicit representative elements in \cite[Table 3]{Dk5}, we can choose a Cartan sublalgebra of $\kf_Q^i$ and follow the strategy outlined in the proof of Theorem \ref{thm:exceptional} to find that $(\tilde{K}_Q^{i})^\circ = (K_Q^{i})^\circ$. By Lemma \ref{Lem:character_integrality} and Lemma \ref{Lem:a2S2_strongly_quantizable}, we conclude that these four $K$-orbits pose no obstruction to the quantization of (twisted) local systems over $\Orb_K$.

We are unable to determine $\tilde{Z}_K$. Continuing the analysis above for $\bar{Z}_K$ shows that either $\tilde{Z}_K \simeq \Z_8 \times \Z_2$ such that the embedding $\Z_4 \simeq \underline{\tilde{K}} \hookrightarrow \tilde{Z}_K \simeq \Z_8 \times \Z_2$ is given by the embedding of $\Z_4$ to the first factor $\Z_8$, or $\tilde{Z}_K \simeq \Z_4 \times \Z_4$, where the second factor is the center of $\tilde{G}_\R$, and the embedding $\Z_4 \simeq \underline{\tilde{K}} \hookrightarrow \tilde{Z}_K \simeq \Z_4 \times \Z_4$ sends the generator $\bar{1}$ of $\Z_4$ to $(\bar{2}, \bar{1})$. The irreducible admissible $\tilde{K}$-equivariant vector bundles are equivalent to irreducible characters of $\tilde{Z}_K$, so in either case there are totally $16$ of them. We now apply Lemma \ref{Lem:sl3_so3_obstructive}. Among the $16$ irreducible admissible $\tilde{K}$-equivariant vector bundles (in fact, local systems), $4$ give no quantization, $4$ quantize to irreducible Harish-Chandra $(\g, K)$-modules and $4$ quantize irreducible genuine Harish-Chandra $(\g, \bar{K})$-modules. The remaining $4$ local systems quantize to genuine Harish-Chandra $(\g, \tilde{K})$-modules.\\


\begin{thebibliography}{99}

\bibitem{Alexeevskii} A. V. Alexeevskii, {\it Component groups of centralizers of unipotent elements in semisimple algebraic groups}, Akad. Nauk Gruzin. SSR Trudy Tbiliss. Mat. Inst. Razmadze 62 (1979), 5–27, Collection of articles on algebra, 2.

\bibitem{Atiyah} M. F. Atiyah. Complex analytic connections in fibre bundles. Trans. Amer. Math. Soc., 85:181–207, 1957.

\bibitem{BC}
V. Baranovsky, T. Chen, {\it Quantization of vector bundles on Lagrangian subvarieties}, International Mathematics Research Notices, Vol. 2017, No. 00, pp. 1–22.

\bibitem{BGKP}
V.~Baranovsky, V.~Ginzburg, D.~Kaledin, and J.~Pecharich.
\newblock Quantization of line bundles on Lagrangian subvarieties.
\newblock {\em Selecta Math. (N.S.)}, 22(1):1--25, 2016.

\bibitem{BMSZ1}
D. Barbasch, J-J. Ma, B. Sun, C.-B. Zhu, {\it Special unipotent representations of real classical groups: construction and unitarity}. arXiv:1712.05552.

\bibitem{BMSZ2} D. Barbasch, J-J. Ma, B. Sun, C.-B. Zhu,
{\it Special unipotent representations of real classical groups: counting and reduction}.
arXiv:2205.05266.

\bibitem{BMSZ3} D. Barbasch, J-J. Ma, B. Sun, C.-B. Zhu,
{\it Genuine special unipotent representations of spin groups}.
arXiv:2303.03571.

\bibitem{BMSZ4} D. Barbasch, J-J. Ma, B. Sun, C.-B. Zhu,
{\it Special unipotent representations of simple linear Lie groups of type A}.
arXiv:2304.04495.

\bibitem{Beauville} A. Beauville, {\it Symplectic singularities}. Invent. Math.  139  (2000),  no. 3, 541-549.	

\bibitem{BB} A. Beilinson, J. Bernstein, {\it A proof of Jantzen conjectures, I.M.Gelfand Seminar}, Adv. Soviet Math., vol. 16, Amer. Math. Soc., Providence, RI, 1–50, 1993.
\bibitem{BBFJLS1}
Gwyn Bellamy, Cédric Bonnafé, Baohua Fu, Daniel Juteau, Paul Levy, Eric Sommers, {\it A new family of isolated symplectic singularities with trivial local fundamental group.} Preprint: arXiv:2112.15494.

%\bibitem{BS} W. M. Beynon and N. Spaltenstein, {\it Green functions of finite Chevalley groups of type En (n = 6, 7, 8)}, J. Algebra, 88 (1984), 584–614.

\bibitem{BK} R. Bezrukavnikov, D. Kaledin, {\it Fedosov quantization in the algebraic context}.
Moscow Math. J. 4 (2004), 559-592.
\bibitem{coh} I. Biswas, P. Chatterjee, {\it On the exactness of Kostant- Kirillov form and the second cohomology of nilpotent orbits}, Internat. J. Math. 23 (2012), no. 8.
\bibitem{Bourbaki} N. Bourbaki,{\it Lie groups and Lie algebras. Chapters 4–6}.
Translated from the 1968 French original by Andrew Pressley. Elements of Mathematics (Berlin). Springer-Verlag, Berlin, 2002. xii+300 pp.
\bibitem{BPW} T. Braden, N. Proudfoot, B. Webster, {\it Quantizations of conical symplectic resolutions I: local and global structure}. Ast\'{e}risque 384 (2016), 1-73.
\bibitem{Deligne} P. Deligne, Équations différentielles à points singuliers réguliers, Lecture Notes in Mathematics, Vol. 163, Springer-Verlag, Berlin, 1970.
\bibitem{Burban_Drozd} I. Burban, Y. Drozd, {\it Maximal Cohen-Macaulay modules over surface singularities}. Trends in representation theory of algebras and related topics, 101–166, EMS Ser. Congr. Rep., Eur. Math. Soc., Zürich, 2008.
\bibitem{Carter} R.W. Carter, {\it Finite groups of Lie type: conjugacy classes and complex characters}. Wiley, 1985.
\bibitem{CM} D. Collingwood, W. McGovern, {\it Nilpotent orbits in semisimple Lie algebras}.
Chapman and Hall, London, 1993.

\bibitem{Dk0}D.~v. \DJ{okovi\'c}. 
\newblock Classification of nilpotent elements in simple exceptional real Lie algebras of inner type and description of their centralizers.
\newblock {\em J. Algebra}. 112 (1987), 503-524.

\bibitem{Dk0.5} D.~v. \DJ{okovi\'c}. 
\newblock Classification of nilpotent elements in simple real lie algebras $E_{6(6)}$ and $E_{6(-26)}$ and description of their centralizers.
\newblock {\em J. Algebra}. 116 (1988), 196-207.

\bibitem{Dk1}D.~v. \DJ{okovi\'c}.
\newblock The closure diagrams for nilpotent orbits of real forms of {$F_4$}
and {$G_2$}.
\newblock {\em J. Lie Theory}, 10(2):491--510, 2000.

\bibitem{Dk2}
D.~v. \DJ{okovi\'c}.
\newblock The closure diagrams for nilpotent orbits of real forms of {$E_6$}.
\newblock {\em J. Lie Theory}, 11(2):381--413, 2001.


\bibitem{Dk3}
D.~v. \DJ{okovi\'c}.
\newblock The closure diagrams for nilpotent orbits of the real forms {E} {VI}
and {E} {VII} of {$E_7$}.
\newblock {\em Represent. Theory}, 5:17--42, 2001.

\bibitem{Dk4}
D.~v. \DJ{okovi\'c}.
\newblock Correction to: ``{T}he closure diagrams for nilpotent orbits of the
real forms {E} {VI} and {E} {VII} of {$E_7$}'' [{R}epresent. {T}heory {\bf 5}
(2001), 17--42; {MR}1826427 (2002b:22033)].
\newblock {\em Represent. Theory}, 5:503, 2001.

\bibitem{Dk5}
D.~v. \DJ{okovi\'c}.
\newblock The closure diagram for nilpotent orbits of the split real form of
{$E_7$}.
\newblock {\em Represent. Theory}, 5:284--316, 2001.

\bibitem{Dk6}
D.~v. \DJ{okovi\'c}.
\newblock The closure diagram for nilpotent orbits of the real form {EIX} of
{$E_8$}.
\newblock {\em Asian J. Math.}, 5(3):561--584, 2001.

\bibitem{Dk7}
D.~v. \DJ{okovi\'c}.
\newblock The closure diagram for nilpotent orbits of the split real form of
{$E_8$}.
\newblock {\em Cent. Eur. J. Math.}, 1(4):573--643, 2003.

\bibitem{Dk8}
D.~v. \DJ{okovi\'c}.
\newblock Corrections for ``{T}he closure diagram for nilpotent orbits of the
split real form of {$E_8$}'' [{C}ent. {E}ur. {J}. {M}ath. {\bf 1} (2003), no.
4, 573--643 (electronic); mr2040654].
\newblock {\em Cent. Eur. J. Math.}, 3(3):578--579, 2005.

\bibitem{Fu} B. Fu, {\it On $\mathbb{Q}$-factorial terminalizations of nilpotent orbits}.
J. Math. Pures Appl. (9) 93 (2010), no. 6, 623–635.

\bibitem{FJLS1} B. Fu, D. Juteau, P. Levy, E. Sommers,
{\it Generic singularities of nilpotent orbit closures}.
Adv. Math. 305 (2017), 1-77.

\bibitem{FJLS2} B. Fu, D. Juteau, P. Levy, E. Sommers, {\it Local geometry of special pieces of nilpotent orbits}. Preprint: arXiv:2308.07398 [math.RT].

\bibitem{GG} W.L. Gan, V. Ginzburg, {\it Quantization of Slodowy slices}. Int. Math. Res. Not. 2002, no. 5, 243–255.

\bibitem{GraafElashvili}
W.A. de Graaf, A.G. Elashvili,
{\it Induced nilpotent orbits of the simple Lie algebras of exceptional type},
Georgian Math. J., 16 (2) (2009), pp. 257-278.

\bibitem{Hesselink} {W. Hesselink},
{\it Cohomologyand the resolutionof the nilpotent variety}. Math. Ann. 223, 249-252(1976).

\bibitem{Hotta} R. Hotta, K. Takeuchi, T. Tanisaki, \emph{$\mathscr{D}$-modules, perverse sheaves, and representation theory}, Progress in Mathematics 236, Birkh\"auser.

\bibitem{Kaledin_symplectic} D. Kaledin, {\it Symplectic singularities from the Poisson point of view}. J. Reine Angew. Math. 600 (2006), 135–156.

\bibitem{King} D. King, {\it The component groups of nilpotents in exceptional simple real Lie algebras}, Comm. Algebra 20 (1992), no. 1, 219–284.
\bibitem{KP_normal} H. Kraft, C. Procesi, {\it Closures of conjugacy classes of matrices are normal}. Invent. Math. 53 (1979), no. 3, 227-247.
\bibitem{KP1} H. Kraft, C. Procesi, {\it Minimal singularities in $\operatorname{GL}_n$}. Invent. Math. 62 (1981), no. 3, 503–515.
\bibitem{KP2} H. Kraft, C. Procesi, {\it On the geometry of conjugacy classes in classical groups}. Comment. Math. Helv. 57 (1982), no. 4, 539-602.

\bibitem{LT}
R. Lawther, D. M. Testerman,
{\it Centres of centralizers of unipotent elements in simple algebraic groups},
 Memoirs of the American Mathematical Society (2011) 210.

\bibitem{LY} N.C. Leung, S. Yu,
{\it Equivariant deformation quantization and coadjoint orbit method}.
arXiv:1809.09082.
\bibitem{Wquant} I.V. Losev. {\it Quantized symplectic actions and
$W$-algebras}.   J. Amer. Math. Soc. 23(2010), 35-59.
\bibitem{1dim_quant} I. Losev, {\it Quantizations of nilpotent orbits vs 1-dimensional representations of W-algebras}. arXiv:1004.1669.
\bibitem{HC} I.V. Losev. {\it Finite dimensional representations of
W-algebras}. Duke Math J. 159(2011), n.1, 99-143.
\bibitem{quant_iso} I. Losev, {\it Isomorphisms of quantizations via quantization of resolutions}. Adv. Math. 231(2012), 1216-1270.
\bibitem{W_prim} I. Losev, {\it Primitive ideals in W-algebras of type A}. J. Algebra, 359 (2012), 80-88.
\bibitem{Wdim} I. Losev, {\it Dimensions of irreducible modules over W-algebras and Goldie ranks}.
Invent. Math. 200 (2015), N3, 849-923.
\bibitem{B_ineq} I. Losev. {\it Bernstein inequality and holonomic modules}. Adv. Math. 308 (2017),  941-963.
\bibitem{Perv}  I. Losev. {\it Wall-crossing functors for quantized symplectic resolutions:
perversity and partial Ringel dualities}.  PAMQ, 13 (2017),
n.2, 247-289.
\bibitem{reg} I. Losev. {\it Quantizations of regular functions on nilpotent orbits}.  Bull. Inst. Math. Acad. Sin. 13 (2018), n.2, 199-225.
\bibitem{HC_symp} I. Losev, {\it Harish-Chandra bimodules over quantized symplectic singularities}.
Transform. Groups. 26 (2021), N2, 565-600.
\bibitem{orbit} I. Losev. {\it Deformations of symplectic singularities and Orbit method for semisimple Lie algebras}. Selecta Math, 28 (2022).
\bibitem{LMM} I. Losev, L. Mason-Brown, D. Matvieievskyi, {\it Unipotent Ideals and Harish-Chandra Bimodules}. arXiv:2108.03453.
\bibitem{MM} L. Mason-Brown, D. Matvieievskyi, {\it Unipotent Ideals for Spin and Exceptional Groups}.
arXiv:2109.09124.
%\bibitem{Namikawa} Y. Namikawa. {\it Induced nilpotent orbits and birational geometry}. Adv. Math. 222 (2009), no. 2, 547-564.
\bibitem{Namikawa} Y. Namikawa, Fundamental groups of symplectic singularities, in Higher Dimensional Algebraic Geometry –
in Honour of Professor Yujiro Kawamata’s Sixtieth Birthday, Adv. Stud. Pure Math., Vol. 74, Editors
K. Oguiso, C. Birkar, S. Ishii, S. Takayama, Math. Soc. Japan, Tokyo, 2017, 321–334, arXiv:1301.1008.

\bibitem{Nevins} M. Nevins, {\it Admissible nilpotent coadjoint orbits in the p-adic reductive groups}. Ph.D. Thesis M.I.T. Cambridge, MA (June 1998).

\bibitem{Noel1}
A.~G. No{\"e}l.
\newblock Classification of admissible nilpotent orbits in simple exceptional
real {L}ie algebras of inner type.
\newblock {\em Represent. Theory}, 5:455--493, 2001.

\bibitem{Noel2}
A.~G. No{\"e}l.
\newblock Classification of admissible nilpotent orbits in simple real {L}ie
algebras {$E_{6(6)}$} and {$E_{6(-26)}$}.
\newblock {\em Represent. Theory}, 5:494--502, 2001.

\bibitem{Ohta1}
T. Ohta, {\it Classification of admissible nilpotent orbits in the classical real Lie algebras}, J. Algebra 136 (1991), 290-333.

\bibitem{Ohta2}
T. Ohta, {\it The closures of nilpotent orbits in classical symmetric pairs and their singularities},
Tohoku Math. J. (2) 43 (1991), no. 2, 161-211.

\bibitem{Premet_slice} A. Premet. {\it Special transverse slices and their enveloping algebras}. Adv. Math. 170(2002), 1-55.

\bibitem{Premet_ideal} Premet, A. {\it Multiplicity-free primitive ideals associated with rigid nilpotent orbits}. Transformation Groups 19, 569–641 (2014).

\bibitem{RawnsleySternberg} J. Rawnsley, S. Sternberg, {\it On representations associated to the minimal nilpotent coadjoint orbit of {${\rm SL}(3,\,{\bf R})$}}, Amer. J. Math. 104 (1982), no. 6, 1153–1180.

\bibitem{Schwartz} J. Schwartz, {\it The Determination of the Admissible Nilpotent Orbits in Real Classical Groups}, Thesis, M.I.T., Cambridge, MA, 1987.

\bibitem{Torasso}  P. Torasso. {\it Quantification g\'{e}om\'{e}trique, op\'{e}rateurs d'entrelacement et
repr\'{e}sentations unitaires de {$\widetilde {\rm SL}_3({\R})$}}. Acta Math. 150 (1983), no. 3-4, 153–242.

\bibitem{realgrp} J. Sekiguchi. {\it Fundamental Groups of Semisimple Symmetric Spaces}. Representations of Lie Groups, Kyoto, Hiroshima, 1986, 519--529, Mathematical Society of Japan, Tokyo, Japan, 1988.
\bibitem{Vogan} D. Vogan, {\it Associated varieties and unipotent representations}.
Harmonic analysis on reductive groups (Brunswick, ME, 1989), 315–388, Progr. Math., 101, Birkh\"{a}user Boston, Boston, MA, 1991.
\bibitem{Yu} S. Yu, {\it The enhanced period map and equivariant deformation quantizations of nilpotent orbits}. To appear in Acta Math Sinica, English Series.
\end{thebibliography}
\end{document} 